% concatenated from BGWX.tex
\documentclass[reqno,12pt]{amsart}
%\usepackage[  left=4.00cm, right=4.00cm, top=4.00cm, bottom=4.00cm]{geometry}
%\renewcommand{\baselinestretch}{2}
\usepackage{amssymb,amsmath,amscd, amsthm,stmaryrd,verbatim,
mathrsfs,tikz,color, amsfonts, enumerate, esvect% times, %txfonts,
}
\usepackage{tikz-cd}
\usepackage[mark=o,edge length=1.5cm,root-radius=0.1cm]{dynkin-diagrams}
\usepackage{graphicx}
\usetikzlibrary{patterns}
\usetikzlibrary{er}
%\usepackage{showlabels}
\usepackage{genyoungtabtikz}
\usepackage[section]{placeins}
%\usepackage[pdfborder={0 0 0}]{hyperref}
%\usepackage[colorlinks,linkcolor=blue,urlcolor=cyan,citecolor=blue,bookmarksnumbered,bookmarksopen,pagebackref,hidelinks]{hyperref}
\usepackage[colorlinks,linkcolor=blue,urlcolor=cyan,citecolor=blue]{hyperref} 	

% \usepackage[style=alphabetic,sorting=nyt]{biblatex}
%\addbibresource{BGWX.bib}



%%% show label %%%
%\usepackage{xcolor}
%\usepackage[color]{showkeys}
%\definecolor{refkey}{gray}{.45}
%\definecolor{labelkey}{gray}{.45}
%\definecolor{citekey}{gray}{.45}
%% 
%%% paper size %%%
\usepackage[a4paper]{geometry}
\renewcommand{\baselinestretch}{1}
\geometry{top=1in, bottom=1in, left=1in, right=1in}
%%%

\numberwithin{equation}{section}
\allowdisplaybreaks
\usepackage{ 
   amsmath}
\usepackage{ 
   autobreak}

\newcommand{\bl}{0}
\newcommand{\dom}{\operatorname{dom}}
\newcommand{\caH}{\mathcal{H}}
\newcommand{\sgn}{\operatorname{sgn}}
%\renewcommand{\sp}{ \mf{sp}}
\newcommand{\Ree}{\mathrm{Re}}
\newcommand{\Imm}{\mathrm{Im}}

\newcommand{\seq}{ \mathrm{Seq}_n(\mathbb{R})}

\theoremstyle{plain}

\newtheorem{Thm}{Theorem}[section]
\newtheorem{Prop}[Thm]{Proposition}
\newtheorem{Lem}[Thm]{Lemma}
\newtheorem{Cor}[Thm]{Corollary}
\newtheorem{Conj}[Thm]{Conjecture}

\theoremstyle{definition}
\newtheorem{Defff}[Thm]{Definition}
\newtheorem{ex}[Thm]{Example}
\newtheorem{Rem}[Thm]{Remark}

\newcommand{\lam}{\lambda}

\newcommand{\Sp}{{Sp(n,\mathbb{R})}}
\newcommand{\mb}{\mathbf}
\newcommand{\bb}{\mathbb}
\newcommand{\ms}{\mathscr}
\newcommand{\mr}{\mathrm}
\newcommand{\frk}{\mathfrak}
\newcommand{\Wsmall}{W_{[\lambda]}}
\newcommand{\ptl}{\partial}
\newcommand{\dlt}{\delta}
\newcommand{\sgm}{\sigma}
\newcommand{\al}{\alpha}
\newcommand{\be}{\beta}
\newcommand{\g}{\gamma}
\newcommand{\gm}{\gamma}
\newcommand{\vs}{\varsigma}
\newcommand{\Lmd}{\Lambda}
\newcommand{\lmd}{\lambda}
\newcommand{\td}{\tilde}
\newcommand{\vf}{\varphi}
\newcommand{\yt}{Y^{\nu}}
\newcommand{\wt}{\mbox{wt}\:}
\newcommand{\der}{\mbox{Der}\:}
\newcommand{\es}{\epsilon}
\newcommand{\dmd}{\diamond}
\newcommand{\clt}{\clubsuit}
\newcommand{\vt}{\vartheta}
\newcommand{\ves}{\varepsilon}
\newcommand{\dg}{\dagger}
\newcommand{\ga}{{\cal G}({\cal A})}
\newcommand{\hga}{\hat{\cal G}({\cal A})}
\newcommand{\Edo}{\mbox{End}\:}
\newcommand{\kn}{\mbox{ker}\:}
\newcommand{\Dlt}{\Delta}
\newcommand{\rad}{\mbox{Rad}}
\newcommand{\rta}{\rightarrow}
\newcommand{\mbb}{\mathbb}
\newcommand{\rd}{\mbox{Res}\:}
\newcommand{\stc}{\stackrel{\circ}}
\newcommand{\stdt}{\stackrel{\bullet}}
\newcommand{\lra}{\Longrightarrow}
\newcommand{\f}{\varphi}
\newcommand{\rw}{\rightarrow}
\newcommand{\op}{\oplus}
\newcommand{\co}{\Omega}
\newcommand{\dar}{\Longleftrightarrow}
%\newcommand{\sta}{\theta}
%\newcommand{\stl}{\stackrel}
\newcommand{\ol}{\overline}
\newcommand{\ul}{\underline}
\newcommand{\ep}{\varepsilon}
\newcommand{\tr}{\mathrm{tr}}
\newcommand{\ad}{\mathrm{ad}}
\newcommand{\diag}{\mathrm{diag}}
\newcommand{\Cl}{\mathrm{Cl}}
\newcommand{\trace}{\mathrm{trace}}
\newcommand{\Char}{\mathrm{Char}}
\newcommand{\Ann}{\mathrm{Ann}}
\newcommand{\eq}{\begin{equation}}
\newcommand{\en}{\end{equation}}
\newcommand{\beqna}[1]{\begin{eqnarray}\label{#1}}
\newcommand{\eeqna}{\end{eqnarray}}
\newcommand{\beqn}[1]{\begin{equation}\label{#1}}
\newcommand{\eeqn}{\end{equation}}
\newcommand{\tabincell}[2]{\begin{tabular}{@{}#1@{}}#2\end{tabular}}
 \newcommand{\minitab}[2][l]{\begin{tabular}{#1}#2\end{tabular}}
\newcommand{\gk}{Gelfand--Kirillov dimension }
\newcommand{\gks}{Gelfand--Kirillov dimensions }
\newcommand{\af}{$ \mathbf{a} $}
\newcommand{\mc}[1]{\mathcal{#1}}
\newcommand{\msc}[1]{\mathscr{#1}}
\newcommand{\mf}[1]{\mathfrak{#1}}
\renewcommand{\subset}{\subseteq}
\newcommand{\rar}{\rightarrow}
\newcommand{\bil}[2]{\langle{#1},{#2}^{\vee} \rangle }
\newcommand{\hs}{ \mathfrak{h}^*}
\newcommand{\Ocat}{\mathscr{O}}
\newcommand{\aff}{ \mathbf{a} }
\newcommand{\bff}{ \mathbf{b} }
\newcommand{\cff}{ \mathbf{c} }
\newcommand{\eff}{ \mathbf{e} }
\newcommand{\gkd}{\operatorname{GKdim}}
\renewcommand{\hom}{\operatorname{Hom}}
\newcommand{\sln}{\mathfrak{sl}(n)}
\newcommand{\lambdarho}{\lambda+\rho=(\lambda_1,\lambda_2,\cdots,\lambda_n)}
\newcommand{\sn}{\mf{S}_n}
\newcommand{\su}{{SU}}
\newcommand{\nodd}{}
\newcommand{\hdom}[3]{\draw (0+#1,0-#2) rectangle (2+#1,-1-#2)++(-1,+0.5) node {$ #3$};}
\newcommand{\vdom}[3]{\draw (0+#1,0-#2) rectangle (1+#1,-2-#2)++(-0.5,+1) node {$ #3$};}
\newcommand{\hobox}[3]{\draw (0+#1,0-#2) rectangle (1+#1,-1-#2)++(-0.5,+0.5) node {$ #3$};}
\newcommand{\hdoms}[3]{\draw (0+#1,0-#2) rectangle (1.4+#1,-1-#2)++(-0.7,+0.5) node {$ #3$};}


\newcommand{\lest}{\leqslant}
\newcommand{\gest}{\geqslant}
\renewcommand{\leq}{\leqslant}
\renewcommand{\geq}{\geqslant}


\newcommand{\domscale}{0.51}
\newcommand{\sh}{\operatorname{sh}}
\newcommand{\ev}{^{\mathrm{ev}}}
\newcommand{\od}{^{\mathrm{odd}}}

\newcommand{\rf}[1]{{#1}^*}
\newcommand{\lf}[1]{^*\!{#1}}

\long\def\bzq#1#2{{\color{purple} #1 \sout{#2}}}







\begin{document}

 %\bibliographystyle{alpha}
%\bibliographystyle{plain}
\title[GK dimensions and annihilator varieties]{Gelfand--Kirillov dimensions and annihilator varieties of highest weight modules of exceptional Lie algebras
}
\author{Zhanqiang Bai}\author{Fan Gao}\author{Yutong Wang}\author{Xun Xie}

\address[Bai]{School of Mathematical Sciences, Soochow University, Suzhou 215006,  China}
\email{zqbai@suda.edu.cn}

 \address[Gao]{School of Mathematical Sciences, Zhejiang University, Hangzhou 310058,  China}
 \email{gaofan@zju.edu.cn}




\address[Wang]{School of Mathematical Sciences,  Peking University, Beijing, 100871, China}
\email{wangyutong25@stu.pku.edu.cn}

\address[Xie]{School of Mathematics and Statistics, Beijing Institute of Technology, Beijing 100081, China}
\email{xieg7@163.com}

% \address[Xiao]{College of Mathematics and statistics, Shenzhen Key Laboratory of Advanced Machine Learning and Applications, Shenzhen University,
% 	Shenzhen 518060, Guangdong, P. R. China}
% \email{xiaow@szu.edu.cn}


\subjclass[2010]{Primary 22E47; Secondary 17B10, 17B08}
\date{\today}
\keywords{Gelfand--Kirillov dimension, annihilator variety,  nilpotent orbit, Springer correspondence, Sommers duality} 



\begin{abstract}
In this paper, we give  an efficient algorithm to compute the Gelfand--Kirillov dimensions of simple highest weight modules of exceptional  Lie algebras. By using the Sommers--Achar duality, we also determine the annihilator varieties of these highest weight modules.
\end{abstract}
\maketitle




%\setcounter{tocdepth}{1}
\tableofcontents

\section{Introduction}

% \subsection{GK dimensions}

 Let $\mathfrak{g}$ be a complex semisimple Lie algebra and $U(\mathfrak{g})$ be its enveloping algebra. The Gelfand--Kirillov (GK) dimension is an important invariant to measure the size of an infinite-dimensional $U(\mathfrak{g})$-module $M$ \cite{Ge-Ki}. 
It plays significant roles in representation theory, see for example \cite{Vo78,Jo78,J81-1}. 

Annihilator variety is another invariant for infinite-dimensional $U(\mathfrak{g})$-modules.  
 For a simple $U(\mathfrak{g})$-module $M$, let $I(M)=\Ann(M)$ be its annihilator ideal in  $U(\mathfrak{g})$ and $J(M)$ be the corresponding graded ideal in $ S(\mathfrak{g}):=\text{gr} (U(\mathfrak{g}))$. The zero set of $J(M)$ in the dual vector space $\mathfrak{g}^*$ of $\mathfrak{g}$ is called {\it the annihilator variety} of $M$, and is denoted by $V(\Ann(M))$. From Joseph \cite{Jo78}, we know that $\dim V(\Ann(M))=2\gkd M$. 
In fact, Borho--Brylinski \cite{BoB1} proved that the annihilator variety of a highest weight module  with a fixed regular integral infinitesimal
	character is the Zariski closure of a nilpotent orbit in $\mathfrak{g}^*$. Joseph \cite{Jo85} extended this result to highest weight modules with  arbitrary  infinitesimal  characters.
In general, the study of Gelfand--Kirillov dimensions and annihilator varieties of  irreducible highest weight modules  is closely related with many research fields, such as representations of Weyl groups, Kazhdan--Lusztig cells, representations of Lie groups and $W$-algebras.  See for example \cite{BarV82,BMSZ-typeC,BMSZ-ABCD,BMSZ-counting,GS,GSK,GT,Mc96,Lo15,LY}.
 
 
However, in general,  these two invariants are not  easy to compute for irreducible highest weight modules. 
Nishiyama--Ochiai--Taniguchi  \cite{NOT} and Enright--Willenbring \cite{EW} independently gave characterizations of the GK dimensions  for unitary highest weight Harish-Chandra modules appearing in the dual pair settings. Bai--Hunziker \cite{Bai-Hu} found a uniform formula for the GK dimensions of   unitary highest weight modules.
By using the theory of Lusztig's $\aff$-function  \cite{lusztig1977symbol, Lu79,Lusztig1982A-class-II, lusztig1984char,lusztig1985cellsI,lusztig1985A_n}, Bai--Xie  \cite{BX} also realized a simple algorithm to compute the GK dimension of a highest weight module of $\mathfrak{sl}(n,\mathbb{C})$, which generalized Joseph's work on  GK dimensions of regular integral highest weight modules for type $A$ in \cite{Jo78}.
 Then Bai--Xiao--Xie \cite{BXX} gave certain explicit formulas of GK dimensions for simple highest weight modules of all classical Lie algebras and highest weight Harish-Chandra modules of exceptional Lie algebras.  
For the trivial infinitesimal 
character case, the computation of the annihilator variety of a highest weight module  has been done in type $A$ by Joseph \cite{Jo-com}, in types $B$ 
and $C$ by Garfinkle \cite{Garfinkle3}, and in type $D$ by McGovern and Pietraho 
\cite{MP23}. 
 Furthermore, building upon the Robinson--Schensted insertion algorithm utilized in \cite{BX,BXX}, Bai--Ma--Wang \cite{BMW} gave a simple description for the nilpotent orbit associated with the annihilator variety of a highest weight module of classical Lie algebras. However, describing these two invariants for a general highest weight module of  exceptional Lie algebras remains notably more intricate and challenging. This is one motivation for the present paper, devoted to this problem. On the other hand, Bai--Xiao \cite{BXiao} recently proved that a scalar generalized Verma module is reducible if and only if the GK dimension of its corresponding highest weight module is strictly less than the dimension of the nilradical of the corresponding parabolic subalgebra.   This also gives us another motivation to study the GK dimensions of arbitrary highest weight modules of exceptional Lie algebras.



From  \cite{BX} and \cite{BXX}, we know that the computation of the GK dimension of a highest weight module $L(\lam)$ with infinitesimal character $\lam$ can be reduced to the computation of Lusztig's $\aff$-function of a certain Weyl group element $w_{\lam}$.  For classical types, we can use the Robinson--Schensted (RS) insertion algorithm to compute the value of $\aff(w_{\lam})$ directly without determining $w_{\lam}$ explicitly. For an exceptional type Lie algebra, the computation of $\aff(w_{\lam})$ is more complicated since there is no RS algorithm and it is not easy to find out the corresponding Weyl group element $w_{\lam}$. In view of this, we first seek to give an efficient algorithm to compute the value of $\aff(w_{\lam})$, which will help us with computing the GK dimensions and annihilator varieties of  highest weight modules of exceptional Lie algebras.


Given an integral  highest weight module $L(\lambda)$ of an exceptional Lie algebra, this is achieved in two steps. First, we use Lemma \ref{find-w-lambda} to obtain the $w_{\lam}$ such that $w_{\lam}$ is minimal and $w_{\lam}^{-1}\lambda$ is antidominant. Second, we use  the computer algebra package PyCox written by Geck \cite{ge} to determine the
value of $\aff(w_\lambda)$, which in turn determines $\gkd L(\lambda)$ for integral $\lambda$.

Our first result in this paper is to generalize the above and to give an efficient algorithm to compute the value of $\aff(w_{\lam})$ for a nonintegral weight $\lam$, which also allows us to obtain $\gkd L(\lambda)$. This we elaborate now.

Fix a Cartan subalgebra $\mathfrak{h}$ of $\mathfrak{g}$. Let $\Phi$ be the root system of $(\mathfrak{g},\mathfrak{h})$ and let $W$ be its Weyl group. For any $\lam\in \mathfrak{h}^*$, the integral Weyl subgroup is denoted by $W_{[\lam]}$, whose root system is denoted by $\Phi_{[\lambda]}$. Suppose  $\Phi_{[\lambda]}=\Phi_{[\lambda]_1}\times \Phi_{[\lambda]_2}\times \cdots \times \Phi_{[\lambda]_k}\simeq \Phi_1\times \Phi_2\times \cdots\times \Phi_k$  is a direct product of irreducible root systems, and we denote such an isomorphism by $\phi$.  Thus we also get an isomorphism between the Cartan subalgebras $\mathfrak{h}_{\Phi_{[\lambda]}}=\mathfrak{h}_{\Phi_{[\lambda]_1}}\times \mathfrak{h}_{\Phi_{[\lambda]_2}}\times \cdots \times \mathfrak{h}_{\Phi_{[\lambda]_k}}$ and $\mathfrak{h}_{\Phi_1}\times \mathfrak{h}_{\Phi_2}\times \cdots\times \mathfrak{h}_{\Phi_k}$, and also between their complex duals, both still denoted by $\phi$.
From the isomorphism  $\phi$, we can write 
$$\phi(\lam|_{\mathfrak{h}_{\Phi_{[\lam]}}})=\prod\limits_{1\lest i\lest k} \phi(\lambda|_{\mathfrak{h}_{\Phi_{[\lambda]_i}}}).$$
 Denote ${\lam}'=\phi(\lam|_{\mathfrak{h}_{\Phi_{[\lam]}}})$.
Then ${\lam}'|_{\Phi_i}:={\lam}'|_{{\mathfrak{h}_{\Phi_{i}}}}=\phi(\lambda|_{\mathfrak{h}_{\Phi_{[\lambda]_i}}})$ is an integral weight of type $\Phi_i$.

\subsection{Main result} With these notations we have the following.

\begin{Thm} \label{T:main1}
    For a given nonintegral weight $\lam$, suppose that $\lambda=w_{\lambda}\mu$, where $\mu$ is antidominant and $w_{\lambda}$ is the unique minimal length element in the integral Weyl subgroup $ W_{[\lam]}$.
%     Then
% \begin{equation*}
% 	\gkd L(\lambda)=|\Phi^+|-\aff_{[\lambda]}(w_{\lambda}),
% \end{equation*}
% 	where $\aff_{[\lambda]}$ is the $\aff$-function on $W_{[\lambda]}$. 
Suppose  $\Phi_{[\lambda]}\simeq \Phi_1\times \Phi_2\times \cdots\times \Phi_k$  is a direct product of some irreducible root systems, and we denote such an isomorphism by $\phi$. Then we have  $$\aff(w_{\lambda})=\aff(w_{{\lam}'})=\sum_{1\lest i\lest k}\aff(w_{{\lam}'|_{\Phi_i}}).$$ 
In particular, when ${\Phi_i}$ is of classical type, we can determine the value of $\aff(w_{{\lam}'|_{\Phi_i}})$ directly from ${\lam}'|_{\Phi_i}$ without finding out ${w_{{\lam}'|_{\Phi_i}}}$, by using the algorithm in \cite{BXX}.
\end{Thm}




From $\dim V(\Ann(L(\lam)))=2\gkd L(\lam)$, we see that the nilpotent orbit $\mathcal{O}_{{\rm Ann}(L(\lambda))}$ associated with $V(\Ann(L(\lam)))$ can be determined by $\gkd L(\lam)$ if there is only one such nilpotent orbit with dimension equal to $2\gkd L(\lam)$. However, if there are more than one nilpotent orbits with dimension equal to $2\gkd L(\lam)$, we need to resort to the result of Joseph \cite{Jo85} to determine $\mathcal{O}_{{\rm Ann}(L(\lambda))}$ explicitly. 

For this purpose, we use $\pi_{w} \in {\rm Irr}(W)$ to denote the special representation of $W$ corresponding to the two-sided cell containing $w$, as in \cite{BMW}. For a special representation $\pi_{w}$, let $\mathcal{O}(\pi_{w})$ be the associated special nilpotent orbit via the Springer correspondence. 


From the isomorphism $\phi$, we see that $\lam'|_{{\Phi_i}}$ is an integral weight of type ${\Phi_i}$. When $\Phi_i$ is of classical type, we use the RS algorithm and H-algorithm given in \cite{BMW} to obtain a special  nilpotent orbit $\mathcal{O}(\pi_{w_{\lam'|_{{\Phi_i}}}})$. When $\Phi_i$ is of exceptional type, we  use  Lemma \ref{find-w-lambda} and PyCox to get  the corresponding character $\chi_i \in {\rm Irr}(W(\Phi_i))$ of $w_{\lam'|_{{\Phi_i}}}$ and the corresponding special  nilpotent orbit. In this way, we see that $\pi_{w_{\lam}}=\prod \pi_{w_{\lam'|_{{\Phi_i}}}} \times \prod \chi_j$ will be the special representation of $W_{[\lam]}$ and we have $\mathcal{O}_{{\rm Ann}(L(\lambda))} = \mathcal{O}_{\rm Spr}^G \circ j_{W_{[\lambda]}}^{W} (\pi_{w_\lambda})$ by Joseph \cite{Jo85}. Granted with $\pi_{w_\lambda}$, the computation of $\mathcal{O}_{{\rm Ann}(L(\lambda))}$ can be further reduced to the Sommers duality as given in \cite{Sommers}. This gives us the following:


\begin{Thm}\label{T:main2}
Keep notations as above.
Let $H^\vee \subseteq G^\vee$ be a pseudo-Levi with root system $\Phi_{[\lambda]}^\vee$. Let $H$ be its Langlands dual group, whose root system is then $\Phi_{[\lambda]}$.
Now $\pi_{w_\lambda} \in {\rm Irr}(W(H))$ is a special representation of $W(H)$. Let $\pi^{\vee}_{w_\lambda}$ be the corresponding special representation of $W(H^{\vee})$ obtained via the canonical isomorphism $W(H) \simeq W(H^\vee)$. Furthermore, let $L^\vee\subset H^\vee$ be a Levi subgroup of $H^\vee$ and $\mathcal{O}_{L^\vee}$ a distinguished nilpotent orbit of $L^\vee$ such that $d_{\rm LS}^{H^\vee}(\mathcal{O}(\pi^{\vee}_{w_\lambda})) = G\cdot \mathcal{O}_{L^\vee}$.
Then 
$$\mathcal{O}_{{\rm Ann}(L(\lambda))} = d_{\rm Som}^{\Phi^\vee}(L^{\vee}, \mathcal{O}_{L^\vee}),$$
where $d_{\rm Som}^{\Phi^\vee}$ is the duality given in \cite{Sommers}.
\end{Thm}

The content and details of Theorem \ref{T:main2} are explained in \S \ref{AVHWM}. We just note here that it essentially follows from applying Sommers' result in \cite{som-98, Sommers} to Joseph's formula above, in order to compute $\mathcal{O}_{{\rm Ann}(L(\lambda))}$ from $\pi_{w_\lambda}$ in a more direct way.

To summarize, Theorem \ref{T:main1}, Theorem \ref{T:main2} and the usage of PyCox (whenever necessary) constitute the main algorithmic steps in our determining $\gkd L(\lambda)$ and $\mathcal{O}_{{\rm Ann}(L(\lambda))}$. Moreover, the same as for the classical case treated in \cite{BMW}, we have incorporated all these in a simple online program. The program takes $\lambda$ as the input and returns $\gkd L(\lambda)$ and $\mathcal{O}_{{\rm Ann}(L(\lambda))}$. It is available at
    \begin{center}
\textcolor{blue}{http://test.slashblade.top:5000/lie/GKdim}
    \end{center}



We remark also that it is an interesting problem to determine all highest weight modules whose annihilator varieties are  equal to a given nilpotent orbit closure, see \cite{Ir85,Jo98,Tamori,BMXX,BMW,BZhang} for example. When $L(\lam)$ is a simple  highest weight module such that $\Phi_{[\lambda]}^\vee$ is pseudo-maximal (see Definition \ref{thmbd}), we tabulate all the possibilities of $\mathcal{O}(\pi_{w_{\lam}})$ or $\pi_{w_{\lam}}$ in the Appendix, when $\mathcal{O}_{{\rm Ann}(L(\lambda))}$ is equal to a fixed given nilpotent orbit.


%Bai-Ma-Xiao-Xie \cite{BMXX} gave a characterization f,or all highest weight modules whose annihilator varieties are  equal to the minimal nilpotent orbit closure. 


  
% The second main result of this paper is a full description of the associated varieties of all the highest weight HC modules. Such a result was expected for a long time. In 1992, Joseph \cite{Jo92} found a recursive relation for associated varieties of unitary highest weight modules in a case-by-case fashion. After that, the GK dimensions of holomorphic representations were obtained for Hermitian symmetric groups in the dual pair settings by Tan-Zhu \cite{tan-zhu}. Nishiyama-Ochiai-Taniguchi  \cite{NOT} and Enright-Willenbring \cite{EW} independently gave characterizations of the GK dimensions and associated varieties for unitary highest weight HC modules appearing in the dual pair settings. In recent years, a uniform formula for the GK dimensions and associated varieties of unitary highest weight HC modules was discovered by Bai-Hunziker \cite{BH}. Almost at the same time, Barchini-Zierau \cite{BZ} and Zierau \cite{Zie18} computed the characteristic cycles of highest weight HC modules with regular integral infinitesimal characters, which can give some information of the corresponding associated varieties. In this paper, we determine associated varieties of the highest weight HC modules  for all the Hermitian types (Theorems 6.2 and 7.1), where the highest weights are not necessarily regular or integral and the modules are not necessarily unitary.  Note that the case of type $ A $ was already known in \cite{BX}.

% \subsection{Lusztig's $ \aff $-functions}\label{subsec:intr2} 

% There is a beautiful formula by Lusztig \cite{lusztig1985A_n} connecting the GK dimensions and the $ \aff $-functions. Using this formula, one can translate the computation of GK dimensions of simple highest weight modules to the computation of Lusztig's $ \aff $-functions on  the integral Weyl subgroups \cite{BX}. Thus, for our purposes we need to find a fine formula of Lusztig's  $ \aff $-functions first.

% Fix a Cartan subalgebra $\mathfrak{h}$ of $\mathfrak{g}$.
% Let $\Phi$ be the root system of $(\mathfrak{g}, \mathfrak{h})$ and $W$ be its Weyl group. The function $\aff: W\rightarrow\mathbb{N}$ first appeared in Lusztig's works on representations of finite groups of Lie type and Weyl groups  \cite{lusztig1977symbol, Lu79,Lusztig1982A-class-II, lusztig1984char}. It is originally defined as a function on the set of irreducible representations of Weyl groups. In  the study of cells of affine Weyl groups, Lusztig \cite{lusztig1985cellsI,lusztig1985A_n} introduced the $ \aff $-functions on Weyl groups defined in terms of Kazhdan-Lusztig basis of Hecke algebras. He found many interesting properties about $ \aff $-functions and cells. In fact, the $ \aff $-functions on the  set of simple modules and on the Weyl group are compatible \cite[Prop. 20.6]{lusztig2003hecke}.

% It is well known that there is an algorithm of the $ \aff $-functions for Weyl groups of type $ A $, using Robinson--Schensted insertion. This gives rise to a combinatorial  algorithm of GK dimensions of simple highest weight modules of $\mathfrak{sl}(n,\mathbb{C})$ \cite[Thm. 4.6]{BX}. From this, we can determine the GK dimensions and the associated varieties of  highest weight HC modules for type $A$ (see \cite{BX} for more details). This paper is a sequel of \cite{BX}, and we will give an answer for other classical types.

% Note that the Weyl groups of type $B$ and $C$ are the same. For Weyl groups of types $ B $ and $ D $, Lusztig gave a formula for the $ \aff $-functions on the set of irreducible representations of the Weyl groups \cite{lusztig1977symbol, Lu79,lusztig1984char}. Since the irreducible representations are parametrized by Lusztig's symbols \cite{lusztig1977symbol}, we need  a certain combinatorial correspondence between  the elements of Weyl groups and the irreducible representations. This was realized in the work \cite{Barbasch-vogan1982classical} by Barbasch and Vogan. In other words, they gave a way of associating an element of the Weyl group to a symbol. Combining these results together, we get combinatorial algorithms for computing $ \aff $-functions on the Weyl groups of types $ B $ and $ D $ (Propositions \ref{prop:afB} and \ref{prop:afD}). 
% %These algorithms are known to experts, but we include a proof for convenience.
% It is quite surprising that we can simplify these algorithms and get a uniform and efficient formula of $ \aff $-functions on classical Weyl groups (Theorem \ref{thm1}).

% \subsection{A uniform formula}

% % Now we describe our formula in detail. Fix a positive integer $N$. We represent a \textit{partition} of $N$ as a decreasing sequence $p=(p_1, p_2, \cdots, p_N)$ of nonnegative integers. For simplicity, we often omit the trailing zeros in $p$. Each partition of $N$ corresponds a \textit{Young diagram}, which is an array of $N$ boxes having left-justified rows with the $i$-th row containing $p_i$ boxes for $1\lest i\lest N$. We will not distinguish a partition and its Young diagram. A \textit{Young tableau} $Y$ of shape $p$ is obtained by filling each box of the Young diagram $p$ with a number (which could be integral, real or complex, as required in context). In this case, we write $p=\sh(Y)$.


% % For a Young diagram $p$, use $ (k,l) $ to denote the box in the $ k $-th row and the $ l $-th column.
% % We say the box $ (k,l) $ is \textit{even} (resp. \textit{odd}) if $ k+l $ is even (resp. odd). Let $ p_i \ev$ (resp. $ p_i\od $) be the number of even (resp. odd) boxes in the $ i $-th row of the Young diagram $ p $. 
% % One can easily check that \begin{equation}\label{eq:ev-od}
% % 	p_i\ev=\begin{cases}
% % 		\left\lceil \frac{p_i}{2} \right\rceil&\text{ if } i \text{ is odd},\\
% % 		\left\lfloor \frac{p_i}{2} \right\rfloor&\text{ if } i \text{ is even},
% % 	\end{cases}
% % 	\quad p_i\od=\begin{cases}
% % 		\left\lfloor \frac{p_i}{2} \right\rfloor&\text{ if } i \text{ is odd},\\
% % 		\left\lceil \frac{p_i}{2} \right\rceil&\text{ if } i \text{ is even}.
% % 	\end{cases}
% % \end{equation}
% % Here for $ a\in \mathbb{R} $, $ \lfloor a \rfloor $ is the largest integer $ n $ such that $ n\lest a $, and $ \lceil a \rceil$ is the smallest integer $n$ such that $ n\geq a $. For convenience, we set
% % \begin{equation*}
% % p\ev=(p_1\ev,p_2\ev,\cdots)\quad\mbox{and}\quad p\od=(p_1\od,p_2\od,\cdots).
% % \end{equation*}

% % %Note that $  \lfloor \frac{k+1}{2} \rfloor = \lceil \frac{k}{2} \rceil$ for any $k\in \mathbb{Z}$. 

% % % \begin{ex}
% % % 	Let $ p=(5,5,4,3,3) $. The even boxes in $p$ are marked as follows.
% % % 	\[
% % % 	\young(\times~\times~\times,~\times~\times~,\times~\times~,~\times~,\times~\times)
% % % 	\]
% % % 	Then $ p\ev=(3,2,2,1,2)$ and $ p\od=(2,3,2,2,1) $.
% % % \end{ex}


% % For  a totally ordered set $ \Gamma $, we  denote by $ \mathrm{Seq}_n (\Gamma)$ the set of sequences $ x=(x_1,x_2,\cdots, x_n) $   of length $ n $ with $ x_i\in\Gamma $.
% % For the aim of this paper, it is enough to take $ \Gamma $ to be $ \mathbb{Z} $ or a coset of $ \mathbb{Z} $ in $ \mathbb{C} $. 

% % \begin{Defff}\label{def:FY}
% % 	Applying the Robinson--Schensted algorithm to $ x\in \mathrm{Seq}_n (\Gamma)$, we  get a Young tableau $ Y(x) $. Denote $ p(x)=\sh(Y(x)) $.
% % \end{Defff}


% % %A convenient reference for the Robinson--Schensted algorithm is \cite{Sagan}. Also see \cite{BX} when equal entries or non-integral entries occurred. 

% % For $ x=(x_1,x_2,\cdots,x_n)\in \seq $, set
% % \begin{equation*}
% % 	\begin{aligned}
% % 		{x}^-=&(x_1,x_2,\cdots,x_{n-1}, x_n,-x_n,-x_{n-1},\cdots,-x_2,-x_1),\\
% % 		{}^-{x}=&(-x_n,-x_{n-1},\cdots, -x_2,-x_1,x_1,x_2,\cdots, x_{n-1}, x_n).
% % 	\end{aligned}
% % \end{equation*}
% % Let $ W $ be a classical Weyl group. For $ w\in W $, denote by $\vv{w}$ the sequence $(w(1), w(2), \cdots, w(n))$.
% % %, where $w=w(1)w(2)\cdots w(n)\in W$ in the one-line notation. 


% % \begin{Thm}\label{thm1}
% % 	For an element $w$ of a classical Weyl group $ W$,
% % 	\begin{equation*}\label{eq:afunc}
% % 		\aff(w)=\begin{cases}
% % 			\sum_{i\geq1}(i-1)p(\vv w)_i &\text{ if }\Phi=A_{n-1},\\
% % 			\sum_{i\geq1}(i-1)p({}^-\vv{w})_i\od &\text{ if }\Phi=B_n,\\
% % 			\sum_{i\geq1}(i-1)p({}^-\vv{w})_i\ev &\text{ if }\Phi=D_n.
% % 		\end{cases}
% % 	\end{equation*}
% % \end{Thm}


% % \begin{ex}
% % 	Let $\Phi=D_5$ and $\vv w=(3, 4, -1, 5, 2)$. We have
% % 	\[
% % 	Y({}^-\vv{w})=\begin{tikzpicture}[scale=\domscale+0.1,baseline=-20pt]
% % 		\hobox{0}{0}{-5}
% % 		\hobox{1}{0}{-4}
% % 		\hobox{2}{0}{-3}
% % 		\hobox{3}{0}{-1}
% % 		\hobox{4}{0}{2}
% % 		\hobox{5}{0}{5}
% % 		\hobox{0}{1}{-2}
% % 		\hobox{1}{1}{1}
% % 		\hobox{2}{1}{3}
% % 		\hobox{3}{1}{4}		
% % 	\end{tikzpicture}.
% % \]
% % Hence  $p({}^-\vv{w})\ev=(3,2)$ and $\aff(w)=0\cdot3+1\cdot2=2$.
% % \end{ex}






% \subsection{GK dimensions of highest weight modules}\label{subsec:gkh}

% Based on  the above result for the $ \aff $-functions, we can deduce an efficient algorithm for the GK dimensions of simple highest weight modules for classical Lie algebras.  

% Let $\Phi^+\subset\Phi$ be the set of positive roots determined by a Borel subalgebra $\mathfrak{b}$ of $\mathfrak{g}$. Denote by $\Delta$ the set of simple roots in $\Phi^+$. For $\lambda\in\mathfrak{h}^*$, define the \textit{Verma module}
% \begin{equation*}
% M(\lambda)=U(\mathfrak{g})\otimes_{U(\mathfrak{b})}\mathbb{C}_{\lambda-\rho},
% \end{equation*}
% where $\mathbb{C}_{\lambda-\rho}$ is a one-dimensional $\mathfrak{b}$-module with weight $\lambda-\rho$ and $\rho$ is  half the sum of positive roots. Denote by $L(\lambda)$ the simple quotient of $M(\lambda)$.

% %We will present the integral case here, while the non-integral case can be found in section 5.

% \begin{Thm}\label{thm2}
% 	Let $\lambda=(\lambda_1, \lambda_2, \cdots, \lambda_n)\in \mathfrak{h}^*$ be an integral weight. Then 
% 	\begin{equation*}
% 		\gkd L(\lambda)=\begin{cases}
% 			\frac{n(n-1)}{2}-\sum_{i\geq1}(i-1)p(\lambda)_i &\text{ if }\Phi=A_{n-1},\\
% 			n^2-\sum_{i\geq1}(i-1)p(\lambda^-)_i\od &\text{ if }\Phi=B_n/C_n,\\
% 			n^2-n-\sum_{i\geq1}(i-1)p(\lambda^-)_i\ev &\text{ if }\Phi=D_n.
% 		\end{cases}
% 	\end{equation*}
% \end{Thm}
% The case of type $A$ is already known (e.g., \cite{BX}). The case of types $B/C$, for which we need to consider upper closures of chambers, is a relatively straightforward consequence of Theorem \ref{thm1}. The case of type $D$, however, is much more difficult than we expected, for which we need to prove a critical property (Theorem \ref{thm:hollow}) about the domino tableaux in \cite{Garfinkle1}. To achieve this, a new invariant, the hollow tableau, is proposed in \S4. Van Leeuwen's result \cite{Leeuwen-elementary}, which is about the connection between Garfinkle's bumping description of domino tableaux and the left-right insertion description given in earlier work of Barbasch-Vogan \cite{Barbasch-vogan1982classical}, is also necessary.

% %It seems that the hollow tableau can also be used in the theory of cells.

% %It was showed by Leeuwen \cite{Leeuwen-elementary} that this bumping description is deeply related to the left-right insertion description given in earlier work of Barbasch-Vogan \cite{Barbasch-vogan1982classical}.

% % \begin{ex}
% % 	Let $\Phi=B_8$ and $\lambda=(3, 4, 1, -2, 0, -3, 5, 6)$. Then 
% % 	\[
% % 	Y(\lambda^-)=\begin{tikzpicture}[scale=\domscale+0.1,baseline=-45pt]
% % 		\hobox{0}{0}{-6}
% % 		\hobox{1}{0}{-5}
% % 		\hobox{2}{0}{-4}
% % 		\hobox{3}{0}{-3}
% % 		\hobox{0}{1}{-3}
% % 		\hobox{1}{1}{-1}
% % 		\hobox{2}{1}{0}
% % 		\hobox{3}{1}{2}	
% % 		\hobox{0}{2}{-2}
% % 		\hobox{1}{2}{0}
% % 		\hobox{2}{2}{5}
% % 		\hobox{3}{2}{6}
% % 		\hobox{0}{3}{1}
% % 		\hobox{1}{3}{3}
% % 		\hobox{0}{4}{3}
% % 		\hobox{1}{4}{4}	
% % 	\end{tikzpicture}.
% % 	\]
% % 	Therefore, $p(\lambda^-)\od=(2, 2, 2, 1, 1)$ and 
% % 	\[
% % 	\gkd L(\lambda)=8^2-(0\cdot 2+1\cdot 2+ 2\cdot 2+ 3\cdot1+4\cdot1)=51.
% % 	\]
% % \end{ex}





% The non-integral case which needs more effort is given in Theorem \ref{thm:algo}.
% %The non-integral case needs more effort.
%  A main step is to decompose of the integral part $ \Phi_{[\lam]} $ of $ \Phi $ relative to $ \lam $ into some orthogonal subsystems. It is proved that these could be of type $ A $ or $ B $ if $ \Phi=B_n $;  $ A $, $ C$ or  $D $  if $ \Phi=C_n $; $ A $ or  $D $  if $ \Phi=D_n $.
 

% \begin{Rem}
% 	It is an easy consequence that the $\aff$-functions and the GK dimensions can also be calculated by  column lengths of the Young diagrams, using Lemma \ref{def:F}. This is useful when the tableaux have many rows and much fewer columns.
% \end{Rem}



%\subsection{Organization}
This paper is organized as follows. In \S\ref{sec:pre}, we give some necessary preliminaries. In \S \ref{a-value-compu}, for a given weight $\lam$, we explain our method of computing $\aff(w_{\lam})$, where we write $\lam=w_{\lam}\mu$ such that  $w_{\lam}$  has the minimal length in the Weyl group $W_{[\lam]}$ and $\mu $ is antidominant. In \S \ref{f4-gkd}--\S\ref{e-gkd}, we give  the GK dimensions of all the simple highest weight modules for exceptional Lie algebras. In \S\ref{AVHWM}--\S\ref{S:gen},   we show how to determine the annihilator varieties of the highest weight  modules for all exceptional Lie algebras, by using our results in \S \ref{a-value-compu} coupled with the Sommers duality. Throughout, we use Bourbaki's notation and coordinates for root systems and the relevant data, as given in \cite{Bour}.



\subsection*{Acknowledgments}
We would like to thank  M. Geck for helpful discussions about PyCox. 
Z. Bai  is partially supported  by NSFC Grant No. 12171344. F. Gao is partially supported by the National Key $\textrm{R}\,\&\,\textrm{D}$ Program of China (No. 2022YFA1005300) and  NSFC Grant No.  12171422. X. Xie is partially supported by 
NSFC Grant No. 12171030 and 12431002.



%
%%%%%%%%%%%%%%%%%%%%%%%%%%%%%%%%%%%%%%%%%%%%%%%%%%%%%%%%%%%%%%%%%%%%
%


\section{Preliminaries}\label{sec:pre}


%
%%%%%%%%%%%%%%%%%%%%%%%%%%%%%%%%%%%%%%%%%%%%%%%%%%%%%%%%%%%%%%%%%%%%
%

In this section, we recall some basic results on Lusztig's $\aff$-functions, Gelfand--Kirillov dimensions and associated varieties. See \cite{lusztig1984char, Vo78, Vo91} for more details.

Now let $\mathfrak{g}$ be a simple complex Lie algebra and $\mathfrak{h}$ be a Cartan subalgebra.
	Let $\Phi^+\subset\Phi$ be the set of positive roots determined by a Borel subalgebra $\mathfrak{b}$ of $\mathfrak{g}$. So we have a Cartan decomposition $\mathfrak{g}=\mathfrak{n}\oplus\mathfrak{h}\oplus\mathfrak{n}^-$ with $\mathfrak{b}=\mathfrak{n}\oplus \mathfrak{h}$. Denote by $\Delta$ the set of simple roots in $\Phi^+$.
% \subsection{GK dimensions and associated varieties}


% Let $\mathfrak{g}$ be a simple complex Lie algebra and $\mathfrak{h}$ be a Cartan subalgebra.
% 	Let $\Phi^+\subset\Phi$ be the set of positive roots determined by a Borel subalgebra $\mathfrak{b}$ of $\mathfrak{g}$. So we have a Cartan decomposition $\mathfrak{g}=\mathfrak{n}\oplus\mathfrak{h}\oplus\mathfrak{n}^-$ with $\mathfrak{b}=\mathfrak{n}\oplus \mathfrak{h}$. Denote by $\Delta$ the set of simple roots in $\Phi^+$.
% %	Choose a subset $I\subset \Delta$, when then generates a subsystem
% %	$\Phi_I\subset\Phi$.
% %	Let $\mathfrak{p}_I$ be the standard parabolic subalgebra corresponding to $I$ with Levi decomposition $\mathfrak{p}_I=\mathfrak{ l}_I\oplus\mathfrak{u}_I$. We frequently drop the
% %	subscript $I$ if there is no confusion.
% %	Let $F(\lambda)$ be a finite-dimensional irreducible $\mathfrak{l}$-module with highest weight $\lambda -\rho\in\mathfrak{h}^*$. It can also be viewed as a
% %	$\mathfrak{p}$-module with trivial $\mathfrak{u}$-action. 
% The Verma module $U(\mathfrak{g})\otimes_{U(\mathfrak{b})} \mathbb{C}_{\lambda-\rho}$
% has a simple quotient of $L(\lambda)$, which is a highest weight module with highest weight $\lambda-\rho$. 
% % We may say that $L(\lambda)$ is a highest weight module induced from the Levi subalgebra $\mathfrak{l}$.









% Let $M$ be a finitely generated $U(\mathfrak{g})$-module. Fix a finite dimensional generating space $M_0$ of $M$. Let $U_{n}(\mathfrak{g})$ be the standard filtration of $U(\mathfrak{g})$. Set $M_n=U_n(\mathfrak{g})\cdot M_0$ and
% \(
% \text{gr} (M)=\bigoplus\limits_{n=0}^{\infty} \text{gr}_n M,
% \)
% where $\text{gr}_n M=M_n/{M_{n-1}}$. Then $\text{gr}(M)$ is a graded module of $\text{gr}(U(\mathfrak{g}))\simeq S(\mathfrak{g})$.
% % The GK dimension of $M$ is the following invariant.
% % \cite[Thm. 1.1]{Vo78}.


% \begin{Defff} The \textit{Gelfand--Kirillov dimension} of $M$  is given by
% \begin{equation*}
% 	\operatorname{GKdim} M = \overline{\lim\limits_{n\rightarrow \infty}}\frac{\log\dim( M_n )}{\log n}.
% \end{equation*}
% \end{Defff}

%  Denote $\varphi_{M,M_0}(n)=\dim( U_n(\mathfrak{g})M_{0})$. By \cite[Chapter VII. Thm. 41]{Za-Sa}, there exists a unique polynomial $\tilde{\varphi}_{M,M_0}(n)$ such that $\varphi_{M,M_0}(n)=\tilde{\varphi}_{M,M_0}(n)$ for large $n$. The leading term
% of $\tilde{\varphi}_{M,M_0}(n)$ is $\frac{c(M)}{(d_{M})!}n^{d_{M}},$ where $c(M)$ is an integer. The integer $d_M$ is the
% Gelfand--Kirillov dimension of $M$, that is, $d_M=\operatorname{GKdim} M$. 
% In particular, $\operatorname{GKdim} M=0$ if and only if $M$ is finite-dimensional.
% %The definition of GK dimension is independent of  the choice of $M_0$.

% \begin{Defff}
% 	The  \textit{associated variety} of $M$ is defined by
% \begin{equation*}
% 	V(M):=\{X\in \mathfrak{g}^* \mid f(X)=0 \text{ for all~} f\in \operatorname{Ann}_{S(\mathfrak{g})}(\operatorname{gr} M)\}.
% \end{equation*}
% \end{Defff}

% The above two definitions are independent of the choice of $M_0$, and $\dim V(M)=\gkd M$ (see e.g. \cite{NOT}). 
% % If $M_0$ is $\mathfrak{a}$-invariant for a subalgebra $\mathfrak{a}\subset\mathfrak{g}$, then 
% % %$ V(M)\subset (\mathfrak{g}/\mathfrak{a})^* $. 
% % %(see \cite[(1.5)(b)]{Vo78})
% % \begin{equation}\label{embed}
% % 	V(M)\subset (\mathfrak{g}/\mathfrak{a})^*.
% % \end{equation}
% 	\begin{Defff} Let $\mathfrak{g}$ be a finite-dimensional semisimple Lie algebra. Let $I$ be a two-sided ideal in $U(\mathfrak{g})$. Then $\text{gr}(U(\mathfrak{g})/I)\simeq S(\mathfrak{g})/\text{gr}I$ is a graded $S(\mathfrak{g})$-module, and its annihilator ideal is $\text{gr}I \subset S(\mathfrak{g})$. We define its associated variety by
% 		$$V(I):=V(U(\mathfrak{g})/I)=\{X\in \mathfrak{g}^* \mid p(X)=0\ \mbox{for all $p\in {\text{gr}}I$}\}.
% 		$$
% 	\end{Defff}
	
% 	Following \cite{GSK}, $V(\Ann (M))$ is called the \textit{annihilator variety} of the $U(\mathfrak{g})$-module $M$. From Joseph \cite[Prop. 2.7]{Jo78}, we have $$\dim V(\Ann L(\lam))=2\gkd L(\lam).$$
	
	
	
% 	% For an irreducible representation $M$ of $\mathfrak{g}$, we usually use $\mathcal{O}(M)$ to denote this nilpotent orbit.
	
	
	
% 	% We use ${\sim_R}$ (resp. ${\sim_L}$ and ${\sim_{LR}}$ ) to denote the Kazhdan-Lusztig right (resp. left and double) cell equivalence relation \cite{KL}.
	
	
	
	
	


%  Let $\langle-, -\rangle: \mathfrak{h} \times \mathfrak{h}^* \to \mathbb{C}$ be the canonical pairing. For $\mu\in\mathfrak{h}^*$, define 
% \begin{equation*}
% \Phi_{[\mu]}:=\{\alpha\in\Phi\mid\langle\mu, \alpha^\vee\rangle\in\mathbb{Z}\},
% \end{equation*}
% where $ \alpha^\vee$ is the coroot associated with the root $\alpha \in \Phi$. 
% Denote $\Phi_{[\mu]}^+=\Phi_{[\mu]}\cap \Phi^+$ and $ \Phi_\mu^+=\{\alpha\in\Phi_{[\mu]}^+\mid\langle\mu, \alpha^\vee\rangle>0\} $. 
% Set 
% \[
% W_{[\mu]}:=\{w\in W\mid w\mu-\mu\in \mathbb{Z}\Phi\}.
% \]
% Here $\Phi_{[\mu]}$ is a root system with Weyl group $W_{[\mu]}$. 
% %(e.g., \cite{Hum08}).
%  Let $\Delta_{[\mu]}$ be the simple system of $\Phi_{[\mu]}$ lying in $\Phi^+$. Set $J=\{\alpha\in\Delta_{[\mu]}\mid\langle\mu, \alpha^\vee\rangle=0\}$. Denote by $W_J$ the Weyl group generated by reflections $s_\alpha$ with $\alpha\in J$. Let $\ell_{[\mu]}$ be the length function on $W_{[\mu]}$ with respect to the chosen $\Delta_{[\mu]}$. In particular, we have $\ell_{[\mu]}=\ell$, the length function on $W$, when $\mu$ is integral. Put
% \begin{equation*}\label{ceq1}
% 	W_{[\mu]}^J:=\{w\in W_{[\mu]}\mid \ell_{[\mu]}(ws_\alpha)=\ell_{[\mu]}(w)+1\ \mbox{for all}\ \alpha\in J\}.
% \end{equation*}
% Thus $W_{[\mu]}^J$ consists of the shortest representatives of the cosets $wW_J$ with $ w\in W_{[\mu]} $. When $\mu$ is integral, we simply write $W^J:=W_{[\mu]}^J$ .

% \begin{Defff} \label{D:ant-d}
% A weight $ \mu\in\hs $ is called \textit{dominant} if $ \bil{\mu}{\al} \notin \mathbb{Z}_{<0}$ for all $ \al\in\Phi^+ $.
% A weight $ \mu\in\hs $ is called \textit{antidominant} if $ \bil{\mu}{\al} \notin\mathbb{Z}_{>0}$ for all $ \al\in\Phi^+ $. 
% \end{Defff}
% We caution the reader that the above notion of dominance and anti-dominance may not be equivalent to the ``standard" one, unless $\lambda$ is integral.

% For any $\lambda\in\mathfrak{h}^*$, there exists a unique antidominant weight $\mu\in\hs$ and a unique $w_{\lambda}\in W_{[\mu]}^J$ such that $\lambda=w_{\lambda}\mu$. 
% %(see for example \cite[Pro. 3.5]{Hum08}). 

% \begin{Prop}[{\cite[Prop. 3.5]{Hum08}}]\label{anti}
%     Let $\lambda\in \mathfrak{h}^*$, with corresponding root system $\Phi_{[\lambda]}$ and Weyl group $W_{[\lambda]}$. Let $\Delta_{[\lambda]}$ be the simple system of  $\Phi_{[\lambda]}$ in $\Phi_{[\lambda]} \cap \Phi^+$.  Then $\lambda$ is antidominant if and only if one of the
% following three equivalent conditions holds:
% \begin{itemize}
%     \item[(1)] $\langle\lambda, \alpha^\vee\rangle\lest 0$ for all $\alpha \in \Delta_{[\lambda]}$;
%     \item[(2)] $\lambda\lest s_{\alpha}\lambda$ for all $\alpha \in \Delta_{[\lambda]}$;
%     \item[(3)]  $\lambda\lest w\lambda$ for all $w \in W_{[\lambda]}$.
% \end{itemize}
% Therefore, there is a unique antidominant weight in the orbit $W_{[\lambda]}\lambda$.
% \end{Prop}






\subsection{Hecke algebra and \texorpdfstring{$\aff$}{}-functions}\label{sec:cell}



Recall that the Weyl group $ W  $ of $ \mf{g} $ is a Coxeter group generated by $S:=\{s_\al\mid\al\in\Delta \}$. Let $\ell(-)$ be the length function on $W$ with respect to $S$. Then we have a Hecke algebra $ \mc{H} $ over $ \mathcal{A} :=\mathbb{Z}[q,q^{-1}]$, which is generated by $ T_w $, $ w\in W $ with relations \[
T_wT_{w'}=T_{ww'} \text{ if }\ell(ww')=\ell(w)+\ell(w'),
\]\[
\text{and }(T_s+q^{-1})(T_s-q)=0 \text{ for any }s\in S.
\]
The Kazhdan--Lusztig basis $C_w, w\in W$ of $ \mc{H} $ are characterized as the unique elements $ C_w \in \mc{H}$ such that
\[
\overline{C_w}=C_w,\qquad C_w\equiv T_w \mod{\mc{H}_{<0}}
\]where $ \bar{\,} :\mc{H}\rar\mc{H}$ is the bar involution such that $ \bar{q}=q^{-1} $, $ \overline{T_w} =T_{w^{-1}}^{-1}$, and $ \mc{H}_{<0}=\bigoplus_{w\in W}\mathcal{A}_{<0}T_w $, $ \mathcal{A}_{<0}=q^{-1}\mathbb{Z}[q^{-1}] $.

If $ C_y $ occurs in the expansion of $ hC_w $ (resp. $C_wh$) with respect to the KL-basis for some $ h\in\mc{H} $, then we write $ y\leftarrow_L w $ (resp. $ y\leftarrow_R w $). Extend $ \leftarrow_L $ (resp. $ \leftarrow_R $) to a preorder $ \lest_L $ (resp. $\lest _R$) on $ W $. For $x, w\in W$, write $x \lest_{LR} w$ if there exist $x=w_1, \dots, w_n=w$ such that for every $1\lest i<n$ we have either $w_i\lest_L w_{i+1}$ or $w_i\lest_R w_{i+1}$. Let $\sim_{L}$, $\sim_{R}$, $\sim_{LR}$ be the equivalence relations associated with $\lest_L$, $\lest_R$, $\lest_{LR}$ (for example, $x\sim_{L}w$ if and only if $x\lest_L w$ and $w\lest_Lx$). The equivalence classes on $W$ for $\sim_L$, $\sim_R$, $\sim_{LR}$ are called \textit{left cells}, \textit{right cells} and \textit{two-sided cells} respectively. 
% $w_1\underset{L}{\sim}w_2$

% Obviously $\{1\}$ (resp. $\{w_0\}$) is the largest (smallest) two sided-cell of $W$. Let $\mathcal{C}$ be the set of elements $w\in W\backslash\{1\}$ such that $w$ has a unique reduced expression. Denote $\mathcal{C}w_0:=\{ww_0\mid w\in\mathcal{C}\}$.

% \begin{Lem}\label{hklem2} Let $w\in W$, then
% 	\begin{itemize}
% 		\item [(1)] The sets $\mathcal{C}$ and $\mathcal{C}w_0$ are two-sided cells of $W$.
% 		\item [(2)] If $w\neq 1, w_0$, then $x\lest_{LR} w\lest_{LR} y$ for any $x\in\mathcal{C}w_0$ and $y\in\mathcal{C}$.
% 	\end{itemize}
% \end{Lem}





We have $ C_xC_y=\sum_{z\in W} h_{x,y,z}C_z $ with $ h_{x,y,x}\in\mathcal{A} $. Then Lusztig's \textit{$\aff$-function} $ \mathbf{a}:W\rar\mathbb{N} $ is defined by\[
\aff(z)=\max\{\deg h_{x,y,z}\mid x,y\in W \} \text{ for } z\in W.
\]
The following lemma is an easy consequence of Lusztig's results \cite[\S 15]{Lus03}.

\begin{Lem}\label{alem1}
Let $x, w\in W$. Then
	\begin{itemize}\item [(1)]$ \aff(w)=\aff(w^{-1}) $.
		\item [(2)] If $x\lest_{LR} w$, then $\aff(x)\geq\aff(w)$. Hence $\aff(x)=\aff(w)$ whenever $x\sim_{LR} w$.
		\item [(3)] If $ w_I $ is the longest element of the parabolic subgroup of $ W $ generated by a subset $ I\subset S $, the $ \aff(w_I)$ is equal to  the length $\ell(w_I) $ of $ w_I $.
		\item [(4)] If $ W $ is a direct product of  Coxeter subgroups $ W_1 $ and $ W_2 $, then\[
		\aff(w)=\aff(w_1)+\aff(w_2)
		\]for  $ w=(w_1,w_2) \in W_1\times W_2=W$.
	\end{itemize}
\end{Lem}


Consider
$$\mathcal{C}:=\{w\in W\backslash \{1\}: \ w \text{ has a unique reduced expression}\}.$$
Denote $\mathcal{C}w_0:=\{ww_0\mid w\in\mathcal{C}\}$, where $w_0$ is the longest element of $W$. From \cite[Prop. 3.8]{Lus83}, we know that $\mathcal{C}$ and $\mathcal{C}w_0$ are two-sided cells of $W$.



% \begin{Prop}[{\cite[Thm. 5.1.2, 5.1.3 $\&$ 5.1.4 and Prop. 6.3.7]{CM}}]\label{bcd}
% The nilpotent orbits of classical types can be identified with some partitions as follows:
% \begin{enumerate}
%     \item Nilpotent orbits in $\mathfrak{so}{(2n+1,\mathbb{C})}$ are in one-to-one correspondence with the set of partitions of $2n+1$ in which even parts occur with even multiplicity. A partition $\bf q$ of type $B$ is special if and only if its dual partition ${\bf q}^t$ is a partition of type $B$.
%     \item Nilpotent orbits in $\mathfrak{sp}{(n,\mathbb{C})}$ are in one-to-one correspondence with the set of partitions of $2n$ in which odd parts occur with even multiplicity. A partition $\bf q$ of type $C$ is special if and only if its dual partition ${\bf q}^t$ is a partition of type $C$.
%     \item Nilpotent orbits in $\mathfrak{so}{(2n,\mathbb{C})}$ are in one-to-one correspondence with the set of partitions of $2n$ in which even parts occur with even multiplicity,
%   except that each ``very even" partition ${\bf d}$ (consisting of only even parts) correspond to two orbits, denoted by $\mathcal{O}^I_{\bf d}$ and $\mathcal{O}^{II}_{\bf d}$. A partition $\bf q$ of type $D$ is special if and only if its dual partition ${\bf q}^t$ is a partition of type $C$.
% \end{enumerate}
% \end{Prop}
Recall that Lusztig singled out a subclass of representations of Weyl groups called \emph{special} in \cite{lusztig1984char}. Roughly speaking, for classical types, special representations are characterized by special symbols (equivalently special partitions). For exceptional types, special representations are listed as the first characters in each family of symbols \cite[Chap. 4]{lusztig1984char}.

Special orbits are then by definitions those corresponding to the special representations via the Springer correspondence, which we briefly recall now. Let $G$ be a connected reductive group over $\mathbb{C}$ with Lie algebra $\mathfrak{g}$. Denote by $\mathcal{N}(G)$ the set of nilpotent orbits in $\mathfrak{g}$. For $\mathcal{O} \in \mathcal{N}(G)$, we write $A_\mathcal{O}:=G_{\rm ad}^x/(G_{\rm ad}^x)^o$ for the connected component group of the centralizer subgroup  of any $x\in \mathcal{O}$ in the adjoint group $G_{\rm ad}$ of $G$. Consider
$$\mathcal{N}^{\rm en}(G):=\{(\mathcal{O}, \eta): \mathcal{O} \in \mathcal{N}(G), \eta \in {\rm Irr}(A_\mathcal{O}) \}.$$
The Springer correspondence gives an injective map
\begin{equation} \label{Spr}
  {\rm Spr}_G: {\rm Irr}(W(G)) \longrightarrow \mathcal{N}^{\rm en}(G)  
\end{equation}
denoted by ${\rm Spr}_G(\sigma) = (\mathcal{O}_{\rm Spr}^G(\sigma), \eta_\sigma)$. In particular ${\rm Spr}_G(\mathbf{1})=(\mathcal{O}_{\rm reg}, \mathbf{1})$ and ${\rm Spr}_G(\varepsilon_W)=(0, \mathbf{1})$. 
For every orbit $\mathcal{O} \in \mathcal{N}(G)$, the pair $(\mathcal{O}, \mathbf{1})$ always lies in the image of ${\rm Spr}_G$, and thus we have a well-defined map
$${\rm Spr}_{G, \mathbf{1}}^{-1}: \mathcal{N}(G) \longrightarrow {\rm Irr}(W(G))$$
given by ${\rm Spr}_{G,\mathbf{1}}^{-1}(\mathcal{O}):= {\rm Spr}_G^{-1}(\mathcal{O}, \mathbf{1})$.  If $\sigma \in {\rm Irr}(W(G))$ is a special representation, then one has ${\rm Spr}_G(\sigma) = (\mathcal{O}_{\rm Spr}^G(\sigma), \mathbf{1})$, and we call $\mathcal{O}_{\rm Spr}^G(\sigma)$ a special representation, as alluded to above.

In fact, it can be seen that ${\rm Spr}_G$ actually only depends on the Lie algebra $\mathfrak{g}$, and thus we may write all the above using $\mathfrak{g}$ as
$${\rm Spr}_\mathfrak{g}:={\rm Spr}_G, \quad \mathcal{O}_{\rm Spr}^\mathfrak{g}:=\mathcal{O}_{\rm Spr}^G, \quad {\rm Spr}_{\mathfrak{g}, \mathbf{1}}^{-1}:={\rm Spr}_{G, \mathbf{1}}^{-1}.$$











On the other hand, in \cite{BXX} or \cite{BMW},  from a classical Weyl group element $w\in W$, we can get a special  partition ${\bf p}$ and a special symbol $\Lambda_{\bf p}$ (representing the irreducible special representation ${\rm Spr}_{\mathfrak{g}, \mathbf{1}}^{-1}({\bf p})$). Then we  can compute the corresponding $\aff(w)$ from this symbol and write $\aff(\Lambda_{\mathfrak{p}})=\aff(w)$ since the $\aff$-function is constant on a two-sided cell by Lemma \ref{alem1}. The following lemma lists all the possible special partitions for a given $\aff$ value. It will be used later.
% \subsection{Algorithms for $\aff$-functions}

% Let us recall the algorithms of $\aff$-functions for classical Weyl groups. Details can be found in \cite{BXX}.




 

% Recall that in \cite{BMW},  from the partition ${\bf p}$ of a special nilpotent orbit $\mathcal{O}$ of classical type, we can get a special symbol $\Lambda_{\bf p}$, then we can compute the corresponding $\aff$ value for this symbol.
% Now we give the following lemma which gives all the possible special partitions for a given $\aff$ value. This lemma will be used later.

\begin{Lem}\label{a-value}
We denote $\aff({\bf p}):=\aff(\Lambda_{{\bf p}})$.
\begin{itemize}\item [(1)]
 The  $\aff$-function of the Weyl group of type $A_1 $ takes the following values:
   $$\aff({\bf p})=\begin{cases}
	0 &\emph{if }{\bf p}=[2],\\
	1 &\emph{if }{\bf p}=[1^2].
	\end{cases}
 $$
    
  %  $$0,1.$$
   
\item [(2)]   The  $\aff$-function of the Weyl group of type $A_2 $ takes the following values:
  %  $$0,1,3.$$
$$\aff({\bf p})=\begin{cases}
	0 &\emph{if }{\bf p}=[3],\\
	1 &\emph{if }{\bf p}=[2,1],\\
 3 &\emph{if }{\bf p}=[1^3].
	\end{cases}
 $$



   
\item [(3)]   The  $\aff$-function of the Weyl group of type $A_3 $ takes the following values:
   % $$0,1,2,3,6.$$
$$\aff({\bf p})=\begin{cases}
	0 &\emph{if }{\bf p}=[4],\\
	1 &\emph{if }{\bf p}=[3,1],\\
 2 &\emph{if }{\bf p}=[2^2],\\
 3 &\emph{if }{\bf p}=[2,1^2],\\
 6 &\emph{if }{\bf p}=[1^4].
	\end{cases}
 $$




\item [(4)]
 The  $\aff$-function of the Weyl group of type $ A_4$  takes the following values:
   % $$0,1,2,3,4,6,10.$$    
  $$\aff({\bf p})=\begin{cases}
	0 &\emph{if }{\bf p}=[5],\\
	1 &\emph{if }{\bf p}=[4,1],\\
 2 &\emph{if }{\bf p}=[3,2],\\
 3 &\emph{if }{\bf p}=[3,1^2],\\
 4 &\emph{if }{\bf p}=[2^2,1],\\
 6 &\emph{if }{\bf p}=[2,1^3],\\
 10 &\emph{if }{\bf p}=[1^5].\\
	\end{cases}
 $$





\item [(5)]
 The  $\aff$-function of the Weyl group of type $ A_5$  takes the following values:
   % $$0,1,2,3,4,6,7,10,15.$$
$$\aff({\bf p})=\begin{cases}
	0 &\emph{if }{\bf p}=[6],\\
	1 &\emph{if }{\bf p}=[5,1],\\
 2 &\emph{if }{\bf p}=[4,2],\\
 3 &\emph{if }{\bf p}=[4,1^2] \emph{~or~} [3^2],\\
 4 &\emph{if }{\bf p}=[3,2,1],\\
 6 &\emph{if }{\bf p}=[3,1^3] \emph{~or~} [2^3], \\
 7 &\emph{if }{\bf p}=[2^2,1^2], \\
10 &\emph{if }{\bf p}=[2,1^4],\\
15 &\emph{if }{\bf p}=[1^5].\\
	\end{cases}
 $$





 \item [(6)]  The  $\aff$-function of the Weyl group of type $ A_7$  takes the following values:
    %$$0,1,2,3,4,5,6,7,8,9,10,11,12,13,15,16,21,28.$$
$$\aff({\bf p})=\begin{cases}
	0 &\emph{if }{\bf p}=[8],\\
	1 &\emph{if }{\bf p}=[7,1],\\
 2 &\emph{if }{\bf p}=[6,2],\\
 3 &\emph{if }{\bf p}=[6,1^2] \emph{~or~} [5,3],\\
 4 &\emph{if }{\bf p}=[5,2,1] \emph{~or~} [4^2],\\
 5 &\emph{if }{\bf p}=[4,3,1],\\
 6 &\emph{if }{\bf p}=[5,1^3] \emph{~or~} [4,2^2], \\
 7 &\emph{if }{\bf p}=[4,2,1^2] \emph{~or~} [3^2,2], \\
8 &\emph{if }{\bf p}=[3^2,1^2],\\
% 9 &\emph{if }{\bf p}=[3,2,2,1],\\
% 10 &\emph{if }{\bf p}=[4,1^4],\\
% 11 &\emph{if }{\bf p}=[3,2,1^3],\\
% 12 &\emph{if }{\bf p}=[2^4],\\
% 13 &\emph{if }{\bf p}=[2^3,1,1],\\
% 15 &\emph{if }{\bf p}=[3,1^5],\\
% 16 &\emph{if }{\bf p}=[2,2,1^4],\\
% 21 &\emph{if }{\bf p}=[2,1^6],\\
% 28 &\emph{if }{\bf p}=[1^8].\\
	\end{cases}
  \emph{~and~} \aff(\bf p)=\begin{cases}
9 &\emph{if }{\bf p}=[3,2^2,1],\\
10 &\emph{if }{\bf p}=[4,1^4],\\
11 &\emph{if }{\bf p}=[3,2,1^3],\\
12 &\emph{if }{\bf p}=[2^4],\\
13 &\emph{if }{\bf p}=[2^3,1^2],\\
15 &\emph{if }{\bf p}=[3,1^5],\\
16 &\emph{if }{\bf p}=[2^2,1^4],\\
21 &\emph{if }{\bf p}=[2,1^6],\\
28 &\emph{if }{\bf p}=[1^8].\\
	\end{cases}
 $$

 
\item [(7)] The  $\aff$-function of the Weyl group of type $B_3 $ takes the following values:
   % $$0,1,2,3,4,9.$$ 
    $$\aff({\bf p})=\begin{cases}
	0 &\emph{if }{\bf p}=[7],\\
	1 &\emph{if }{\bf p}=[5,1^2],\\
 2 &\emph{if }{\bf p}=[3^2,1],\\
 3 &\emph{if }{\bf p}=[3,2^2],\\
 4 &\emph{if }{\bf p}=[3,1^4],\\
9 &\emph{if }{\bf p}=[1^7].\\
	\end{cases}
 $$

  \item [(8)]  The  $\aff$-function of the Weyl group of type $C_3 $ takes the following values:
   % $$0,1,2,3,4,9.$$ 
    $$\aff({\bf p})=\begin{cases}
	0 &\emph{if }{\bf p}=[6],\\
	1 &\emph{if }{\bf p}=[4,2],\\
 2 &\emph{if }{\bf p}=[3^2],\\
 3 &\emph{if }{\bf p}=[2^3],\\
 4 &\emph{if }{\bf p}=[2^2,1^2],\\
9 &\emph{if }{\bf p}=[1^6].\\
	\end{cases}
 $$
 
  \item [(9)]  The  $\aff$-function of the Weyl group of type $C_4$  takes the following values:
   % $$0,1,2,3,4,5,6,9,16.$$
$$\aff({\bf p})=\begin{cases}
	0 &\emph{if }{\bf p}=[8],\\
	1 &\emph{if }{\bf p}=[6,2],\\
 2 &\emph{if }{\bf p}=[4^2],\\
 3 &\emph{if }{\bf p}=[4,2^2],\\
 4 &\emph{if }{\bf p}=[4,2,1^2],\\
 5 &\emph{if }{\bf p}=[3^2,1^2], \\
 6 &\emph{if }{\bf p}=[2^4], \\
9 &\emph{if }{\bf p}=[2^2,1^4],\\
16 &\emph{if }{\bf p}=[1^8].\\
	\end{cases}
 $$




 % The  $\aff$-function of the Weyl group of type $B_3 \times A_1$ (or $C_3\times A_1$) takes the following values:
 %    $$0,1,2,3,4,5,9,10.$$


\item [(10)]   The  $\aff$-function of the Weyl group of type $ A_8$  takes the following values:
    % $$0,1,2,3,4,5,6,7,8,9,10,11,12,13,15,16,18,21,22,28,36.$$

$$\aff({\bf p})=\begin{cases}
	0 &\emph{if }{\bf p}=[9],\\
	1 &\emph{if }{\bf p}=[8,1],\\
 2 &\emph{if }{\bf p}=[7,2],\\
 3 &\emph{if }{\bf p}=[7,1^2] \emph{~or~} [6,3],\\
 4 &\emph{if }{\bf p}=[6,2,1] \emph{~or~} [5,4],\\
 5 &\emph{if }{\bf p}=[5,3,1],\\
 6 &\emph{if }{\bf p}=[5,2^2] \emph{~or~} [4^2,1] \emph{~or~} [6,1^3], \\
 7 &\emph{if }{\bf p}=[5,2,1^2] \emph{~or~} [4,3,2], \\
8 &\emph{if }{\bf p}=[4,3,1^2],\\
9 &\emph{if }{\bf p}=[3^3] \emph{~or~} [4,2^2,1],\\
\end{cases}$$
 and $$\aff({\bf p})=\begin{cases}
% 10 &\emph{if }{\bf p}=[3,3,2,1] \emph{~or~} [5,1^{4}],\\
% 11 &\emph{if }{\bf p}=[4,2,1^3],\\
% 12 &\emph{if }{\bf p}=[3,2^3] \emph{~or~} [3,3,1^3],\\
% 13 &\emph{if }{\bf p}=[3,2,2,1,1],\\
% 15 &\emph{if }{\bf p}=[4,1^5],\\
% 16 &\emph{if }{\bf p}=[3,2,1^4] \emph{~or~} [2^4,1],\\

% 18 &\emph{if }{\bf p}=[2^3,1^3],\\
% 21 &\emph{if }{\bf p}=[3,1^6],\\
% 22 &\emph{if }{\bf p}=[2,2,1^5],\\
% 28 &\emph{if }{\bf p}=[2,1^7],\\
% 36 &\emph{if }{\bf p}=[1^9].\\  
% \end{cases}$$
% $$
% \emph{~and~} \aff({\bf p})=\begin{cases}
% 	0 &\emph{if }{\bf p}=[9],\\
% 	1 &\emph{if }{\bf p}=[8,1],\\
%  2 &\emph{if }{\bf p}=[7,2],\\
%  3 &\emph{if }{\bf p}=[7,1,1] \emph{~or~} [6,3],\\
%  4 &\emph{if }{\bf p}=[6,2,1] \emph{~or~} [5,4],\\
%  5 &\emph{if }{\bf p}=[5,3,1],\\
%  6 &\emph{if }{\bf p}=[5,2,2] \emph{~or~} [4,4,1] \emph{~or~} [6,1^3], \\
%  7 &\emph{if }{\bf p}=[5,2,1,1] \emph{~or~} [4,3,2], \\
% 8 &\emph{if }{\bf p}=[4,3,1,1],\\
% 9 &\emph{if }{\bf p}=[3^3] \emph{~or~} %[4,2,2,1],\\
10 &\emph{if }{\bf p}=[3^2,2,1] \emph{~or~} [5,1^{4}],\\
11 &\emph{if }{\bf p}=[4,2,1^3],\\
12 &\emph{if }{\bf p}=[3,2^3] \emph{~or~} [3^2,1^3],\\
13 &\emph{if }{\bf p}=[3,2^2,1^2],\\
15 &\emph{if }{\bf p}=[4,1^5],\\
16 &\emph{if }{\bf p}=[3,2,1^4] \emph{~or~} [2^4,1],\\

18 &\emph{if }{\bf p}=[2^3,1^3],\\
21 &\emph{if }{\bf p}=[3,1^6],\\
22 &\emph{if }{\bf p}=[2^2,1^5],\\
28 &\emph{if }{\bf p}=[2,1^7],\\
36 &\emph{if }{\bf p}=[1^9].\\
	\end{cases}
 $$


% \begin{table}[htbp]
% 	\centering
% 	\renewcommand{\arraystretch}{1.5}
% 	\setlength\tabcolsep{5pt}
	
% 	\begin{tabular}{|c|c|c|c|c|c|}
% 		%\begin{tabular}{|c|c|c|c|}
% 		\hline 
% $\aff({\bf p})$ &  0  & $1$  & $ 2$  &    $3$       \\\hline
% ${\bf p}$ & $[1^9]$ & $[8,1]$ & $[7,2]$ & $[6,3]$ or $[7,1,1]$\\
% \hline


% \end{tabular}
% 	\bigskip
% 	\caption{ GKdim and AV for $H:=A_8$ (continued)}
% 	\label{taba7a1-2}
% \end{table}

  



  % The  $\aff$-function of the Weyl group of type $ D_4$  takes the following values:
  %   $$0,1,2,3,6,7,12.$$

  \item [(11)]   The  $\aff$-function of the Weyl group of type $ D_5$  takes the following values:
   % $$0,1,2,3,4,5,6,7,10,12,13,20.$$
$$\aff({\bf p})=\begin{cases}
	0 &\emph{if }{\bf p}=[9,1],\\
	1 &\emph{if }{\bf p}=[7,3],\\
 2 &\emph{if }{\bf p}=[7,1^3] \emph{~or~} [5^2],\\
 3 &\emph{if }{\bf p}=[5,3,1^2],\\
 4 &\emph{if }{\bf p}=[4^2,1^2],\\
 5 &\emph{if }{\bf p}=[3^3,1], \\
 6 &\emph{if }{\bf p}=[5, 1^5] \emph{~or~} [3^2,2^2], \\
7 &\emph{if }{\bf p}=[3^2, 1^4],\\
10 &\emph{if }{\bf p}=[2^4,1^2],\\
12 &\emph{if }{\bf p}=[3,1^7],\\
13 &\emph{if }{\bf p}=[2^2,1^6],\\
20 &\emph{if }{\bf p}=[1^{10}].\\
	\end{cases}
 $$






  \item [(12)]  The  $\aff$-function of the Weyl group of type $ D_6$  takes the following values:
  %  $$0,1,2,3,4,5,6,7,8,9,10,12,13,15,16,20,21,30.$$
$$\aff({\bf p})=\begin{cases}
	0 &\emph{if }{\bf p}=[11,1],\\
	1 &\emph{if }{\bf p}=[9,3],\\
 2 &\emph{if }{\bf p}=[9,1^3] \emph{~or~} [7,5],\\
 3 &\emph{if }{\bf p}=[7,3,1^2] \emph{~or~} [6^2],\\
 4 &\emph{if }{\bf p}=[5^2,1^2],\\
 5 &\emph{if }{\bf p}=[5,3^2,1], \\
 6 &\emph{if }{\bf p}=[7, 1^5] \emph{~or~} [5,3,2^2] \emph{~or~} [4^2,3,1], \\
7 &\emph{if }{\bf p}=[5,3, 1^4] \emph{~or~} [4^2,2^2],\\

8 &\emph{if }{\bf p}=[3^4] \emph{~or~} [4^2,1^4],\\
\end{cases}$$
 and $$\aff({\bf p})=\begin{cases}
9 &\emph{if }{\bf p}=[3^3, 1^3],\\
10 &\emph{if }{\bf p}=[3^2,2^2,1^2],\\
12 &\emph{if }{\bf p}=[5,1^7],\\
13 &\emph{if }{\bf p}=[3,3,1^6],\\
15 &\emph{if }{\bf p}=[2^6],\\
16 &\emph{if }{\bf p}=[2^4,1^4],\\
20 &\emph{if }{\bf p}=[3,1^{9}],\\
21 &\emph{if }{\bf p}=[2^2,1^8],\\
30 &\emph{if }{\bf p}=[1^{12}].\\
	\end{cases}
 $$





   \item [(13)] The  $\aff$-function of the Weyl group of type $ D_8$  takes the following values:
%     \begin{align*}
% &0,1,2,3,4,5,6,7,8,9,10,11,12,13,14,15,16,17,\\
% &18,20,21,22,23,24,28,29,30,31,34,42,43,56.
%     \end{align*}
$$\aff({\bf p})=\begin{cases}
	0 &\emph{if }{\bf p}=[15,1],\\
	1 &\emph{if }{\bf p}=[13,3],\\
 2 &\emph{if }{\bf p}=[13,1^3] \emph{~or~} [11,5],\\
 3 &\emph{if }{\bf p}=[11,3,1^2] \emph{~or~} [9,7],\\
 4 &\emph{if }{\bf p}=[9,5,1^2] \emph{~or~} [8^2],\\
 5 &\emph{if }{\bf p}=[9,3^2,1] \emph{~or~} [7^2,1^2], \\
 6 &\emph{if }{\bf p}=[11, 1^5] \emph{~or~} [9,3,2^2] \emph{~or~} [7,5,3,1], \\
7 &\emph{if }{\bf p}=[9,3, 1^4] \emph{~or~} [7,5,2^2] \emph{~or~} [6^2,3,1],\\

8 &\emph{if }{\bf p}=[7,3^3] \emph{~or~} [7,5,1^4] \emph{~or~} [6,6,2,2] \emph{~or~} [5^3,1],\\
9 &\emph{if }{\bf p}=[7,3^2, 1^3] \emph{~or~} [6^2,1^4] \emph{~or~} [5^2,3^2],\\
\end{cases}$$
 and $$\aff({\bf p})=\begin{cases}
10 &\emph{if }{\bf p}=[7,3,2^2,1^2] \emph{~or~} [5^2,3,1^3],\\
11 &\emph{if }{\bf p}=[5^2,2^2,1^2],\\
12 &\emph{if }{\bf p}=[9,1^7]  \emph{~or~ }[5,3^3,1^2] \emph{~or~}[4^4],\\
13 &\emph{if }{\bf p}=[7,3,1^6]\emph{~or~}[4^2,3^2,1^2],\\
14 &\emph{if }{\bf p}=[5^2,1^6],\\
 
15 &\emph{if }{\bf p}=[5,3,2^4] \emph{~or~}[5,3^2,1^5],\\
16 &\emph{if }{\bf p}=[4^2,2^4] \emph{~or~}[4,4,3,1^5]\emph{~or~}[5,3,2^2,1^4] \emph{~or~}[3^5,1],\\
17 &\emph{if }{\bf p}=[4^2,2^2,1^4] \emph{~or~}[3^4,2^2],\\
18 &\emph{if }{\bf p}=[3^4,1^4],\\

20 &\emph{if }{\bf p}=[7,1^{9}],\\

21 &\emph{if }{\bf p}=[5,3,1^8] \emph{~or~}[3^2,2^4,1^2],\\
\end{cases}$$
 and $$\aff({\bf p})=\begin{cases}
22 &\emph{if }{\bf p}=[4^2,1^8],\\
23 &\emph{if }{\bf p}=[3^3,1^7],\\
24 &\emph{if }{\bf p}=[3^2,2^2,1^6],\\
28 &\emph{if }{\bf p}=[2^8],\\
29 &\emph{if }{\bf p}=[2^6,1^4],\\
30 &\emph{if }{\bf p}=[5,1^{11}],\\
31 &\emph{if }{\bf p}=[3^2,1^{10}],\\
34 &\emph{if }{\bf p}=[2^4,1^8],\\
42 &\emph{if }{\bf p}=[3,1^{13}],\\
43 &\emph{if }{\bf p}=[2^2,1^{12}],\\
56 &\emph{if }{\bf p}=[1^{16}].\\
	\end{cases}
 $$
  
\end{itemize}
%  The  $\aff$-function of the Weyl group of type $ E_6$  takes the following values:
% $$0,1,2,3,4,5,6,7,10,11,12,13,15,20,25,36.$$

%     The  $\aff$-function of the Weyl group of type $ E_7$  takes the following values:
% $$0,1,2,3,4,5,6,7,8,9,10,11,12,13,15,16,20,21,22,25,30,36,37,46,63.$$
 
\end{Lem}

The computation of $\aff$-functions of classical Weyl groups can be found in \cite{BXX}.  

\subsection{GK dimensions and associated varieties}



% We may say that $L(\lambda)$ is a highest weight module induced from the Levi subalgebra $\mathfrak{l}$.









Let $M$ be a finitely generated $U(\mathfrak{g})$-module. Fix a finite dimensional generating space $M_0$ of $M$. Let $U_{n}(\mathfrak{g})$ be the standard filtration of $U(\mathfrak{g})$. Set $M_n=U_n(\mathfrak{g})\cdot M_0$ and
\(
\text{gr} (M)=\bigoplus\limits_{n=0}^{\infty} \text{gr}_n M,
\)
where $\text{gr}_n M=M_n/{M_{n-1}}$. Then $\text{gr}(M)$ is a graded module of $\text{gr}(U(\mathfrak{g}))\simeq S(\mathfrak{g})$.
% The GK dimension of $M$ is the following invariant.
% \cite[Thm. 1.1]{Vo78}.


\begin{Defff} The \textit{Gelfand--Kirillov dimension} of $M$  is given by
\begin{equation*}
	\operatorname{GKdim} M = \overline{\lim\limits_{n\rightarrow \infty}}\frac{\log\dim( M_n )}{\log n}.
\end{equation*}
\end{Defff}

 Denote $\varphi_{M,M_0}(n)=\dim( U_n(\mathfrak{g})M_{0})$. By \cite[Chapter VII. Thm. 41]{Za-Sa}, there exists a unique polynomial $\tilde{\varphi}_{M,M_0}(n)$ such that $\varphi_{M,M_0}(n)=\tilde{\varphi}_{M,M_0}(n)$ for large $n$. The leading term
of $\tilde{\varphi}_{M,M_0}(n)$ is $\frac{c(M)}{(d_{M})!}n^{d_{M}},$ where $c(M)$ is an integer. The integer $d_M$ is the
Gelfand--Kirillov dimension of $M$, that is, $d_M=\operatorname{GKdim} M$. 
In particular, $\operatorname{GKdim} M=0$ if and only if $M$ is finite-dimensional.
%The definition of GK dimension is independent of  the choice of $M_0$.

\begin{Defff}
	The  \textit{associated variety} of $M$ is defined by
\begin{equation*}
	V(M):=\{X\in \mathfrak{g}^* \mid f(X)=0 \text{ for all~} f\in \operatorname{Ann}_{S(\mathfrak{g})}(\operatorname{gr} M)\}.
\end{equation*}
\end{Defff}

The above two definitions are independent of the choice of $M_0$, and $\dim V(M)=\gkd M$ (see e.g. \cite{NOT}). 
% If $M_0$ is $\mathfrak{a}$-invariant for a subalgebra $\mathfrak{a}\subset\mathfrak{g}$, then 
% %$ V(M)\subset (\mathfrak{g}/\mathfrak{a})^* $. 
% %(see \cite[(1.5)(b)]{Vo78})
% \begin{equation}\label{embed}
% 	V(M)\subset (\mathfrak{g}/\mathfrak{a})^*.
% \end{equation}
	\begin{Defff} Let $\mathfrak{g}$ be a finite-dimensional semisimple Lie algebra. Let $I$ be a two-sided ideal in $U(\mathfrak{g})$. Then $\text{gr}(U(\mathfrak{g})/I)\simeq S(\mathfrak{g})/\text{gr}I$ is a graded $S(\mathfrak{g})$-module, and its annihilator ideal is $\text{gr}I \subset S(\mathfrak{g})$. We define its associated variety by
		$$V(I):=V(U(\mathfrak{g})/I)=\{X\in \mathfrak{g}^* \mid p(X)=0\ \mbox{for all $p\in {\text{gr}}I$}\}.
		$$
	\end{Defff}
	
	Following \cite{GSK}, $V(\Ann (M))$ is called the \textit{annihilator variety} of the $U(\mathfrak{g})$-module $M$. From Joseph \cite[Prop. 2.7]{Jo78}, we have $$\dim V(\Ann L(\lam))=2\gkd L(\lam).$$
	
	
	
	% For an irreducible representation $M$ of $\mathfrak{g}$, we usually use $\mathcal{O}(M)$ to denote this nilpotent orbit.
	
	
	
	% We use ${\sim_R}$ (resp. ${\sim_L}$ and ${\sim_{LR}}$ ) to denote the Kazhdan-Lusztig right (resp. left and double) cell equivalence relation \cite{KL}.
	
	
	
	
	


 Let $\langle-, -\rangle: \mathfrak{h} \times \mathfrak{h}^* \to \mathbb{C}$ be the canonical pairing. For $\mu\in\mathfrak{h}^*$, define 
\begin{equation*}
\Phi_{[\mu]}:=\{\alpha\in\Phi\mid\langle\mu, \alpha^\vee\rangle\in\mathbb{Z}\},
\end{equation*}
where $ \alpha^\vee$ is the coroot associated with the root $\alpha \in \Phi$. 
Denote $\Phi_{[\mu]}^+=\Phi_{[\mu]}\cap \Phi^+$ and $ \Phi_\mu^+=\{\alpha\in\Phi_{[\mu]}^+\mid\langle\mu, \alpha^\vee\rangle>0\} $. 
Set 
\[
W_{[\mu]}:=\{w\in W\mid w\mu-\mu\in \mathbb{Z}\Phi\}.
\]
Here $\Phi_{[\mu]}$ is a root system with Weyl group $W_{[\mu]}$. 
%(e.g., \cite{Hum08}).
 Let $\Delta_{[\mu]}$ be the simple system of $\Phi_{[\mu]}$ lying in $\Phi^+$. Set $J=\{\alpha\in\Delta_{[\mu]}\mid\langle\mu, \alpha^\vee\rangle=0\}$. Denote by $W_J$ the Weyl group generated by reflections $s_\alpha$ with $\alpha\in J$. Let $\ell_{[\mu]}$ be the length function on $W_{[\mu]}$ with respect to the chosen $\Delta_{[\mu]}$. In particular, we have $\ell_{[\mu]}=\ell$, the length function on $W$, when $\mu$ is integral. Put
\begin{equation*}\label{ceq1}
	W_{[\mu]}^J:=\{w\in W_{[\mu]}\mid \ell_{[\mu]}(ws_\alpha)=\ell_{[\mu]}(w)+1\ \mbox{for all}\ \alpha\in J\}.
\end{equation*}
Thus $W_{[\mu]}^J$ consists of the shortest representatives of the cosets $wW_J$ with $ w\in W_{[\mu]} $. When $\mu$ is integral, we simply write $W^J:=W_{[\mu]}^J$ .

\begin{Defff} \label{D:ant-d}
A weight $ \mu\in\hs $ is called \textit{dominant} if $ \bil{\mu}{\al} \notin \mathbb{Z}_{<0}$ for all $ \al\in\Phi^+ $.
A weight $ \mu\in\hs $ is called \textit{antidominant} if $ \bil{\mu}{\al} \notin\mathbb{Z}_{>0}$ for all $ \al\in\Phi^+ $. 
\end{Defff}
We caution the reader that the above notion of dominance and anti-dominance may not be equivalent to the ``standard" one, unless $\lambda$ is integral.

For any $\lambda\in\mathfrak{h}^*$, there exists a unique antidominant weight $\mu\in\hs$ and a unique $w_{\lambda}\in W_{[\mu]}^J$ such that $\lambda=w_{\lambda}\mu$. 
%(see for example \cite[Pro. 3.5]{Hum08}). 

\begin{Prop}[{\cite[Prop. 3.5]{Hum08}}]\label{anti}
    Let $\lambda\in \mathfrak{h}^*$, with corresponding root system $\Phi_{[\lambda]}$ and Weyl group $W_{[\lambda]}$. Let $\Delta_{[\lambda]}$ be the simple system of  $\Phi_{[\lambda]}$ in $\Phi_{[\lambda]} \cap \Phi^+$.  Then $\lambda$ is antidominant if and only if one of the
following three equivalent conditions holds:
\begin{itemize}
    \item[(1)] $\langle\lambda, \alpha^\vee\rangle\lest 0$ for all $\alpha \in \Delta_{[\lambda]}$.
    \item[(2)] $\lambda\lest s_{\alpha}\lambda$ for all $\alpha \in \Delta_{[\lambda]}$.
    \item[(3)]  $\lambda\lest w\lambda$ for all $w \in W_{[\lambda]}$.
\end{itemize}
Therefore, there is a unique antidominant weight in the orbit $W_{[\lambda]}\lambda$.
\end{Prop}

\begin{Prop}[{\cite[Prop. 3.8]{BX}}]\label{pr:main1}
	Let $ \lam\in\hs $. Suppose that $\lambda=w_{\lambda}\mu$, where $\mu$ is antidominant and $w_{\lambda}\in W_{[\mu]}^J$. Then
\begin{equation*}
	\gkd L(\lambda)=|\Phi^+|-\aff_{[\lambda]}(w_{\lambda}),
\end{equation*}
	where $\aff_{[\lambda]}$ is the $\aff$-function on $W_{[\lambda]}=W_{[\mu]}$.
\end{Prop}


For  a totally ordered set $ \Gamma $, we  denote by $ \mathrm{Seq}_n (\Gamma)$ the set of sequences $ x=(x_1,x_2,\cdots, x_n) $   of length $ n $ with $ x_i\in\Gamma $. In our paper, we usually take $\Gamma$ to be $\mathbb{Z}$ or a coset of $\mathbb{Z}$ in $\mathbb{C}$.
 \begin{Defff}[Robinson--Schensted insertion algorithm]
For an element  $ x \in  \mathrm{Seq}_n (\Gamma)$, we write  $x=(x_1,\dots,x_n)$. We associate to $x $ a  Young tableau  $ P(x) $ as follows. Let $ P_0 $ be an empty Young tableau. Assume that we have constructed Young tableau $ P_k $ associated to $ (x_1,\dots,x_k) $, $ 0\leq k<n $. Then $ P_{k+1} $ is obtained by adding $ x_{k+1} $ to $ P_k $ as follows. Firstly we add $ x_{k+1} $ to the first row of $ P_k $ by replacing the leftmost entry $ x_j $ in the first row which is \textit{strictly} bigger than $ x_{k+1} $.  (If there is no such an entry $ x_j $, we just add a box with entry $x_{k+1}  $ to the right side of the first row, and end this process). Then add this $ x_j $ to the next row as the same way of adding $x_{k+1} $ to the first row.  Finally we put $P(x)=P_n$.

\end{Defff}

Note that this  algorithm  was  originally
discovered by Robinson
 \cite{Rob38} and then rediscovered independently in a different form by Schensted \cite{Sc61}. 
 
We use $p(x)=(p_1,\dots, p_k)$ to denote the shape of $P(x)$, where $p_i$ is the number of boxes in the $i$-th row of  $P(x)$.
When $\sum\limits_{1\leq i\leq k} p_i=N$, $p(x)$ will be a partition of $N$.







	% Let $p(x)=(p_1,p_2,\cdots)$ be the shape of the Young tableau $P(x)$ obtained by applying the Robinson--Schensted algorithm (\cite{BX}) to $x\in \mathrm{Seq}_n (\Gamma)$, where $p_i$ is the number of boxes in the $i$-th row of the Young tableau $P(x)$. 




	For a Young diagram $P$, use $ (k,l) $ to denote the box in the $ k $-th row and the $ l $-th column.
	We say that the box $ (k,l) $ is \textit{even} (resp. \textit{odd}) if $ k+l $ is even (resp. odd). Let $ p_i ^{\rm ev}$ (resp. $ p_i^{\rm od} $) be the numbers of even (resp. odd) boxes in the $ i $-th row of the Young diagram $ P $.
	One can easily check that
	\begin{equation}\label{eq:ev-od}
	p_i^{\rm ev}=\begin{cases}
	\left\lceil \frac{p_i}{2} \right\rceil,&\text{ if } i \text{ is odd},\\
	\left\lfloor \frac{p_i}{2} \right\rfloor,&\text{ if } i \text{ is even},
	\end{cases}
	\quad p_i^{\rm od}=\begin{cases}
	\left\lfloor \frac{p_i}{2} \right\rfloor,&\text{ if } i \text{ is odd},\\
	\left\lceil \frac{p_i}{2} \right\rceil,&\text{ if } i \text{ is even}.
	\end{cases}
	\end{equation}
	Here for $ a\in \mathbb{R} $, $ \lfloor a \rfloor $ is the largest integer $ n $ such that $ n\leq a $, and $ \lceil a \rceil$ is the smallest integer $n$ such that $ n\geq a $. For convenience, we set
	\begin{equation*}
	p^{\rm ev}=(p_1^{\rm ev},p_2^{\rm ev},\cdots)\quad\mbox{and}\quad p^{\rm od}=(p_1^{\rm od},p_2^{\rm od},\cdots).
	\end{equation*}
	
	For $ x=(x_1,x_2,\cdots,x_n)\in \mathrm{Seq}_n (\Gamma) $, set
	\begin{equation*}
	\begin{aligned}
	{x}^-=&(x_1,x_2,\cdots,x_{n-1}, x_n,-x_n,-x_{n-1},\cdots,-x_2,-x_1).
	%{}^-{x}=&(-x_n,-x_{n-1},\cdots, -x_2,-x_1,x_1,x_2,\cdots, x_{n-1}, x_n).
	\end{aligned}
	\end{equation*}
	
	\begin{Prop}[{\cite[Thm. 1.5]{BXX}}]\label{integral}
		Let $\lambda=(\lambda_1, \lambda_2, \cdots, \lambda_n)\in \mathfrak{h}^*$ be  integral. Then
		
		\[	\aff(w_{\lambda})=\left\{
		\begin{array}{ll}
		\sum\limits_{i\gest 1}(i-1)p(\lambda)_i, &\textnormal{if}\:\Phi =A_{n-1},\\
        \sum\limits_{i\gest 1}(i-1)p((\lambda)^-)_i^{\rm od}, &\textnormal{if}\:\Phi=B_{n}/C_n,\\
		\sum\limits_{i\gest 1}(i-1)p((\lambda)^-)_i^{\rm ev}, &\textnormal{if}\:\Phi=D_{n}.
		\end{array}	
		\right.
		\]	

  
	\end{Prop}



\subsection{Annihilator variety}

 The Verma module $U(\mathfrak{g})\otimes_{U(\mathfrak{b})} \mathbb{C}_{\lambda-\rho}$
has a simple quotient  $L(\lambda)$, which is a highest weight module with highest weight $\lambda-\rho$, where $\rho$ is half the sum of positive roots of $\Phi^+$. 
Let $W$ be the Weyl group of $(\mathfrak{g}, \mathfrak{h})$.  We use $L_w, w\in W$ to denote the simple highest weight $\mathfrak{g}$-module of highest weight $-w\rho-\rho$.  We denote $I_w=\Ann(L_w)$.  Borho--Brylinski \cite{BoB1} proved that the annihilator variety of $L_w$ is irreducible, and in fact is the closure of a single nilpotent orbit $\mathcal{O}_w$. Thus, we also write  $V(I_w):= V(\Ann (L_w))=\overline{\mathcal{O}}_w$.
	

Moreover, it is known that the map $w\mapsto \mathcal{O}_w$ above induces a bijection between the two-sided cells in the Weyl group $W$ and special nilpotent orbits, see for example \cite{BarV82, BoB1, Ta}. In other words, $\mathcal{O}_w=\mathcal{O}_y$ if and only if $w\sim_{LR}y$. 
  
We summarize the above as follows.

\begin{Prop} \label{intehral-W}
The map $w\mapsto \mathcal{O}_w$ and the Springer correspondence induce a bijective correspondence between the set of two-sided cells in $W_{[\lambda]}$ and the set  of special representations of $W_{[\lambda]}$.
	\end{Prop}
    
% Let $ P $ be a Young tableau. We denote by $ (a,b) $ the box in the $ a $-th row and $ b $-th column of $ P$. We call $ (a,b) $  an even (resp. odd) box if $ a+b $ is \textit{even} (resp. \textit{odd}). Let $ p=\sh(P) $ be the shape of $ P $, which by definition is a partition $ [p_1,p_2,\dots] $ so that $ p_i $ is the length of the $ i $-th row of $ P $. We define
% \[
% p\ev=(p_1\ev,p_2\ev, \cdots) \quad(\text{resp. }p\od=(p_1\od,p_2\od, \cdots) )
% \]
% such that $ p_i\ev $ (resp. $ p_i\od $) is the number of even (resp. odd) boxes in the $ i $-th row of $ P$.




% Let $ x=(x_1,\cdots, x_m) $ be a sequence of elements from a totally ordered set. One can obtain a Young tableau $ P(x) $ by applying Robinson--Schensted algorithm to $ x $ (see \cite[3.1]{Sagan} and \cite[4.1]{BX}).  Let $ p(x) $ be the shape of $ P(x) $. Set
% \begin{align}
% 	F_a(x)&=\sum_{k\geq 1} (k-1) p_k,\\
% 	F_b(x)&=\sum_{k\geq 1} (k-1) p_k\ev,\\
% 	F_d(x)&=\sum_{k\geq 1} (k-1) p_k\od,
% \end{align}
% where $ p=p(x)=(p_1,p_2,\cdots)$.

% \begin{Thm}[{\cite[Thm. 1.3]{BXX}}] Let $w=(w(1),w(2),\dots, w(n))\in W$ in the one-line notation. Then $\aff(w)=F_a( w)$ $($resp. $F_b({}^- w)$, $F_d({}^- w))$ for $\Phi=A_{n-1}$ $($resp. $\Phi=B_n$, $\Phi=D_n)$, where $ w=(w(1), w(2), \cdots, w(n))$ and
% 	\[
% 	{}^- w=(-w(n), -w(n-1), \cdots, -w(1), w(1), w(2), \cdots, w(n)).
% 	\]
% \end{Thm}



	
	
	
	% Then we have the following result.
	
	
	% \begin{Prop}[ \cite{BoB1}]
	% 	$V(I_w)=\overline{\mathcal{O}}_w$.
		
		
		
		
	% \end{Prop}




% \subsection{GK dimensions of highest weight modules for classical types}
% For  a totally ordered set $ \Gamma $, we  denote by $ \mathrm{Seq}_n (\Gamma)$ the set of sequences $ x=(x_1,x_2,\cdots, x_n) $   of length $ n $ with $ x_i\in\Gamma $. In our paper, we usually take $\Gamma$ to be $\mathbb{Z}$ or a coset of $\mathbb{Z}$ in $\mathbb{C}$.
% 	Let $p(x)$ be the shape of the Young tableau $P(x)$ obtained by applying Robinson--Schensted algorithm (\cite{BX}) to $x\in \mathrm{Seq}_n (\Gamma)$. 




% 	For $ x=(x_1,x_2,\cdots,x_n)\in \mathrm{Seq}_n (\Gamma) $, set
% 	\begin{equation*}
% 	\begin{aligned}
% 	{x}^-=&(x_1,x_2,\cdots,x_{n-1}, x_n,-x_n,-x_{n-1},\cdots,-x_2,-x_1).
% 	%{}^-{x}=&(-x_n,-x_{n-1},\cdots, -x_2,-x_1,x_1,x_2,\cdots, x_{n-1}, x_n).
% 	\end{aligned}
% 	\end{equation*}
	
% 	\begin{Prop}[{\cite[Theorem 1.5]{BXX}}]\label{integral}
% 		Let $\lambda=(\lambda_1, \lambda_2, \cdots, \lambda_n)\in \mathfrak{h}^*$ be an integral weight. Then 
% 	\begin{equation*}
% 		\gkd L(\lambda)=\begin{cases}
% 			\frac{n(n-1)}{2}-\sum_{i\geq1}(i-1)p(\lambda)_i &\text{ if }\Phi=A_{n-1},\\
% 			n^2-\sum_{i\geq1}(i-1)p(\lambda^-)_i\od &\text{ if }\Phi=B_n/C_n,\\
% 			n^2-n-\sum_{i\geq1}(i-1)p(\lambda^-)_i\ev &\text{ if }\Phi=D_n.
% 		\end{cases}
% 	\end{equation*}

  
% 	\end{Prop}
	
% 	When $\lambda$ is non-integral,  we need some more notations before we give the algorithm of GK dimension.
	
	
		
% 		For $ \mathfrak g= \mathfrak{sl}(n, \mathbb{C})$,  we define $[\lambda]$ to be the set of  maximal subsequence $ x $ of  $ \lambda+\rho $ such that any two entries of $ x $ has an integral difference.
		
% 		For $\mathfrak{g}= \mathfrak{so}(2n+1, \mathbb{C}) $,  $\mathfrak{sp}(n, \mathbb{C}) $ or $\mathfrak{so}(2n, \mathbb{C}) $, we define $[\lambda] $ to be the set of  maximal subsequence $ x $ of  $ \lambda+\rho $ such that any two entries of $ x $ have an integral  difference or sum. In this case, we set $ [\lambda]_1 $ (resp. $ [\lambda]_2 $)  to be the subset of $ [\lambda] $ consisting of sequences with  all entries belonging to $ \mathbb{Z} $ (resp. $ \frac12+\mathbb{Z} $).
% 		Since there is at most one element in $[\lambda]_1 $ and $[\lambda]_2 $, we denote them by  $(\lambda)_{(0)}$ and $(\lambda)_{(\frac{1}{2})}$.
% 		We set $[\lambda]_{1,2}=[\lambda]_1\cup [\lambda]_2, \quad [\lambda]_3=[\lambda]\setminus[\lambda]_{1,2}$.
		
% 		%	and set $$ [\lambda]_{1,2}=[\lambda]_1\cup [\lambda]_2, \quad [\lambda]_3=[\lambda]\setminus[\lambda]_{1,2}  .$$
	
	
	
	

% 		Let $\lambda\in\mathfrak{h}^{\ast}$ and $ x=(\lambda_{i_1}, \lambda_{i_2},\cdots \lambda_{i_r})\in[\lambda]_3 $. Let $  y=(\lambda_{j_1}, \lambda_{j_2},\cdots, \lambda_{j_p}) $ be the maximal subsequence of $ x $ such that $ j_1=i_1 $ and the difference of any two entries of $ y$ is an integer. Let $ z= (\lambda_{k_1}, \lambda_{k_2},\cdots, \lambda_{k_q}) $ be the subsequence obtained by deleting $ y$ from $ x $, which is possible empty.
% 		Define
% 		$$  \tilde{x}=(\lambda_{j_1}, \lambda_{j_2},\cdots, \lambda_{j_p}, -\lambda_{k_q}, -\lambda_{k_{q-1}},\cdots ,-\lambda_{k_1}).  $$
	
	
	
	
% 	\begin{Prop}[{\cite[Theorem 4.6]{BX}} and {\cite[Theorem 5.7]{BXX}} ]\label{GKdim}
% 		The GK dimension of  $ L(\lambda) $  can be computed as follows.
% 		\begin{enumerate}
% 			\item If $  \mathfrak{g}= \mathfrak{sl}(n,\mathbb{C})$,
% 			\[
% 			\gkd L(\lambda)=\frac{n(n-1)}{2}-\sum _{x\in [\lambda]} F_a(x).
% 			\]
			
% 			\item  If $ \mathfrak{g}= \mathfrak{so}(2n+1,\mathbb{C})$, then
% 		\[
% 		\gkd L(\lambda)=n^2-F_b(\lambda_{(0)}^-)-F_b(\lambda_{(\frac{1}{2})}^-)-\sum _{x\in [\lambda]_3} F_a(\tilde{x}).
% 		\]
% 		\item  If $ \mathfrak{g}= \mathfrak{sp}(n,\mathbb{C}) $, then
% 		\[
% 		\gkd L(\lambda)=n^2-F_b(\lambda_{(0)}^-)-F_d(\lambda_{(\frac{1}{2})}^-)-\sum _{x\in [\lambda]_3} F_a(\tilde{x}).
% 		\]
		
% 			\item  If $  \mathfrak{g} = \mathfrak{so}(2n,\mathbb{C}) $,
% 			\[
% 			\gkd L(\lambda)=n^2-n-F_d((\lambda)_{(0)}^-)-F_d((\lambda)_{(\frac{1}{2})}^-)-\sum _{x\in [\lambda]_3} F_a(\tilde{x}).
% 			\]
% 		\end{enumerate}
% 	\end{Prop}





% \subsection{Partitions and dual partitions} 
% %All the above invariants can be calculated by partitions. 
% We say $q=(q_1, \cdots, q_N)$ is the \textit{dual partition} of a partition $p=(p_1, \cdots, p_N)$ and write $q={}^tp$ if $q_i$ is the length of $i$-th column of the Young diagram $p$. 
% Recall that $p(x)$ the shape of the Young tableau $Y(x)$ obtained by applying Robinson--Schensted algorithm to $x\in \mathrm{Seq}_n (\Gamma)$. For convenience, we set $q(x)={}^tp(x)$.

% \begin{Lem}[\textbf{Definition}]\label{def:F}
% 	 Let $x\in \mathrm{Seq}_n (\Gamma)$. Then
% 	\begin{equation*}
% 	\begin{aligned}
% 		F_a(x):&=\sum_{i\geq1}(i-1)p_i=\sum_{i\geq1}\frac{q_i(q_i-1)}{2},\\
% 		F_b(x):&=\sum_{i\geq1}(i-1)p_i\od=\sum_{2\nmid i}(q_i\od)^2+\sum_{2\mid i}q_i\od(q_i\od-1),\\
% 		F_d(x):&=\sum_{i\geq1}(i-1)p_i\ev=\sum_{2\nmid i}q_i\ev(q_i\ev-1)+\sum_{2\mid i}(q_i\ev)^2,
% 	\end{aligned}
%    \end{equation*}
% where $p=p(x)=(p_1,p_2,\cdots, p_n) $ and $q=q(x)={}^tp(x)=(q_1, q_2, \cdots, q_n)$.
% \end{Lem}





%\begin{Defff}
%	For $x\in \mathrm{Seq}_n (\Gamma)$, an \emph{increasing} (respectively, \emph{decreasing}) \emph{subsequence} of $x$ is a sequence $(x_{i_1}, x_{i_2}, \cdots, x_{i_m})$ with $i_1<i_2<\cdots<i_m$ such that $x_{i_j}-x_{i_k}\in\mathbb{Z}_{\lest 0}$ (resp. $x_{i_j}-x_{i_k}\in\mathbb{Z}_{>0}$) for $1\lest j<k\lest m$. %The integer $m$ is the length of the subsequence.
%\end{Defff}

% For $x\in \mathrm{Seq}_n (\Gamma)$ and $k\lest n$, denote by $c_k(x)$ (resp. $d_k(x)$) the maximal length of
% subsequences which can be written as a union of $k$ increasing (resp. decreasing) subsequences of $x$. The following result is due to \cite{greene1974}.

% \begin{Thm}[Greene's Theorem]\label{gthm}
% 	Let $x\in \mathrm{Seq}_n (\Gamma)$ and $k\lest n$. Then
% 	\begin{equation*}
% 	c_k(x)=p_1+p_2+...+p_k\mbox{ and } d_k(x)=q_1+q_2+...+q_k,
% 	\end{equation*}
% 	where $p=p(x)=(p_1,p_2,\cdots, p_n)$ and $q=q(x)={}^tp(x)=(q_1,q_2,\cdots, q_n)$.
% \end{Thm}
Let $\mathfrak{g}$ be a finite dimensional semisimple Lie algebra with a fixed Cartan subalgebra $\mathfrak{h}$.
% Let $W$ be the Weyl group of $(\mathfrak{g}, \mathfrak{h})$.  We use $L_w, w\in W$ to denote the simple highest weight $\mathfrak{g}$-module of highest weight $-w\rho-\rho$.  We denote $I_w=\Ann(L_w)$.  Borho--Brylinski \cite{BoB1} proved that the annihilator variety of $L_w$ is irreducible, and in fact is the closure of a single nilpotent orbit $\mathcal{O}_w$. Thus, we also write  $V(I_w):= V(\Ann (L_w))=\overline{\mathcal{O}}_w$.
	%  Following this, one has the special nilpotent orbits, which correspond to the special representations of the Weyl group via the Springer correspondence.
% Moreover, it is known that the map $w\mapsto \mathcal{O}_w$ above induces a bijection between the two-sided cells in the Weyl group $W$ and special nilpotent orbits, see for example \cite{BarV82, BoB1, Ta}. In other words, $\mathcal{O}_w=\mathcal{O}_y$ if and only if $w\sim_{LR}y$.  
% We summarize the above as follows.
% \begin{Prop} \label{intehral-W}
% The map $w\mapsto \mathcal{O}_w$ and the Springer correspondence induce a bijective correspondence between the set of two-sided cells in $W_{[\lambda]}$ and the set  of special representations of $W_{[\lambda]}$.
% 	\end{Prop}
Let $W$ be the Weyl group attached to a root system $\Phi$ (with a fixed subset of simple roots $\Delta$) in a finite dimensional real vector space $V$. Then $W$ acts on the space $P_k(V)$ of degree $k$ homogeneous polynomials on $V$.  
Following Joseph \cite{J80-2}, a representation $\sigma\in {\rm Irr}(W)$ is called \emph{univalent} if it
occurs with multiplicity one in  $P_{b(\sigma)}(V)$ where $b(\sigma)$ is the minimal degree  such that $\sigma$ occurs in $P_{b(\sigma)}(V)$.
The number $b(\sigma)$ is called the \emph{fake degree} of $\sigma$.
Let ${\rm Irr}(W)^{\rm uv} \subset {\rm Irr}(W)$ be the set of univalent representations. 
Under the Springer correspondence, the representation associated with a nilpotent orbit and its trivial local system is always univalent \cite[Corollary~4]{BM}, i.e., ${\rm Im}({\rm Spr}_{G,\mathbf{1}}^{-1}) \subseteq {\rm Irr}(W)^{\rm uv}$ using the notation from \ref{Spr}.

Suppose $W'$ is a subgroup of $W$ generated by reflections in a root subsystem of $\Phi$.   
The $j$-induction from $W'$ to $W$ is a well defined map 
\[
\begin{array}{rcl}
j_{W'}^{W} \colon {\rm Irr}(W')^{\rm uv}& \rightarrow &{\rm Irr}(W)^{\rm uv} \\
\sigma' & \mapsto & j_{W'}^{W}\sigma'
\end{array}
\]
where $j_{W'}^{W}\sigma'$ is the representation generated by the $\sigma'$ 
isotypical component in $P_{b(\sigma')}(V)$. Furthermore, $b(j_{W'}^{W}\sigma') =b(\sigma')$ and $j_{W'}^{W}\sigma'$ occurs in $P_{b(\sigma')}(V)$ with multiplicity one.  
The definitions of ``univalent'' and ``$j$-induction'' are independent of the choice of $V$. 

The $j$-induction satisfies the property of induction by stages as follows. 
\begin{Lem}[{See \cite[Theorem~11.2.4]{Ca85}}]\label{j.1}
Let $W''\subset W'$ be two subgroups of $W$ generated by two root subsystems (which are not necessary the parabolic subgroups). 
Then 
\[
j^W_{W''}=j^W_{W'}\circ j^{W'}_{W''}.
\]
\end{Lem}
Some details can be found in \cite[\S 11.2]{Ca85} and \cite[Chapter~4]{lusztig1984char}.


Let $\lambda\in \mathfrak{h}^*$, with associated root system $\Phi_{[\lambda]}$ and Weyl subgroup $W_{[\lambda]}$. Then $\lam$ is an integral weight for $\Phi_{[\lambda]}$. Suppose that $\lambda=w_{\lambda}\mu$, where $\mu$ is antidominant and $w_{\lambda}\in W_{[\mu]}^J$. Then we use $\pi_{w_{\lam}}$ to denote the special representation of $W_{[\lambda]}$ corresponding to the two-sided cell of $W_{[\lambda]}$ containing $w_{\lam}$. 
By using the Springer correspondence and $j$-induction operator, Joseph extended the result in \cite{BoB1} to arbitrary infinitesimal character, which gives the following result.

 % We denote $\tilde{\pi}_w:=j_{W_{[\lambda]}}^W(\pi_w)$, which is a representation of $W$.  
 % Then we have the following result.



 
\begin{Prop}[{\cite[Theorem~3.10]{Jo85}}]\label{2dim}
Let $\mu\in \mathfrak{h}^*$ be an anti-dominant element and   $w\in W_{[\mu]}^J$. 
    Then $\tilde{\pi}_w:=j_{W_{[\mu]}}^W(\pi_w)$ is an
    irreducible $W$-module.   This $W$-module $\tilde{\pi}_w$ corresponds to  a nilpotent orbit $\mathcal{O}_{\tilde{\pi}_w}$
    with trivial local system via the Springer correspondence. Furthermore,
$$V(\Ann (L(w\mu)))=\overline{\mathcal{O}_{\rm Spr}^\mathfrak{g} ( \tilde{\pi}_w)}.$$
\end{Prop}




    
	% \begin{Prop}[ \cite{Jo85}]\label{2dim}
	% 	Let $\mathfrak{g}$ be a reductive Lie algebra and $I$ be a primitive ideal in $U(\mathfrak{g})$. Then $V(I)$ is the closure of a single nilpotent coadjoint orbit $\mathcal{O}_I$ in $\mathfrak{g}^*$. 
 %  % Also we have $$2\gkd L(\lam)=\dim \mathcal{O}_{Ann(L(\lambda))}.$$
	% \end{Prop}
 Henceforth, for a highest weight module $L(\lambda)$, Proposition \ref{2dim} gives 
 $$V(\Ann (L(\lambda)))=\overline{\mathcal{O}}_{\Ann(L(\lambda))}$$
 for a nilpotent orbit $\mathcal{O}_{{\rm Ann}(L(\lambda))}$ that is uniquely determined by $L(\lambda)$. We call $\mathcal{O}_{{\rm Ann}(L(\lambda))}$ the annihilator variety or orbit associated with $L(\lambda)$.
 
\subsection{PyCox}
In this subsection, we  illustrate how to use the computer algebra package PyCox  to find the nilpotent orbit $\mathcal{O}_w$ for a given element $w\in W$. More details about PyCox can be found in \cite{ge}.

From now on, for a simple root $\alpha_i\in \Delta$ we may simply write the simple reflection $s_{\alpha_i}$ by $s_i$. 
 When there are too many factors in an element $w$, we  adopt Geck's notation in \cite{ge} and use it to compute the value of $\aff(w)$. For example, we use $[i_1-1,i_2-1,\cdots,i_k-1]$ to denote
$ w=s_{i_1}s_{i_2}\cdots s_{i_k}$ for type $A_n$, $E_n$, $F_4$ or $G_2$, and use $[n-i_1,n-i_2,\cdots,n-i_k]$ to denote
$ w=s_{i_1}s_{i_2}\cdots s_{i_k}$ for type $B_n$, $C_n$ and $D_n$.




In PyCox, the function ``klcellrepem''   returns us the 
 $\aff$ value of $w \in W$ and the character of the left cell to which $w$ belongs. Note that the notation for characters in PyCox is different with Carter \cite[p. 428]{Ca85}. 

\begin{ex}
    In the case of $F_4$, we use $w=[1,3,2,1]$ to denote the element $w=s_2s_4s_3s_2$. Using PyCox, we have:
   \begin{verbatim}
    >>> W = coxeter("F", 4)
    >>> print(klcellrepelm(W, [1,3,2,1]))
    {'size': 9, 'character': [['9_1', 1]], 'a': 2,
    'special': '9_1', 'index': 22, 'elms': False, 
    'distinv': False}
\end{verbatim}

From this we know that $\aff(w)=2$ and the corresponding character is $9_1$. By comparing \cite[Table C.3]{GP} and \cite[p. 428]{Ca85}, we will know that the corresponding character is  $\phi_{9,2}$ in the notation of \cite[p. 428]{Ca85}. Also, the corresponding special nilpotent orbit is $\mathcal{O}_w=F_4(a_2)$.
\end{ex}
 
 
 
 
 % Then from \cite[Table C.3]{GP}, we will know the corresponding character.
% \begin{figure}[!t]
% \includegraphics[scale=0.3]{a-value.pdf}
% \end{figure}
	
% \begin{figure}
%     \centering
%     \includegraphics[width=0.5\linewidth]{a-value.pdf}
%     \caption{Enter Caption}
%     \label{fig:PyCox}
% \end{figure}
% \begin{figure}[h]
%     \centering
%     \includegraphics[width=1.0\linewidth]{a-fun-value.png}
%     \caption{PyCox example}
%     \label{fig:1}
% \end{figure}

% In \cite[Chapter 13.3]{Ca85}, we can find the correspondence between characters of $W$ and Bala--Carter labels of unipotent classes (nilpotent orbits). 

Thus for exceptional types, by using PyCox we can determine the nilpotent orbit $\mathcal{O}_w$ for a given element $w\in W$.


\subsection{Pseudo-maximal root subsystem}
%Now we recall some definitions and properties about closed subsystems and maximal subsystems. More details  can be found in \cite[\S12]{Kan01}. 


%\begin{Defff}\label{defcl}
%	We say a subsystem $\Phi'\subset\Phi$ is \textit{closed} if $\alpha+\beta\in\Phi'$ whenever $\alpha, \beta\in\Phi'$ and $\alpha+\beta\in\Phi$.
%\end{Defff}



%From Definition \ref{defcl}, we have the following result.

%\begin{Lem}\label{lemni1}
%	Let $\lambda\in\mathfrak{h}^*$. Then $\Phi_{[\lambda]}^\vee$ is a closed subsystem of $\Phi^\vee$.
%\end{Lem}

Let $\Phi$ be a root system and let
\[
\alpha_0=\sum_{i=1}^nh_i\alpha_i
\]
be the highest root of $\Phi$, where $h_i$ are non-negative integers and $\Delta=\{\alpha_i\mid 1\leq i\leq n\}$ is the set of simple roots in $\Phi^+$. 
%Denote $\beta:=\alpha_0$.
For every $1\lest i, j \lest n$, we set
$$\tilde{\Phi}(i):=\Phi \cup \{-\alpha_0\} - \{\alpha_i\}, \quad \Phi(j):=\Phi - \{\alpha_j\}.$$

\begin{Defff}\label{thmbd}
	Let $\Phi$ be irreducible. A root subsystem of $\Phi$ is called pseudo-maximal if it is a proper root subsystem and is equal (up to the action of $W$) to a maximal element of the set
$\{\tilde{\Phi}(i), \Phi(j): \ 1\lest i, j \lest n \}$.
\end{Defff}

Note that a pseudo-maximal root subsystem of $\Phi$ has rank $n-1 $ or $n$. 


Since $j$-induction satisfies the property of induction by stages as in Lemma \ref{j.1}, we only need to consider the case of $\Phi_{[\lam]}^\vee$ which is a pseudo-maximal subsystem of $\Phi^\vee$, as specified in Definition \ref{thmbd}. This is what we will focus in the sections \S \ref{g2-gkd}--\S \ref{e-gkd} when computing the Gelfand--Kirillov dimension of $L(\lambda)$. For general case we discuss in \S \ref{S:gen}.





%Such a subsystem of $\Phi$ is simply called a \emph{maximal subsystem of $\Phi$}.

% For example, let $\Phi=B_n$.
%   In this case $\Phi^\vee\simeq C_n$ with
%   simple roots $\{\ep_1-\ep_2, \cdots, \ep_{n-1}-\ep_n, 2\ep_n\}$.
%    By Lemma \ref{lemni1}, $\Phi_{[\lambda]}^\vee$ is a proper closed subsystem of $\Phi^{\vee}$. Since the rank of $\Phi_{[\lambda]}$ (also equal to the rank of $\Phi_{[\lambda]}^\vee$) is maximal,   $\Phi_{[\lambda]}^\vee$ must be a closed root subsystem of maximal rank of $\Phi^\vee \cong  C_i\times C_{n-i}$ in view of Lemma \ref{thmbd}. 
%   %either $\{\ep_1-\ep_2, \cdots, \ep_{n-1}-\ep_n\}$ (type $A_{n-1}$) or
%   %$\{-2\ep_1, \ep_1-\ep_2, \cdots, \ep_{i-1}-\ep_i, \ep_{i+1}-\ep_{i+2}, \cdots, \ep_{n-1}-\ep_n, 2\ep_n\}$ for $1\leq i<n$ (type $C_{i}\times C_{n-i}$).
%   %It follows that $\Phi_{[\lambda]}^\vee$ is isomorphic to
%   %$A_{n-1}$ or $C_i\times C_{n-i}$ for $1\leq i<n$.
%   Hence the rank of $\Phi_{[\lambda]}$ is maximal if and only if
%   $\Phi_{[\lambda]}$ has type $ C_i^\vee\times C_{n-i}^\vee =  B_i\times B_{n-i}$.

% Now we  recall some definitions and properties about  subsystems and maximal subsystems. Most of which can be found in \cite[\S12]{Kan01} and \cite{Wa68}. 


% \begin{Defff}\label{defcl}
% 	A subset $\Phi'\subset\Phi$ is called a \textit{subsystem of $\Phi$} if $\Phi'$ satisfies the following two conditions:
%  \begin{enumerate}
%      \item if $\alpha\in \Phi'$, then $-\alpha\in \Phi'$.
%      \item if $\alpha, \beta\in\Phi'$ and $\alpha+\beta\in\Phi$, then $\alpha+\beta\in\Phi'$.
%  \end{enumerate}

%  A subsystem $\Phi'$ is called maximal if $\Phi'$ is a proper subset of $\Phi$ and $\Phi'$ is not 
% properly contained in any proper subsystem of $\Phi$.
% \end{Defff}






% Let $\alpha_0$ be the highest root of $\Phi$ and write
% \[
% \alpha_0=\sum_{i=1}^nh_i\alpha_i.
% \]

% \begin{Thm}[Borel-de Siebenthal]\label{thmbd}
% 	Let $\Phi$ be irreducible. Then the maximal proper closed subsystem of $\Phi$ $($up to the action of $W$$)$ are those with simple systems
% 	\begin{itemize}
% 		\item [(1)] $\{\alpha_1, \alpha_2, \cdots, \hat\alpha_i, \cdots, \alpha_n\}$ with $h_i=1$.
% 		\item [(2)] $\{-\alpha_0, \alpha_1, \cdots, \hat\alpha_i, \cdots, \alpha_n\}$ with $h_i=p$ being prime.
% 	\end{itemize}
% where $\alpha_0=\sum_{i=1}^nh_i\alpha_i$ is the highest root of $\Phi$.
% \end{Thm}

% \begin{Prop}\label{max-integral}
% 	Let $\Phi$ be irreducible and $\lambda\in \mathfrak{h}^*$ be nonintegral. Then $\Phi_{[\lambda]}^{\vee}$ is maximal if and only if it is conjugate to one of the maximal subsystems given in \cite[Table 1]{Wa68}.

% \end{Prop}


%%%%%%%%%%%%%%%%%%%%%%%%%%%%%%%%%%%%%%%%%%%%%%%%%%%%%%%%%%%%%%%%%%%%
%






% \section{ Sommers-Achar  duality}
% In this section, we recall  the Sommers-Achar  duality from \cite{Sommers} and \cite{ac03}.

% \subsection{Bala--Carter label}


\section{The computation of \texorpdfstring{$\aff(w_{\lam})$}{} for exceptional types}\label{a-value-compu}

In view of Proposition \ref{pr:main1}, in order to compute the GKdim of  highest weight modules, it is essential to find the required $w_{\lam}$  or the value of $\aff(w_{\lam})$. In this section, we will give a simple method to find a $w_{\lam}$ which has minimal length in the Weyl group $W_{[\lam]}$ such that $\lam=w_{\lam}\mu$, where $\mu $ is antidominant. Furthermore, we explain how to compute $\aff(w_\lambda)$ by utilizing a natural decomposition of $\Phi_{[\lambda]}$.

\subsection{Integral highest weight modules}\label{inte-a}

Let $\mathfrak{g}$ be a finite-dimensional complex simple Lie algebra $\mathfrak{g}$ with a fixed Cartan subalgebra $\mathfrak{h}$. Recall that $\langle \cdot, \cdot\rangle$ denotes the canonical pairing between  $\mathfrak{h}$ and its dual $\mathfrak{h}^*$. Let $\Delta=\{\alpha_i \mid  1\lest i\lest n\}$ be the set of simple roots of $\mathfrak{g}$ with corresponding fundamental weights $\{\omega_i\mid 1\lest i\lest n\}$. 
Recall that the Cartan matrix is defined by
$A=(A_{ij})_{n\times n}$ with $A_{ij}=\langle \alpha_i, \alpha_j^\vee \rangle$ and $A_{ij}\in \{0,-1,-2,-3\}$ for $i \neq j$.

%Then we have the following lemma.

\begin{Lem}\label{rho}
   We have $\rho=\sum\limits_{k=1}^n \omega_k$ and 
   $\alpha_j=\sum\limits_{k=1}^n A_{jk}\omega_k$. For any $\lam\in \mathfrak{h}^*$, we write $\lam=\sum\limits_{1\lest i\lest n}k_i\omega_i$, where $k_i=\langle \lam, \alpha_i^{\vee}\rangle$. Then $\lam$ is an integral weight if and only if $k_i\in \mathbb{Z}$ for all $1\lest i\lest n$. Also, $\lam$ is an integral and antidominant weight if and only if   $k_i\in \mathbb{Z}_{\lest 0}$ for  all $1\lest i\lest n$.
\end{Lem}

Now we want to find the aforementioned  $w_{\lam}$ for an integral weight $\lambda$. When $\lam$ is antidominant, $w_{\lam}=\mathrm{id}$. Suppose $\lam$ is not antidominant, then there are some $k_i\in \mathbb{Z}_{>0}$. Suppose the largest index is $i_{\lam}={i_1}$ such that $k_{i_1}\in \mathbb{Z}_{>0}$. Then we have $$s_{i_1}\lam=\lam-k_{i_1}\alpha_{i_1}=\lam-k_{i_1}\sum\limits_{j=1}^n A_{i_1j}\omega_j=\lam-2k_{i_1}\omega_{i_1}-k_{i_1}\sum\limits_{j=1,j\neq i_1}^n A_{i_1j}\omega_j.$$
 Then the coefficient of $\omega_{i_1}$ in $s_{i_1}\lam$ becomes $-k_{i_1}$. We continue this process until all the coefficients of $\omega_{i}$ are in $\mathbb{Z}_{\lest 0}$. This process will stop after finite steps. Multiplying all the $s_i$ appeared in this process will give us the desired $w_{\lam}$.

We call the above process \emph{positive index reduction} algorithm.  In fact, we have the following result.
% After this algorithm, we can get a new weight, denoted it by $\pi(\lambda)$. Then we have $\sum \pi(\lam) <\sum \lam$ and $i_{\pi(\lambda)}\lest i_{\lam}
%  $.
% Then we repeat this process until $i_{\pi^k(\lambda)}< i_{\lam}$ where $\pi^k(\lambda)$ means that we do $k$ times of this positive index reduction algorithm to $\lambda$. Finally we can get $i_{\pi^m(\lambda)}=0$ for some $m$, which means that we can get an antidominant weight $\pi^m(\lambda)$.



\begin{Lem}\label{find-w-lambda}
    For any integral weight $\lam$ (regular or singular),  we can get an antidominant weight $\mu$ by applying the positive index reduction algorithm. We multiply all the $s_i, 1\lest i \lest k$ that appeared in this process and get a $w_{\lambda}:=s_1 s_2 \cdots s_k$. Then this $w_{\lambda}$ has the minimal length in $W$ such that $w_{\lambda}^{-1}\lam$ is antidominant.
\end{Lem}
 \begin{proof}
     This is an easy consequence of \cite[Lem. 3.1]{do23}, which is a reformulation of the results in 
     \cite[Thm. 4.3.1(iv)]{BB05} and \cite[Prop. 4.1]{Er95}.
 \end{proof}

\begin{Rem}
    From \cite[Lem. 3.1]{do23},  we can choose any positive index (whose corresponding coefficient is positive) during the positive induction algorithm. 
For convenience, we have designed   the following web page for an interested reader to apply our algorithm:
    \begin{center}
\textcolor{blue}{http://test.slashblade.top:5000/lie/antidominant}
    \end{center}
We write $\lam=\sum\limits_{0\lest i\lest n-1}k_i\omega_i:=[k_0,\cdots,k_{n-1}]$, where $k_i=\langle \lam, \alpha_i^{\vee}\rangle$. Here  the order of simple roots is the same with that in \cite{ge}. In our web page, we only need to input $``k_0,\cdots,k_{n-1}"$, and the corresponding $w_{\lambda}$ will appear as the output.

    
\end{Rem}



 

% Suppose $\lam$ is regular. Then we have $k_i\neq 0$ for all $1\lest i\lest n$. First we  consider the case of type $E_6$. Then $n=6$.
% We write $\lam=(k_1,\dots,k_6)$ (see Figure \ref{d678} for the numbering of simple roots). We have $\alpha_1=2\omega_1-\omega_3$,  $\alpha_2=2\omega_2-\omega_4$, $\alpha_3=-\omega_1+2\omega_3-\omega_4$, $\alpha_4=-\omega_2-\omega_3+2\omega_4-\omega_5$, $\alpha_5=-\omega_4+2\omega_5-\omega_6$, $\alpha_6=-\omega_5+2\omega_6$.  Now we have $k_{i_1}>0$ and $k_{i_1+j}<0$ for all $0<j\lest 6-i_1$. Then 
% $$s_{i_1}\lam=(\dots,k_{i_1-1}+k_{i_1},-k_{i_1},k_{i_1}+k_{i_1+1},\dots ) \text{~if~} \alpha_{i_1}=-\omega_{i_1-1}+2\omega_{i_1}-\omega_{i_1+1}.$$
% When $k_{i_1}+k_{i_1+1}\lest 0$, the largest index of $s_{i_1}\lam$ with a positive element will be smaller than $i_1$. When $k_{i_1}+k_{i_1+1}> 0$, we have $s_{i_1+1}s_{i_1}\lam=(\dots,k_{i_1-1}+k_{i_1},k_{i_1+1},-k_{i_1}-k_{i_1+1},k_{i_1+1}+k_{i_1+2},\dots ) \text{~if~} \alpha_{i_1+1}=-\omega_{i_1}+2\omega_{i_1+1}-\omega_{i_1+2}.$
% Then the largest index of $s_{i_1+1}s_{i_1}\lam$ with a positive element will be smaller than $i_1$. For other forms of
% $\alpha_{i_1}$ and $\alpha_{i_1+1}$, we can get similar result. Since $i_1\lest 6$, the process will stop after finite steps.

% We will give the proof by induction on the value of $\sum \lam:=\sum_{1\lest i\lest 6} k_i$ and $i_{\lam}$.
% \begin{enumerate}
%     \item If $i_{\lam}=1$, then  $k_1>0$ and $k_i\lest 0$ for all $1< i\lest n$. We have $s_1\lam=\lam-k_1(2\omega_1-\omega_3)=(-k_1,k_2,k_3+k_1,k_4,k_5,k_6)$. Then  $\sum s_1\lam=\sum_{2\lest i\lest 6} k_i<k_1+\sum_{2\lest i\lest 6} k_i=\sum \lam$. 
%     \item If $k_3+k_1>0$, then $s_3s_1\lam=(k_3,k_2,-k_3-k_1,k_4+k_3+k_1,k_5,k_6)$. 
%     \item If $k_4+k_3+k_1>0$, we have $s_4s_3s_1\lam=(k_3,k_2+k_4+k_3+k_1,k_4,-k_4-k_3-k_1,k_5+k_4+k_3+k_1,k_6)$.
%     \item \begin{enumerate}
%         \item If $k_5+k_4+k_3+k_1\lest 0$ and $k_2+k_4+k_3+k_1>0$, we have $s_2s_4s_3s_1\lam=(k_3,-(k_2+k_4+k_3+k_1),k_4,k_2,k_5+k_4+k_3+k_1,k_6)$, which is already antidominant.
%         \item If $k_5+k_4+k_3+k_1> 0$, we have $s_5s_4s_3s_1\lam=(k_3,k_2+k_4+k_3+k_1,k_4,k_5,-(k_5+k_4+k_3+k_1),k_6+k_5+k_4+k_3+k_1)$.
        
%       {\bf (A).}  If $k_6+k_5+k_4+k_3+k_1> 0$, we have $s_6s_5s_4s_3s_1\lam=(k_3,k_2+k_4+k_3+k_1,k_4,k_5,k_6,-(k_6+k_5+k_4+k_3+k_1))$.
            
%              If $k_2+k_4+k_3+k_1> 0$, we have $s_2s_6s_5s_4s_3s_1\lam=(k_3,-(k_2+k_4+k_3+k_1),k_4,k_5+k_4+k_3+k_2+k_1,k_6,-(k_6+k_5+k_4+k_3+k_1))$.
          
%            If $k_5+k_4+k_3+k_2+k_1>0$, we have $s_4s_2s_6s_5s_4s_3s_1\lam=(k_3,k_5,k_5+2k_4+k_3+k_2+k_1,-(k_5+k_4+k_3+k_2+k_1),k_6+k_5+k_4+k_3+k_2+k_1,-(k_6+k_5+k_4+k_3+k_1))$.
%             \begin{enumerate}
%            \item If $k_6+k_5+k_4+k_3+k_2+k_1>0$, we have $s_5s_4s_2s_6s_5s_4s_3s_1\lam=(k_3,k_5,k_5+2k_4+k_3+k_2+k_1,k_6,-(k_6+k_5+k_4+k_3+k_2+k_1),k_2)$.

%            If $k_5+2k_4+k_3+k_2+k_1>0$, we have $s_3s_5s_4s_2s_6s_5s_4s_3s_1\lam=(k_5+2k_4+2k_3+k_2+k_1,k_5,-(k_5+2k_4+k_3+k_2+k_1),k_6+k_5+2k_4+k_3+k_2+k_1,-(k_6+k_5+k_4+k_3+k_2+k_1),k_2)$. Then  $\sum s_3s_5s_4s_2s_6s_5s_4s_3s_1 \lam <\sum \lam$ and $i_{s_3s_5s_4s_2s_6s_5s_4s_3s_1\lam}\lest  1$.
           
%            \item If $k_6+k_5+k_4+k_3+k_2+k_1\lest 0$ and $k_5+2k_4+k_3+k_2+k_1>0$, we have $s_3s_4s_2s_6s_5s_4s_3s_1\lam=(k_5+2k_4+2k_3+k_2+k_1,k_5,-(k_5+2k_4+k_3+k_2+k_1),k_4,k_6+k_5+k_4+k_3+k_2+k_1,-(k_6+k_5+k_4+k_3+k_1))$. Then  $\sum s_3s_4s_2s_6s_5s_4s_3s_1 \lam <\sum \lam$ and $i_{s_3s_4s_2s_6s_5s_4s_3s_1 \lam}\lest  1$.
%             \end{enumerate}

%        {\bf (B).}
%         If $k_6+k_5+k_4+k_3+k_1\lest 0$ and $k_2+k_4+k_3+k_1>0$, we have $s_2s_5s_4s_3s_1\lam=(k_3,-(k_4+k_3+k_2+k_1),k_4,k_5+k_4+k_3+k_2+k_1,-(k_5+k_4+k_3+k_1),k_6+k_5+k_4+k_3+k_1)$. 
           
%             If $k_5+k_4+k_3+k_2+k_1>0$, we have $s_4s_2s_5s_4s_3s_1\lam=(k_3,k_5,k_5+2k_4+k_3+k_2+k_1,-(k_5+k_4+k_3+k_2+k_1),k_2,k_6+k_5+k_4+k_3+k_1)$.
           
%             If $k_5+2k_4+k_3+k_2+k_1>0$, we have $s_3s_4s_2s_5s_4s_3s_1\lam=(k_5+2k_4+2k_3+k_2+k_1,k_5,-(k_5+2k_4+k_3+k_2+k_1),k_4,k_2,k_6+k_5+k_4+k_3+k_1)$. Then $\sum s_3s_4s_2s_5s_4s_3s_1\lam <\sum \lam$ and $i_{s_3s_4s_2s_5s_4s_3s_1\lam}\lest 1$.
          
%          \end{enumerate}
    
 
    
% \end{enumerate}
% When $i_{\lam}>1$, the process is similar and we can get a similar result.

% The case of type $E_7$ and $E_8$ is similar with type $E_6.$

 









\begin{ex}
  Let  $\lam=(2,-1,0,-5,-6,-8,12,-12)$ be  an integral weight of type $E_7$. We can write $\lam=-\omega_1+\omega_2-3\omega_3+\omega_4-5\omega_5-\omega_6-2\omega_7:=[-1,1,-3,1,-5,-1,-2]$. The largest index of $\lam$ with a positive element is $i_1=4$. Recall that $\alpha_4=-\omega_{2}-\omega_{3}+2\omega_{4}-\omega_5$. Thus $s_{4}\lam=\lam-\alpha_4=[-1,2,-2,-1,-4,-1,-2]$. The largest index of $s_4\lam$ with a positive element is $i_2=2$. Thus $s_2s_{4}\lam=s_{4}\lam-2\alpha_2=[-1,-2,-2,1,-4,-1,-2]$ since $\alpha_2=2\omega_{2}-\omega_4$. The largest index of $s_2s_4\lam$ with a positive element is $i_4=4$. Thus $s_4s_2s_{4}\lam=s_2s_{4}\lam-\alpha_4=[-1,-1,-1,-1,-3,-1,-2]$, which is already antidominant. Thus we have $w_{\lam}=s_4s_2s_{4}$. By using PyCox, we have $\aff(w_{\lam})=3$.
\end{ex}

\subsection{Nonintegral highest weight modules}\label{nonint}
To highlight the root system $\Phi$, we use $\mathfrak{h}_{\Phi}$ to denote the corresponding Cartan subalgebra of the Lie algebra $\mathfrak{g}_\Phi$. 
% We use $E_{\Phi}$ to denote the $\mathbb{R}$ linear span of $\Phi$ in $\mathfrak{h}_{\Phi}^*$. 
% Then $V_{\Phi}$ is a subspace of some Euclidean space $E_{\Phi}=\mathbb{R}^{n}=\mathrm{span}_{\mathbb{R}}\{\ep_1,\dots,\ep_n\}$. It may happen that $\dim V_{\Phi}=\mathrm{rank} (\Phi)< n$ (for example, $\dim V_{A_n}=n<n+1=\dim E_{A_n}$) or $\dim V_{\Phi}=\mathrm{rank} (\Phi)= n$. When $\dim V_{\Phi}=\mathrm{rank} (\Phi)<n= \dim E_{\Phi}$, we can choose another root system $\Phi^{\perp}$ such that the roots in $\Phi$ and $ \Phi^{\perp}$ are orthogonal to each other and $E_{\Phi}=\mathrm{span}_{\mathbb{R}}\{\Phi, \Phi^{\perp} \}$.  

Recall that we use $\Phi_{[\lam]}$ to denote the integral root system for $\lam\in \mathfrak{h}_{\Phi}^*$ and $\Delta_{[\lambda]}$ to denote the simple system of  $\Phi_{[\lambda]}$ lying in $\Phi_{[\lambda]} \cap \Phi^+$. 
The following enables us to determine  $\Delta_{[\lambda]}$ more explicitly.


% \begin{Lem}
%     For $\lam\in \mathfrak{h}_{\Phi}^*$, we decompose $\Phi_{[\lam]}$ into several orthogonal components:
%     $$\Phi_{[\lam]}=\Phi_{[\lam]_1}\cup \Phi_{[\lam]_2}\cup \cdots \cup \Phi_{[\lam]_k}\simeq \Phi_1\times \Phi_2\times \cdots\times \Phi_k,$$
%     such that $\Phi_{[\lam]_i}\simeq \Phi_i$ for $1\leq i\leq k$. Then we have the following way to find the simple system $\Delta_{[\lambda]_i}$ of $\Phi_{[\lam]_i}$:
% \begin{enumerate}
%     \item If $|\Phi^+_{[\lam]_i}|=\frac{n(n+1)}{2}$, then $\Phi_{[\lam]_i}\simeq A_n$. Then we have a longest chain of positive roots $I_n=\{\alpha_{i_1},\alpha_{i_2},\dots,\alpha_{i_n}\}$ such that $\langle \alpha_{i_j}, \alpha_{i_{j+1}}\rangle=-1$ for $1\leq j\leq n-1$ and $\langle \alpha_{i_l}, \alpha_{i_{k}}\rangle=0$ for all $1\leq l, k\leq n$ with  $|l-k|> 1$. Thus we have $\Delta_{[\lambda]_i}=I_n$.
%     \item If $|\Phi^+_{[\lam]_i}|=n^2$ and the number of short roots is $n$ ($n\geq 2$), then $\Phi_{[\lam]_i}\simeq B_n$. Then we have a longest chain of positive long roots $I_{n-1}=\{\alpha_{i_1},\alpha_{i_2},\dots,\alpha_{i_{n-1}}\}$ such that $\langle \alpha_{i_j}, \alpha_{i_{j+1}}\rangle=-1$ for $1\leq j\leq n-2$ and $\langle \alpha_{i_l}, \alpha_{i_{k}}\rangle=0$ for all $1\leq l, k\leq n-1$ with  $|l-k|> 1$. Also we have another positive short root $\alpha_{i_n}\in \Phi^+_{[\lam]_i}$ such that $\langle \alpha_{i_{n-1}}, \alpha_{i_{n}}\rangle=-1$ and $\langle \alpha_{i_l}, \alpha_{i_{n}}\rangle=0$ for all $1\leq l\leq n-2$.
%     Thus we have $\Delta_{[\lambda]_i}=I_{n-1}\cup\{\alpha_{i_n}\}$.
%     \item If $|\Phi^+_{[\lam]_i}|=n^2$ and the number of long roots is $n$ ($n\geq 3$), then $\Phi_{[\lam]_i}\simeq C_n$. Then similar to the case in (2), we have $\Delta_{[\lambda]_i}=I_{n-1}\cup\{\alpha_{i_n}\}$, where $I_{n-1}$ is the longest chain of positive short roots  and $\alpha_{i_n}$ is a long root.
%        \item If $|\Phi^+_{[\lam]_i}|=n^2-n$ ($n\geq 4$), then $\Phi_{[\lam]_i}\simeq D_n$.  Then we have a longest chain of positive  roots $I_{n-1}=\{\alpha_{i_1},\alpha_{i_2},\dots,\alpha_{i_{n-1}}\}$ such that $\langle \alpha_{i_j}, \alpha_{i_{j+1}}\rangle=-1$ for $1\leq j\leq n-2$ and $\langle \alpha_{i_l}, \alpha_{i_{k}}\rangle=0$ for all $1\leq l, k\leq n-1$ with  $|l-k|> 1$. Also we have another positive  root $\alpha_{i_n}\in \Phi^+_{[\lam]_i}$ such that $\langle \alpha_{i_{n-1}}, \alpha_{i_{n}}\rangle=0$, $\langle \alpha_{i_{n-2}}, \alpha_{i_{n}}\rangle=-1$ and $\langle \alpha_{i_{l}}, \alpha_{i_{n}}\rangle=0$ for all $1\leq l\leq n-3$.
%     Thus we have $\Delta_{[\lambda]_i}=I_{n-1}\cup\{\alpha_{i_n}\}$.
% \item If $|\Phi^+_{[\lam]_i}|=36$, then $\Phi_{[\lam]_i}\simeq E_6$.  Then we have a longest chain of positive  roots $I_{5}=\{\alpha_{i_1},\alpha_{i_2},\dots,\alpha_{i_{5}}\}$ such that $\langle \alpha_{i_j}, \alpha_{i_{j+1}}\rangle=-1$ for $1\leq j\leq 4$ and $\langle \alpha_{i_l}, \alpha_{i_{k}}\rangle=0$ for all $1\leq l, k\leq 5$ with  $|l-k|> 1$. Also we have another positive  root $\alpha_{i_6}\in \Phi^+_{[\lam]_i}$ such that $\langle \alpha_{i_{k}}, \alpha_{i_{6}}\rangle=0$ for $k\neq 3$ and  $\langle \alpha_{i_{3}}, \alpha_{i_{6}}\rangle=-1$.
%     Thus we have $\Delta_{[\lambda]_i}=I_{5}\cup\{\alpha_{i_6}\}$.

% \item If $|\Phi^+_{[\lam]_i}|=63$, then $\Phi_{[\lam]_i}\simeq E_7$. Then we have a longest chain of positive  roots $I_{6}=\{\alpha_{i_1},\alpha_{i_2},\dots,\alpha_{i_{6}}\}$ such that $\langle \alpha_{i_j}, \alpha_{i_{j+1}}\rangle=-1$ for $1\leq j\leq 5$ and $\langle \alpha_{i_l}, \alpha_{i_{k}}\rangle=0$ for all $1\leq l, k\leq 6$ with  $|l-k|> 1$. Also we have another positive  root $\alpha_{i_7}\in \Phi^+_{[\lam]_i}$ such that $\langle \alpha_{i_{k}}, \alpha_{i_{7}}\rangle=0$ for $k\neq 3$ and  $\langle \alpha_{i_{3}}, \alpha_{i_{7}}\rangle=-1$.
%     Thus we have $\Delta_{[\lambda]_i}=I_{6}\cup\{\alpha_{i_7}\}$.


    
% \end{enumerate}




    
% \end{Lem}

\begin{Lem}
    For $\lam\in \mathfrak{h}_{\Phi}^*$, we decompose the root system $\Phi_{[\lam]}$ into several orthogonal irreducible subsystems
    $$\Phi_{[\lam]}=\Phi_{[\lam]_1}\sqcup \Phi_{[\lam]_2}\sqcup \cdots \sqcup \Phi_{[\lam]_k}.$$
 We can determine the simple root subsystem $\Delta_{[\lambda]_i}$ of $\Phi_{[\lam]_i}$ as follows:
\begin{itemize}
    \item[(1)] If $|\Phi^+_{[\lam]_i}|=\frac{n(n+1)}{2}$, then $\Phi_{[\lam]_i}\simeq A_n$. We compute $\rho_i=\frac{1}{2}\sum_{\alpha\in \Phi^+_{[\lam]_i}}\alpha$. We filter out all the positive roots  $\alpha \in \Phi^+_{[\lam]_i}$ such that $\langle \rho_i, \alpha^{\vee}\rangle=1$ and denote them by $I_n=\{\alpha_{i_1},\alpha_{i_2},\dots,\alpha_{i_n}\}$. Then we have $\Delta_{[\lambda]_i}=I_n$.
    \item[(2)] If $|\Phi^+_{[\lam]_i}|=n^2$ and the number of short roots is $n$ ($n\geq 2$), then $\Phi_{[\lam]_i}\simeq B_n$.
    We compute $\rho_i=\frac{1}{2}\sum_{\alpha\in \Phi^+_{[\lam]_i}}\alpha$. We filter out all the positive roots  $\alpha \in \Phi^+_{[\lam]_i}$ such that $\langle \rho_i, \alpha^{\vee}\rangle=1$ and denote them by $I_{n}=\{\alpha_{i_1},\alpha_{i_2},\dots,\alpha_{i_{n}}\}$.  Then we have $\Delta_{[\lambda]_i}=I_{n}$.
      
    \item[(3)] If $|\Phi^+_{[\lam]_i}|=n^2$ and the number of long roots is $n$ ($n\geq 3$), then $\Phi_{[\lam]_i}\simeq C_n$. Similar to the case in (2), we have $\Delta_{[\lambda]_i}=I_{n}$.
    
       \item[(4)] If $|\Phi^+_{[\lam]_i}|=n^2-n$ ($n\geq 4$), then $\Phi_{[\lam]_i}\simeq D_n$.  Similar to the case in (1), we have $\Delta_{[\lambda]_i}=I_{n}$.
       
\item[(5)] If $|\Phi^+_{[\lam]_i}|=36$ and the roots are of the same length, then $\Phi_{[\lam]_i}\simeq E_6$. Similar to the case in (1), we have $\Delta_{[\lambda]_i}=I_{6}$.

\item[(6)] If $|\Phi^+_{[\lam]_i}|=63$, then $\Phi_{[\lam]_i}\simeq E_7$. Then similar to the case in (1), we have $\Delta_{[\lambda]_i}=I_{7}$.


    
\end{itemize}




    
\end{Lem}
\begin{proof}
  Recall that $\Phi^+_{[\lam]}=\Phi_{[\lam]}\cap \Phi^+$. Let $\Delta=\{\alpha_i \mid  1\lest i\lest n\}$ be the simple roots of  $\Phi$ with corresponding fundamental weights $\{\omega_i\mid 1\lest i\lest n\}$. Then we  have $\rho=\sum\limits_{k=1}^n \omega_k$. Note that any positive root in $\Phi^+$ is a linear combination of simple roots with nonnegative integral coefficients. Thus we have $\langle \rho, \alpha^{\vee}\rangle=1$ if and only if $\alpha\in \Phi^+$ is a simple root. Therefore our results follow from this observation. 
\end{proof}



\begin{Thm}\label{transf}
    Suppose  $\Phi$ and $  \Psi$ are two isomorphic root systems. Let $\phi: \Phi \rightarrow \Psi$ be an isomorphism. By abuse of notation, still denote by $\phi$ the associated isomorphism  from $\mathfrak{h}_{\Phi}$ to $\mathfrak{h}_{\Psi}$, and from $\mathfrak{h}_{\Phi}^*$ to $\mathfrak{h}_{\Psi}^*$. Suppose for $\lam\in \mathfrak{h}_{\Phi}^*$, we have  an isomorphism  $\phi|_{\Phi_{[\lam]}}: \Phi_{[\lam]} \rightarrow \Psi_0 $, where $\Psi_0$ is a root subsystem of $\Psi$ and $\mathrm{rank} {(\Psi_0)}=\mathrm{rank} {(\Psi)}=n$. Then $\lam'=\phi(\lam)$ is an integral weight of type $\Psi_0$ and  we have: $$\aff(w_{\lam})=\aff(w_{\lam'}).$$

\end{Thm}

\begin{proof}
Note that if $\phi: \Phi_1 \rightarrow \Psi_1$  is an isomorphism between two root systems, then their Weyl groups $W_{\Phi_1}$  and $W_{\Psi_1}$ are isomorphic (we still denote this isomorphism by $\phi$) and we have $\phi s_{\alpha}\phi ^{-1}=s_{\phi(\alpha)}$ for any root $\alpha \in \Phi_1$  and simple reflection $s_{\alpha}$. For any $w\in W_{\Phi_1}$, we will have 
$\aff(w)=\aff(\phi(w))$. 

Recall that the Cartan matrix of a root system $\Phi$ (its simple system is $\Delta=\{\alpha_i \mid  1\lest i\lest n\})$ is defined by
$A=(A_{ij})_{n\times n}$ with $A_{ij}={\langle \alpha_i, \alpha_j^\vee \rangle}$ and $A_{ij}\in \{0,-1,-2,-3\}$ for $i \neq j$. Then we have $\alpha_j=\sum\limits_{k=1}^n A_{jk}\omega_k$. 
Suppose the simple system of $\Phi_{[\lam]}$ is $\Delta_{[\lam]}=\{\beta_1,\cdots,\beta_n\}$. For $\lam \in \mathfrak{h}_{\Phi}^*$, we can write $$\lam=\sum_{1\lest i\lest n}x_i\beta_i=\sum_{1\lest i\lest n}x'_i\omega_{\beta_i}\in \mathfrak{h}_{\Phi_{[\lam]}}^*,$$ since $\mathrm{rank} {(\Psi_0)}=\mathrm{rank} {(\Psi)}$ and $$\mathrm{span}_{\mathbb{R}}\{\omega_{\beta_1},\cdots,\omega_{\beta_n}\}=\mathrm{span}_{\mathbb{R}}\{\beta_1,\cdots,\beta_n\}.$$
Since $\lam$ is an integral weight in $\mathfrak{h}_{\Phi_{[\lam]}}^*$, we must have $x'_i\in \mathbb{Z}$ and thus
$$\lam'=\phi(\lam)=\sum_{1\lest i\lest n}x_i\phi(\beta_i)=\sum_{1\lest i\lest n}x'_i\omega_{\phi({\beta_i})}$$ is an integral weight in $\mathfrak{h}_{\Psi_{0}}^*$.


% Thus, we must have $\mathfrak{h}_{\Phi_{[\lam]}}^*=\mathfrak{h}_{\Phi}^*$ and $\mathfrak{h}_{\Psi_0}^*=\mathfrak{h}_{\Psi}^*$. 
% In other words,
% when $\lam\in \mathfrak{h}_{\Phi}^*$,  $\lam'=F(\lam)\in \mathfrak{h}_{\Psi}^*=\mathfrak{h}_{\Psi_0}^*$ will be  a weight of type $\Psi_0$.



% Then we have the following lemma.
% \begin{lemma}\label{rho}
% Let $\Pi=\{\alpha_i \mid  1\lest i\lest n\}$ be the simple roots of a finite-dimensional complex simple Lie algebra $\mathfrak{g}$. 
%    Then we have $\rho=\sum\limits_{k=1}^n \omega_k$ and 
%    $\alpha_j=\sum\limits_{k=1}^n A_{jk}\omega_k$.
% \end{lemma}

% Let $K=(k_{ij})_{n\times n}$ be the inverse of the Cartan matrix $A$. Then from Lemma \ref{rho}, we have $\omega_j=\sum\limits_{i=1}^n k_{ji}\alpha_i$. The explicit formulas for the inverses of the Cartan matrices of the simple Lie algebras are given in \cite{WZ}.





Now  we write $\lambda=w_{\lambda}\mu$, where $\mu$ is antidominant and $w_{\lambda}\in W_{[\mu]}^J$. Then $\lam'=\phi(\lam)=\phi(w_{\lambda}\mu)=\phi(w_{\lambda})(\mu)=\phi w_{\lambda}\phi ^{-1}\phi(\mu)$. From $\langle \lam, \alpha\rangle=\langle \phi(\lam),\phi(\alpha)\rangle$, we can see that $\phi(\mu)$ is still antidominant with respect to $\Psi$. If we write $w_{\lam}=s_{\beta_{i_1}}s_{\beta_{i_2}}\cdots s_{\beta_{i_k}}$ for some simple roots $\{\beta_{i_j}\mid 1\lest j\lest k\}\subset \Phi_{[\lam]}^+$, then \begin{align*}
   \phi(w_{\lam})&= \phi w_{\lam}\phi^{-1}=(\phi s_{\beta_{i_1}}\phi ^{-1})(\phi s_{\beta_{i_2}}\phi ^{-1})\cdots (\phi s_{\beta_{i_k}}\phi ^{-1})\\
    &=s_{\phi(\beta_{i_1})}s_{\phi(\beta_{i_2})}\cdots s_{\phi(\beta_{i_k})}.
\end{align*}
Thus $\aff(w_{\lam})=\aff(\phi(w_{\lam}))=\aff(w_{\lam'})$ since $\phi(w_{\lam})=w_{\lam'}\in W_{[\phi(\mu)]}^{\phi(J)}$.
\end{proof}

From the argument in the above proof, we have the following corollary.

\begin{Cor}\label{transf-1}
    Keep the notation as above. When $\mathrm{rank} {(\Psi_0)}<\mathrm{rank} {(\Psi)}$, 
  %   we can choose two  subspaces  $E_{\Phi _{[\lam]}}^{\perp}$ and  $E_{\Psi _{0}}^{\perp}$  such that $E_{\Phi}=E_{\Phi_{[\lam]}}\oplus E_{\Phi_{[\lam]}}^{\perp}$ and $E_{\Psi}=E_{\Psi_{0}}\oplus E_{\Psi_{0}}^{\perp}$. Choose two bases 
  % $\Phi_{[\lam]}^{\perp}$ and $\Psi_{0}^{\perp}$ such that $E_{\Phi}=\mathrm{span}_{\mathbb{R}}\{\Phi_{[\lam]}, \Phi_{[\lam]}^{\perp} \}$ and $E_{\Psi}=\mathrm{span}_{\mathbb{R}}\{\Psi_{0}, \Psi_{0}^{\perp} \}$. Then $\lam'=\phi(\lam)$ is  a  weight of type $\Psi_0\times\Psi_0^{\perp}$. If we write $\lam'=(\lam',\lam'^{\perp})$,
  we denote ${\lam}'=\phi(\lam|_{\mathfrak{h}_{\Phi_{[\lam]}}})$.
% Then $\lam'|_{{\Phi_i}}=\phi(\lambda|_{\mathfrak{h}_{\Phi_{[\lam]}^*}})|_{{\Phi_i}}$ is an integral weight of type $\Phi_i$.
  Then we have $$\aff(w_{\lam})=\aff(w_{{\lam}'}).$$
\end{Cor}
\begin{proof}
    Suppose the simple root system of $\Phi_{[\lam]}$ is $\Delta_{[\lam]}=\{\alpha_1,\cdots,\alpha_n\}$.
    % and the simple root system of $\Phi_{[\lam]}^{\perp}$ is $\Delta_{[\lam]}^{\perp}=\{\beta_1,\cdots,\beta_m\}$. 
    For $\lam \in \mathfrak{h}_{\Phi}^*$, we have 
$\mathfrak{h}_{\Phi}=\mathfrak{h}_{\Phi_{[\lam]}}\oplus \mathfrak{h}_{\Phi_{[\lam]}}^{\perp}$, and 
\begin{align*}\lam|_{\mathfrak{h}_{\Phi_{[\lam]}}}&=\sum_{1\lest i\lest n}x_i\alpha_i=\sum_{1\lest i\lest n}x'_i\omega_{\alpha_i},
         \end{align*} 
         since $\mathrm{span}_{\mathbb{R}}\{\omega_{\alpha_1},\cdots,\omega_{\alpha_n}\}=\mathrm{span}_{\mathbb{R}}\{\alpha_1,\cdots,\alpha_n\}$. 
As $\lam$ is an integral weight of ${\Phi_{[\lam]}}$, we must have $x'_i\in \mathbb{Z}$. 
So ${\lam}'=\phi(\lam|_{\mathfrak{h}_{\Phi_{[\lam]}}})=\sum_{1\lest i\lest n}x_i\phi(\alpha_i)=\sum_{1\lest i\lest n}x'_i\omega_{\phi({\alpha_i})}$
is an integral weight of ${\Psi_{0}}$.









Now  we write $\lambda=w_{\lambda}\mu$, where $\mu$ is antidominant and $w_{\lambda}\in W_{[\mu]}^J$. 
% Then $\lam'=\phi(\lam)=\phi(w_{\lambda}\mu)=\phi(w_{\lambda})(\mu)=\phi w_{\lambda}\phi^{-1}\phi(\mu)$. 
If we write $w_{\lam}=s_{\alpha_{i_1}}s_{\alpha_{i_2}}\cdots s_{\alpha_{i_k}}$ for some simple roots $\{\alpha_{i_j}\mid 1\lest j\lest k\}\subset \Phi_{[\lam]}^+$, then \begin{align*}
   \phi(w_{\lam})&= \phi w_{\lam}\phi^{-1}=(\phi s_{\alpha_{i_1}}\phi^{-1})(\phi s_{\alpha_{i_2}}\phi^{-1})\cdots (\phi s_{\alpha_{i_k}}\phi^{-1})\\
    &=s_{\phi(\alpha_{i_1})}s_{\phi(\alpha_{i_2})}\cdots s_{\phi(\alpha_{i_k})}.
\end{align*}
Thus $\aff(w_{\lam})=\aff(\phi(w_{\lam}))=\aff(w_{{\lam}'})$ since $\phi(w_{\lam})=w_{{\lam}'}\in W_{[\phi(\mu)]}^{\phi(J)}$.
\end{proof}

% \begin{Cor}\label{transf-2}
%     Keep the notation as above. When $\mathrm{rank} {(\Psi_0)}<\mathrm{rank} {(\Psi)}<\dim E_{\phi}$ or $\mathrm{rank} {(\Psi_0)}=\mathrm{rank} {(\Psi)}<\dim E_{\Phi}$, we can choose $\Phi_{[\lam]}^{\perp}$ and $\Psi_{0}^{\perp}$ such that $E_{\Phi}=\mathrm{span}_{\mathbb{R}}\{\Phi_{[\lam]}, \Phi_{[\lam]}^{\perp} \}$ and $E_{\Psi}=\mathrm{span}_{\mathbb{R}}\{\Psi_{0}, \Psi_{0}^{\perp} \}$. Then $\lam'=\phi(\lam)$ is  a weight of type $\Psi_0\times \Psi_0^{\perp}$. If we write $\lam'=(\lam',\lam'^{\perp})$, then we have $$\aff(w_{\lam})=\aff(w_{\lam'}).$$
% \end{Cor}

Now Corollary \ref{transf-1} entails the following algorithm of computing $\aff(w_{\lam})$:
\begin{enumerate}
    \item suppose $\Phi_{[\lam]}=\Phi_{[\lambda]_1}\times \Phi_{[\lambda]_2}\times \cdots \times \Phi_{[\lambda]_k}\simeq \Phi_1\times \Phi_2\times \cdots\times \Phi_k$ is a direct product of irreducible root systems. We use $\phi$ to denote such an isomorphism.  Then $\phi$ induces an isomorphism from $\mathfrak{h}_{\Phi_{[\lam]}}$ to $\prod_{1\leq i\leq k} \mathfrak{h}_{\Phi_i}$. Suppose the simple system of $\Phi_{i}$ is $\Delta_i$ for $1\leq i\leq k$ and   the simple system of $\Phi_{[\lam]}$ is $\Delta_{[\lam]}=\{\alpha_1,\cdots,\alpha_n\}=\prod_{1\leq i\leq k}\Delta_{[\lam]_i}$, where $\Delta_{[\lam]_i}$ is the simple system of $\Phi_{[\lambda]_i}$ and isomorphic  to $\Delta_i$.  Write $\lambda|_{\mathfrak{h}_{\Phi_{[\lambda]}}}=\sum_{\alpha_i\in \Delta_{[\lam]}}k_i\omega_i$, where $k_i=\langle \lam, \alpha_i^{\vee}\rangle$;
    % \item[2)] $\lam'=\phi(\lam)=\sum_{\alpha_i\in \Delta_{[\lam]}}k_i\phi(\omega_i)+ \sum_{\beta_j\in {\Delta_{[\lam]}^{\perp}}}t_j\phi(\omega_j)$;
    % \item[3)] $\lam'=\sum_{\alpha_i\in \Phi_{[\lam]}}k_i\phi(\omega_i)=\sum_{1\lest i\lest m}\lam'_i\ep_i=(\lam'_1,\cdots,\lam'_{m})$, where $m=\dim{E}_{\Phi_{[\lam]}}$.
    \item denote ${\lam}'=\phi(\lam|_{\mathfrak{h}_{\Phi_{[\lambda]}}})=\sum_{\alpha_i\in \Delta_{[\lam]}}k_i\phi(\omega_i)$;
    \item  when $\Phi_i$ is of classical type, denote $${\lam}'|_{\Phi_i}=\sum_{\alpha_{i_t}\in \Delta_{[\lam]_i}}k_{i_t}\phi(\omega_{i_t})=(t_{1},t_{2},\cdots,t_{d_i}).$$ 
    Then we can get the value of $\aff(w_{\lam'|_{{\Phi_i}}})$ without finding out $w_{\lam'|_{{\Phi_i}}}$ by using the RS algorithm in \cite{BXX}. When some $\Phi_i$ is of exceptional type, we can find $w_{\lam'|_{{\Phi_i}}}$ by using the algorithm in Lemma \ref{find-w-lambda}, and then $\aff(w_{\lam'|_{{\Phi_i}}})$ is given by using PyCox;
     \item$\aff(w_{\lam})=\sum_{1\leq i\leq k} \aff(w_{\lam'|_{{\Phi_i}}})$.
\end{enumerate}

\begin{ex}
    Suppose $\lam=(2,1,1.1,3,0.9,1.9,4,2.1)$ is a weight for $\Phi=D_8$. Then $\Phi^+_{[\lam]}=\{\ep_i\pm \ep_{j}\mid i<j \in \{1,2,4,7\}\}\cup \{\ep_3+\ep_5,\ep_3+\ep_6,\ep_3-\ep_8,\ep_5-\ep_6,\ep_5+\ep_8,\ep_6+\ep_8\}$ and $\Phi_{[\lambda]}=\Phi_{[\lambda]_1}\times \Phi_{[\lambda]_2}\simeq D_4\times A_3$.  The simple system of $\Phi_{[\lambda]_1}$ is $\Delta_{[\lambda]_1}=\{\gamma_1=\ep_1- \ep_{2},\gamma_2=\ep_2- \ep_{4},\gamma_3=\ep_4- \ep_{7},\gamma_4=\ep_4+ \ep_{7}\}$ and the simple system of $\Phi_{[\lambda]_2}$ is $\Delta_{[\lambda]_2}=\{\beta_1=\ep_3-\ep_8,\beta_2=\ep_8+\ep_6,\beta_3=\ep_5-\ep_6\}$. Suppose the simple system of $D_4\times A_3$ is $\{\alpha_1=\ep_1-\ep_2,\alpha_2=\ep_2-\ep_3,\alpha_3=\ep_3-\ep_4,\alpha_4=\ep_3+\ep_4\}\times \{\alpha'_1=\ep_1-\ep_2,\alpha'_2=\ep_2-\ep_3,\alpha'_3=\ep_3-\ep_4\}$. We define an isomorphism
     $\phi: \Phi_{[\lam]}\rightarrow D_4\times A_3$ such that   $\phi(\gamma_i)=\alpha_i$ (for $1\leq i\leq 4$) and $\phi(\beta_i)=\alpha'_i$ (for $1\leq i\leq 3$).
    Then we have 
   $$\lambda|_{\mathfrak{h}_{\Phi_{[\lambda]}}}=\omega_1-2\omega_2-\omega_3+7\omega_4-\omega'_1+4\omega'_2-\omega'_3,$$
    where $\{\omega_i\mid 1\leq i\leq 4\}$ are the fundamental weights for $\Delta_{[\lam]_1}$ and $\{\omega'_i\mid 1\leq i\leq 3\}$ are the fundamental weights for $\Delta_{[\lam]_2}$. 
Denote $\lam'=\phi(\lam|_{\mathfrak{h}_{\Phi_{[\lam]}}})$.  Thus  
  $$\lam'|_{D_4}=\phi(\omega_1)-2\phi(\omega_2)-\phi(\omega_3)+7\phi(\omega_4)=(2,1,3,4),$$
 which is an integral weight of $D_4$ with $\aff(w_{\lam'|_{{D_4}}})=3$ by using the RS algorithm in \cite{BXX}, and $$\lam'|_{A_3}=-\phi(\omega'_1)+4\phi(\omega'_2)-\phi(\omega'_3)=(1,2,-2,-1),$$
 which is an integral weight of $A_3$ with $\aff(w_{\lam'|_{{A_3}}})=2$ by using the RS algorithm in \cite{BXX}. 
 Therefore we have $\aff(w_{\lambda})=3+2=5$.

Note that \begin{align*}
    \lam'|_{A_3}+(0.1,0.1,0.1,0.1)=(1.1,2.1,-1.9,-0.9):=x.
\end{align*}
Thus the Young tableaux $P(\lam'|_{A_3})$ and $P(x)$ have the same shape. This is just the algorithm used in \cite{BXX} for the case of type $D$. 

\end{ex}


%\begin{proof}
%	The second statement of (2) follows from Greene's Theorem (Theorem \ref{gthm}), while the other assertions are obvious.
%\end{proof}
%For convenience, we set $F_a(x)=\sum_i(i-1)p(x)_i$, $F_b(x)=\sum_i(i-1)p(x)_i\od$ and $F_d(x)=\sum_i(i-1)p(x)_i\ev$ for $x\in\mathrm{Seq}_n (\Gamma)$, where $\Gamma=c+\mathbb{Z}$ for a $c\in\mathbb{C}$.



%

% \section{Associated varieties for Exceptional groups} \label{sec:exp}

% %
% %%%%%%%%%%%%%%%%%%%%%%%%%%%%%%%%%%%%%%%%%%%%%%%%%%%%%%%%%%%%%%%%%%%%
% %

% In this section, we give the associated varieties of highest weight HC modules for the two exceptional groups. Now $\Phi$ is of type $E_6$ or $E_7$. As usual $\Phi$ can be realized as a subset of $\mathbb{R}^8$. The simple roots are 
% \begin{align*}
% 	\alpha_1&=\frac{1}{2}(\ep_1-\ep_2-\ep_3-\ep_4-\ep_5-\ep_6-\ep_7+\ep_8) ,\\
% 	\alpha_2&=\ep_1+\ep_2,\\
% 	\alpha_k&=\ep_{k-1}-\ep_{k-2} \text{ with }3\lest k\lest n,
% \end{align*}
% where $n=6, 7$. 
% Here we adopt notation in \cite{EHW}. 
% In particular, if $\Phi=E_6$, then $\Delta_c=\Delta\backslash\{\alpha_1\}$; if $\Phi=E_7$, then $\Delta_c=\Delta\backslash\{\alpha_7\}$. 

% \begin{figure}[htpb]
% 	\begin{tikzpicture}[scale=1.5,baseline=0]
% 		\filldraw (0,0.05) node[above=1pt]{$ 1 $} circle [radius=0.05];
% 		\draw (0.05,0.05)--++(0.65,0); 
% 		\draw (0.75,0.05) node[above=1pt]{$ 3 $} circle[radius=0.05];  
% 		\draw (0.8,0.05)--++(0.65,0);
% 		\draw  (1.45,0) ++(0.05,0.05) node[above=1pt]{$ 4$} circle[radius=0.05];
% 		\draw (1.55,0.05)--++(0.65,0);
% 		\draw  (2.2,0) ++(0.05,0.05) node[above=1pt]{$ 5$} circle[radius=0.05];
% 		\draw (2.3,0.05)--++(0.65,0);
% 		\draw  (2.95,0) ++(0.05,0.05) node[above=1pt]{$ 6$} circle[radius=0.05];
% 		\draw (1.5,0)--++(0,-0.5);
% 		\draw  (1.45,-0.6) ++(0.05,0.05) node[right=1pt]{$ 2$} circle[radius=0.05];
% 	\end{tikzpicture}
% 	\quad
% 	\begin{tikzpicture}[scale=1.5,baseline=0]
% 		\draw (0,0.05) node[above=1pt]{$ 1 $} circle [radius=0.05];
% 		\draw (0.05,0.05)--++(0.65,0); 
% 		\draw (0.75,0.05) node[above=1pt]{$ 3 $} circle[radius=0.05];  
% 		\draw (0.8,0.05)--++(0.65,0);
% 		\draw  (1.45,0) ++(0.05,0.05) node[above=1pt]{$ 4$} circle[radius=0.05];
% 		\draw (1.55,0.05)--++(0.65,0);
% 		\draw  (2.2,0) ++(0.05,0.05) node[above=1pt]{$ 5$} circle[radius=0.05];
% 		\draw (2.3,0.05)--++(0.65,0);
% 		\draw  (2.95,0) ++(0.05,0.05) node[above=1pt]{$ 6$} circle[radius=0.05];
% 		\draw (3.05,0.05)--++(0.65,0);
% 		\filldraw  (3.7,0) ++(0.05,0.05) node[above=1pt]{$ 7$} circle[radius=0.05];
% 		\draw (1.5,0)--++(0,-0.5);
% 		\draw  (1.45,-0.6) ++(0.05,0.05) node[right=1pt]{$ 2$} circle[radius=0.05];
% 	\end{tikzpicture}
% 	\caption{Diagrams of $ E_6 $ and $ E_7 $ }
% 	\label{d67}
% \end{figure}

% \subsection{The associated varieties} %First we state our criteria for exceptional groups.

% Recall  that, for an integral weight $ \lam\in\hs $, we have $ \Phi_\lam^+=\{\alpha\in\Phi^+\mid\langle\lambda, \alpha^\vee\rangle>0\} $.
% \begin{Thm}\label{67}
% 	Let $L(\lambda)$ be a HC module of $G$ of type $E_6/E_7$. If $\lambda\in\hs$ is not integral, then $V(L(\lambda))=\overline{\mathcal{O}}_{k(\lambda)}$ with $k(\lambda)=2$ $($for $E_6$$)$ or $3$ $($for $E_7$$)$. In the case when $\lambda$ is integral, $ k(\lam) $ is given as follows.
% 	\begin{itemize}
% 		\item [(1)] If $G$ is of type $E_6$, then
% 		\[
% 		k(\lambda)=\begin{cases}
% 			0& \text{ if } \Psi_{\lam}^+\cap S_2\neq\emptyset,\\
% 			1& \text{ if } \Psi_{\lam}^+\cap S_1\neq\emptyset \text{ and } \Psi_{\lam}^+\cap S_2=\emptyset ,\\
% 			2& \text{ if } \Psi_{\lam}^+\cap S_1=\emptyset,
% 		\end{cases}
% 		\]
% 		where $S_1=\{\frac{1}{2}(\pm(\ep_1+\ep_2+\ep_3-\ep_4)-\ep_5-\ep_6-\ep_7+\ep_8)\}$ and $S_2=\{\alpha_1\}$.
	
% 	\item [(2)] If $G$ is of type $E_7$, then
% 		\[
% 	k(\lambda)=\begin{cases}
% 		0& \text{ if } \Psi_{\lam}^+\cap S_3\neq\emptyset,\\
% 		1& \text{ if } \Psi_{\lam}^+\cap S_2\neq\emptyset \text{ and } \Psi_{\lam}^+\cap S_3=\emptyset ,\\
% 		2& \text{ if } \Psi_{\lam}^+\cap S_1\neq\emptyset \text{ and } \Psi_{\lam}^+\cap S_2=\emptyset ,\\
% 		3& \text{ if } \Psi_{\lam}^+\cap S_1=\emptyset,
% 	\end{cases}
% 	\]
% 	where $S_1=\{\frac{1}{2}(\pm(\ep_1-\ep_2-\ep_3+\ep_4)-\ep_5+\ep_6-\ep_7+\ep_8), \ep_5+\ep_6\}$, $S_2=\{\pm\ep_1+\ep_6\}$ and $S_3=\{-\ep_5+\ep_6\}$.
% 	\end{itemize}

% \end{Thm}



% \begin{Rem}
% 	When $L(\lambda)$ is a HC module of exceptional types, we know that its associated variety are determined by the corresponding $\aff$-functions, in view of Propositions \ref{C: dimYk} and \ref{pr:main1}. The computation of $\aff$-functions for exceptional Weyl groups is by no means an easy problem even using computers. However, as shown in the above example, with Theorem \ref{67}, the associated variety can be easily obtained by hand. 
% \end{Rem}

% %Then $ w^c:=w_cw_\circ $ is the longest element in $ W^c $, and $ \Phi^+\backslash\Phi_c^+=\Phi_{w^c}^- $. 

% \subsection{Some facts} 
% Let $W_c$ be the Weyl group of the compact root system $\Phi_c$. Denote by $W^c$ the set of minimal length coset representatives of $W_c\backslash W$. Then $W=W_cW^c$. Let $ w\in W$. For simplicity, we write (see \S2.5)
% \[
% \Phi_w^+:=\Phi_{wC_\circ}^+=\{\alpha\in \Phi^+\mid w^{-1}\alpha<0 \}.
% \]
% Let $ w=s_{\gamma_1}s_{\gamma_2} \cdots s_{\gamma_k}$ be a reduced expression of $ w $ with $\gamma_1, \cdots, \gamma_k\in\Delta$. For $ 1\lest r\lest k $, define positive roots
% \[
% \beta_r=s_{\gamma_1}s_{\gamma_2} \cdots s_{\gamma_{r-1}} \gamma_r.
% \]
% Then $ \beta_r $'s are distinct and (see for example \cite[\S1.7]{Hum90}) 
% \begin{equation}\label{gfeq1}
% \Phi_w^+=\{\beta_{r}\mid   1\lest r\lest k \}.
% \end{equation}
% In particular, this set is independent of the choice of the reduced expressions of $ w $. Recall the Bruhat ordering ``$\lest$'' on $W$ (e.g., \cite[\S5.9]{Hum90}).

% \begin{Lem}\label{hermp}
% 	Let $G$ be a Lie group of Hermitian type. Suppose that $y\lest w$ for $y, w\in W^c$. Write $z=y^{-1}w$. Then $ \ell(w)=\ell(y)+\ell(z) $, and  \begin{equation}\label{eq:Phi-w}
% 	\Phi_{w_cw}^+=\Phi_{w_cy}^+\sqcup  w_cy(\Phi_z^+) .
% 	\end{equation}
% \end{Lem}




% \subsection{A decomposition of $W^c$}

% Denote by $w_\circ$ (resp. $w_c$) the longest element in $W$ (resp. $W_c$). All the $\Phi_c^+$-dominant weights with infinitesimal character $-\rho$ can be written as $w_cw(-\rho)$ with $w\in W^c$. According to  \cite[Lem. 7.1]{CIS}, the set $ W^c $ can be partitioned into some subsets such that $ \gkd L(-w_cw\rho)$ is constant when $ w$ varies in each subset. 
 
% Indeed, we label elements of $ W^c $ in the same way as \cite[Figs. 7.1, 7.2]{CIS}, where  
% $ W^c=\{w_0,w_1,\cdots, w_{26}\} $ for type $ E_6 $ and $ W^c=\{w_0,w_1,\cdots, w_{55}\} $
% for type $ E_7 $.
% Reduced expressions of these elements can be easily read from  \cite[Figs. 7.1, 7.2]{CIS}, keeping in mind that our labels on simple roots (see Figure \ref{d67}) are different from those
% in \cite[page 81]{CIS}. For example, in type $ E_6 $, a reduced expression of $ w_9 $  is given by $ s_{\alpha_1}s_{\alpha_3}s_{\alpha_4}s_{\alpha_5}s_{\alpha_2}s_{\alpha_6} $.

% %One can also see that the longest element $ w^c $ of $ W^c $ is equal to $ w_{26} $  in type $ E_6 $ (or $ w_{55} $ in type $ E_7 $).
 
 
% Denote by $ W^c_{\lest w}  $ the subset $ \{ y\in W^c\mid y\lest w \}$ for any $ w\in W^c$. In \cite[Fig. 7.1]{CIS}, the set $ W^c $  for type $ E_6 $ is decomposed into three subsets $  \msc{C}_0=W^c_{\lest w_{19}}$, $  \msc{C}_1=W^c_{\lest w_{25}}\backslash W^c_{\lest w_{19}} $ and $  \msc{C}_2=W^c_{\lest w_{26}}\backslash W^c_{\lest w_{25}}=\{w_{26}\}$. Similarly, in \cite[Fig. 7.2]{CIS}, the set $ W^c $  for type $ E_7 $ is decomposed into 4 disjoint subsets $ \msc{C}_0=W^c_{\lest w_{23}}$, $\msc{C}_1=W^c_{\lest w_{48}}\backslash W^c_{\lest w_{23}}$, $\msc{C}_2=W^c_{\lest w_{54}}\backslash   W^c_{\lest w_{48}} $ and $\msc{C}_3=W^c_{\lest w_{55}}\backslash   W^c_{\lest w_{54}}=\{w_{55}\}$.






% \begin{Lem}\label{lemlast}
% 	Recall the sets $S_k$ defined in Theorem \ref{67}.
% 	\begin{itemize}
% 		\item [(1)] Let $\Phi=E_6$. Fix $0\lest k\lest 2$. The following conditions are equivalent:
% 		\begin{itemize}
% 			\item [(1a)] $z\in\msc{C}_k$;
% 			\item [(1b)] $\gkd L(-w_cz\rho)=(2-k)(8+3k)=|\Phi^+|-\aff(w_cz)$;
% 			\item [(1c)] $\Phi_{w_cz}^+\cap S_k\neq\emptyset$ when $k\geq1$ and $\Phi_{w_cz}^+\cap S_{k+1}=\emptyset$ when $k\lest1$.
% 		\end{itemize}
% 		\item [(2)] Let $\Phi=E_7$. Fix $0\lest k\lest 3$. The following conditions are equivalent:
% 		\begin{itemize}
% 			\item [(2a)] $z\in\msc{C}_k$;
% 			\item [(2b)] $\gkd L(-w_cz\rho)=(3-k)(9+4k)=|\Phi^+|-\aff(w_cz)$;
% 			\item [(2c)] $\Phi_{w_cz}^+\cap S_k\neq\emptyset$ when $k\geq1$ and $\Phi_{w_cz}^+\cap S_{k+1}=\emptyset$ when $k\lest2$.
% 		\end{itemize}
% 	\end{itemize}
% \end{Lem}


\subsection{Longest element of a Weyl group}
Recall that we use $\Delta $ to denote the simple root system of a simple Lie algebra $\mathfrak{g}$. 
Suppose $ w_I $ is the longest element of the parabolic subgroup of $ W $ generated by simple reflections arising from $ I\subset \Delta $, then $ \aff(w_I)$ is equal to  the length $\ell(w_I) $ of $ w_I $ by Lemma \ref{alem1}.
Let $\Phi_{I}$ be the root subsystem 
generated by simple roots in $I$, and we denote $\Phi_{I}^+=\Phi_I \cap \Phi^+ $.
% From now on, for a simple root $\alpha_i\in \Delta$ we simply write the simple reflection $s_{\alpha_i}$ by $s_i$.

%  When there are too many factors in an element $w$, we  adopt Geck's notation in \cite{ge}, which will be used to compute the value of $\aff(w)$. For example, we use $[i_1-1,i_2-1,\cdots,i_k-1]$ to denote
% $ w=s_{i_1}s_{i_2}\cdots s_{i_k}$ for type $A_n$, $F_4$ or $E_n$, and use $[n-i_1,n-i_2,\cdots,n-i_k]$ to denote
% $ w=s_{i_1}s_{i_2}\cdots s_{i_k}$ for type $B_n$, $C_n$ and $D_n$.



From \cite[Table 1]{BKOP}, we reproduce the following result, which will be used in our characterization of annihilator varieties.

\begin{Lem}\label{long}
    The longest element $w_0$ of a Weyl group can be expressed as the following:
    \begin{itemize}
        \item for type $A_n$,
        \begin{align*}
            w_{A_{n}}=&[n-1,n-2,n-1,\cdots,0,1,\cdots,n-1]\\
            =&s_n(s_{n-1}s_n)\cdots (s_1s_2\cdots s_n);
        \end{align*}
        \item for type $B_n$ or $C_n$,
        \begin{align*}
w_{B_{n}}=w_{C_{n}}=&[0,1,0,1,\cdots,n-k,\cdots,1,0,1,\cdots,n-k,\\
&\cdots,n-1,n-2,\cdots,1,0,1,\cdots,n-2,n-1]\\
=&s_n(s_{n-1}s_ns_{n-1})\cdots (s_ks_{k+1}\cdots s_{n-1} s_ns_{n-1}\cdots s_{k+1}s_k)\\
&\cdots (s_1s_2\cdots s_{n-1} s_ns_{n-1}\cdots s_2s_1);
        \end{align*}
        
        % $w_{B_{n}}=s_n(s_{n-1}s_n)\cdots (s_1s_2\cdots s_n)(s_1s_2\cdots s_{n-1})(s_1s_2\cdots s_{n-2})\cdots (s_1s_2)s_1$;
       
        \item for type $D_n$, 
        \begin{align*}
w_{D_{n}}=&[0,1,2,0,1,2,\cdots,n-k,\cdots,2,0,1,2,\cdots,n-k,\\
&\cdots,n-1,\cdots,2,0,1,\cdots,n-1]\\
=& s_ns_{n-1}(s_{n-2}s_ns_{n-1}s_{n-2})\cdots(s_k\cdots s_{n-2}s_ns_{n-1}s_{n-2}\cdots s_k)\\
&\cdots(s_1\cdots s_{n-2}s_n s_{n-1}\cdots s_{1});
% =&s_n(s_{n-2}s_{n-1}s_{n-3}s_{n-2}s_n)(s_{n-4}s_{n-3}s_{n-2}s_{n-1}s_{n-5}s_{n-4}s_{n-3}s_{n-2}s_n)\\
% &\cdots(s_2\cdots s_{n-1}s_1\cdots s_{n-2}s_n)(s_1s_2\cdots s_{n-1})\\
% &(s_1s_2\cdots s_{n-2})\cdots (s_1s_2)s_1;
        \end{align*}
% for type $D_n$ with $n$ odd, 
% \begin{align*}
% w_{D_{n}}=&[1,2,0,3,2,1,4,3,2,0,\\
% &5,4,3,2,1,6,5,4,3,2,0,\\
% &\cdots,n-2,\cdots,1,n-1,\cdots,2,0,n-1,n-2,\cdots,1,\\
% &n-1,n-2,\cdots,2,\cdots,n-1,n-2,n-1]\\
% =&s_{n-1}s_{n-2}s_{n}(s_{n-3}s_{n-2}s_{n-1}s_{n-4}s_{n-3}s_{n-2}s_n)\\
% & \cdot (s_{n-5}s_{n-4}s_{n-3}s_{n-2}s_{n-1}s_{n-6}s_{n-5}s_{n-4}s_{n-3}s_{n-2}s_n)\\
% &\cdots(s_2\cdots s_{n-1}s_1\cdots s_{n-2}s_n)(s_1s_2\cdots s_{n-1})\\
% & \cdot (s_1s_2\cdots s_{n-2})\cdots (s_1s_2)s_1;
%         \end{align*}


        
   
        \item for type $E_6$,
        \begin{align*}
 w_{E_{6}}=&[4,3,2,4,3,2,1,2,4,3,2,1,0,1,2,4,3,2,\\
 &1,0,5,4,2,3,1,0,2,1,4,2,3,5,4,2,1,0] \\
=&s_5s_{4}s_{3}s_5s_{4}s_{3}s_2 s_{3}s_5s_{4}s_{3} s_2s_1s_2 s_{3}s_5s_{4}s_3 \\
& \cdot 
s_2s_{1}s_6s_5s_3s_4s_2s_1s_3s_2s_5s_3s_4s_6s_5s_3s_2s_1;
% w_{E_{6}}=&[0,2,3,4,5,1,3,4,2,3,0,2,1,3,4,5,1,3,4,\\
% &2,3,0,2,1,3,4,2,3,0,2,1,3,0,2,1,0]\\
% =&s_1s_3s_4s_5s_6s_2s_4s_5s_3s_4s_1s_3s_2s_4s_5s_6s_2s_4s_5\\
% & \cdot s_3s_4s_1s_3s_2s_4s_5s_3s_4s_1s_3s_2s_4s_1s_3s_2s_1;
        \end{align*}
     \item for type $E_7$,
     \begin{align*}
w_{E_{7}}
=&[4,3,2,4,3,2,1,2,4,3,2,1,0,1,2,4,3,2,1,0,5,\\
 &4,2,3,1,0,2,1,4,2,3,5,4,2,1,0,6, 5, 4, 2, 3, 1,\\
 & 0, 2, 1, 4, 2, 3, 5, 4, 2, 1, 0, 6, 5, 4, 2, 3, 1, 2, 4, 5, 6] \\
=&s_5s_{4}s_{3}s_5s_{4}s_{3}s_2 s_{3}s_5s_{4}s_{3} s_2s_1s_2 s_{3}s_5s_{4}s_3 
s_2s_{1}s_6\\
 \cdot & s_5s_3s_4s_2s_1s_3s_2s_5s_3s_4s_6s_5s_3s_2s_1 s_7s_6s_5s_3s_4s_2\\
 \cdot & s_1s_3s_2s_5s_3s_4s_6s_5s_3s_2s_1s_7s_6s_5s_3s_4s_2s_3s_5s_6s_7;
% &[6,5,4,3,2,1,3,4,5,6,0,2,3,4,5,1,3,4,2,3,0,2,1,3,4,5,6,0,2,3,4,\\
% &5,1,3,4,2,3,0,2,1,3,4,5,1,3,4,2,3,0,2,1,3,4,2,3,0,2,1,3,0,2,1,0]\\
% =&s_7s_6s_5s_4s_3s_2s_4s_5s_6s_7s_1s_3s_4s_5s_6s_2s_4s_5s_3s_4s_1s_3s_2s_4s_5s_6s_7s_1s_3s_4s_5\\
% & \cdot s_6s_2s_4s_5s_3s_4s_1s_3s_2s_4s_5s_6s_2s_4s_5s_3s_4s_1s_3s_2s_4s_5s_3s_4s_1s_3s_2s_4s_1s_3s_2s_1;
        \end{align*}
        
        \item for type $E_8$, 
        \begin{align*}
w_{E_{8}}=&[4,3,2,4,3,2,1,2,4,3,2,1,0,1,2,4,3,2,1,0,5,\\
 &4,2,3,1,0,2,1,4,2,3,5,4,2,1,0,6, 5, 4, 2, 3, 1,\\
 & 0, 2, 1, 4, 2, 3, 5, 4, 2, 1, 0, 6, 5, 4, 2, 3, 1, 2, 4, 5, 6,\\
 &7, 6, 5, 4, 2, 3, 1, 0, 2, 1, 4, 2, 3, 5, 4, 2, 1, 0, 6, 5, 4,\\
 &2, 3, 1, 2, 4, 5, 6, 7, 6, 5, 4, 2, 3, 1, 0, 2, 1, 4, 2, 3, 5,\\
 &4, 2, 1, 0, 6, 5, 4, 2, 3, 1, 2, 4, 5, 6, 7] \\
=&s_5s_{4}s_{3}s_5s_{4}s_{3}s_2 s_{3}s_5s_{4}s_{3} s_2s_1s_2 s_{3}s_5s_{4}s_3 
s_2s_{1}s_6\\
 \cdot & s_5s_3s_4s_2s_1s_3s_2s_5s_3s_4s_6s_5s_3s_2s_1 s_7s_6s_5s_3s_4s_2\\
 \cdot & s_1s_3s_2s_5s_3s_4s_6s_5s_3s_2s_1s_7s_6s_5s_3s_4s_2s_3s_5s_6s_7\\
 \cdot &  
 s_8s_7s_6s_5s_3s_4s_2s_1s_3s_2s_5s_3s_4s_6s_5s_3s_2s_1s_7s_6s_5\\ \cdot & s_3s_4s_2s_3s_5s_6s_7s_8s_7s_6s_5s_3s_4s_2s_1s_3s_2s_5s_3s_4s_6\\ \cdot &s_5s_3s_2s_1s_7s_6s_5s_3s_4s_2s_3s_5s_6s_7s_8;
% =&[s_6s_,3,2,1,3,4,5,6,0,2,3,4,5,1,3,4,2,3,0,2,1,3,4,5,6,7,6,5,\\
% &4,3,2,1,3,4,5,6,0,2,3,4,5,1,3,4,2,3,0,2,1,3,4,5,6,7,6,5,4,3,2,\\
% &1,3,4,5,6,0,2,3,4,5,1,3,4,2,3,0,2,1,3,4,5,6,0,2,3,4,5,1,3,4,2,\\
% &3,0,2,1,3,4,5,1,3,4,2,3,0,2,1,3,4,2,3,0,2,1,3,0,2,1,0]\\
% =&s_8s_7s_6s_5s_4s_3s_2s_4s_5s_6s_7s_1s_3s_4s_5s_6s_2s_4s_5s_3s_4s_1s_3s_2s_4s_5s_6s_7s_8s_7s_6\\
% &\cdot s_5s_4s_3s_2s_4s_5s_6s_7s_1s_3s_4s_5s_6s_2s_4s_5s_3s_4s_1s_3s_2s_4s_5s_6s_7s_8s_7s_6s_5s_4s_3\\
% &\cdot s_2s_4s_5s_6s_7s_1s_3s_4s_5s_6s_2s_4s_5s_3s_4s_1s_3s_2s_4s_5s_6s_7s_1s_3s_4s_5s_6s_2s_4s_5s_3\\
% & \cdot s_4s_1s_3s_2s_4s_5s_6s_2s_4s_5s_3s_4s_1s_3s_2s_4s_5s_3s_4s_1s_3s_2s_4s_1s_3s_2s_1.
        \end{align*}
    \end{itemize}
\end{Lem}








\section{GK dimensions of  highest weight modules of type \texorpdfstring{$G_2$}{}}\label{g2-gkd}




In this section, $ \mf{g} $ is of type $ G_2 $. The positive roots are \[
\alpha_1=\ep_1-\ep_2, \alpha_2=-2\ep_1+\ep_2+\ep_3
\]
and $ \alpha_1+\alpha_2=-\ep_1+\ep_3 $, $ 2\alpha_1+\alpha_2=-\ep_2+\ep_3 $, $ 3\alpha_1+\alpha_2=\ep_1-2\ep_2 +\ep_3 $, $ 3\alpha_1+2\alpha_2=-\ep_1-\ep_2+2\ep_3 $. The simple roots are $ \alpha_1 $ and $ \alpha_2 $. The highest root is $\beta=3\alpha_1+2\alpha_2=-\ep_1-\ep_2+2\ep_3.$ 



\begin{figure}[htpb]
	\centering 
	\begin{tikzpicture}
 	% \node at(0,0) {$G_2$};
  \draw(6,0)circle[radius=0.1];
 	\draw(2,0)circle[radius=0.1];
 	\draw(4,0)circle[radius=0.1];
 	\draw (4.13,0)--(5.87,0) (2.12,0)--(3.87,0)  (2.11,0.07)--(3.88,0.07) (2.11,-0.07)--(3.88,-0.07);
 	\node at (2,0.5) {1};
 	\node at (4,0.5) {2};
 \end{tikzpicture}
	\caption{Extended Dynkin diagram of $ G_2 $ }
	\label{dg2}
\end{figure}

% \begin{center}
% $\dynkin[labels={\alpha_3,\alpha_2,\alpha_1},extended]G2$
%   \end{center}








We have $ \hs=\{ (\lam_1,\lam_2,\lam_3)\in\mathbb{C}^3\mid \lam_1+\lam_2+\lam_3 =0\} $. Fix $ \lam\in\hs $. We have \[
\bil{\lam}{(\ep_1- \ep_2)}=\lam_1- \lam_2,
\]
\[
\bil{\lam}{(-2\ep_1+\ep_2+\ep_3)}=\frac{-2\lam_1+\lam_2+\lam_3}{3}=-\lam_1.
\]
Then $ \lam=(\lambda_1,\lambda_2,\lambda_3) \in \mathfrak{h}^*$ is an integral weight if and only if $ (\lam_1,\lam_2)\in\mathbb{Z}^2 $.

By \cite[p. 128]{CM}, we know that all nilpotent orbits of $G_2$ have different dimensions.

\begin{Prop}\label{g2}
Let $L(\lam)$ be a simple integral highest weight module of $G_2$. Then
  the following holds:
    \begin{itemize}
        \item[(1)] $\gkd L(\lambda)=0$ if and only if $0>\lam_1 > \lam_2 $, i.e., $\lambda$ is dominant.
        \item[(2)] $\gkd L(\lambda)=5$ if and only if $\lambda$ is neither dominant nor antidominant.
         \item[(3)] $\gkd L(\lambda)=6$ if and only if $0\leq \lam_1\lest \lam_2 $, i.e., $\lambda$ is antidominant.  
    \end{itemize}
\end{Prop}
\begin{proof}
    From \cite[\S 4.8]{lusztig1984char}, we know that there are three two-sided cells in the Weyl group $W$ of $G_2$: $\mathcal{C}_1$, $\mathcal{C}_2$, $\mathcal{C}_3$. We denote the corresponding special nilpotent orbits via the Springer correspondence by $\mathcal{O}_1=\{0\}$, $\mathcal{O}_2=G_2(a_1)$, and $\mathcal{O}_3=G_2$. Note that $\mathcal{O}_3$ has the maximal dimension $12$ and $\mathcal{O}_1$ is the zero orbit. Also we have $ \mathcal{C}_3=\{e\}$, $s_{1}\in \mathcal{C}_2$ and $\mathcal{C}_1=\{s_{1}s_{2}s_{1}s_{2}s_{1}s_{2}\}$. We write $\lambda=w_{\lambda}\mu$, where $\mu$ is antidominant and $w_{\lambda}\in W_{[\mu]}^J$. When $\lambda$ is antidominant integral, we have $w_{\lambda}=e$; from Proposition \ref{pr:main1} we get $\gkd L(\lambda)=6$ since $\aff(e)=0$.  When $\lambda$ is dominant integral, $L(\lambda)$ is finite-dimensional; thus we have  $\gkd L(\lambda)=0$. The case $(2)$ is obvious since $w_{\lambda}\neq e$ or $s_{1}s_{2}s_{1}s_{2}s_{1}s_{2}$ for this case.  
\end{proof}


From \cite[\S 8.4]{CM}, we know that there are $5$ nilpotent orbits for $G_2$, of which there are only $3$ special ones, given as above. Thus, if we consider general $\lambda \in \mathfrak{h}^*$, then a-priori there are $5$ possibilities for the values of $\gkd L(\lambda)$.

As usual, the label $`\tilde{~}$' as in $A_1\times \tilde{A}_1\subset  G_2$  is attached to the connected component
corresponding to short roots.
Now we have the following result for nonintegral highest weight modules of $G_2$.

\begin{Prop}\label{g2-nonint}
Let $L(\lam)$ be a simple nonintegral highest weight module of $G_2$  such that $\Phi_{[\lam]}^\vee$ is pseudo-maximal. Then one of the followings holds:
\begin{itemize}
    \item[(1)] $\Phi_{[\lam]}\simeq A_2$. Then
$ \gkd L(\lam) $ can achieve the following values:
$$3,5,6.$$

\item[(2)] $\Phi_{[\lam]}\simeq A_1\times \tilde{A}_1$. Then
$ \gkd L(\lam) $ can achieve the following values:
$$4,5,6.$$
     
% \item $\Phi_{[\lam]}\simeq \tilde{A}_2$. Then
% $ \gkd L(\lam) $ can achieve the following values:
% $$3,5,6.$$
   
 
  \end{itemize}

\end{Prop}


\begin{proof}
 % The argument is similar with that in Theorem \ref{f4-nonint}. 
    
 From the extended Dynkin diagram of the root system of $G_2$ and Definition \ref{thmbd}, we can see that $\Phi_{[\lam]}$ with $\Phi_{[\lam]}^\vee$ being pseudo-maximal takes the following cases:
$$ A_2, A_1\times \tilde{A}_1.$$
    
    From Proposition \ref{pr:main1} we have \begin{equation*}
	\gkd L(\lambda)=|\Phi^+|-\aff(w_{\lambda}).
\end{equation*}
When $\Phi_{[\lam]}\simeq A_2$, from Lemma \ref{a-value} we know that $\aff(w_{\lambda})$ takes the following values: $$0,1,3.$$ 
Thus $ \gkd L(\lam) $ can achieve the following values:
$$3,5,6,$$
since $|\Phi^+|=6$ for $G_2$.

The proof of (2) is similar.
\end{proof}


\section{GK dimensions of  highest weight modules of type \texorpdfstring{$F_4$}{}}\label{f4-gkd}

%
%%%%%%%%%%%%%%%%%%%%%%%%%%%%%%%%%%%%%%%%%%%%%%%%%%%%%%%%%%%%%%%%%%%%
%


 In this section, $ \mf{g}$ is of type $ F_4 $.
The set $ \Phi^+ $ of positive roots is
\[
\{\ep_i\mid 1\lest i \lest 4\}\cup \{  \ep_i\pm \ep_j\mid 1\lest i<j \lest 4\} \cup\left \{ \frac12(\ep_1\pm \ep_2\pm\ep_3\pm \ep_4) \right\}.
\]
The  simple roots are $ \alpha_1=\ep_2-\ep_3,\alpha_2=\ep_3-\ep_4,\alpha_3=\ep_4, \alpha_4=\frac12(\ep_1- \ep_2-\ep_3- \ep_4)$.  The highest root is $\beta=2\alpha_1+3\alpha_2+4\alpha_3+2\alpha_4=\ep_1+\ep_2.$



\begin{figure}[htpb]
	\centering 
	\begin{tikzpicture}
 	% \node at(0,0) {$F_4$};
 	\foreach \i in {0,2,4,...,8}{\draw(\i,0)circle[radius=0.1];};
 		\draw (0.14,0)--(1.9,0)(2.14,0)--(3.86,0)  (4.13,0.05)--(5.865,0.05) (4.13,-0.05)--(5.865,-0.05)  (6.14,0)--(7.855,0) ;
 	\node at (0,0.5) {};
  \node at (2,0.5) {1};
 	\node at (4,0.5) {2};
 	\node at (6,0.5) {3};
 	\node at (8,0.5) {4};
 \end{tikzpicture}
	\caption{Extended Dynkin diagram of $ F_4 $ }
	\label{df4}
\end{figure}








\hspace{2cm}










For $ \lam=(\lam_1,\lam_2,\lam_3,\lam_4)\in\hs $, we have \[
\bil{\lam}{(\ep_i)}=2\lam_i,
\]
\[
\bil{\lam}{(\ep_i\pm \ep_j)}=\lam_i\pm \lam_j,
\]
\[
\left\langle{\lam},{\left(\frac12(\ep_1\pm \ep_2\pm\ep_3\pm \ep_4) \right)^{\vee}}\right\rangle=\lam_1\pm\lam_2\pm \lam_3\pm \lam_4.
\]

% Recall that we use $\Delta $ to denote the simple system of a simple Lie algebra $\mathfrak{g}$. 
% Suppose $ w_I $ is the longest element of the parabolic subgroup of $ W $ generated by simply reflections of $ I\subset \Delta $, then $ \aff(w_I)$ is equal to  the length $\ell(w_I) $ of $ w_I $ by Lemma \ref{alem1}.
% Let $\Phi_{I}$ be a sub root system 
% generated by simple roots in $I$, and denote $\Phi_{I}^+=\Phi_I \cap \Phi^+ $.









To avoid confusion, we use $\mathcal{O}(L)$ to denote the nilpotent orbit with Bala--Carter label $L$ in \cite{Ca85}. For many cases, the nilpotent orbit is uniquely determined by its dimension. So we only need to classify those highest weight modules whose annihilator varieties have the same dimension. 

For any $\lambda\in\mathfrak{h}^*$, recall that we can write $\lambda=w_{\lambda}\mu$, where $\mu$ is antidominant and $w_{\lambda}\in W_{[\mu]}^J$.  Then we have the following result.



\begin{Prop}\label{f4}
Let $L(\lam)$ be a simple integral highest weight module of $F_4$. Then
  the following holds:
    \begin{itemize}
     \item[(1)] $\gkd L(\lambda)=24$ if and only if $\lam_1\lest \lam_2\lest \lam_3\lest \lam_4\lest 0$ and $\lam_1\lest \lam_2+\lam_3+\lam_4$, i.e., $\lambda$ is antidominant.
      \item[(2)] $\gkd L(\lambda)=0$ if and only if $\lambda$ is dominant.
      \item[(3)] $\gkd L(\lambda)=11$ if and only if $w_{\lambda}\in \mathcal{C}w_0 $.
       \item[(4)] $\gkd L(\lambda)=23$ if and only if $w_{\lambda}\in \mathcal{C} $.
        \item[(5)] $\gkd L(\lambda)=22$ if and only if $w_{\lambda}\sim_{LR} w_{\{\alpha_1,\alpha_3\}}$.
         \item[(6)] $\gkd L(\lambda)=20$ if and only if $w_{\lambda}\sim_{LR} w_{\{\alpha_1,\alpha_3,\alpha_4\}}$.
         \item[(7)] $\gkd L(\lambda)=15$ if and only if $$w_{\lambda}\sim_{LR} w_{\{\alpha_1,\alpha_2,\alpha_3\}}=[2,1,2,0,1,2,0,1,0]
         %s_3s_{2}s_{3}s_{1}s_{2}s_{3}s_{1}s_{2}s_{1} 
         \text{~and~} V(\Ann (L(\lam)))=\overline{\mathcal{O}}(\tilde{A}_2),$$ or $$w_{\lambda}\sim_{LR} w_{\{\alpha_2,\alpha_3,\alpha_4\}}=[1,2,1,3,2,1,3,2,3]       
        % s_{2}s_{3}s_{2}s_{4}s_{3}s_{2}s_{4}s_{3}s_{4}  
          \text{~and~} V(\Ann (L(\lam)))=\overline{\mathcal{O}}(A_2).$$
         \item[(8)] $\gkd L(\lambda)=21$ if and only if $$w_{\lambda}\sim_{LR} w_{\{\alpha_3,\alpha_4\}} 
  \text{~and~} V(\Ann (L(\lam)))=\overline{\mathcal{O}}(B_3)$$ or $$w_{\lambda}\sim_{LR} w_{\{\alpha_1,\alpha_2\}}  \text{ ~and~} V(\Ann (L(\lam)))=\overline{\mathcal{O}}(C_3).$$
          \item[(9)] $\gkd L(\lambda)=14$ if and only if 
          $$w_{\lambda}\sim_{LR}[1, 2, 1, 2, 3, 2, 1, 0, 2, 1, 3, 2, 1, 0].$$
          %s_2s_3s_2s_3s_4s_3s_2s_1s_3s_2s_4s_3s_2s_1.$$
%           \begin{align*}&w_{\lambda}\sim_{LR} w_{\{\alpha_1,\alpha_2,\alpha_3,\ep_1\}}=w_{\{\alpha_1,\alpha_2,\alpha_3\}}s_{\ep_1}\\
%           =&s_{3}s_{2}s_{3}s_{1}s_{2}s_{3}s_{1}s_{2}s_{1}*
% s_{3}s_{4}s_{3}s_{4}s_{2}s_{3}s_{2}s_{4}s_1s_2s_3s_2s_1s_{3}s_{4}s_{3}s_{4}s_{2}s_{3}s_{2}s_{4}\\
% =&s_1s_2s_1s_3s_2s_1s_3s_2s_3s_4s_3s_2s_3s_4.              
%  \end{align*}
           \end{itemize}
\end{Prop}
\begin{proof}
    
  From \cite[\S 8.4]{CM}, we know that there are $11$ two-sided cells in the Weyl group $W$ of $F_4$ corresponding to the $11$ special nilpotent orbits listed in \cite[Table in page 128]{CM}. We denote these cells by by  $\mathcal{C}_i, 1\lest i\lest 11$, and the corresponding special nilpotent orbits via the Springer correspondence by $\mathcal{O}_i, 1\lest i\lest 11$. Here $\mathcal{O}_{11}$ is the regular orbit of maximal dimension $48$, $\mathcal{O}_{10}$ is the subregular orbit with $\dim \mathcal{O}_{10}=46$, and $\mathcal{O}_2$ is the minimal special orbit with $\dim \mathcal{O}_2=22$ and $\mathcal{O}_1$ is the zero orbit. 

By Lemma \ref{alem1}, we have $ \aff(w_I)=\ell(w_I)=|\Phi_I^+|$, where $ w_I $ is the longest element of the parabolic subgroup of $ W $ generated by simply reflections of some $ I\subset \Delta $. Thus our results in (1)--(8)  will follow from the equation 
$$\gkd L(\lambda)=24-\aff(w_{\lambda})=24-\aff(w_I)=24-|\Phi_I^+|,$$ 
if we know $w_{\lambda}\sim_{LR} w_I$ for some $w_{I}$ since there is only one or two two-sided cells with the given $\aff$ value, see \cite[Chapter 8]{CM} and \cite[\S 13.3]{Ca85}.



  
 Note that $ \mathcal{C}_1=\{e\}$.  We write $\lambda=w_{\lambda}\mu$, where $\mu$ is antidominant and $w_{\lambda}\in W_{[\mu]}^J$. When $\lambda$ is antidominant integral, we have $w_{\lambda}=e$.  Then from Proposition \ref{pr:main1}, we can get $\gkd L(\lambda)=24$ since $\aff(e)=0$. 
  
  When $\lambda$ is dominant integral, $L(\lambda)$ is finite-dimensional. We we have  $\gkd L(\lambda)=0$ in this case. 
  
   Note that  $ \mathcal{C}_2=\mathcal{C}w_0$ by \cite[Lem. 3.5]{BMXX}. Thus from Proposition \ref{2dim}, we have $\gkd L(\lambda)=11$ if and only if $w_{\lambda}\in \mathcal{C}w_0 $. Similarly  we have $\gkd L(\lambda)=23$ if and only if $w_{\lambda}\in \mathcal{C} $ since  $ \mathcal{C}_{10}=\mathcal{C}$ by \cite[Lem. 3.5]{BMXX}.

From \cite{BC76} and \cite[\S 8.4]{CM}, we know that $\mathcal{O}_{9}=F_4(a_2)$ is the distinguished nilpotent orbit corresponding to $I=\{\alpha_1,\alpha_3\}$  with $\dim \mathcal{O}_{9}=44$ and $w_{\{\alpha_1,\alpha_3\}}\in \mathcal{C}_9$ is the longest element in $W_{I}$. Thus we have $\gkd L(\lambda)=22$ if and only if $w_{\lambda}\sim_{LR} w_{\{\alpha_1,\alpha_3\}} $.

From \cite{BC76} and \cite[\S 8.4]{CM}, we know that $\mathcal{O}_{6}=F_4(a_3)$ is the distinguished nilpotent orbit corresponding to $I=\{\alpha_1,\alpha_3,\alpha_4\}$  with $\dim \mathcal{O}_{6}=40$ and $w_{\{\alpha_1,\alpha_3,\alpha_4\}}=s_{1}s_{4}s_{3}s_{4}\in \mathcal{C}_6$ is the longest element in $W_{I}$. Thus we have $\gkd L(\lambda)=20$ if and only if $w_{\lambda}\sim_{LR} w_{\{\alpha_1,\alpha_3,\alpha_4\}}=s_{1}s_{4}s_{3}s_{4}$.



By using PyCox, we find that $w_{\{\alpha_1,\alpha_2\}}$ (or the nilpotent orbit $\mathcal{O}_{w_{\{\alpha_1,\alpha_2\}}}$) corresponds to the character of $F_4$ with label $8_3$. By comparing the table in \cite[Table C.3]{GP}
and \cite[p. 428]{Ca85}, we find that the corresponding nilpotent orbit is $C_3$. Similarly, the  nilpotent orbit $\mathcal{O}_{w_{\{\alpha_3,\alpha_4\}}}$ corresponding to $w_{\alpha_3, \alpha_4}$ is $B_3$.

By noting that $s_3s_{2}s_{3}s_{1}s_{2}s_{3}s_{1}s_{2}s_{1}$ is denoted by $[2,1,2,0,1,2,0,1,0]$ and $s_{2}s_{3}s_{2}s_{4}s_{3}$ $
\cdot s_{2}s_{4}s_{3}s_{4}$ denoted by $[1,2,1,3,2,1,3,2,3]$ in PyCox, we find that the nilpotent orbit $\mathcal{O}_{w_{\{\alpha_1,\alpha_2,\alpha_3\}}}$ corresponding to $\{\alpha_1, \alpha_2, \alpha_3\}$ is $\tilde{A}_2$, and the  nilpotent orbit $\mathcal{O}_{w_{\{\alpha_2,\alpha_3,\alpha_4\}}}$ is $A_2$.

When $\gkd L(\lam)=14$, from Lemma \ref{alem1} we know that $\aff(w_{\lam})=10$. By \cite[Table in page 128]{CM}, there is only one special nilpotent orbit $\mathcal{O}_3$ with dimension equal to $28$. So we only need to find a representative element in the corresponding two-sided cell $\mathcal{C}_3$.
By using PyCox, we find that $\aff([1, 2, 1, 2, 3, 2, 1, 0, 2, 1, 3, 2, 1, 0])=10$. Thus $\gkd L(\lam)=14$ if and only if $w_{\lambda}\sim_{LR}[1, 2, 1, 2, 3, 2, 1, 0, 2, 1, 3, 2, 1, 0]$.
\end{proof}

\begin{Rem}\label{represent}
    In PyCox, we can find a representative element for a given $\bf a$ value $k$ in the following way:
    \begin{verbatim}
>>>W = coxeter("F", 4)
>>>left_cell_infos = klcellreps(W)
... for left_cell_info in left_cell_infos:
... if left_cell_info['a'] == 10:
>>>print(f"Index: {left_cell_info['index']}")
Index: 19
Index: 21
Index: 23
...for elm in left_cell_infos[19]['elms']:
>>>print(W.coxelmtoword(elm))
[1, 2, 1, 2, 3, 2, 1, 0, 2, 1, 3, 2, 1, 0]
[0, 1, 2, 1, 2, 3, 2, 1, 0, 2, 1, 3, 2, 1, 0]
[1, 0, 2, 1, 0, 3, 2, 1, 0, 2, 1, 2, 3, 2, 1, 0, 2, 1, 2]
[0, 2, 1, 0, 3, 2, 1, 0, 2, 1, 2, 3, 2, 1, 0, 2, 1, 2]
[0, 1, 0, 3, 2, 1, 0, 2, 1, 2, 3, 2, 1, 0, 2, 1, 2]
[0, 1, 0, 2, 1, 2, 3, 2, 1, 0, 2, 1, 3, 2, 1, 0]
[0, 2, 1, 0, 2, 1, 2, 3, 2, 1, 0, 2, 1, 3, 2, 1, 0]
[0, 1, 0, 2, 1, 0, 2, 1, 2, 3, 2, 1, 0, 2, 1, 2]
[1, 0, 2, 1, 0, 2, 1, 2, 3, 2, 1, 0, 2, 1, 3, 2, 1, 0]    
\end{verbatim}
In the above process, from the first step, we can find out all the left cells (with different Index numbers) with the same $\bf a$ value $10$. Then from the second step, we can find out all the Weyl group elements in the left cell with Index $19$. Then we choose a shortest one as a representative with the given $\bf a$ value $10$: $w_{\lambda}=[1, 2, 1, 2, 3, 2, 1, 0, 2, 1, 3, 2, 1, 0]$.

In the case of type $E_n$ (for $n=6,7,8$), it may happen that more than one two-sided cell has the same given $\bf a$ value. Then from the above process, we can get all the left cells with the  given $\bf a$ value. For each Index number, we choose an element in the corresponding left cell and find its corresponding special nilpotent orbit. Since there are finite Index numbers, we can give a complete set of representatives for the two sided cells  with the given $\bf a$ value. 
\end{Rem}





\begin{Rem}
    In Proposition \ref{f4}, we omit the expression of the longest element $w_I$ when $\Phi_I$ contains only type $A$ components. By Lemma \ref{long}, we know that its expression is very simple in this case.
\end{Rem}

From \cite[\S 8.4]{CM}, we know that there are $16$ nilpotent orbits for $F_4$, among which there are $11$ special ones. So there are $16$ possibilities for the values of $\gkd L(\lambda)$ for general $\lambda$.

We use $[m,n]$ to denote the set $\{m, m+1,\cdots,n\}$ for any positive integers $m<n$. Now we have the following result for nonintegral highest weight modules of $F_4$.





\begin{Prop}\label{f4-nonint}
Let $L(\lam)$ be a simple nonintegral highest weight module of $F_4$  such that $\Phi_{[\lam]}^\vee$ is pseudo-maximal.
Then one of the following holds:
    \begin{itemize}
	\item[(1)]  $\Phi_{[\lam]}\simeq C_4$. Then 
 $ \gkd L(\lam) $ can achieve the following values:
$$8,15,[18,24].$$

% When $ \gkd L(\lam)=15 $, we have $V(\Ann (L(\lam)))=\overline{\mathcal{O}}(A_2)$ if and only if $\pi_{w_{\lam}}=([1^3],[1])$ $(\text{equivalently~} \mathcal{O}(\pi_{w_{\lam}})=[2,2,1^4])$.



% When $ \gkd L(\lam)=18 $, we have $V(\Ann (L(\lam)))=\overline{\mathcal{O}}(B_2)$ if and only if $\pi_{w_{\lam}}=([1,1],[1,1])$ $(\text{equivalently~} \mathcal{O}(\pi_{w_{\lam}})=[2^4])$.

% When $ \gkd L(\lam)=21 $, we have $V(\Ann (L(\lam)))=\overline{\mathcal{O}}(B_3)$ if and only if $\pi_{w_{\lam}}=([2,1],[1])$ $(\text{equivalently~} \mathcal{O}(\pi_{w_{\lam}})=[4,2,2])$.
% \begin{enumerate}
%      \item $ \gkd L(\lam) =24$ 
%     % if and only if  	$ [\lam]_{1,2}=\{ (\lam_1,\lam_2),(\lam_3,\lam_4) \} $, $ \lam_1>\lam_2>0$, $\lam_3>\lam_4>0$,  $ \lam_1-\lam_2-\lam_3-\lam_4>0$. 
% if and only if $\lam$ is antidominant.
   
    
%     \item $ \gkd L(\lam) =23$ if and only if 
%     \item $ \gkd L(\lam) =22$ if and only if 
%     \item $ \gkd L(\lam) =21$ if and only if 
%     \item $ \gkd L(\lam) =20$ if and only if 
%      \item $ \gkd L(\lam) =8$ 
%     % if and only if  	$ [\lam]_{1,2}=\{ (\lam_1,\lam_2),(\lam_3,\lam_4) \} $, $ \lam_1>\lam_2>0$, $\lam_3>\lam_4>0$,  $ \lam_1-\lam_2-\lam_3-\lam_4>0$. 
% if and only if $\lam_2 \in \frac{1}{2}\mathbb{Z}_{>0}$,  $\lam_4 \in \frac{1}{2}\mathbb{Z}_{>0}$, $\lam_2 -\lam_3\notin \mathbb{Z}$, $\lam_1-\lam_2 \in \mathbb{Z}_{>0}$,$\lam_3-\lam_4 \in \mathbb{Z}_{>0}$, 
%     $ \lam_1-\lam_2-\lam_3-\lam_4\in \mathbb{Z}_{>0}$.

%     if and only if  	$ [\lam]_{1,2}=\{ (\lam_1,\lam_2),(\lam_3,\lam_4) \} $, $ \lam_1>\lam_2>0$, $\lam_3>\lam_4>0$,  $ \lam_1-\lam_2-\lam_3-\lam_4>0$. 
% \end{enumerate} 


  \item[(2)] $\Phi_{[\lam]}\simeq A_2\times \tilde{A}_2$. Then
$ \gkd L(\lam) $ can achieve the following values:
$$18,[20,24].$$
% When $ \gkd L(\lam)=18 $, we have $V(\Ann (L(\lam)))=\overline{\mathcal{O}}(\tilde{A}_2+A_1)$ if and only if $\mathcal{O}(\pi_{w_{\lam}})=[1^3]\times [1^3]$.

  \item[(3)] $\Phi_{[\lam]}\simeq A_1\times \tilde{A}_3$. Then
$ \gkd L(\lam) $ can achieve the following values:
$$17, 18, [20, 24].$$







% When $ \gkd L(\lam)=21 $, we have $V(\Ann (L(\lam)))=\overline{\mathcal{O}}(B_3)$ if and only if $\mathcal{O}(\pi_{w_{\lam}})=[1^3]\times [3]$, and $V(\Ann (L(\lam)))=\overline{\mathcal{O}}(C_3)$ if and only if $\mathcal{O}(\pi_{w_{\lam}})=[3]\times [1^3]$.


%   \item $\Phi_{[\lam]}\simeq A_1\times \tilde{A}_3$. Then
% $ \gkd L(\lam) $ can achieve the following values:
% $$17, 18, 20, 21, 22, 23, 24.$$



% When $ \gkd L(\lam)=18 $, we have $V(\Ann (L(\lam)))=\overline{\mathcal{O}}(B_2)$ if and only if $\mathcal{O}(\pi_{w_{\lam}})=[2]\times [1^4]$.

% When $ \gkd L(\lam)=21 $, we have  $V(\Ann (L(\lam)))=\overline{\mathcal{O}}(C_3)$ if and only if $\mathcal{O}(\pi_{w_{\lam}})=[2]\times [2,1,1]$.

    
%     \item $\Phi_{[\lam]}\simeq C_4$ and $\Delta_{[\lambda]}=\{\alpha_3,\alpha_4,\varepsilon_3, \alpha_1\}$. Then
% $ \gkd L(\lam) $ can achieve the following values:
% $$8,15,18,19,20,21,22,23,24.$$

% When $ \gkd L(\lam)=15 $, we have $V(\Ann (L(\lam)))=\overline{\mathcal{O}}({A}_2).$

% When $ \gkd L(\lam)=21 $, we have $V(\Ann (L(\lam)))=\overline{\mathcal{O}}(B_3).$

% When $ \gkd L(\lam)=18 $, we have $V(\Ann (L(\lam)))=\overline{\mathcal{O}}(B_2).$


% \begin{enumerate}
%     \item $ \gkd L(\lam) =24$ 
%     % if and only if  	$ [\lam]_{1,2}=\{ (\lam_1,\lam_2),(\lam_3,\lam_4) \} $, $ \lam_1>\lam_2>0$, $\lam_3>\lam_4>0$,  $ \lam_1-\lam_2-\lam_3-\lam_4>0$. 
% if and only if $\lam$ is antidominant.
    
%     \item $ \gkd L(\lam) =23$ if and only if 
%     \item $ \gkd L(\lam) =22$ if and only if 
%     \item $ \gkd L(\lam) =21$ if and only if 
%     \item $ \gkd L(\lam) =20$ if and only if 
% \end{enumerate} 

%     \item $\Phi_{[\lam]}\simeq A_2\times \tilde{A}_2$ and
% $\Delta_{[\lam]}=\{\alpha_1,\alpha_2,\alpha_4, \frac{1}{2}(\varepsilon_1+\varepsilon_2+\varepsilon_3+\varepsilon_4)\}$. Then
% $ \gkd L(\lam) $ can achieve the following values:
% $$18, 20, 21, 22, 23, 24.$$
%   When $ \gkd L(\lam)=21 $, we have $V(\Ann (L(\lam)))=\overline{\mathcal{O}}(B_3)$ if $w_{\lam}\sim_{LR}w_{\{\alpha_1,\alpha_2\}}=w_{A_2}$ and $V(\Ann (L(\lam)))=\overline{\mathcal{O}}(C_3)$ if $w_{\lam}\sim_{LR}w_{\{\alpha_4, \frac{1}{2}(\varepsilon_1+\varepsilon_2+\varepsilon_3+\varepsilon_4)\}}=w_{A^*_2}$.

% When $ \gkd L(\lam)=18 $, we have $V(\Ann (L(\lam)))=\overline{\mathcal{O}}(\tilde{A}_2+A_1).$  
    
    
    \item[(4)] $\Phi_{[\lam]}\simeq B_3\times \tilde{A}_1$. Then
$ \gkd L(\lam) $ can achieve the following values:
$$14, 15, [19,  24].$$

% When $ \gkd L(\lam)=15 $, we have $V(\Ann (L(\lam)))=\overline{\mathcal{O}}(\tilde{A}_2)$ if and only if $\pi_{w_{\lam}}=(\emptyset,[1^3])\times [2]$ $(\text{equivalently~}\mathcal{O}(\pi_{w_{\lam}})=[1^7]\times [2])$.

% When $ \gkd L(\lam)=21 $, we have $V(\Ann (L(\lam)))=\overline{\mathcal{O}}(C_3)$ if and only if $\pi_{w_{\lam}}=([1,1],[1])\times [2]$ $(\text{equivalently~}\mathcal{O}(\pi_{w_{\lam}})=[3,2,2]\times [2])$, and $V(\Ann (L(\lam)))=\overline{\mathcal{O}}(B_3)$ if and only if $\pi_{w_{\lam}}=([1],[2])\times [1,1]$ $(\text{equivalently~}\mathcal{O}(\pi_{w_{\lam}})=[3,3,1]\times [1,1])$.


% \begin{enumerate}
%     \item $ \gkd L(\lam) =24$ 
%     % if and only if  	$ [\lam]_{1,2}=\{ (\lam_1,\lam_2),(\lam_3,\lam_4) \} $, $ \lam_1>\lam_2>0$, $\lam_3>\lam_4>0$,  $ \lam_1-\lam_2-\lam_3-\lam_4>0$. 
% if and only if $\lam$ is antidominant.
    
%     \item $ \gkd L(\lam) =23$ if and only if 
%     \item $ \gkd L(\lam) =22$ if and only if 
%     \item $ \gkd L(\lam) =21$ if and only if 
%     \item $ \gkd L(\lam) =20$ if and only if 
% \end{enumerate} 
% \item $\Phi_{[\lam]}\simeq B_3\times \tilde{A}_1$ and
% $\Delta_{[\lam]}=\{\alpha_1,\alpha_2,\alpha_3,\varepsilon_1\}$. Then
% $ \gkd L(\lam) $ can achieve the following values:
% $$14, 15, 19, 20, 21, 22, 23, 24.$$

% When $ \gkd L(\lam)=15 $, we have $V(\Ann (L(\lam)))=\overline{\mathcal{O}}(\tilde{A}_2).$

% When $ \gkd L(\lam)=21 $, we have $V(\Ann (L(\lam)))=\overline{\mathcal{O}}(C_3).$
% \begin{enumerate}
%     \item $ \gkd L(\lam) =24$ 
%     % if and only if  	$ [\lam]_{1,2}=\{ (\lam_1,\lam_2),(\lam_3,\lam_4) \} $, $ \lam_1>\lam_2>0$, $\lam_3>\lam_4>0$,  $ \lam_1-\lam_2-\lam_3-\lam_4>0$. 
% if and only if $\lam$ is antidominant.
    
%     \item $ \gkd L(\lam) =23$ if and only if 
%     \item $ \gkd L(\lam) =22$ if and only if 
%     \item $ \gkd L(\lam) =21$ if and only if 
%     \item $ \gkd L(\lam) =20$ if and only if 
% \end{enumerate} 

	% 	\item 	$ [\lam]_{1,2}=\{ (\lam_1,\lam_3),(\lam_2,\lam_4) \} $, $ \lam_1>\lam_3>0$, $\lam_2>\lam_4>0$,  $ \lam_1-\lam_2-\lam_3-\lam_4>0$. In this case, $ V(L(\lam))=\overline{Be_{\alpha_1}} \cup \overline{Be_{\alpha_2}} $.
		
	% 	\item 	$ [\lam]_{1,2}=\{ (\lam_1,\lam_4),(\lam_2,\lam_3) \} $, $ \lam_1>\lam_4>0$, $\lam_2>\lam_3>0$,  $ \lam_1-\lam_2-\lam_3-\lam_4>0$. In this case, $ V(L(\lam))=  \overline{Be_{\alpha_2}} $.
	% \end{itemize}
\end{itemize}
\end{Prop}

\begin{proof}
% It is easy to see that  $ \Phi_{[\lam]}=\Phi  $ (i.e., $\lam$ is an integral weight) if and only if $2\lam_i\in \mathbb{Z}$ and $\lam_i-\lam_j\in \mathbb{Z}$ for all $1\lest i,j\lest 4$. 


From the extended Dynkin diagram of the root system of $F_4$ and Definition \ref{thmbd}, we can see that $\Phi_{[\lam]}$ with $\Phi_{[\lam]}^\vee$ being pseudo-maximal takes the following cases:
$$C_4, A_2\times \tilde{A}_2, A_1\times \tilde{A}_3, B_3\times \tilde{A}_1.$$


 From Proposition \ref{pr:main1} we have \begin{equation*}
	\gkd L(\lambda)=|\Phi^+|-\aff(w_{\lambda}).
\end{equation*}
When $\Phi_{[\lam]}\simeq C_4$,
  from Lemma \ref{a-value} we know that $\aff(w_{\lambda})$ takes the following values: $$0,1,2,3,4,5,6,9,16.$$ 
 Thus $ \gkd L(\lam) $ can achieve the following values:
$$8,15,18,19,20,21,22,23,24.$$
The proof of other cases is similar.




\end{proof}




% From Corollary \ref{transf-2}, we can see that for any $\lam\in \mathfrak{h}_{\Phi}^*$, it  is an integral weight in $\mathfrak{h}_{\Phi_{[\lam]}}^*$. Thus $\lam'=F(\lam)$ (or $\lam'$) is an integral weight of $\Psi_0$.

% \begin{Rem} From now on we denote $\beta_1=\ep_1-\ep_2, \beta_2=\ep_2-\ep_3, \beta_3=\ep_3$, and $\beta_4=\ep_4$. Then the simple root system of $C_4$ is $\Delta=\{\gamma_1=\beta_1=\ep_1-\ep_2, \gamma_2=\beta_2=\ep_2-\ep_3, \gamma_3=\beta_3-\beta_4, \gamma_4=2\beta_4\}$.
%     In general, by Theorem \ref{transf} we have the following:
%     \begin{enumerate}
%         \item In the case of $\mathfrak{l}_1$, we define a map $\phi: \Phi_{[\lambda]} \rightarrow C_4$ such that $\phi(\ep_2)=\gamma_1, \phi(\alpha_4)=\gamma_2, \phi(\alpha_3)=\gamma_3$ 
%  and $\phi(\alpha_2)=\gamma_4$. Then we can extend it to an isomorphism and still denote it by $\phi: \mathfrak{h}_{\Phi_{[\lambda]}}^*\rightarrow \mathfrak{h}_{C_4}^*$.  Thus for $\lam\in \mathfrak{h}_{F_4}^* $, it will be an integral weight in  $\mathfrak{h}_{\Phi_{[\lambda]}}^*$ and 
%  $$\phi(\lam)=\lambda'=(\lam_1+\lam_2,\lam_1-\lam_2,\lam_3+\lam_4,\lam_3-\lam_4)$$ will be an integral  weight of type $C_4$.
%         \item In the case of $\mathfrak{l}_2$, we have three cases:
%         \begin{enumerate}
%             \item we define a map $\phi: \Phi_{[\lambda]} \rightarrow A_2\times A^*_2$ such that  
%             $\phi(\varepsilon_1+\varepsilon_3)=\beta_1, \phi(\alpha_1)=\beta_2, \phi(\alpha_3)=\alpha_3$ and $ \phi(\alpha_4)=\frac{1}{2}(\ep_1+\ep_2+\ep_3-\ep_4)$. Then we can extend $\phi$ to an isomorphism and still denote it by $\phi: \mathfrak{h}_{\Phi_{[\lambda]}}^*\rightarrow \mathfrak{h}_{A_2\times A^*_2}^*$. We define another map $\psi: A^*_2\rightarrow A_2$ such that $\psi(\alpha_3)=\beta_1$ and $\psi(\frac{1}{2}(\ep_1+\ep_2+\ep_3-\ep_4))=\beta_2$.
%            Thus $$\phi(\lam)|_{\mathfrak{h}_{A_2}^*}=(\frac{2\lam_1+\lam_2+\lam_3}{3},\frac{-\lam_1+\lam_2-2\lam_3}{3},\frac{-\lam_1-2\lam_2+\lam_3}{3})$$ will be an integral weight of type $A_2$
%            and $$\psi\circ \phi(\lam)|_{\mathfrak{h}_{A^*_2}^*}=(\frac{\lam_1-\lam_2-\lam_3-3\lam_4}{3},\frac{\lam_1-\lam_2-\lam_3+3\lam_4}{3},\frac{-2\lam_1+2\lam_2+2\lam_3}{3})$$ will be an integral weight of type $A_2$ (or $ A^*_2$).
            
            
            
          
            
            
           
%             \item we define a map $\phi: \Phi_{[\lambda]} \rightarrow A_3\times A_1$ such that  
%             $\phi(\frac{1}{2}(\varepsilon_1+\varepsilon_2+\varepsilon_3-\varepsilon_4))=\beta_1, \phi(\alpha_3)=\beta_2, \phi(\alpha_4)=\ep_3-\ep_4$ and $ \phi(\alpha_1)=\ep_1+\ep_2+\ep_3+\ep_4$. Then we can extend $\phi$ to an isomorphism and still denote it by $\phi: \mathfrak{h}_{\Phi_{[\lambda]}}^*\rightarrow \mathfrak{h}_{A_3\times A_1}^*$. 
%            Thus $$\phi(\lam)|_{\mathfrak{h}_{A_3}^*}=(\frac{2\lam_1+\lam_2+\lam_3}{2},\frac{-\lam_2-\lam_3+2\lam_4}{2},\frac{-\lam_2-\lam_3-2\lam_4}{2}, \frac{-2\lam_1+\lam_2+\lam_3}{2})$$ will be an integral weight of type $A_3$
%            and $$\phi(\lam)|_{\mathfrak{h}_{A^*_1}^*}=\frac{\lam_2-\lam_3}{2}(\ep_1+\ep_2+\ep_3+\ep_4)$$ will be an integral weight of type $ A_1$. The value of  $\aff(\phi(\lam)|_{\mathfrak{h}_{A^*_1}^*})$  is $1$ if $\lam_2-\lam_3>0$  and $0$ if $\lam_2-\lam_3\lest 0$.

% % $\Phi_{[\lam]}\simeq A_1\times A_3$ and $\Delta_{[\lambda]}=\{ \alpha_1,\frac{1}{2}(\varepsilon_1+\varepsilon_2+\varepsilon_3-\varepsilon_4),\alpha_3,\alpha_4\}$.



            
%             \item we define a map $\phi: \Phi_{[\lambda]} \rightarrow C_4$ such that $\phi(\alpha_3)=\gamma_1, \phi(\alpha_4)=\gamma_2, \phi(\ep_3)=\gamma_3$ and $ \phi(\alpha_1)=\gamma_4$. Then we can extend it to an isomorphism and still denote it by $\phi: \mathfrak{h}_{\Phi_{[\lambda]}}^*\rightarrow \mathfrak{h}_{C_4}^*$.  Thus $$\phi(\lam)=\lambda'=(\lam_1+\lam_4,\lam_1-\lam_4,\lam_2+\lam_3,\lam_2-\lam_3)$$ will be an integral weight of type $C_4$.
%         \end{enumerate}




        
%        \item In the case of $\mathfrak{l}_3$, we have two cases:
%         \begin{enumerate}
%             \item we define a map $\phi: \Phi_{[\lambda]} \rightarrow A_2\times A^*_2$ such that  
%             $\phi(\alpha_1)=\beta_1, \phi(\alpha_2)=\beta_2, \phi(\alpha_4)=\alpha_3$ and $ \phi(\frac{1}{2}(\ep_1+\ep_2+\ep_3+\ep_4))=\frac{1}{2}(\ep_1+\ep_2+\ep_3-\ep_4)$. Then we can extend $\phi$ to an isomorphism and still denote it by $\phi: \mathfrak{h}_{\Phi_{[\lambda]}}^*\rightarrow \mathfrak{h}_{A_2\times A^*_2}^*$. We define another map $\psi: A^*_2\rightarrow A_2$ such that $\psi(\alpha_3)=\beta_1$ and $\psi(\frac{1}{2}(\ep_1+\ep_2+\ep_3-\ep_4))=\beta_2$.
%            Thus $$\phi(\lam)|_{\mathfrak{h}_{A_2}^*}=(\frac{2\lam_1-\lam_3-\lam_4}{3},\frac{-\lam_2+2\lam_3-\lam_4}{3},\frac{-\lam_2-\lam_3+2\lam_4}{3})$$ will be an integral weight of type $A_2$
%            and $$\psi\circ \phi(\lam)|_{\mathfrak{h}_{A^*_2}^*}=(\frac{3\lam_1-\lam_2-\lam_3-\lam_4}{3},\frac{2\lam_2+2\lam_3+2\lam_4}{3},-\frac{3\lam_1+\lam_2+\lam_3+\lam_4}{3})$$ will be an integral weight of type $A_2$ (or $ A^*_2$).
            
            
            
            
% %             $\Phi_{[\lam]}\simeq A_2\times A^*_2$ and
% % $\Delta_{[\lam]}=\{\alpha_1,\alpha_2,\alpha_4, \frac{1}{2}(\varepsilon_1+\varepsilon_2+\varepsilon_3+\varepsilon_4)\}$





%             \item we define a map $\phi: \Phi_{[\lambda]} \rightarrow B_3\times A_1$ such that $\phi(\alpha_2)=\beta_1, \phi(\alpha_1)=\beta_2, \phi(\frac{1}{2}(\ep_1-\ep_2+\ep_3+\ep_4))=\beta_3, \phi(\alpha_4)=\beta_4$. Then we can extend it to an isomorphism and still denote it by $\phi: \mathfrak{h}_{\Phi_{[\lambda]}}^*\rightarrow \mathfrak{h}_{B_3\times A_1}^*$.  Thus  \begin{align*}
%            &\phi(\lam)=\lambda'\\
%            =& (\frac{\lam_1+\lam_2+\lam_3-\lam_4}{2},\frac{\lam_1+\lam_2-\lam_3+\lam_4}{2}, \frac{\lam_1-\lam_2+\lam_3+\lam_4}{2},\\
%            &\frac{\lam_1-\lam_2-\lam_3-\lam_4}{2})
%        \end{align*}
%        will be a weight of type $B_3\times A_1$.
%             \end{enumerate}
       
       
     
%        \item In the case of $\mathfrak{l}_4$, we define a map $\phi: \Phi_{[\lambda]} \rightarrow B_3\times A_1$ such that $\phi(\alpha_1)=\beta_1, \phi(\alpha_2)=\beta_2, \phi(\alpha_3)=\beta_3, \phi(\ep_1)=\beta_4$. 
% Then we can extend it to an isomorphism and still denote it by $\phi: \mathfrak{h}_{\Phi_{[\lambda]}}^*\rightarrow \mathfrak{h}_{B_3\times A_1}^*$.  Thus $$\phi(\lam)=\lambda'=(\lam_2,\lam_3,\lam_4,\lam_1)$$  will be a weight of type $B_3\times A_1$.
%     \end{enumerate}
% Then the value of $\aff(w_{\lam})=\aff(w_{\lam'})$ can be computed by the algorithm in \cite{BXX}.    
% \end{Rem}

\begin{ex}\label{exf4}
	Let $\mathfrak{g}=F_4$ and $L(\lam)$ be the highest weight module of $\mathfrak{g}$ with  $\lambda=(4, 5, \frac{3}{2}, \frac{1}{2})$. It is easy to verify that $\Phi_{[\lambda]}$ is a subsystem with simple roots $\{\alpha_2, \alpha_3, \alpha_4, \ep_2\}$. Therefore $\Phi_{[\lambda]}\simeq C_4$. Suppose the simple root system of $C_4$ is $\Delta=\{\beta_1=\ep_1-\ep_2, \beta_2=\ep_2-\ep_3, \beta_3=\ep_3-\ep_4,\beta_4=2\ep_4\}$. We consider the isomorphism $\phi: \Phi_{[\lambda]} \rightarrow C_4$ such that $\phi(\varepsilon_2)=\beta_1, \phi(\alpha_4)=\beta_2, \phi(\alpha_3)=\beta_3, \phi(\alpha_2)=\beta_4$. We have 
 $$\lambda=\omega_1+\omega_2-3\omega_3+10\omega_{4},$$
 where $\{\omega_i\mid 1\leq i\leq 4\}$ are the fundamental weights for the simple system of $\Phi_{[\lambda]}$. Thus we can get  
 $$\phi(\lam)=\lambda'=\phi(\omega_1)+\phi(\omega_2)-3\phi(\omega_3)+10\phi(\omega_{4})=(9,-1,2,1),$$ which is an integral weight of type $C_4$. So $\aff(w_{\lambda})=\aff(w_{\lambda'})=9$ by the RS algorithm in \cite{BXX} and Proposition \ref{integral}.
 By Proposition \ref{pr:main1},  we have $\gkd L(\lambda)=|\Phi^+|-\aff(w_{\lambda})=24-9=15$.
% and $V(\Ann (L(\lam)))=\overline{\mathcal{O}}(A_2)$.

% Note that $\lambda'=(9,-1,2,1)$ will give us a $B$-symbol
% $\Lambda =\begin{pmatrix}
% 	0~1~2~3\\1~2~4
% 	\end{pmatrix} $. The corresponding special symbol is $\Lambda^s =\begin{pmatrix}
% 	0~1~2~4\\1~2~3
% 	\end{pmatrix} $, which gives  a bipartition
%  $({\bf p},{\bf q})=([1]],[1,1,1])$.
%  From \cite[ (4.5.4)]{lusztig1984char}, we have $\Lambda=\pi_{({\bf p},{\bf q})}=j_{W_3}^{W_4}\Lambda'$ with $\Lambda' =\begin{pmatrix}
% 	0~1~2~3\\1~2~3
% 	\end{pmatrix}.$ The corresponding bipartition of $\Lambda'$ is $({\bf p}',{\bf q}')=(\emptyset,[1,1,1])$. Thus from \cite[Table C.3 \& Table D.3]{GP} and \cite[p.428]{Ca85}, we have $V(\Ann (L(\lam)))=\overline{\mathcal{O}}(A_2)$.
\end{ex}
 
\begin{ex}
	Let $\mathfrak{g}=F_4$ and $L(\lam)$ be the highest weight module of $\mathfrak{g}$ with  $\lambda=(\frac{7}{4}, \frac{1}{4}, \frac{5}{4}, -\frac{3}{4})$. It is easy to verify that $\Phi_{[\lambda]}$ is a subsystem with simple roots $\{ \alpha_2,\alpha_1, \frac{1}{2}(\ep_1-\ep_2+\ep_3+\ep_4), \alpha_4\}$. Therefore $\Phi_{[\lambda]}\simeq B_3\times A_1$. Suppose the simple root system of $B_3\times A_1$ is $\Delta=\{\beta_1=\ep_1-\ep_2, \beta_2=\ep_2-\ep_3, \beta_3=\ep_3\}\times \{\beta_1'=\ep_1-\ep_2\}$. We define a map $\phi: \Phi_{[\lambda]} \rightarrow B_3\times A_1$ such that $\phi(\alpha_2)=\beta_1, \phi(\alpha_1)=\beta_2, \phi(\frac{1}{2}(\ep_1-\ep_2+\ep_3+\ep_4))=\beta_3, \phi(\alpha_4)=\beta_1'$. 
 Then we have 
$$\lambda=2\omega_1-\omega_2+2\omega_3+\omega'_1$$
    where $\{\omega_i\mid 1\leq i\leq 3\}\cup \{\omega'_1\}$ are the fundamental weights for $\Delta$. 

We get  
 \begin{align*}
     \phi(\lam)=\lambda'=& 2\phi(\omega_1)-\phi(\omega_2)+2\phi(\omega_3)+\phi(\omega'_1)
 \end{align*}
and thus $$\lam'|_{B_3}=2\phi(\omega_1)-\phi(\omega_2)+2\phi(\omega_3)=(2,0,1)$$
 which is an integral weight of type  $B_3$ with $\aff(w_{\lam'|_{{B_3}}})=4$, and $$\lam'|_{A_1}=\phi(\omega'_1)=(\frac{1}{2},-\frac{1}{2}),$$
 which is an integral weight of type  $A_1$ with $\aff(w_{\lam'|_{{A_1}}})=1$ by using  the RS algorithm in \cite{BXX} and Proposition \ref{integral}.

 
 Therefore we have  $\aff(w_{\lambda})=\aff(w_{\lambda'})=4+1=5$ by the RS algorithm in \cite{BXX}.
 By Proposition \ref{pr:main1},  we have $\gkd L(\lambda)=|\Phi^+|-\aff(w_{\lambda})=24-5=19$. 
% Thus from \cite[p.128]{CM}, we have $V(\Ann (L(\lam)))=\overline{\mathcal{O}}(C_3(a_1))$.
\end{ex}

% \section{GK dimensions of  highest weight modules of type \texorpdfstring{$G_2$}{}}\label{g2-gkd}




% In this section, $ \mf{g} $ is of type $ G_2 $. The positive roots are \[
% \alpha_1=\ep_1-\ep_2, \alpha_2=-2\ep_1+\ep_2+\ep_3
% \]
% and $ \alpha_1+\alpha_2=-\ep_1+\ep_3 $, $ 2\alpha_1+\alpha_2=-\ep_2+\ep_3 $, $ 3\alpha_1+\alpha_2=\ep_1-2\ep_2 +\ep_3 $, $ 3\alpha_1+2\alpha_2=-\ep_1-\ep_2+2\ep_3 $. The simple roots are $ \alpha_1 $ and $ \alpha_2 $. The highest root is $$\beta=3\alpha_1+2\alpha_2.$$ In this case, we have $E_{\Phi}=\mathbb{R}^3$.



% \begin{figure}[htpb]
% 	\centering 
% 	\begin{tikzpicture}
%  	% \node at(0,0) {$G_2$};
%   \draw(6,0)circle[radius=0.1];
%  	\draw(2,0)circle[radius=0.1];
%  	\draw(4,0)circle[radius=0.1];
%  	\draw (4.2,0)--(5.87,0) (2.2,0)--(3.8,0)  (2.18,0.1)--(3.82,0.1) (2.18,-0.1)--(3.82,-0.1);
%  	\node at (2,0.5) {1};
%  	\node at (4,0.5) {2};
%  \end{tikzpicture}
% 	\caption{Extended Dynkin diagram of $ G_2 $ }
% 	\label{dg2}
% \end{figure}










% We have $ \hs=\{ (\lam_1,\lam_2,\lam_3)\in\mathbb{C}^3\mid \lam_1+\lam_2+\lam_3 =0\} $. Fix $ \lam\in\hs $. We have \[
% \bil{\lam}{(\ep_1- \ep_2)}=\lam_1- \lam_2,
% \]
% \[
% \bil{\lam}{(-2\ep_1+\ep_2+\ep_3)}=\frac{-2\lam_1+\lam_2+\lam_3}{3}=-\lam_1.
% \]
% Then $ \lam=(\lambda_1,\lambda_2,\lambda_3) $ is an integral weight if and only if $ (\lam_1,\lam_2)\in\mathbb{Z}^2 $.

% By \cite[p. 128]{CM}, we know that all nilpotent orbits of $G_2$ have different dimensions.

% \begin{Thm}\label{g2}
% Let $L(\lam)$ be a simple integral highest weight module of $G_2$. Then
%   the following holds:
%     \begin{enumerate}
%         \item $\gkd L(\lambda)=0$ if and only if $0>\lam_1 > \lam_2 $, i.e., $\lambda$ is dominant integral.
%         \item $\gkd L(\lambda)=5$ if and only if $\lambda$ is neither dominant nor antidominant.
%          \item $\gkd L(\lambda)=6$ if and only if $0\leq \lam_1\lest \lam_2 $, i.e., $\lambda$ is antidominant integral.  
%     \end{enumerate}
% \end{Thm}
% \begin{proof}
%     From \cite[\S 4.8]{lusztig1984char}, we know that there are three two-sided cells in the Weyl group $W$ of $G_2$: $\mathcal{C}_1$, $\mathcal{C}_2$, $\mathcal{C}_3$. We denote the corresponding special nilpotent orbits via the Springer correspondence by $\mathcal{O}_1=\{0\}$, $\mathcal{O}_2=G_2(a_1)$, and $\mathcal{O}_3=G_2$. Note that $\mathcal{O}_3$ has the maximal dimension $12$ and $\mathcal{O}_1$ is the zero orbit. Also we have $ \mathcal{C}_3=\{e\}$, $s_{1}\in \mathcal{C}_2$ and $\mathcal{C}_1=\{s_{1}s_{2}s_{1}s_{2}s_{1}s_{2}\}$. We write $\lambda=w_{\lambda}\mu$, where $\mu$ is antidominant and $w_{\lambda}\in W_{[\mu]}^J$. When $\lambda$ is antidominant integral, we have $w_{\lambda}=e$; from Proposition \ref{pr:main1} we get $\gkd L(\lambda)=6$ since $\aff(e)=0$.  When $\lambda$ is dominant integral, $L(\lambda)$ is finite-dimensional; thus we have  $\gkd L(\lambda)=0$. The case $(2)$ is obvious since $w_{\lambda}\neq e$ or $s_{1}s_{2}s_{1}s_{2}s_{1}s_{2}$ for this case.  
% \end{proof}


% From \cite[\S 8.4]{CM}, we know that there are $5$ nilpotent orbits for $G_2$, of which there are $3$ special ones. Thus, if we consider general $\lambda \in \mathfrak{h}^*$, then a-priori there are $5$ possibilities for the values of $\gkd L(\lambda)$.



% % Denote $\Phi_{[\lam]}^+=\Phi_{[\lam]}\cap \Phi^+$ and $ \Phi_\lam^+=\{\alpha\in\Phi_{[\lam]}^+\mid\langle\lambda, \alpha^\vee\rangle>0\} $.

% %Then we have the following result for nonintegral highest weight modules of $G_2$.

% \begin{Thm}\label{g2-nonint}
% Let $L(\lam)$ be a simple nonintegral highest weight module of $G_2$  such that $\Phi_{[\lam]}^\vee$ is pseudo-maximal. Then one of the followings holds:
% \begin{enumerate}
%     \item $\Phi_{[\lam]}\simeq A_2$. Then
% $ \gkd L(\lam) $ can achieve the following values:
% $$3,5,6.$$

% \item $\Phi_{[\lam]}\simeq A_1\times \tilde{A}_1$. Then
% $ \gkd L(\lam) $ can achieve the following values:
% $$4,5,6.$$
     
% % \item $\Phi_{[\lam]}\simeq \tilde{A}_2$. Then
% % $ \gkd L(\lam) $ can achieve the following values:
% % $$3,5,6.$$
   
 
%   \end{enumerate}

% \end{Thm}


% \begin{proof}
%  The argument is similar with that in Theorem \ref{f4-nonint}. 
    
% %     From Proposition \ref{pr:main1} we have \begin{equation*}
% % 	\gkd L(\lambda)=|\Phi^+|-\aff(w_{\lambda}).
% % \end{equation*}
% % When $\Phi_{[\lam]}\simeq A_2$, from Lemma \ref{a-value} we know that $\aff(w_{\lambda})$ takes the following values: $$0,1,3.$$ 
% % Thus $ \gkd L(\lam) $ can achieve the following values:
% % $$3,5,6.$$
% % The proof of other cases is similar.
% \end{proof}




% \begin{Thm}Suppose $L(\lam)$ is a simple nonintegral highest weight module of $G_2$.
% 	Then
%     the following holds:
%     \begin{enumerate}
%         \item $\gkd L(\lambda)=3$  if and only if
% 	$ \lam_1\equiv \lam_2\equiv \pm \frac13 ~(\mathrm{mod}~ \mathbb{Z})$, $ \lam_1- \lam_2>0 $, and $ 2\lam_1+\lam_2<0 $. 
%         \item $\gkd L(\lambda)=4$ if and only if $|\Phi_{[\lam]}^+|=|\Phi_{\lam}^+|=2$.
%         \item $\gkd L(\lambda)=5$ if and only if $|\Phi_{\lam}^+|=1$ or $2$.
%          \item $\gkd L(\lambda)=6$ if and only if $|\Phi_{\lam}^+|=0$.
%     \end{enumerate}
 
 

% \end{Thm}
% \begin{proof}
% There are $6$ possible subsystems $ \Psi $ of $ \Phi $ with $ \Psi^+=\Psi\cap\Phi^+ $ given by
% \begin{enumerate}
% 	\item[(a)] $ \{ \alpha,\beta+\alpha,\beta+2\alpha \} $,
% 	\item[(b)] $ \{ \beta,\beta+3\alpha,2\beta+3\alpha \} $,
% 	\item[(c)] $ \{ \beta,\beta+2\alpha\} $,
% 	\item[(d)] $ \{\beta+\alpha,\beta+3\alpha \} $,
%  \item[(e)] $ \{\alpha,2\beta+3\alpha \} $,
% 	\item[(f)] $ |\Psi^+|=1 $.
% \end{enumerate}

% In case (a), $ \Psi\subset\Phi_{{[\lam]}} $ if and only if $ \lam_1-\lam_2,-2\lam_1-\lam_2\in\mathbb{Z} $, if and only if $ 3\lam_1\in2\mathbb{Z},\lam_2\in \lam_1+\mathbb{Z} $. Thus $ \Psi\subset \Phi_{[\lam]}\neq\Phi $ if and only if $ \lam_1\in\pm \frac13+\mathbb{Z} $, $ \lam_2\in \lam_1+\mathbb{Z} $. So we have  $W_{[\lam]}\simeq A_2$.

% In case (b),  $ \Psi\subset\Phi_{{[\lam]}} $ if and only if $- \lam_1,-\lam_2\in\mathbb{Z} $, if and only if $ \Phi_{[\lam]}=\Phi  $.

% In case (c), $\Phi_{[\lam]}=\{ \pm\beta,\pm(\beta+2\alpha)\}$. We will have $-\lam_1\in \mathbb{Z}$ and  $-\lam_2+\lam_3\in \mathbb{Z}$. From $\lam_1+\lam_2+\lam_3=0$, we can get $\lam_2,\lam_3\in \frac{1}{2}+\mathbb{Z}$. 

% In case (d), $ \Phi_{[\lam]}=\{\pm(\beta+\alpha),\pm(\beta+3\alpha) \}$. We will have  $-\lam_1+\lam_3\in \mathbb{Z}$
% and $-\lam_2\in \mathbb{Z}$. From $\lam_1+\lam_2+\lam_3=0$, we can get $\lam_1,\lam_3\in \frac{1}{2}+\mathbb{Z}$. 

% In case (e), $ \Phi_{[\lam]}=\{\pm\alpha, \pm(2\beta+3\alpha) \}$. We will have  $\lam_1-\lam_2\in \mathbb{Z}$
% and $\lam_3\in \mathbb{Z}$. From $\lam_1+\lam_2+\lam_3=0$, we can get $\lam_1,\lam_2\in \frac{1}{2}+\mathbb{Z}$. 

% For the above three cases, we have $W_{[\lam]}\simeq A_1\times A_1$.

% In case (f), $ |\Phi_{[\lam]}^+|=1$ and  one of $\{\lam_i\mid 1\lest i\lest 3\}$ or $\{\lam_i-\lam_j\mid 1\lest i< j\lest 3\}$ will belong to $\mathbb{Z}$.
% Thus we have  $W_{[\lam]}\simeq A_1$.

% Recall that we  write  $\lambda=w_{\lambda}\mu$, where $\mu$ is antidominant and $w_{\lambda}\in W_{[\mu]}^J$. Then $w_{\lambda}$ is a product of elements in the subgroups of $W_{[\lam]}$. So $\aff(w_{\lambda})$ is the summation of the values of $\aff$-functions of these elements. 

% For type (a), from Lemma \ref{a-value} we have the following: 
% \begin{enumerate}
%     \item $\aff(w_{\lambda})=3$ if and only if 
% $ \lam_1\equiv \lam_2\equiv \pm \frac13 
%   ~(\mathrm{mod}~ \mathbb{Z})$, $ \lam_1- \lam_2>0 $, and $ 2\lam_1+\lam_2<0 $, equivalently, $|\Phi_{\lam}^+|=3$;
% \item  $\aff(w_{\lambda})=0$ if and only if 
% $ \lam_1\equiv \lam_2\equiv \pm \frac13 ~(\mathrm{mod}~ \mathbb{Z})$, $ \lam_1\lest \lam_2 $, and $ 2\lam_1+\lam_2\geq 0 $, equivalently, $|\Phi_{\lam}^+|=0$;
% \item  $\aff(w_{\lambda})=1$ if and only if 
% $ \lam_1\equiv \lam_2\equiv \pm \frac13 ~(\mathrm{mod}~ \mathbb{Z})$ and $|\Phi_{\lam}^+|=1$ or $2$.
% \end{enumerate}

% For types (b), (c) and (d), we have $\aff(w_{\lambda})=0$ if and only if $|\Phi_{\lam}^+|=0$, $\aff(w_{\lambda})=1$ if and only if $|\Phi_{\lam}^+|=1$, and $\aff(w_{\lambda})=2$ if and only if $|\Phi_{\lam}^+|=2$.

% For type (f), we have $\aff(w_{\lambda})=0$ if and only if $|\Phi_{\lam}^+|=0$, and $\aff(w_{\lambda})=1$ if and only if $|\Phi_{\lam}^+|=1$.

% Recall that from Proposition \ref{pr:main1}, we have \begin{equation*}
% 	\gkd L(\lambda)=|\Phi^+|-\aff(w_{\lambda}).
% \end{equation*}
% Thus we  complete the proof.
  




% %   must 
% % have $\gk\Phi_{[\lam]}d L(\lambda) = 4$ since there is only one nilpotent orbit corresponding to it.
% \end{proof}

% \begin{Cor}
%     Suppose $L(\lam)$ is a simple nonintegral highest weight module of $G_2$.
% 	Then $\gkd L(\lambda)=6-k$ with $k=0,2,3$ if and only if $|\Phi_{\lam}^+|=k$.
% \end{Cor}

\section{GK dimensions of  highest weight modules of type \texorpdfstring{$E$}{} }\label{e-gkd}



In this section, $\mf{g}$ is of type $ E $. The root system $ \Phi $ can be realized as a subset of $\mathbb{R}^8$.  The simple roots are 
\begin{align*}
\alpha_1&=\frac{1}{2}(\ep_1-\ep_2-\ep_3-\ep_4-\ep_5-\ep_6-\ep_7+\ep_8) ,\\
\alpha_2&=\ep_1+\ep_2,\\
\alpha_k&=\ep_{k-1}-\ep_{k-2} \text{ with }3\lest k\lest n,
\end{align*}
where $n=6, 7,8$. 




\subsection{Type \texorpdfstring{$E_6$}{}}
For type $E_6$, we have $$\Phi^+=
\{\ep_i \pm \ep_j \mid 5\geq i > j\geq 1\}\cup 
\Big\{\frac{1}{2}(\sum_{i=1}^5 (-1)^{n(i)}\ep_i -\ep_6-\ep_7+\ep_8)\mid \sum_{i=1}^5n(i) \text{ even}\Big\}$$
and $ \hs=\{ (\lam_1,\lam_2,\cdots,\lam_8)\in\mathbb{C}^8\mid \lam_6=\lam_7=-\lam_8\} $.  The highest root is $$\beta=\alpha_1+2\alpha_2+2\alpha_3+3\alpha_4+2\alpha_5+\alpha_6=\frac{1}{2}(\ep_1+\ep_2+\ep_3+\ep_4+\ep_5-\ep_6-\ep_7+\ep_8).$$

\begin{figure}[htpb]
	\centering 
	
	\begin{tikzpicture}[scale=1.5,baseline=0]
	
 
 \draw (0,0.05) node[above=1pt]{$ 1 $} circle [radius=0.05];
	\draw (0.05,0.05)--++(0.65,0); 
	\draw (0.75,0.05) node[above=1pt]{$ 3 $} circle[radius=0.05];  
	\draw (0.8,0.05)--++(0.65,0);
	\draw  (1.45,0) ++(0.05,0.05) node[above=1pt]{$ 4$} circle[radius=0.05];
	\draw (1.55,0.05)--++(0.65,0);
	\draw  (2.2,0) ++(0.05,0.05) node[above=1pt]{$ 5$} circle[radius=0.05];
	\draw (2.3,0.05)--++(0.65,0);
	\draw  (2.95,0) ++(0.05,0.05) node[above=1pt]{$ 6$} circle[radius=0.05];
	\draw (1.5,0)--++(0,-0.5);
	\draw  (1.45,-0.6) ++(0.05,0.05) node[right=1pt]{$ 2$} circle[radius=0.05];

\draw (1.49,-1.2) node[above=1pt]{} circle [radius=0.05];
	\draw (1.5,-0.6)--++(0,-0.53); 
 
	\end{tikzpicture}
	
	
	
	
	\caption{Extended Dynkin diagram of $ E_6 $ }
	\label{E6-dy}
\end{figure}







Fix $ \lam\in\hs $. We have \[
\bil{\lam}{(\alpha_1)}=\frac{1}{2}(\lam_1- \lam_2-\cdots-\lam_7+\lam_8),
\]
\[
\bil{\lam}{(\ep_i\pm \ep_j)}=\lam_i\pm \lam_j.
\]
Then $ \lam=(\lam_1,\lam_2,\cdots,\lam_8) $ is an integral weight if and only if $\lam_1- \lam_2-\cdots-\lam_7+\lam_8\in 2\mathbb{Z}$,  $ \lam_1-\lam_2\in\mathbb{Z}, \lam_2-\lam_3\in\mathbb{Z},\lam_3-\lam_4\in\mathbb{Z},\lam_4-\lam_5\in\mathbb{Z}$ and $2\lam_i\in \mathbb{Z} $ for $1\lest i\lest 5$.

  By similar arguments as in Proposition \ref{f4}, we have the following result.


\begin{Prop}\label{E6}
Let $L(\lam)$ be a simple integral highest weight module of $E_6$. Then
  the following holds:
    \begin{itemize}
        \item[(1)] $\gkd L(\lambda)=0$ if and only if $\lambda$ is dominant.
         \item[(2)] $\gkd L(\lambda)=36$ if and only if  $\lambda$ is antidominant.  
        \item[(3)] $\gkd L(\lambda)=11$ if and only if $w_{\lambda}\in\mathcal{C}w_0$.
       \item[(4)] $\gkd L(\lambda)=35$ if and only if $w_{\lambda}\in \mathcal{C}$.
        \item[(5)] $\gkd L(\lambda)=16$ if and only if $$w_{\lambda}\sim_{LR}w_{\{\alpha_2,\alpha_3,\alpha_4,\alpha_5,\alpha_6\}}=[1,2,3,1,2,3,4,3,1,2,3,4,5,4,3,1,2,3,4,5].$$
        %=s_3s_4s_2s_5s_4s_3s_6s_5s_4s_2s_6s_5s_4s_3s_6s_5s_4s_6s_5s_6.$$
        \item[(6)] $\gkd L(\lambda)=21$ if and only if $$w_{\lambda}\sim_{LR}w_{\{\alpha_1,\alpha_3,\alpha_4,\alpha_5,\alpha_6\}}.$$
        \item[(7)] $\gkd L(\lambda)=23$ if and only if $$w_{\lambda}\sim_{LR}[3,1,4,3,2,5,4,3,1,5,4,3,2,5,4,3,5,4,5].
   % s_4s_2s_5s_4s_3s_6s_5s_4s_2s_6s_5s_4s_3s_6s_5s_4s_6s_5s_6
    $$
        \item[(8)] $\gkd L(\lambda)=24$ if and only if $$w_{\lambda}\sim_{LR}w_{\{\alpha_2,\alpha_3,\alpha_4,\alpha_5\}}=[1,2,3,1,2,3,4,3,1,2,3,4].$$
%s_2s_4s_3s_5s_4s_2s_5s_4s_3s_5s_4s_5.$$
        \item[(9)] $\gkd L(\lambda)=25$ if and only if $$w_{\lambda}\sim_{LR}w_{\{\alpha_1,\alpha_2,\alpha_4,\alpha_5,\alpha_6\}}.$$
        \item[(10)] $\gkd L(\lambda)=26$ if and only if $$w_{\lambda}\sim_{LR}w_{\{\alpha_3,\alpha_4,\alpha_5,\alpha_6\}}.$$
        \item[(11)] $\gkd L(\lambda)=29$ if and only if $$w_{\lambda}\sim_{LR}w_{\{\alpha_1,\alpha_4,\alpha_5,\alpha_6\}}.$$
        \item[(12)] $\gkd L(\lambda)=30$ if and only if $$w_{\lambda}\sim_{LR}w_{\{\alpha_4,\alpha_5,\alpha_6\}}\text{~and~} V(\Ann (L(\lam)))=\overline{\mathcal{O}}({A}_4),$$ or $$w_{\lambda}\sim_{LR}w_{\{\alpha_1,\alpha_3,\alpha_5,\alpha_6\}} \text{~and~} V(\Ann (L(\lam)))=\overline{\mathcal{O}}({D}_4).$$
        \item[(13)] $\gkd L(\lambda)=31$ if and only if $$w_{\lambda}\sim_{LR}w_{\{\alpha_1,\alpha_2,\alpha_5,\alpha_6\}}.$$
        \item[(14)] $\gkd L(\lambda)=32$ if and only if $$w_{\lambda}\sim_{LR}w_{\{\alpha_1,\alpha_5,\alpha_6\}} \text{~and~} V(\Ann (L(\lam)))=\overline{\mathcal{O}}({D}_5(a_1)).$$
        \item[(15)] $\gkd L(\lambda)=33$ if and only if $w_{\lambda}\sim_{LR}w_{\{\alpha_5,\alpha_6\}}$.
        \item[(16)] $\gkd L(\lambda)=34$ if and only if $w_{\lambda}\sim_{LR}w_{\{\alpha_1,\alpha_6\}}$.
    \end{itemize}
\end{Prop}








In the above proposition, we give the expression of $w_I$ when $\Phi_I $ contains one component of type $D$. When the two-sided cell does not contain any $w_I$, we use PyCox to find a representative element in that two-sided cell with  a given $\aff$ value. When there are two special nilpotent orbits with the same dimension, we use Remark \ref{represent} to find  a
representative element for each corresponding two-sided cell.









% For the cases of nonintegral weights, we need the following lemma.

% \begin{Lem}\label{a-value-d}
 
%     The  $\aff$-function of the Weyl group of type $ A_4$  takes the following values:
%     $$0,1,2,3,4,6,10.$$    
  

%  The  $\aff$-function of the Weyl group of type $ A_5$  takes the following values:
%     $$0,1,2,3,4,6,7,10,15.$$

%   The  $\aff$-function of the Weyl group of type $ D_4$  takes the following values:
%     $$0,1,2,3,6,7,12.$$

%      The  $\aff$-function of the Weyl group of type $ D_5$  takes the following values:
%     $$0,1,2,3,4,5,6,7,10,12,13,20.$$

%     The  $\aff$-function of the Weyl group of type $ D_6$  takes the following values:
%     $$0,1,2,3,4,5,6,7,8,9,10,12,13,15,16,20,21,30.$$

%     The  $\aff$-function of the Weyl group of type $ D_8$  takes the following values:
%     \begin{align*}
% &0,1,2,3,4,5,6,7,8,9,10,11,12,13,14,15,16,17,\\
% &18,20,21,22,23,24,28,29,30,31,34,42,43,56.
%     \end{align*}

%      The  $\aff$-function of the Weyl group of type $ A_7$  takes the following values:
%     $$0,1,2,3,4,5,6,7,8,9,10,11,12,13,15,16,21,28.$$

%  The  $\aff$-function of the Weyl group of type $ E_6$  takes the following values:
% $$0,1,2,3,4,5,6,7,10,11,12,13,15,20,25,36.$$

%     The  $\aff$-function of the Weyl group of type $ E_7$  takes the following values:
% $$0,1,2,3,4,5,6,7,8,9,10,11,12,13,15,16,20,21,22,25,30,36,37,46,63.$$
% \end{Lem}

By similar arguments as in Proposition \ref{f4-nonint}, we have the following result for nonintegral highest weight modules of $E_6$.
%     \item $ \gkd L(\lam) =21$ if and only if 
%     \item $ \gkd L(\lam) =20$ if and only if 
% \end{enumerate} 
\begin{Prop}\label{e6-non}
Let $L(\lam)$ be a simple nonintegral highest weight module of $E_6$  such that $\Phi_{[\lam]}^\vee$ is pseudo-maximal.
	Then one of 
    the following holds:
    \begin{itemize}
        
	\item[(1)] $\Phi_{[\lam]}\simeq D_5$. Then 
 $ \gkd L(\lam) $ can achieve the following values:
$$16, 23, 24, 26, [29, 36].$$


\item[(2)] $\Phi_{[\lam]}\simeq A_5\times A_1$. Then
$ \gkd L(\lam) $ can achieve the following values:
$$20, 21, 25, 26, [28, 36].$$




            \item[(3)] $\Phi_{[\lam]}\simeq A_2\times A_2\times A_2$. Then
$ \gkd L(\lam) $ can achieve the following values:
$$27,  [29, 36].$$







\end{itemize}
\end{Prop}

% \begin{proof}




% Recall that we have the highest root $\beta=\frac{1}{2}(\ep_1+\ep_2+\ep_3+\ep_4+\ep_5-\ep_6-\ep_7+\ep_8)$ of $\Phi$.

% $(1).$ In the case of $\mathfrak{l}_1$, there is  one   maximal subsystem $ \Psi_1 $ of $ \Phi $ with $ \Psi_1^+$ spanned by  $\Delta_{1}=\{ \alpha_2,\alpha_3,\alpha_4,\alpha_5,\alpha_6\}$.
% So $\Psi_{1}\simeq D_5$.






% If $\Psi_{1}\subset \Phi_{[\lam]}$, we have $\lam_i\pm \lam_j\in \mathbb{Z}$ for $5\geq i>j\geq 1$. We can see that $|\Phi^+_{[\lam]}|=20$. From Lemma \ref{a-value},  one can see that $ \gkd L(\lam) $ can achieve the following values:
% $$16,23,24,26,29,30,31,32,33,34,35,36.$$
% % We have
% % $ \gkd L(\lam) = 16 $ if and only if 
% %  $\lam$ is $\Phi^+_{[\lam]}$-dominant integral.

% % From  \cite[Table D.4]{GP}, \cite[p.429]{Ca85} and \cite[p.129]{CM}, we can get the conclusions about $V(\Ann (L(\lam)))$ when $\gkd L(\lam)=30$ or $32$.




% $(2).$ In the case of $\mathfrak{l}_2$, there is  one   maximal subsystem $ \Psi_2 $ of $ \Phi $ with $ \Psi_2^+$ spanned by  $\Delta_{2}=\{ \alpha_1,\alpha_3,\alpha_4,\alpha_5,\alpha_6,\beta\}$.
% So $\Psi_{2}\simeq A_5\times A_1$.






% If $\Psi_{2}\subset \Phi_{[\lam]}$, we have $\lam_i- \lam_j\in \mathbb{Z}$ for $5\geq i>j\geq 1$ and $\lam_1+ \lam_2+\cdots+\lam_5-\lam_6-\lam_7+\lam_8\in 2\mathbb{Z}$. We can see that $|\Phi^+_{[\lam]}|=16$. From Lemma \ref{a-value},  one can see that $ \gkd L(\lam) $ can achieve the following values:
% $$20,21,25,26,28,29,30,31,32,33,34,35,36.$$
% % We have
% % $ \gkd L(\lam) = 20 $ if and only if 
% %  $\lam$ is $\Phi^+_{[\lam]}$-dominant integral.


% % From  \cite[p.100]{lusztig1984char}, \cite[p.429]{Ca85} and \cite[p.129]{CM}, we can get the conclusions about $V(\Ann (L(\lam)))$ when $\gkd L(\lam)=30$ or $32$.


% $(3).$ In the case of $\mathfrak{l}_3$, there is  one   maximal subsystem $ \Psi_3 $ of $ \Phi $ with $ \Psi_3^+$ spanned by  $\Delta_{3}=\{ -\beta, \alpha_2,\alpha_4,\alpha_5,\alpha_6,\alpha_1\}$.
% So $\Psi_{3}\simeq A_5\times A_1$.






% If $\Psi_{3}\subset \Phi_{[\lam]}$, we have $\lam_1+\lam_2\in \mathbb{Z}$, $\lam_i- \lam_j\in \mathbb{Z}$ for $5\geq i>j\geq 2$ and $\lam_1+ \lam_2+\cdots+\lam_5-\lam_6-\lam_7+\lam_8\in 2\mathbb{Z}$. Thus $-\lam_1+ \lam_2+\lam_3+\lam_4-\lam_5-\lam_6-\lam_7+\lam_8\in 2\mathbb{Z}$. We can choose the simple system of $\Phi_{[\lam]}$ by $\Delta_{[\lam]}=\{ \alpha_2,\alpha_4,\alpha_5,\alpha_6,\frac{1}{2}(-\ep_1+\ep_2+\ep_3+\ep_4-\ep_5-\ep_6-\ep_7+\ep_8),\alpha_1\}$.

% We can see that $|\Phi^+_{[\lam]}|=16$. From Lemma \ref{a-value},  one can see that $ \gkd L(\lam) $ can achieve the following values:
% $$20,21,25,26,28,29,30,31,32,33,34,35,36.$$
% % We have
% % $ \gkd L(\lam) = 20 $ if and only if 
% %  $\lam$ is $\Phi^+_{[\lam]}$-dominant integral.

% % From  \cite[p.100]{lusztig1984char}, \cite[p.429]{Ca85} and \cite[p.129]{CM}, we can get the conclusions about $V(\Ann (L(\lam)))$ when $\gkd L(\lam)=30$ or $32$.


% $(4).$ In the case $(a)$ of case $\mathfrak{l}_4$, there is  one   maximal subsystem $ \Psi_4 $ of $ \Phi $ with $ \Psi_4^+$ spanned by  $\Delta_{4}=\{ \alpha_3,\alpha_1,\ep_3+\ep_4,\alpha_6,\alpha_5,\alpha_2\}$.
% So $\Psi_{4}\simeq A_5\times A_1$.

% If $\Psi_{4}\subset \Phi_{[\lam]}$, we have $\lam_1\pm \lam_2\in \mathbb{Z},\lam_3+ \lam_4\in \mathbb{Z}$, $\lam_i- \lam_j\in \mathbb{Z}$ for $5\geq i>j\geq 3$ and $\lam_1-\lam_2-\cdots -\lam_7+\lam_8\in 2\mathbb{Z}$.  

% We can see that $|\Phi^+_{[\lam]}|=16$. From Lemma \ref{a-value},  one can see that $ \gkd L(\lam) $ can achieve the following values:
% $$20,21,25,26,28,29,30,31,32,33,34,35,36.$$
% % We have
% % $ \gkd L(\lam) = 20 $ if and only if 
% %  $\lam$ is $\Phi^+_{[\lam]}$-dominant integral.


% Case $(b)$ is similar to case $(a)$.
% % From  \cite[p.100]{lusztig1984char}, \cite[p.429]{Ca85} and \cite[p.129]{CM}, we can get the conclusions about $V(\Ann (L(\lam)))$ when $\gkd L(\lam)=30$ or $32$.

% $(5).$ In the case of $\mathfrak{l}_5$, there is  one   maximal subsystem $ \Psi_5 $ of $ \Phi $ with $ \Psi_5^+$ spanned by $\Delta_{5}=\{ \ep_4+\ep_5,\alpha_1,\alpha_3,\alpha_4,\alpha_2,\alpha_6\}$.
% So $\Psi_{5}\simeq A_5\times A_1$.

% If $\Psi_{5}\subset \Phi_{[\lam]}$, we have $\lam_1\pm \lam_2\in \mathbb{Z},\lam_4+ \lam_5\in \mathbb{Z}$, $\lam_3- \lam_2\in \mathbb{Z},\lam_5- \lam_4\in \mathbb{Z}$ and $\lam_1-\lam_2-\cdots -\lam_7+\lam_8\in 2\mathbb{Z}$.

% We can see that $|\Phi^+_{[\lam]}|=16$. From Lemma \ref{a-value},  one can see that $ \gkd L(\lam) $ can achieve the following values:
% $$20,21,25,26,28,29,30,31,32,33,34,35,36.$$
% % We have
% % $ \gkd L(\lam) = 20 $ if and only if 
% %  $\lam$ is $\Phi^+_{[\lam]}$-dominant integral.

% % From  \cite[p.100]{lusztig1984char}, \cite[p.429]{Ca85} and \cite[p.129]{CM}, we can get the conclusions about $V(\Ann (L(\lam)))$ when $\gkd L(\lam)=30$ or $32$.



% $(6).$ In the case of $\mathfrak{l}_6$, it is similar to case $(1)$.

% % For all cases, the Bala--Carter label of the nilpotent orbit in $V(\Ann (L(\lam)))$ is determined by  Alvis \cite{Alv}.


 
% \end{proof}


% \begin{Rem}\label{e6-gouzao} From now on we denote $\beta_1=\ep_1-\ep_2, \beta_2=\ep_2-\ep_3, \beta_3=\ep_3-\ep_4$,  $\beta_4=\ep_4-\ep_5$ and $\beta_5=\ep_4+\ep_5$. Then the simple root system of $D_5$ is $\Delta=\{\beta_1, \beta_2, \beta_3,\beta_4, \beta_5\}$.
%     In general, by Theorem \ref{transf} and Corollary \ref{transf-2} we have the following:
%     \begin{enumerate}
%         \item In the case of $\mathfrak{l}_1$, we define an isomorphism map $\phi: \mathbb{C}^8\rightarrow \mathbb{C}^8$ such that 
%  $\phi|_{\Phi_{[\lambda]}}: \Phi_{[\lambda]} \rightarrow D_5$ is an isomorphism and  $\phi(\alpha_6)=\beta_1, \phi(\alpha_5)=\beta_2, \phi(\alpha_4)=\beta_3, \phi(\alpha_3)=\beta_4,\phi(\alpha_2)=\beta_5,\phi(\ep_i)=\ep_i$ for $6\lest i\lest 8$. Thus for $\lam\in \mathfrak{h}_{E_6}^* $, it will be a weight in  $\mathfrak{h}_{\Phi_{[\lambda]}}^*\times \mathfrak{h}_{\Phi^{\perp}_{[\lambda]}}^*$ and $$\phi({\lam})=\lam'=(\lam_5,\lam_4,\lam_3,\lam_2,\lam_1,-\lam_8,-\lam_8,\lam_8).$$  So  $\bar{\lambda}'=(\lam_5,\lam_4,\lam_3,\lam_2,\lam_1)=F(\lam|_{\mathfrak{h}_{\Phi_{[\lambda]}}^*})$ will  be an integral weight of type $D_5$.
      
%         \item In the case of $\mathfrak{l}_2$, we define an isomorphism map $\phi: \mathbb{C}^8\rightarrow \mathbb{C}^8$ such that 
%  $\phi|_{\Phi_{[\lambda]}}: \Phi_{[\lambda]} \rightarrow A_5\times A_1$ is an isomorphism and $\phi(\alpha_1)=\beta_1, \phi(\alpha_3)=\beta_2, \phi(\alpha_4)=\beta_3, \phi(\alpha_5)=\beta_4,\phi(\alpha_6)=\ep_5-\ep_6,\phi(\beta)=\ep_7-\ep_8$, $\phi(-2\ep_6+\ep_7-\ep_8)=\ep_1+\cdots+\ep_6$, $\phi(\ep_7+\ep_8)=\ep_7+\ep_8$. 
%        Then we can extend it to an isomorphism and still denote it by $\phi: \mathfrak{h}_{\Phi_{[\lambda]}}^*\rightarrow \mathfrak{h}_{A_5\times A_1}^*$.   Thus \begin{align*}
%            &\bar{\lambda}'=\phi(\lam)=\lambda'\\
%            =&(\frac{-\lam_1-\lam_2-\lam_3-\lam_4-\lam_5+5\lam_8}{4},
%                \frac{-3\lam_1+\lam_2+\lam_3+\lam_4+\lam_5-\lam_8}{4},\\
%            & \frac{\lam_1-3\lam_2+\lam_3+\lam_4+\lam_5-\lam_8}{4},\frac{\lam_1+\lam_2-3\lam_3+\lam_4+\lam_5-\lam_8}{4},\\
%            &\frac{\lam_1+\lam_2+\lam_3-3\lam_4+\lam_5-\lam_8}{4},\frac{\lam_1+\lam_2+\lam_3+\lam_4-3\lam_5-\lam_8}{4},\\
%            &\frac{\lam_1+\lam_2+\lam_3+\lam_4+\lam_5+3\lam_8}{4},
%            -\frac{\lam_1+\lam_2+\lam_3+\lam_4+\lam_5+3\lam_8}{4}) 
%         \end{align*} will be a weight of type $A_5\times A_1$.


        
%        \item In the case of $\mathfrak{l}_3$, we define an isomorphism map $\phi: \mathbb{C}^8\rightarrow \mathbb{C}^8$ such that 
%  $\phi|_{\Phi_{[\lambda]}}: \Phi_{[\lambda]} \rightarrow A_5\times A_1$ is an isomorphism and $\phi(\alpha_2)=\beta_1, \phi(\alpha_4)=\beta_2, \phi(\alpha_5)=\beta_3,\phi(\alpha_6)=\beta_4,\phi(\frac{1}{2}(-\ep_1+\ep_2+\ep_3+\ep_4-\ep_5-\ep_6-\ep_7+\ep_8))=\ep_5-\ep_6, \phi(\alpha_1)=\ep_7-\ep_8,$ $\phi(-2\ep_6+\ep_7-\ep_8)=\ep_1+\cdots+\ep_6$, $\phi(\ep_7+\ep_8)=\ep_7+\ep_8$. Then we can extend it to an isomorphism and still denote it by $\phi: \mathfrak{h}_{\Phi_{[\lambda]}}^*\rightarrow \mathfrak{h}_{A_5\times A_1}^*$.  Thus \begin{align*}
%            &\bar{\lambda}'=\phi(\lam)=\lambda'\\
%            =&(\frac{3\lam_1+\lam_2+\lam_3+\lam_4+\lam_5+\lam_8}{4},
%                \frac{-\lam_1-3\lam_2+\lam_3+\lam_4+\lam_5+\lam_8}{4},\\
%            & \frac{-\lam_1+\lam_2-3\lam_3+\lam_4+\lam_5+\lam_8}{4},\frac{-\lam_1+\lam_2+\lam_3-3\lam_4+\lam_5+\lam_8}{4},\\
%            &\frac{-\lam_1+\lam_2+\lam_3+\lam_4-3\lam_5+\lam_8}{4},\frac{\lam_1-\lam_2-\lam_3-\lam_4-\lam_5-5\lam_8}{4},\\
%            &\frac{\lam_1-\lam_2-\lam_3-\lam_4-\lam_5+3\lam_8}{4},          -\frac{\lam_1-\lam_2-\lam_3-\lam_4-\lam_5+3\lam_8}{4})
%         \end{align*} will be a weight of type $A_5\times A_1$.

        
%        \item In the case of $\mathfrak{l}_4$, we have two cases:
%        \begin{enumerate}
%            \item   we define an isomorphism map $\phi: \mathbb{C}^8\rightarrow \mathbb{C}^8$ such that 
%  $\phi|_{\Phi_{[\lambda]}}: \Phi_{[\lambda]} \rightarrow A_5\times A_1$ is an isomorphism and $\phi(\alpha_3)=\beta_1, \phi(\alpha_1)=\beta_2, \phi(\ep_3+\ep_4)=\beta_3, \phi(\alpha_6)=\beta_4,\phi(\alpha_5)=\ep_5-\ep_6,\phi(\alpha_2)=\ep_7-\ep_8$, $\phi(-2\ep_6+\ep_7-\ep_8)=\ep_1+\cdots+\ep_6$, $\phi(\ep_7+\ep_8)=\ep_7+\ep_8$. Then we can extend it to an isomorphism and still denote it by $\phi: \mathfrak{h}_{\Phi_{[\lambda]}}^*\rightarrow \mathfrak{h}_{A_5\times A_1}^*$. 
%        Thus \begin{align*}
%            &\bar{\lambda}'=\phi(\lam)=\lambda'\\
%            =&(\frac{-\lam_1+\lam_2+2\lam_8}{2},
%                \frac{\lam_1-\lam_2+2\lam_8}{2},\frac{\lam_3+\lam_4+\lam_5-\lam_8}{2},\\
%            & \frac{-\lam_3-\lam_4+\lam_5-\lam_8}{2},\frac{-\lam_3+\lam_4-\lam_5-\lam_8}{2},\frac{\lam_3-\lam_4-\lam_5-\lam_8}{2},\\
%            &\frac{\lam_1+\lam_2}{2},
%            -\frac{\lam_1+\lam_2}{2})
%         \end{align*} will be a weight of type $A_5\times A_1$.
%         \item   we define an isomorphism map $\phi: \mathbb{C}^8\rightarrow \mathbb{C}^8$ such that
% $\phi|_{\Phi_{[\lambda]}}: \Phi_{[\lambda]} \rightarrow A_2\times A_2\times A_2$ is an isomorphism and $\phi(\alpha_1)=\beta_1, \phi(\alpha_3)=\beta_2, \phi(\frac{1}{2}(-\ep_1-\ep_2+\ep_3+\ep_4+\ep_5-\ep_6-\ep_7+\ep_8))=\beta_4, \phi(\alpha_2)=\ep_5-\ep_6,\phi(\alpha_5)=\ep_7-\ep_8,\phi(\alpha_6)=\frac{1}{2}(\ep_1+\ep_2+\ep_3+\ep_4+\ep_5+\ep_6-\ep_7+\ep_8)$, $\phi(\ep_1+\ep_2+\ep_3-\ep_4-\ep_5-\ep_6)=\ep_1+\ep_2+\ep_3-\ep_4-\ep_5-\ep_6$, $\phi(\ep_7+\ep_8)=\ep_7+\ep_8$. We define another map $\psi: A_2\rightarrow A_2$ such that $\psi(\ep_7-\ep_8)=\ep_7-\ep_8$ and $\psi(\frac{1}{2}(\ep_1+\ep_2+\ep_3+\ep_4+\ep_5+\ep_6-\ep_7+\ep_8))=\ep_8-\ep_9$.



% Then we can extend $\phi$ to an isomorphism and still denote it by $\phi: \mathfrak{h}_{\Phi_{[\lambda]}}^*\rightarrow \mathfrak{h}_{A_2\times A_2\times A_2}^*$. 
      
%       % Thus $$\psi\circ \phi(\lam)=(\frac{2\lam_1+\lam_2+\lam_3}{2},\frac{-\lam_2-\lam_3+2\lam_4}{2},\frac{-\lam_2-\lam_3-2\lam_4}{2}, \frac{-2\lam_1+\lam_2+\lam_3}{2})$$ will be an integral weight of type $A_3$ 
       
       
       
       
%        Thus \begin{align*}
%            &\psi\circ \phi(\lam)=\\
%            =&(\frac{-\lam_3-\lam_4-\lam_5+3\lam_8}{3},
%                \frac{-3\lam_1+3\lam_2+\lam_3+\lam_4+\lam_5-3\lam_8}{6},\\
%                &\frac{3\lam_1-3\lam_2+\lam_3+\lam_4+\lam_5-3\lam_8}{6},\\
%            & \frac{\lam_3+\lam_4+\lam_5+3\lam_8}{3},\frac{3\lam_1+3\lam_2-\lam_3-\lam_4-\lam_5-3\lam_8}{6},\\
%            &-\frac{3\lam_1+3\lam_2+\lam_3+\lam_4+\lam_5+3\lam_8}{6},\\
%            &\frac{-2\lam_3+\lam_4+\lam_5}{3},
%            \frac{\lam_3-2\lam_4+\lam_5}{3},\frac{\lam_3+\lam_4-2\lam_5}{3})
%         \end{align*} will be a weight of type $A_2\times A_2\times A_2$.






        
%        \end{enumerate}
       
       
       
       
     

%   \item In the case of $\mathfrak{l}_5$, we define an isomorphism map $\phi: \mathbb{C}^8\rightarrow \mathbb{C}^8$ such that 
%  $\phi|_{\Phi_{[\lambda]}}: \Phi_{[\lambda]} \rightarrow A_5\times A_1$ is an isomorphism and $\phi(\ep_4+\ep_5)=\beta_1, \phi(\alpha_1)=\beta_2, \phi(\alpha_3)=\beta_3,\phi(\alpha_4)=\beta_4,\phi(\alpha_2)=\ep_5-\ep_6, \phi(\alpha_6)=\ep_7-\ep_8$, $\phi(-2\ep_6+\ep_7-\ep_8)=\ep_1+\cdots+\ep_6$, $\phi(\ep_7+\ep_8)=\ep_7+\ep_8$. Then we can extend it to an isomorphism and still denote it by $\phi: \mathfrak{h}_{\Phi_{[\lambda]}}^*\rightarrow \mathfrak{h}_{A_5\times A_1}^*$. Thus \begin{align*}
%            &\bar{\lambda}'=\phi(\lam)=\lambda'\\
%            =&(\frac{\lam_4+\lam_5+2\lam_8}{2},
%                \frac{-\lam_4-\lam_5+2\lam_8}{2}, \frac{-\lam_1+\lam_2+\lam_3-\lam_8}{2},\\
%            &\frac{\lam_1-\lam_2+\lam_3-\lam_8}{2},\frac{\lam_1+\lam_2-\lam_3-\lam_8}{2},-\frac{\lam_1+\lam_2+\lam_3+\lam_8}{2},\\
%            &\frac{\lam_5-\lam_4}{2},
%            \frac{\lam_4-\lam_5}{2})
%         \end{align*} will be a weight of type $A_5\times A_1$.


%  \item In the case of $\mathfrak{l}_6$, we define an isomorphism map $\phi: \mathbb{C}^8\rightarrow \mathbb{C}^8$ such that 
%  $\phi|_{\Phi_{[\lambda]}}: \Phi_{[\lambda]} \rightarrow D_5$ is an isomorphism and $\phi(\alpha_1)=\beta_1, \phi(\alpha_3)=\beta_2, \phi(\alpha_4)=\beta_3, \phi(\alpha_5)=\beta_4,\phi(\alpha_2)=\beta_5,\phi(\ep_5-\ep_6)=\ep_6-\ep_7, \phi(\ep_6-\ep_7)=\ep_7-\ep_8, \phi(\frac{1}{2}(\ep_5-\ep_6+\ep_7+\ep_8))=\ep_8$.  Thus for $\lam\in \mathfrak{h}_{E_6}^* $, it will be a weight in  $\mathfrak{h}_{\Phi_{[\lambda]}}^*\times \mathfrak{h}_{\Phi^{\perp}_{[\lambda]}}^*$ and
% \begin{align*}
%            &\phi({\lam})={\lambda}'=\\
%            =&(\frac{-\lam_5+3\lam_8}{2},
%                \frac{-\lam_1+\lam_2+\lam_3+\lam_4}{2}, \\
%            &\frac{\lam_1-\lam_2+\lam_3+\lam_4}{2},\frac{\lam_1+\lam_2-\lam_3+\lam_4},\\
%            &\frac{\lam_1+\lam_2+\lam_3-\lam_4}{2},\frac{\lam_5+\lam_8}{2},0,0)
%         \end{align*} will be a weight of type $D_5\times B_3$, and \begin{align*}
%           \lam'= &(\frac{-\lam_5+3\lam_8}{2},
%                \frac{-\lam_1+\lam_2+\lam_3+\lam_4}{2}, 
%            \frac{\lam_1-\lam_2+\lam_3+\lam_4}{2},\\
%            &\frac{\lam_1+\lam_2-\lam_3+\lam_4},
%            \frac{\lam_1+\lam_2+\lam_3-\lam_4}{2})
%         \end{align*} will be an integral weight of type $D_5$.
% % \begin{align*}
%  %           &\phi(\lam)=\lambda'\\
%  %           =&(2\lam_8,\frac{-\lam_1+\lam_2+\lam_3+\lam_4+\lam_5+\lam_8}{2},
%  %               \frac{\lam_1-\lam_2+\lam_3+\lam_4+\lam_5+\lam_8}{2}, \\
%  %           &\frac{\lam_1+\lam_2-\lam_3+\lam_4+\lam_5+\lam_8}{2},\frac{\lam_1+\lam_2+\lam_3-\lam_4-\lam_5-\lam_8}{2})
%  %        \end{align*} 
 
%       %  will be a weight of type $D_5$.
        
%     \end{enumerate}
% Then the value of $\aff(w_{\lam})=\aff(w_{\lam'})=\aff(w_{\lam'})$ can be computed by the algorithm in \cite{BXX}.    
% \end{Rem}


\begin{ex}
	Let $\mathfrak{g}=E_6$ and $L(\lam)$ be the highest weight module of $\mathfrak{g}$ with  $\lambda=(1,2,1,4,4.5,0.5, 0.5, -0.5)$. It is easy to verify that $\Phi_{[\lambda]}$ is a subsystem with simple roots $\{ \alpha_1,\alpha_2,\alpha_3,  \alpha_4, \alpha_5\}$. Therefore $\Phi_{[\lambda]}\simeq D_5$. Suppose the simple root system of $D_5$ is $\Delta=\{\beta_1=\ep_1-\ep_2, \beta_2=\ep_2-\ep_3, \beta_3=\ep_3-\ep_4,\beta_4=\ep_4-\ep_5,\beta_5=\ep_4+\ep_5\}$. We define an isomorphism map $\phi: \Phi_{[\lambda]} \rightarrow D_5$  such that  $\phi(\alpha_1)=\beta_1, \phi(\alpha_3)=\beta_2, \phi(\alpha_4)=\beta_3, \phi(\alpha_5)=\beta_4,\phi(\alpha_2)=\beta_5$.

Denote ${\lam}'=\phi(\lam|_{\mathfrak{h}_{\Phi_{[\lambda]}}})$. 
Thus 
% $$\lambda=-6\omega_1+3\omega_2+\omega_3-\omega_4+3\omega_5+\lambda|_{\Phi_{[\lambda]}\perp},$$
%     where $\{\omega_i\mid 1\leq i\leq 5\}$ are the fundamental weights for the simple system of $\Phi_{[\lambda]}$. 
% Thus we get  
 \begin{align*}
     \lam'|_{D_5}=& -6\phi(\omega_1)+\phi(\omega_2)-\phi(\omega_3)+3\phi(\omega_4)+3\phi(\omega_5)=(-3,3,2,3,0),
 \end{align*}
which is an integral weight of type $D_5$.
 So $\aff(w_{\lambda})=\aff(w_{\lam'})=7$ by the RS algorithm in \cite{BXX} and Proposition \ref{integral}.
 By Proposition \ref{pr:main1},  we have $\gkd L(\lambda)=|\Phi^+|-\aff(w_{\lambda})=36-7=29$. 
% By \cite[p.129]{CM}, we have $V(\Ann (L(\lam)))=\overline{\mathcal{O}}({D}_4(a_1))$.
\end{ex}


\subsection{Type \texorpdfstring{$E_7$}{}}
For type $E_7$, we have \begin{align*}
  \Phi^+=&
\{\ep_i \pm \ep_j \mid 6\geq i > j\geq 1\}\cup\{\ep_8-\ep_7\} \\
&\cup
\Big\{ \frac{1}{2}(\sum_{i=1}^6 (-1)^{n(i)}\ep_i -\ep_7+\ep_8)\mid \sum_{i=1}^6n(i) \text{ odd} \Big\}  
\end{align*}
and $ \hs=\{ (\lam_1,\lam_2,\cdots,\lam_8)\in\mathbb{C}^8\mid \lam_7=-\lam_8\} $. The highest root is $$\beta=2\alpha_1+2\alpha_2+3\alpha_3+4\alpha_4+3\alpha_5+2\alpha_6+\alpha_7=-\ep_7+\ep_8.$$

\begin{figure}[htpb]
	\centering 
	
	
	\begin{tikzpicture}[scale=1.5,baseline=0]
	\draw (-0.75,0.05) node[above=1pt]{} circle [radius=0.05];
	\draw (-0.7,0.05)--++(0.63,0); 
 
 
 \draw (0,0.05) node[above=1pt]{$ 1 $} circle [radius=0.05];
	\draw (0.05,0.05)--++(0.65,0); 
	\draw (0.75,0.05) node[above=1pt]{$ 3 $} circle[radius=0.05];  
	\draw (0.8,0.05)--++(0.65,0);
	\draw  (1.45,0) ++(0.05,0.05) node[above=1pt]{$ 4$} circle[radius=0.05];
	\draw (1.55,0.05)--++(0.65,0);
	\draw  (2.2,0) ++(0.05,0.05) node[above=1pt]{$ 5$} circle[radius=0.05];
	\draw (2.3,0.05)--++(0.65,0);
	\draw  (2.95,0) ++(0.05,0.05) node[above=1pt]{$ 6$} circle[radius=0.05];
	\draw (3.05,0.05)--++(0.65,0);
	\draw  (3.7,0) ++(0.05,0.05) node[above=1pt]{$ 7$} circle[radius=0.05];
	\draw (1.5,0)--++(0,-0.5);
	\draw  (1.45,-0.6) ++(0.05,0.05) node[right=1pt]{$ 2$} circle[radius=0.05];
	\end{tikzpicture}
	
	\caption{Extended Dynkin diagram of $ E_7 $ }
	\label{E7-dy}
\end{figure}






Fix $ \lam\in\hs $. We have \[
\bil{\lam}{(\alpha_1)}=\frac{1}{2}(\lam_1- \lam_2-\cdots-\lam_7+\lam_8),
\]
\[
\bil{\lam}{(\ep_i\pm \ep_j)}=\lam_i\pm \lam_j.
\]
Then $ \lam=(\lam_1,\lam_2,\cdots,\lam_8) $ is an integral weight if and only if $\lam_1- \lam_2-\cdots-\lam_7+\lam_8\in 2\mathbb{Z}$,  $ \lam_1-\lam_2\in\mathbb{Z}, \lam_2-\lam_3\in\mathbb{Z},\lam_3-\lam_4\in\mathbb{Z},\lam_4-\lam_5\in\mathbb{Z}, \lam_5-\lam_6\in\mathbb{Z}$ and $2\lam_i\in \mathbb{Z} $ for $1\lest i\lest 6$.

%From now on,  we will write $w_I\simeq w_{\Phi}$ if $\Phi_I$ is isomorphic to a standard root system $\Phi$.

By similar arguments as in Proposition \ref{f4}, we have the following result.

\begin{Prop}\label{E7}
Let $L(\lam)$ be a simple integral highest weight module of $E_7$. Then
    the following holds:
    \begin{itemize}
        \item[(1)] $\gkd L(\lambda)=0$ if and only if $\lambda$ is dominant integral.
         \item[(2)] $\gkd L(\lambda)=63$ if and only if  $\lambda$ is antidominant integral.  
        \item[(3)] $\gkd L(\lambda)=17$ if and only if $w_{\lambda}\in\mathcal{C}w_0$.
       \item[(4)] $\gkd L(\lambda)=62$ if and only if $w_{\lambda}\in \mathcal{C}$.
     \item[(5)] $\gkd L(\lambda)=26$ if and only if \begin{align*}
          w_{\lambda}\sim_{LR}&[0, 1, 2, 0, 3, 1, 2, 0, 3, 2, 4, 3, 1, 2, 0, 3, 2, 4, 3, 1, 5, 4, 3,\\
          & 1, 2, 0, 3, 2, 4, 3, 1, 5, 4, 3, 2, 0, 6, 5, 4, 3, 1, 2, 3, 4, 5, 6].
      \end{align*}
      \item[(6)] $\gkd L(\lambda)=27$ if and only if \begin{align*}
            w_{\lambda}&\sim_{LR}w_{\{\alpha_1,\alpha_2,\alpha_3,\alpha_4,\alpha_5,\alpha_6\}}=w_{E_6}.   
      \end{align*}
      \item[(7)] $\gkd L(\lambda)=33$ if and only if \begin{align*}
           w_{\lambda}&\sim_{LR}w_{\{\alpha_2,\alpha_3,\alpha_4,\alpha_5,\alpha_6,\alpha_7\}}\\ 
           =&[1, 2, 3, 1, 2, 3, 4, 3, 1, 2, 3, 4, 5, 4, 3, 1, 2, 3, 4, 5, 6, 5, 4, 3, 1, 2, 3, 4, 5, 6].
          % &=s_2s_3s_4s_2s_3s_4s_5s_4s_2s_3s_4s_5s_6s_5s_4s_2s_3s_4s_5s_6s_7s_6s_5s_4s_2s_3s_4s_5s_6s_7.
      \end{align*}
\item[(8)] $\gkd L(\lambda)=38$ if and only if 
\begin{align*}
    w_{\lambda}\sim_{LR}&[1, 2, 3, 1, 2, 0, 3, 4, 3, 1, 2, 0,3, 2, 4,3,\\
    &  5, 4, 3, 1, 2, 0, 3, 2, 4, 3, 1, 5, 4, 3, 2].
\end{align*}
% $$w_{\lambda}\sim_{LR}s_2s_3s_4s_2s_3s_1s_4s_5s_4s_2s_3s_1s_4s_3s_5s_4s_6s_5s_4s_2s_3s_1s_4s_3s_5s_4s_2s_6s_5s_4s_3.$$

\item[(9)] $\gkd L(\lambda)=41$ if and only if 
\begin{align*}
    w_{\lambda}\sim_{LR}&[0, 2, 0, 3, 2, 0, 4, 3, 2, 0, 5, 4, 3, 2, 0, 6, \\
    & 5, 4, 3, 1, 2, 0, 3, 2, 4, 3, 1, 5, 4, 3, 2].
\end{align*}
% $$w_{\lambda}\sim_{LR}s_1s_3s_1s_4s_3s_1s_5s_4s_3s_1s_6s_5s_4s_3s_1s_7s_6s_5s_4s_2s_3s_1s_4s_3s_5s_4s_2s_6s_5s_4s_3.$$


\item[(10)] $\gkd L(\lambda)=42$ if and only if
% one of the followings holds:
% \begin{enumerate}
%     \item \begin{align*}w_{\lambda}&\sim_{LR}w_{\{\alpha_1,\alpha_2,\alpha_3,\alpha_4,\alpha_5,\alpha_7\}}=[4,3,1,2,3,4,0,2,3,1,0,2,3,4,0,2,3,0,2,0,6]\\
% %=s_5s_4s_2s_3s_4s_5s_1s_3s_4s_2s_1s_3s_4s_5s_1s_3s_4s_1s_3s_1s_7\\
% &\text{and~} V(\Ann (L(\lam)))=\overline{\mathcal{O}}(2A_2),  \end{align*}
% \item \begin{align*}
% w_{\lambda}&\sim_{LR}w_{\{\alpha_1,\alpha_3,\alpha_4,\alpha_5,\alpha_6,\alpha_7\}} \text{and~} V(\Ann (L(\lam)))=\overline{\mathcal{O}}(A_2+3A_1),
%         \end{align*}
%         \item \begin{align*}
% w_{\lambda}&\sim_{LR}=[1, 3, 1, 2, 4, 3, 1, 2, 3, 5, 4, 3, 1, 2, 3, 4, 6, 5, 4, 3, 1, 2, 3, 4, 5]\\
% % s_2s_4s_2s_3s_5s_4s_2s_3s_4s_6s_5s_4s_2s_3s_4s_5s_7s_6s_5s_4s_2s_3s_4s_5s_6\\
% & \text{and~} V(\Ann (L(\lam)))=\overline{\mathcal{O}}(A_3).
% \end{align*}
% \end{enumerate}
\begin{align*}w_{\lambda}&\sim_{LR}w_{\{\alpha_1,\alpha_2,\alpha_3,\alpha_4,\alpha_5,\alpha_7\}}=[4,1,3,4,1,3,2,3,4,1,3,2,0,2,3,4,1,3,2,0,6]\\
%=s_5s_4s_2s_3s_4s_5s_1s_3s_4s_2s_1s_3s_4s_5s_1s_3s_4s_1s_3s_1s_7\\
&\text{~and~} V(\Ann (L(\lam)))=\overline{\mathcal{O}}(2A_2),  \end{align*}
% $$w_{\lambda}\sim_{LR}s_5s_4s_2s_3s_4s_5s_1s_3s_4s_2s_1s_3s_4s_5s_1s_3s_4s_1s_3s_1s_7=w_{\{\alpha_1,\alpha_2,\alpha_3,\alpha_4,\alpha_5,\alpha_7\}}$$
        or
        \begin{align*}
w_{\lambda}&\sim_{LR}w_{\{\alpha_1,\alpha_3,\alpha_4,\alpha_5,\alpha_6,\alpha_7\}} \text{~and~} V(\Ann (L(\lam)))=\overline{\mathcal{O}}(A_2+3A_1),
        \end{align*}
% $$w_{\lambda}\sim_{LR}s_7s_6s_7s_5s_6s_7s_4s_5s_6s_7s_3s_4s_5s_6s_7s_1s_3s_4s_5s_6s_7=w_{\{\alpha_1,\alpha_3,\alpha_4,\alpha_5,\alpha_6,\alpha_7\}}=w_{A_6}$$
or 
\begin{align*}
w_{\lambda}&\sim_{LR} [1, 3, 1, 2, 4, 3, 1, 2, 3, 5, 4, 3, 1, 2, 3, 4, 6, 5, 4, 3, 1, 2, 3, 4, 5]\\
% s_2s_4s_2s_3s_5s_4s_2s_3s_4s_6s_5s_4s_2s_3s_4s_5s_7s_6s_5s_4s_2s_3s_4s_5s_6\\
& \text{~and~} V(\Ann (L(\lam)))=\overline{\mathcal{O}}(A_3).
\end{align*}
% $$w_{\lambda}\sim_{LR}s_2s_4s_2s_3s_5s_4s_2s_3s_4s_6s_5s_4s_2s_3s_4s_5s_7s_6s_5s_4s_2s_3s_4s_5s_6.$$
       
        \item[(11)] $\gkd L(\lambda)=43$ if and only if $$w_{\lambda}\sim_{LR}w_{\{\alpha_2,\alpha_3,\alpha_4,\alpha_5,\alpha_6\}}=[1,2,3,1,2,3,4,3,1,2,3,4,5,4,3,1,2,3,4,5].$$
%=s_3s_4s_2s_5s_4s_3s_6s_5s_4s_2s_6s_5s_4s_3s_6s_5s_4s_6s_5s_6.$$
    \item[(12)] $\gkd L(\lambda)=47$ if and only if \begin{align*}
          w_{\lambda}&\sim_{LR}w_{\{\alpha_1,\alpha_2,\alpha_4,\alpha_5,\alpha_6,\alpha_7\}} \text{~and~} V(\Ann (L(\lam)))=\overline{\mathcal{O}}({D}_4(a_1)).
      \end{align*}
        \item[(13)] $\gkd L(\lambda)=48$ if and only if \begin{align*}w_{\lambda}&\sim_{LR}w_{\{\alpha_1,\alpha_3,\alpha_4,\alpha_5,\alpha_6\}}\text{~and~} V(\Ann (L(\lam)))=\overline{\mathcal{O}}({D}_4(a_1)+A_1),
\end{align*}
        or
        \begin{align*}w_{\lambda}&\sim_{LR}w_{\{\alpha_2,\alpha_4,\alpha_5,\alpha_6,\alpha_7\}}\text{~and~} V(\Ann (L(\lam)))=\overline{\mathcal{O}}({D}_4).
            \end{align*}
% $$w_{\lambda}\sim_{LR}s_7s_6s_7s_5s_6s_7s_4s_5s_6s_7s_2s_4s_5s_6s_7=w_{\{\alpha_2,\alpha_4,\alpha_5,\alpha_6,\alpha_7\}},\text{and~} V(\Ann (L(\lam)))=\overline{\mathcal{O}}({D}_4).$$
   
      \item[(14)] $\gkd L(\lambda)=49$ if and only if 
%$$w_{\lambda}\sim_{LR}s_1s_2s_3s_1s_4s_2s_3s_1s_4s_3s_5s_4s_2s_3s_1s_4s_5s_6s_5s_4s_7s_6s_5s_4s_2s_3s_1s_4.$$
     $$w_{\lambda}\sim_{LR} [0, 1, 2, 0, 3, 1, 2, 0, 3, 2, 4, 3, 1, 2, 0, 3, 4, 5, 4, 3, 6, 5, 4, 3, 1, 2, 0, 3].$$
         \item[(15)] $\gkd L(\lambda)=50$ if and only if \begin{align*}
w_{\lambda}&\sim_{LR}w_{\{\alpha_2,\alpha_3,\alpha_4,\alpha_5,\alpha_7\}}=[1,2,3,1,2,3,4,3,1,2,3,4,6]\\
%=s_2s_4s_3s_5s_4s_2s_5s_4s_3s_5s_4s_5s_7\\
&\text{~and~} V(\Ann (L(\lam)))=\overline{\mathcal{O}}({A}_4),
         \end{align*}
        or
\begin{align*}
w_{\lambda}&\sim_{LR}w_{\{\alpha_1,\alpha_2,\alpha_3,\alpha_4,\alpha_6,\alpha_7\}} \text{~and~} V(\Ann (L(\lam)))=\overline{\mathcal{O}}(A_3+A_2+A_1).
         \end{align*}
        
% $$w_{\lambda}\sim_{LR}s_2s_4s_2s_3s_4s_2s_1s_3s_4s_2s_7s_6s_7=w_{\{\alpha_1,\alpha_2,\alpha_3,\alpha_4,\alpha_6,\alpha_7\}},\text{and~} V(\Ann (L(\lam)))=\overline{\mathcal{O}}(A_3+A_2+A_1).$$
        \item[(16)] $\gkd L(\lambda)=51$ if and only if 
        \begin{align*}
w_{\lambda}&\sim_{LR}w_{\{\alpha_2,\alpha_3,\alpha_4,\alpha_5\}}=[1,2,3,1,2,3,4,3,1,2,3,4]\\
%=s_2s_4s_3s_5s_4s_2s_5s_4s_3s_5s_4s_5 \\
&\text{~and~} V(\Ann (L(\lam)))=\overline{\mathcal{O}}(({A}_5)'').
        \end{align*}
        
        % $$w_{\lambda}\sim_{LR}s_2s_4s_3s_5s_4s_2s_5s_4s_3s_5s_4s_5=w_{\{\alpha_2,\alpha_3,\alpha_4,\alpha_5\}},\text{and~} V(\Ann (L(\lam)))=\overline{\mathcal{O}}(({A}_5)'').$$
        
        \item[(17)] $\gkd L(\lambda)=52$ if and only if $$w_{\lambda}\sim_{LR}w_{\{\alpha_1,\alpha_2,\alpha_4,\alpha_5,\alpha_6\}}.$$
        \item[(18)] $\gkd L(\lambda)=53$ if and only if \begin{align*}
w_{\lambda}&\sim_{LR}w_{\{\alpha_3,\alpha_4,\alpha_5,\alpha_6\}}\text{~and~} V(\Ann (L(\lam)))=\overline{\mathcal{O}}({D}_5(a_1)),
\end{align*}
% $$w_{\lambda}\sim_{LR}s_6s_5s_4s_3s_6s_5s_4s_6s_5s_6=w_{\{\alpha_3,\alpha_4,\alpha_5,\alpha_6\}},\text{and~} V(\Ann (L(\lam)))=\overline{\mathcal{O}}({D}_5(a_1))$$
        or
\begin{align*}
w_{\lambda}&\sim_{LR}w_{\{\alpha_1,\alpha_3,\alpha_2,\alpha_5,\alpha_6,\alpha_7\}} \text{~and~} V(\Ann (L(\lam)))=\overline{\mathcal{O}}(A_4+A_2).
\end{align*}

% $$w_{\lambda}\sim_{LR}s_3s_1s_3s_2s_7s_6s_7s_5s_6s_7=w_{\{\alpha_1,\alpha_3,\alpha_2,\alpha_5,\alpha_6,\alpha_7\}},\text{and~} V(\Ann (L(\lam)))=\overline{\mathcal{O}}(A_4+A_2).$$
  \item[(19)] $\gkd L(\lambda)=54$ if and only if \begin{align*}
    w_{\lambda}&\sim_{LR}w_{\{\alpha_1,\alpha_3,\alpha_5,\alpha_6,\alpha_7\}}\text{~and~} V(\Ann (L(\lam)))=\overline{\mathcal{O}}({D}_5(a_1)+A_1).
  \end{align*}
  % $$w_{\lambda}\sim_{LR}s_3s_1s_3s_7s_6s_7s_5s_6s_7=w_{\{\alpha_1,\alpha_3,\alpha_5,\alpha_6,\alpha_7\}},\text{and~} V(\Ann (L(\lam)))=\overline{\mathcal{O}}({D}_5(a_1)+A_1).$$
  
    \item[(20)] $\gkd L(\lambda)=55$ if and only if $$w_{\lambda}\sim_{LR}w_{\{\alpha_1,\alpha_2,\alpha_5,\alpha_6,\alpha_7\}}\text{~and~} V(\Ann (L(\lam)))=\overline{\mathcal{O}}({E}_6(a_3)).$$ 
      
        \item[(21)] $\gkd L(\lambda)=56$ if and only if $$w_{\lambda}\sim_{LR}w_{\{\alpha_1,\alpha_4,\alpha_5,\alpha_6\}}\text{~and~} V(\Ann (L(\lam)))=\overline{\mathcal{O}}({E}_7(a_5)),$$
        or
$$w_{\lambda}\sim_{LR}w_{\{\alpha_7,\alpha_2,\alpha_4,\alpha_5\}}\text{~and~} V(\Ann (L(\lam)))=\overline{\mathcal{O}}({D}_5).$$
% $$w_{\lambda}\sim_{LR}s_2s_3s_4s_3s_5s_4s_2s_3s_1s_4s_3s_6s_5s_4s_2s_3s_1s_7s_6s_5s_4s_2s_3s_1s_4s_5s_6s_7.$$
        \item[(22)] $\gkd L(\lambda)=57$ if and only if $$w_{\lambda}\sim_{LR}w_{\{\alpha_4,\alpha_5,\alpha_6\}} \text{~and~} V(\Ann (L(\lam)))=\overline{\mathcal{O}}({D}_6(a_1)),$$ or $$w_{\lambda}\sim_{LR}w_{\{\alpha_1,\alpha_3,\alpha_5,\alpha_6\}} \text{~and~} V(\Ann (L(\lam)))=\overline{\mathcal{O}}({D}_5+A_1),$$ or 
$$w_{\lambda}\sim_{LR}w_{\{\alpha_1,\alpha_3,\alpha_2,\alpha_5,\alpha_7\}} \text{~and~} V(\Ann (L(\lam)))=\overline{\mathcal{O}}(A_6).$$
        \item[(23)] $\gkd L(\lambda)=58$ if and only if $$w_{\lambda}\sim_{LR}w_{\{\alpha_1,\alpha_2,\alpha_5,\alpha_6\}}.$$
        \item[(24)] $\gkd L(\lambda)=59$ if and only if $$w_{\lambda}\sim_{LR}w_{\{\alpha_1,\alpha_5,\alpha_6\}} \text{~and~} V(\Ann (L(\lam)))=\overline{\mathcal{O}}({E}_6(a_1)).$$
        \item[(25)] $\gkd L(\lambda)=60$ if and only if $$w_{\lambda}\sim_{LR}w_{\{\alpha_5,\alpha_6\}} \text{~and~} V(\Ann (L(\lam)))=\overline{\mathcal{O}}({E}_7(a_3)),$$ or 
$$w_{\lambda}\sim_{LR}w_{\{\alpha_2,\alpha_5,\alpha_7\}}  \text{~and~} V(\Ann (L(\lam)))=\overline{\mathcal{O}}({E}_6).$$
% $$w_{\lambda}\sim_{LR}s_2s_5s_4s_3s_1s_7s_6s_5s_4s_2s_3s_4s_5s_6s_7.$$
        \item[(26)] $\gkd L(\lambda)=61$ if and only if $w_{\lambda}\sim_{LR}w_{\{\alpha_1,\alpha_6\}}$.
    \end{itemize}
\end{Prop}





In the above proposition, we still use PyCox to find a representative element in that two-sided cell with  a given $\aff$ value.

For non-integral highest weight modules of $E_7$, we have the following.

\begin{Prop}
Let $L(\lam)$ be a simple nonintegral highest weight module of $E_7$ such that $\Phi_{[\lam]}^\vee$ is pseudo-maximal. Then one of the following holds:
    \begin{itemize}
	\item[(1)] $\Phi_{[\lam]}\simeq D_6\times A_1$. Then 
 $ \gkd L(\lam) $ can achieve the following values:
$$32,33,41,42,43,[46,63].$$

\item[(2)]  $\Phi_{[\lam]}\simeq A_7$. Then
 $ \gkd L(\lam) $ can achieve the following values:
$$35,42,47,48,[50,63].$$


  
   
    \item[(3)] $\Phi_{[\lam]}\simeq A_5\times A_2$. 
 Then 
 $ \gkd L(\lam) $ can achieve the following values:
$$45,47,48,50,[52,63].$$ 



    \item[(4)] $\Phi_{[\lam]}\simeq A_3\times A_3\times A_1$. Then 
 $ \gkd L(\lam) $ can achieve the following values:
$$50,51,[53,63].$$








\item[(5)]  $\Phi_{[\lam]}\simeq E_6$. Then 
 $ \gkd L(\lam) $ can achieve the following values:
$$27,38,43,48,50,51,52,53,[56,63].$$



\end{itemize}
\end{Prop}

\begin{proof}
The argument is similar to that in the proof of Proposition \ref{e6-non}.


\end{proof}


% \begin{Rem}\label{e7-gou} From now on we denote $\beta_1=\ep_1-\ep_2, \beta_2=\ep_2-\ep_3, \beta_3=\ep_3-\ep_4$,  $\beta_4=\ep_4-\ep_5$, $\beta_5=\ep_5-\ep_6$  and $\beta_6=\ep_5+\ep_6$. Then the simple root system of $D_6$ is $\Delta=\{\beta_1, \beta_2, \beta_3,\beta_4, \beta_5,\beta_6\}$. In general, by Theorem \ref{transf} and Corollary \ref{transf-2} we have the following:
%     \begin{enumerate}
%         \item In the case of $\mathfrak{l}_1$, we define an isomorphism map $\phi: \mathbb{C}^8\rightarrow \mathbb{C}^8$ such that 
%  $\phi|_{\Phi_{[\lambda]}}: \Phi_{[\lambda]} \rightarrow D_6\times A_1$ is an isomorphism and  $\phi(\alpha_7)=\beta_1, \phi(\alpha_6)=\beta_2, \phi(\alpha_5)=\beta_3, \phi(\alpha_4)=\beta_4,\phi(\alpha_3)=\beta_5,\phi(\alpha_2)=\beta_6,\phi(\ep_7-\ep_8)=\ep_7-\ep_8, \phi(\ep_7+\ep_8)=\ep_7+\ep_8$. Then we can extend it to an isomorphism and still denote it by $\phi: \mathfrak{h}_{\Phi_{[\lambda]}}^*\rightarrow \mathfrak{h}_{D_6\times A_1}^*$.
%  Thus for $\lam\in \mathfrak{h}_{E_7}^* $, it will be an integral weight  in  $\mathfrak{h}_{\Phi_{[\lambda]}}^*$ and $$\lam'=\phi({\lam})=\lam'=(\lam_6,\lam_5,\lam_4,\lam_3,\lam_2,\lam_1,-\lam_8,\lam_8)$$  will  be an integral weight of type $D_6\times A_1$.
      
%         \item In the case of $\mathfrak{l}_2$, we define an isomorphism map $\phi: \mathbb{C}^8\rightarrow \mathbb{C}^8$ such that 
%  $\phi|_{\Phi_{[\lambda]}}: \Phi_{[\lambda]} \rightarrow A_7$ is an isomorphism and $\phi(\alpha_1)=\beta_1, \phi(\alpha_3)=\beta_2, \phi(\alpha_4)=\beta_3, \phi(\alpha_5)=\beta_4,\phi(\alpha_6)=\beta_5,$
%         $\phi(\alpha_7)=\ep_6-\ep_7,\phi(\gamma_2)=\ep_7-\ep_8,\phi(2\ep_7+2\ep_8)=\ep_1+\cdots+\ep_8$. 
%        Then we can extend it to an isomorphism and still denote it by $\phi: \mathfrak{h}_{\Phi_{[\lambda]}}^*\rightarrow \mathfrak{h}_{A_7}^*$.   Thus \begin{align*}
%            &\lam'=\phi(\lam)=\lambda'\\
%            =&(\frac{-\lam_1-\lam_2-\lam_3-\lam_4-\lam_5-\lam_6+4\lam_8}{4},
%                \frac{-3\lam_1+\lam_2+\lam_3+\lam_4+\lam_5+\lam_6}{4},\\
%            & \frac{\lam_1-3\lam_2+\lam_3+\lam_4+\lam_5+\lam_6}{4},\frac{\lam_1+\lam_2-3\lam_3+\lam_4+\lam_5+\lam_6}{4},\\
%            &\frac{\lam_1+\lam_2+\lam_3-3\lam_4+\lam_5+\lam_6}{4},\frac{\lam_1+\lam_2+\lam_3+\lam_4-3\lam_5+\lam_6}{4},\\
%            &\frac{\lam_1+\lam_2+\lam_3+\lam_4+\lam_5-3\lam_6}{4},
%            \frac{-\lam_1-\lam_2-\lam_3-\lam_4-\lam_5-\lam_6-4\lam_8}{4}) 
%         \end{align*} will be a weight of type $A_7$.
 
%  \item In the case of $\mathfrak{l}_3$, we have two cases:
%  \begin{enumerate}
%      \item we define an isomorphism map $\phi: \mathbb{C}^8\rightarrow \mathbb{C}^8$ such that 
%  $\phi|_{\Phi_{[\lambda]}}: \Phi_{[\lambda]} \rightarrow D_6\times A_1$ is an isomorphism and  $\phi(\alpha_2)=\beta_1, \phi(\alpha_4)=\beta_2, \phi(\alpha_5)=\beta_3, \phi(\alpha_6)=\beta_4,\phi(\alpha_7)=\beta_5,\phi(\gamma_3)=\beta_6,\phi(\alpha_1)=\ep_7-\ep_8, \phi(2\ep_7+2\ep_8)=\ep_1+\cdots+\ep_8$. Then we can extend it to an isomorphism and still denote it by $\phi: \mathfrak{h}_{\Phi_{[\lambda]}}^*\rightarrow \mathfrak{h}_{D_6\times A_1}^*$.
%  Thus for $\lam\in \mathfrak{h}_{E_7}^* $, it will be an integral weight  in  $\mathfrak{h}_{\Phi_{[\lambda]}}^*$ and \begin{align*}
%            &\lam'=\phi(\lam)=\lambda'\\
%            =&(\frac{3\lam_1+\lam_2+\cdots+\lam_6+2\lam_8}{4},
%                \frac{-\lam_1-3\lam_2+\lam_3+\cdots+\lam_6+2\lam_8}{4},\\
%            & \frac{-\lam_1+\lam_2-3\lam_3+\lam_4+\lam_5+\lam_6+2\lam_8}{4},\\
%            &\frac{-\lam_1+\lam_2+\lam_3-3\lam_4+\lam_5+\lam_6+2\lam_8}{4},\\
%            &-\lam_5-\lam_8,\frac{-\lam_1+\lam_2+\cdots+\lam_5-3\lam_6+2\lam_8}{4}\\
%            &\frac{\lam_1-\lam_2-\cdots-\lam_6+2\lam_8}{4},
%            -\frac{\lam_1-\lam_2-\cdots-\lam_6+2\lam_8}{4}) 
%         \end{align*} will  be an integral weight of type $D_6\times A_1$.

%         \item we define an isomorphism map $\phi: \mathbb{C}^8\rightarrow \mathbb{C}^8$ such that 
%  $\phi|_{\Phi_{[\lambda]}}: \Phi_{[\lambda]} \rightarrow A_5\times A_2$ is an isomorphism and  $\phi(\alpha_2)=\beta_1, \phi(\alpha_4)=\beta_2, \phi(\alpha_5)=\beta_3, \phi(\alpha_6)=\beta_4,\phi(\alpha_7)=\beta_5,\phi(\gamma'_3)=\ep_7-\ep_8,\phi(\alpha_1)=\frac{1}{2}(\ep_1+\ep_2+\ep_3+\ep_4+\ep_5+\ep_6-\ep_7+\ep_8), \phi(\ep_7+\ep_8)=\ep_7+\ep_8$. We define another map $\psi: A_2\rightarrow A_2$ such that $\psi(\ep_7-\ep_8)=\ep_7-\ep_8$ and $\psi(\frac{1}{2}(\ep_1+\ep_2+\ep_3+\ep_4+\ep_5+\ep_6-\ep_7+\ep_8))=\ep_8-\ep_9$.
 
 
 
%  Then we can extend it to an isomorphism and still denote it by $\phi: \mathfrak{h}_{\Phi_{[\lambda]}}^*\rightarrow \mathfrak{h}_{A_5\times A_2}^*$.
%  Thus for $\lam\in \mathfrak{h}_{E_7}^* $, it will be an integral weight  in  $\mathfrak{h}_{\Phi_{[\lambda]}}^*$ and \begin{align*}
%            &\psi\circ \phi(\lam)=\lambda'\\
%            =&(\frac{5\lam_1+\lam_2+\cdots+\lam_6}{6},
%                \frac{-\lam_1-5\lam_2+\lam_3+\cdots+\lam_6}{6},\\
%            & \frac{-\lam_1+\lam_2-5\lam_3+\lam_4+\lam_5+\lam_6}{6},
%            \frac{-\lam_1+\lam_2+\lam_3-5\lam_4+\lam_5+\lam_6}{6},\\
%            &\frac{-\lam_1+\lam_2+\lam_3+\lam_4-5\lam_5+\lam_6}{6},\frac{-\lam_1+\lam_2+\lam_3+\lam_4+\lam_5-5\lam_6}{6},\\
%            &\frac{-\lam_1+\lam_2+\cdots+\lam_6+6\lam_8}{6},\frac{\lam_1-\lam_2-\cdots-\lam_6}{3},\frac{-\lam_1+\lam_2+\cdots+\lam_6-6\lam_8}{6}) 
%         \end{align*} will  be an integral weight of type $A_5\times A_2$.
%  \end{enumerate}
 
 
 
 
 

%  \item In the case of $\mathfrak{l}_4$, we have two cases:
%  \begin{enumerate}
%      \item  we define an isomorphism map $\phi: \mathbb{C}^8\rightarrow \mathbb{C}^8$ such that 
%  $\phi|_{\Phi_{[\lambda]}}: \Phi_{[\lambda]} \rightarrow D_6\times A_1$ is an isomorphism and  $\phi(\alpha_3)=\beta_1, \phi(\alpha_1)=\beta_2, \phi(\ep_3+\ep_4)=\beta_3, \phi(\alpha_6)=\beta_4,\phi(\alpha_5)=\beta_5,\phi(\alpha_7)=\beta_6,\phi(\alpha_2)=\ep_7-\ep_8, \phi(2\ep_7+2\ep_8)=\ep_1+\cdots+\ep_8$. Then we can extend it to an isomorphism and still denote it by $\phi: \mathfrak{h}_{\Phi_{[\lambda]}}^*\rightarrow \mathfrak{h}_{D_6\times A_1}^*$.
%  Thus for $\lam\in \mathfrak{h}_{E_7}^* $, it will be an integral weight  in  $\mathfrak{h}_{\Phi_{[\lambda]}}^*$ and \begin{align*}
%            &\lam'=\phi(\lam)=\lambda'\\
%            =&(\frac{-\lam_1+\lam_2+2\lam_8}{2},
%                \frac{\lam_1-\lam_2+2\lam_8}{2}, \frac{\lam_3+\lam_4+\lam_5+\lam_6}{2},\frac{-\lam_3-\lam_4+\lam_5+\lam_6}{2},\\
%            &\frac{-\lam_3+\lam_4-\lam_5-\lam_6-2\lam_8}{2},\frac{\lam_3-\lam_4-\lam_5+\lam_6}{2},\frac{\lam_1+\lam_2}{2},
%            \frac{-\lam_1-\lam_2}{2}) 
%         \end{align*}  will  be an integral weight of type $D_6\times A_1$.
   
%      \item  we define an isomorphism map $\phi: \mathbb{C}^8\rightarrow \mathbb{C}^8$ such that 
%  $\phi|_{\Phi_{[\lambda]}}: \Phi_{[\lambda]} \rightarrow A_3\times A_3\times A_1$ is an isomorphism and  $\phi(\alpha_5)=\beta_1, \phi(\alpha_6)=\beta_2, \phi(\alpha_7)=\beta_3, \phi(\alpha_1)=\beta_5,\phi(\alpha_3)=\ep_6-\ep_7,\phi(\gamma_4)=\ep_7-\ep_8,\phi(\alpha_2)=\ep_1+\ep_2, \phi(\ep_7+\ep_8)=\frac{1}{\sqrt{2}}(\ep_5+\cdots+\ep_8)$. Then we can extend it to an isomorphism and still denote it by $\phi: \mathfrak{h}_{\Phi_{[\lambda]}}^*\rightarrow \mathfrak{h}_{A_3\times A_3\times A_1}^*$.
%  Thus for $\lam\in \mathfrak{h}_{E_7}^* $, it will be an integral weight  in  $\mathfrak{h}_{\Phi_{[\lambda]}}^*$ and \begin{align*}
%            &\lam'=\phi(\lam)|_{\mathfrak{h}_{A_3\times A_3}^*}=\\
%            =&(\frac{-3\lam_3+\lam_4+\lam_5+\lam_6}{4},
%                \frac{\lam_3-3\lam_4+\lam_5+\lam_6}{4}, \frac{\lam_3+\lam_4-3\lam_5+\lam_6}{4},\frac{\lam_3+\lam_4+\lam_5-3\lam_6}{4},\\
%            &\frac{-\lam_3-\lam_4-\lam_5-\lam_6+4\lam_8}{4},\frac{-2\lam_1+2\lam_2+\lam_3+\lam_4+\lam_5+\lam_6}{4},\\
%            &\frac{2\lam_1-2\lam_2+\lam_3+\lam_4+\lam_5+\lam_6}{4},-\frac{\lam_3+\lam_4+\lam_5+\lam_6+4\lam_8}{4}) 
%         \end{align*}  will  be an integral weight of type $A_3\times A_3$ and 
%     $$\phi(\lam)|_{\mathfrak{h}_{A_1}^*}=\frac{\lam_1+\lam_2}{2}(\ep_1+\ep_2)$$ will  be an integral weight of type $A_1$.
%  \end{enumerate}
 
 



%         \item In the case of $\mathfrak{l}_5$, we define an isomorphism map $\phi: \mathbb{C}^8\rightarrow \mathbb{C}^8$ such that 
%  $\phi|_{\Phi_{[\lambda]}}: \Phi_{[\lambda]} \rightarrow  A_7$ is an isomorphism and $\phi(\alpha_2)=\beta_1, \phi(\alpha_4)=\beta_2, \phi(\alpha_3)=\beta_3, \phi(\alpha_1)=\beta_4,\phi(\ep_4+\ep_5)=\beta_5,$
%         $\phi(\alpha_7)=\ep_6-\ep_7,\phi(\alpha_6)=\ep_7-\ep_8, \phi(2\ep_7+2\ep_8)=\ep_1+\cdots+\ep_8$. 
%        Then we can extend it to an isomorphism and still denote it by $\phi: \mathfrak{h}_{\Phi_{[\lambda]}}^*\rightarrow \mathfrak{h}_{A_7}^*$.   Thus \begin{align*}
%            &\lam'=\phi(\lam)=\lambda'\\
%            =&(\frac{\lam_1+\lam_2+\lam_3+\lam_8}{2},
%                \frac{-\lam_1-\lam_2+\lam_3+\lam_8}{2},\frac{-\lam_1+\lam_2-\lam_3+\lam_8}{2},\\
%            & \frac{\lam_1-\lam_2-\lam_3+\lam_8}{2},\frac{\lam_4+\lam_5+\lam_6-\lam_8}{2},\frac{-\lam_4-\lam_5+\lam_6-\lam_8}{2},\\
%            &\frac{-\lam_4+\lam_5-\lam_6-\lam_8}{2},
%            \frac{\lam_4-\lam_5-\lam_6-\lam_8}{2}) 
%         \end{align*} will be a weight of type $A_7$.

% \item In the case of $\mathfrak{l}_6$, we define an isomorphism map $\phi: \mathbb{C}^8\rightarrow \mathbb{C}^8$ such that 
%  $\phi|_{\Phi_{[\lambda]}}: \Phi_{[\lambda]} \rightarrow D_6\times A_1$ is an isomorphism and
%    $\phi(\ep_5+\ep_6)=\beta_1, \phi(\alpha_1)=\beta_2, \phi(\alpha_3)=\beta_3, \phi(\alpha_4)=\beta_4,\phi(\alpha_2)=\beta_5,\phi(\alpha_5)=\beta_6,\phi(\alpha_7)=\ep_7-\ep_8, \phi(2\ep_7+2\ep_8)=\ep_1+\cdots+\ep_8$. Then we can extend it to an isomorphism and still denote it by $\phi: \mathfrak{h}_{\Phi_{[\lambda]}}^*\rightarrow \mathfrak{h}_{D_6\times A_1}^*$.
%  Thus for $\lam\in \mathfrak{h}_{E_7}^* $, it will be an integral weight  in  $\mathfrak{h}_{\Phi_{[\lambda]}}^*$ and \begin{align*}
%            &\lam'=\phi(\lam)=\lambda'\\
%            =&(\frac{\lam_5+\lam_6+2\lam_8}{2},
%                \frac{-\lam_5-\lam_6+2\lam_8}{2},\frac{-\lam_1+\lam_2+\lam_3+\lam_4}{2},\\
%            & \frac{\lam_1-\lam_2+\lam_3+\lam_4}{2},\frac{\lam_1+\lam_2-\lam_3-\lam_4-2\lam_8}{2},\frac{-\lam_1-\lam_2-\lam_3+\lam_4}{2},\\
%            &\frac{-\lam_5+\lam_6}{2},
%            \frac{\lam_5-\lam_6}{2}) 
%         \end{align*}  will  be an integral weight of type $D_6\times A_1$.

%  \item In the case of $\mathfrak{l}_7$, we define an isomorphism map $\phi: \mathbb{C}^8\rightarrow \mathbb{C}^8$ such that 
%  $\phi|_{\Phi_{[\lambda]}}: \Phi_{[\lambda]} \rightarrow E_6$ is an isomorphism and  $\phi(\alpha_1)=\alpha_1, \phi(\alpha_2)=\alpha_2, \phi(\alpha_3)=\alpha_3, \phi(\alpha_4)=\alpha_4,\phi(\alpha_5)=\alpha_5,\phi(\alpha_6)=\alpha_6,\phi(\ep_6-\ep_7)=\ep_6-\ep_7,\phi(\ep_7+\ep_8)=\ep_7+\ep_8$. 
%  Thus for $\lam\in \mathfrak{h}_{E_7}^* $, it will be a weight in  $\mathfrak{h}_{\Phi_{[\lambda]}}^*\times \mathfrak{h}_{\Phi^{\perp}_{[\lambda]}}^*$ and also an integral weight in $\mathfrak{h}_{\Phi_{[\lambda]}}^*$. 
% So  $\bar{\lambda}'=F(\lam|_{\mathfrak{h}_{\Phi_{[\lambda]}}^*})=\lam|_{\mathfrak{h}_{\Phi_{[\lambda]}}^*}=(\lam_1,\cdots,\lam_5,\frac{\lam_6-2\lam_8}{3},\frac{\lam_6-2\lam_8}{3},\frac{-\lam_6+2\lam_8}{3})$ will  be an integral weight of type $E_6$.



 
%     \end{enumerate}
% Then the value of $\aff(w_{\lam})=\aff(w_{\lam'})$ (or $\aff(w_{\lam'})$) can be computed by the algorithm in \cite{BXX} (or Theorem \ref{E6}).    
% \end{Rem}


\begin{ex}
	Let $\mathfrak{g}=E_7$ and $L(\lam)$ be the highest weight module of $\mathfrak{g}$ with  $\lambda=(\frac{1}{4},\frac{1}{4},\frac{1}{4},\frac{1}{4},\frac{1}{4},-\frac{3}{4},-1,1)$. It is easy to verify that $\Phi_{[\lambda]}$ is a subsystem with simple roots $\{ \alpha_1,\alpha_3,\alpha_4,\alpha_5,\alpha_6,\alpha_7,\gamma\}$ where $\gamma=\frac{1}{2}(\ep_1+\ep_2+\ep_3+\ep_4+\ep_5-\ep_6-\ep_7+\ep_8)$. Therefore $\Phi_{[\lambda]}\simeq A_7$.  We define the  isomorphism $\phi:  \Phi_{[\lambda]} \rightarrow A_7$ such that  $\phi(\alpha_1)=\ep_1-\ep_2, \phi(\alpha_3)=\ep_2-\ep_3, \phi(\alpha_4)=\ep_3-\ep_4, \phi(\alpha_5)=\ep_4-\ep_5,\phi(\alpha_6)=\ep_5-\ep_6,$  $\phi(\alpha_7)=\ep_6-\ep_7,\phi(\gamma)=\ep_7-\ep_8$.
Thus by the algorithm in \S \ref{nonint}, we can get
\begin{align*}
            {\lam}'={\lam}'|_{A_7}=\phi(\omega_1)-\phi(\omega_6)+2\phi(\omega_7)
           =\left(\frac{7}{8},-\frac{1}{8},-\frac{1}{8},-\frac{1}{8},-\frac{1}{8},-\frac{1}{8},\frac{7}{8},-\frac{9}{8}\right),
        \end{align*} 
      which  is an integral weight of type $A_7$.
 It gives us a partition ${\bf p}=[6,1,1]$. So $\aff(w_{\lambda})=\aff(w_{\lambda'})=3$ by the RS algorithm in \cite{BXX} and Proposition \ref{integral}. 
 By Proposition \ref{pr:main1},  we have $\gkd L(\lambda)=|\Phi^+|-\aff(w_{\lambda})=63-3=60$.
% From Alvis \cite{Alv}, we have  $V(\Ann (L(\lam)))=\overline{\mathcal{O}}(E_7(a_3))$.
\end{ex}

\subsection{Type \texorpdfstring{$E_8$}{}}
For type $E_8$, we have \begin{align*}
  \Phi^+=&
\{\ep_i \pm \ep_j \mid 8\geq i > j\geq 1\} \cup
\Big\{\frac{1}{2}(\sum_{i=1}^7 (-1)^{n(i)}\ep_i+\ep_8)\mid \sum_{i=1}^7n(i) \text{ even}\Big\}
\end{align*}
and $ \hs=\mathbb{C}^8$. The highest root is $$\beta=2\alpha_1+3\alpha_2+4\alpha_3+6\alpha_4+5\alpha_5+4\alpha_6+3\alpha_7+2\alpha_8=\ep_7+\ep_8.$$

\begin{figure}[htpb]
	\centering 
	

	
	\begin{tikzpicture}[scale=1.5,baseline=0]
	
 
 \draw (0,0.05) node[above=1pt]{$ 1 $} circle [radius=0.05];
	\draw (0.05,0.05)--++(0.65,0); 
	\draw (0.75,0.05) node[above=1pt]{$ 3 $} circle[radius=0.05];  
	\draw (0.8,0.05)--++(0.65,0);
	\draw  (1.45,0) ++(0.05,0.05) node[above=1pt]{$ 4$} circle[radius=0.05];
	\draw (1.55,0.05)--++(0.65,0);
	\draw  (2.2,0) ++(0.05,0.05) node[above=1pt]{$ 5$} circle[radius=0.05];
	\draw (2.3,0.05)--++(0.65,0);
	\draw  (2.95,0) ++(0.05,0.05) node[above=1pt]{$ 6$} circle[radius=0.05];
	\draw (3.05,0.05)--++(0.65,0);
	\draw  (3.7,0) ++(0.05,0.05) node[above=1pt]{$ 7$} circle[radius=0.05];
	\draw (3.8,0.05)--++(0.65,0);
	\draw  (4.45,0) ++(0.05,0.05) node[above=1pt]{$ 8$} circle[radius=0.05];
	\draw (1.5,0)--++(0,-0.5);
	\draw  (1.45,-0.6) ++(0.05,0.05) node[right=1pt]{$ 2$} circle[radius=0.05];

\draw (5.2,0.05) node[above=1pt]{} circle [radius=0.05];
	\draw (4.53,0.05)--++(0.65,0); 

 
	\end{tikzpicture}
	\caption{Extended Dynkin Diagram of $ E_8 $ }
	\label{E8-dy}
\end{figure}






Fix $ \lam=\sum_{i=1}^8 \lambda_i e_i \in\hs $. We have \[
\bil{\lam}{(\alpha_1)}=\frac{1}{2}(\lam_1- \lam_2-\cdots-\lam_7+\lam_8),
\]
\[
\bil{\lam}{(\ep_i\pm \ep_j)}=\lam_i\pm \lam_j.
\]
Then $ \lam=(\lam_1,\lam_2,\cdots,\lam_8) $ is an integral weight if and only if $\lam_1- \lam_2-\cdots-\lam_7+\lam_8\in 2\mathbb{Z}$,  $ \lam_1-\lam_2\in\mathbb{Z}, \lam_2-\lam_3\in\mathbb{Z},\lam_3-\lam_4\in\mathbb{Z},\lam_4-\lam_5\in\mathbb{Z}, \lam_5-\lam_6\in\mathbb{Z},\lam_6-\lam_7\in\mathbb{Z}$ and $2\lam_i\in \mathbb{Z} $ for $1\lest i\lest 7$.





  By similar arguments as in Proposition \ref{f4}, we have the following result.

\begin{Prop}\label{E8}
Let $L(\lam)$ be a simple integral highest weight module of $E_8$. Then
    the following holds:
    \begin{itemize}
        \item[(1)] $\gkd L(\lambda)=0$ if and only if $\lambda$ is dominant integral.
         \item[(2)] $\gkd L(\lambda)=120$ if and only if  $\lambda$ is antidominant integral.  
        \item[(3)] $\gkd L(\lambda)=29$ if and only if $w_{\lambda}\in\mathcal{C}w_0$.
       \item[(4)] $\gkd L(\lambda)=119$ if and only if $w_{\lambda}\in \mathcal{C}$.

\item[(5)] $\gkd L(\lambda)=46$ if and only if \begin{align*}
          w_{\lambda}\sim_{LR} &  
      [0, 1, 2, 0, 3, 1, 2, 0, 3, 2, 4, 3, 1, 2, 0, 3, 2, 4, 3, 1,5,\\
      & 4, 3, 1, 2, 0, 3, 2, 4, 3, 1, 5, 4, 3, 2, 0, 6, 5, 4, 3, 1, 2,\\
      & 0, 3, 2, 4, 3, 1, 5, 4, 3, 2, 6, 5, 4, 3, 1, 7, 6, 5, 4, 3, 1, \\
      &2, 0, 3, 2, 4, 3, 1, 5, 4, 3, 2, 0, 6, 5, 4, 3, 1, 2, 3, 4, 5, \\
      &7, 6, 5, 4, 3, 1, 2, 0, 3, 2, 4, 3, 1, 5, 4, 3, 2, 0].
\end{align*}







    
 \item[(6)] $\gkd L(\lambda)=57$ if and only if \begin{align*}
          w_{\lambda}\sim_{LR}&w_{\{\alpha_1,\alpha_2,\alpha_3,\alpha_4,\alpha_5,\alpha_6,\alpha_7\}}=w_{E_7}.
          % \\    =&[6,5,4,3,2,1,3,4,5,6,0,2,3,4,5,1,3,4,2,3,0,2,1,3,4,5,6,\\
          % &0,2,3,4,5,1,3,4,2,3,0,2,1,3,4,5,1,3,4,2,3,0,2,1,3,4,2,\\
          % &3,0,2,1,3,0,2,1,0].
%           &s_7s_6s_5s_4s_3s_2s_4s_5s_6s_7s_1s_3s_4s_5s_6s_2s_4s_5s_3s_4s_1s_3s_2s_4s_5s_6s_7s_1s_3s_4s_5s_6\\
% &s_2s_4s_5s_3s_4s_1s_3s_2s_4s_5s_6s_2s_4s_5s_3s_4s_1s_3s_2s_4s_5s_3s_4s_1s_3s_2s_4s_1s_3s_2s_1.
      \end{align*}    

\item[(7)] $\gkd L(\lambda)=68$ if and only if \begin{align*}
          w_{\lambda}\sim_{LR} &  
      [1, 2, 3, 1, 2, 0, 3, 4, 3, 1, 2, 0, 3, 2, 4, 3, 5, 4, 3, 1, 2, 0,  \\& 3, 2, 4, 3, 1, 5, 4, 3, 2, 6, 5, 4, 3, 1, 2, 0, 3, 2, 4, 3, 1, 5,  \\& 4, 3, 2, 0, 6, 5, 4, 3, 1, 2, 3, 4, 5, 6, 7, 6, 5, 4, 3, 1, 2, 0,  \\&3, 2, 4, 3, 1, 5, 4, 3, 2, 0, 6, 5, 4, 3, 1, 2, 3, 4, 5, 6].
\end{align*}
\item[(8)] $\gkd L(\lambda)=73$ if and only if \begin{align*}
          w_{\lambda}\sim_{LR}   &
      [0, 1, 2, 0, 3, 1, 2, 0, 3, 2, 4, 3, 1, 2, 0, 3, 2, 4, 3, 1, 5, 4, 3, \\& 1,  2, 0, 3, 2, 4, 3, 1, 5, 4, 3, 2, 0, 6, 5, 4, 3, 1, 2, 0, 3, 2, 4,  \\&3, 1, 5, 4, 3, 2, 0, 6, 5, 4, 3, 1, 7, 6, 5, 4, 3, 1, 2, 0, 3, 2, 4, \\&3, 1, 5,  4, 3, 2, 6, 5, 4, 7, 6, 5].
\end{align*}

\item[(9)] $\gkd L(\lambda)=74$ if and only if \begin{align*}
          w_{\lambda}\sim_{LR}   &
      [0, 1, 2, 0, 3, 1, 2, 0, 3, 2, 4, 3, 1, 2, 0, 3, 2, 4, 3, 1, \\& 5, 4, 3, 1, 2, 0, 3, 2, 4, 3, 1, 5, 4, 3, 2, 0, 6, 5, 4, 3,   \\&1, 2, 0, 3, 2, 4, 3, 1, 5, 4, 3, 2, 6, 5, 4, 3, 1].
\end{align*}

\item[(10)] $\gkd L(\lambda)=78$ if and only if \begin{align*}
          w_{\lambda}&\sim_{LR} w_{\{\alpha_2,\alpha_3,\alpha_4,\alpha_5,\alpha_6,\alpha_7,\alpha_8\}}\\  
           &=[1, 2, 3, 1, 2, 3, 4, 3, 1, 2, 3, 4, 5, 4, 3, 1, 2, 3, 4, 5, 6,\\
           & \ \quad 5, 4, 3, 1, 2, 3, 4, 5, 6, 7, 6, 5, 4, 3, 1, 2, 3, 4, 5, 6, 7].\\          
%           \\
%           &s_3s_4s_2s_5s_4s_3s_6s_5s_4s_2s_7s_6s_5s_4s_3s_8s_7s_6s_5s_4\\
% &s_2s_8s_7s_6s_5s_4s_3s_8s_7s_6s_5s_4s_8s_7s_6s_5s_8s_7s_6s_8s_7s_8.
      \end{align*}
      

        \item[(11)] $\gkd L(\lambda)=83$ if and only if 
    \begin{align*}
           &w_{\lambda}\sim_{LR}w_{\{\alpha_1,\alpha_2,\alpha_3,\alpha_4,\alpha_5,\alpha_6,\alpha_8\}}=w_{E_6}s_8.  
          % =&s_1s_3s_4s_5s_6s_2s_4s_5s_3s_4s_1s_3s_2s_4s_5s_6s_2s_4s_5s_3s_4s_1s_3s_2s_4s_5s_3s_4s_1s_3s_2s_4s_1s_3s_2s_1s_8.
      \end{align*}
     
     
     
    %  \begin{align*}
    %       w_{\lambda}\sim_{LR}&s_1s_2s_3s_1s_4s_2s_3s_1s_4s_3s_5s_4s_2s_3s_1s_4s_3s_5s_4s_2s_6\\
    % &s_5s_4s_2s_3s_1s_4s_3s_5s_4s_2s_6s_5s_4
    %       s_3s_1s_7s_6s_5s_2s_2s_3s_4s_5s_6s_7.
    %   \end{align*}



     
 \item[(12)] $\gkd L(\lambda)=84$ if and only if \begin{align*}
          &w_{\lambda}\sim_{LR}w_{\{\alpha_1,\alpha_2,\alpha_3,\alpha_4,\alpha_5,\alpha_6\}}=w_{E_6}
          % \\    =&s_1s_3s_4s_5s_6s_2s_4s_5s_3s_4s_1s_3s_2s_4s_5s_6s_2s_4s_5s_3s_4s_1s_3s_2s_4s_5s_3s_4s_1s_3s_2s_4s_1s_3s_2s_1
          \text{~and~} V(\Ann (L(\lam)))=\overline{\mathcal{O}}(D_4).
      \end{align*}


\item[(13)] $\gkd L(\lambda)=88$ if and only if \begin{align*}
          w_{\lambda}\sim_{LR}   &
     [1, 2, 3, 1, 2, 3, 4, 3, 1, 2, 3, 4, 5, 4, 3, 1, 2, 0, 3, 4, 5, 6, 5, 4,\\& 3, 1, 2, 0, 3, 2, 4, 3, 1, 5, 4, 3, 2, 0, 6, 5, 4, 7, 6, 5, 4, 3, 1, 2,\\& 0, 3, 2, 4, 3, 1, 5, 4, 3, 2, 0, 6, 5, 4, 3, 1, 2, 3, 4, 5].
\end{align*}
\item[(14)] $\gkd L(\lambda)=89$ if and only if \begin{align*}
          w_{\lambda}\sim_{LR}   &
      [1, 3, 1, 2, 4, 3, 1, 2, 3, 5, 4, 3, 1, 2, 3, 4, 5, 6, 5, 4, 3, 1, 2,\\& 0, 3, 2, 4, 3, 5, 6, 7, 6, 5, 4, 3, 1, 2, 0, 3, 2, 4, 3, 1, 5, 4, 3,\\& 2, 0, 6, 5, 4, 3, 1, 2, 3, 4, 5, 6, 7, 6, 5, 4, 3, 1, 2, 0, 3].
\end{align*}

     

  

\item[(15)] $\gkd L(\lambda)=90$ if and only if \begin{align*}
           w_{\lambda}&\sim_{LR}w_{\{\alpha_2,\alpha_3,\alpha_4,\alpha_5,\alpha_6,\alpha_7\}}\\  &=[1, 2, 3, 1, 2, 3, 4, 3, 1, 2, 3, 4, 5, 4, 3, 1, 2, 3, 4, 5, 6, 5, 4, 3, 1, 2, 3, 4, 5, 6]. 
           % =&s_2s_3s_4s_2s_3s_4s_5s_4s_2s_3s_4s_5s_6s_5s_4s_2s_3s_4s_5s_6s_7s_6s_5s_4s_2s_3s_4s_5s_6s_7.
      \end{align*}
\item[(16)] $\gkd L(\lambda)=92$ if and only if \begin{align*}
          &w_{\lambda}\sim_{LR}w_{\{\alpha_1,\alpha_3,\alpha_4,\alpha_5,\alpha_6,\alpha_7,\alpha_8\}} \text{~and~} V(\Ann (L(\lam)))=\overline{\mathcal{O}}({D}_4(a_1)+A_2).
          % \\    =&s_8s_7s_8s_6s_7s_8s_5s_6s_7s_8s_4s_5s_6s_7s_8s_3s_4s_5s_6s_7s_8s_1s_3s_4s_5s_6s_7s_8
      \end{align*}
  
      
\item[(17)] $\gkd L(\lambda)=94$ if and only if \begin{align*}
          w_{\lambda}&\sim_{LR}   
      [0, 1, 2, 0, 3, 1, 2, 0, 3, 2, 4, 3, 1, 2, 3, 4, 6, 5, 4, 3, 1, 2,\\
      &\quad \quad \quad  0, 3, 2, 4, 3, 1, 5, 4, 3, 2, 0, 6, 7, 6, 5, 4, 3, 1, 2, 3, 4, 5]\\
    &\text{~and~}  V(\Ann (L(\lam)))=\overline{\mathcal{O}}({D}_4(a_1)+A_2).
\end{align*}



\item[(18)] $\gkd L(\lambda)=95$ if and only if \begin{align*}
    w_{\lambda}\sim_{LR}&[1, 2, 3, 1, 2, 0, 3, 4, 3, 1, 2, 0, 3, 2, 4,3,\\
    & 5, 4, 3, 1, 2, 0, 3, 2, 4, 3, 1, 5,4,3,2].
\end{align*}
% s_2s_3s_4s_2s_3s_1s_4s_5s_4s_2s_3s_1s_4s_3s_5s_4s_6s_5s_4s_2s_3s_1s_4s_3s_5s_4s_2s_6s_5s_4s_3



\item[(19)] $\gkd L(\lambda)=96$ if and only if \begin{align*}
          w_{\lambda}\sim_{LR} &  
      [0, 1, 2, 3, 1, 2, 3, 4, 3, 1, 2, 3, 4, 5, 4, 3, 1, 6, 5, 4, 3, 1,\\& 2, 0, 3, 2, 4, 3, 5, 4, 6, 5, 7, 6, 5, 4, 3, 1, 2, 0, 3, 2, 4, 5].
\end{align*}



\item[(20)] $\gkd L(\lambda)=97$ if and only if \begin{align*}
           &w_{\lambda}\sim_{LR}w_{\{\alpha_1,\alpha_2,\alpha_3,\alpha_4,\alpha_5,\alpha_7,\alpha_8\}}=w_{\{\alpha_1,\alpha_2,\alpha_3,\alpha_4,\alpha_5\}} \cdot w_{\{\alpha_7,\alpha_8\}}\\ 
           =& [4,1,3,4,1,3,2,3,4,1,3,2,0,2,3,4,1,3,2,0,7,6,7].          %=&s_5s_4s_2s_3s_4s_5s_1s_3s_4s_2s_1s_3s_4s_5s_1s_3s_4s_1s_3s_1          
          % s_8s_7s_8.
      \end{align*}






\item[(21)] $\gkd L(\lambda)=98$ if and only if $$w_{\lambda}\sim_{LR}
%s_1s_8s_7s_8s_6s_7s_8s_5s_6s_7s_8s_4s_5s_6s_7s_8s_2s_4s_5s_6s_7s_8=
w_{\{\alpha_1,\alpha_2,\alpha_4,\alpha_5,\alpha_6,\alpha_7,\alpha_8\}} \text{~and~} V(\Ann (L(\lam)))=\overline{\mathcal{O}}(A_4+A_2+A_1),$$ 
or \begin{align*}
    w_{\lambda}&\sim_{LR}[0,2,0,3,2,0,4,3,2,0,5,4,3,2,0,6,5,4,3,1,2,0,3,2,4,3,1,5,4,3,2]\\
%s_1s_3s_1s_4s_3s_1s_5s_4s_3s_1s_6s_5s_4s_3s_1s_7s_6s_5s_4s_2s_3s_1s_4s_3s_5s_4s_2s_6s_5s_4s_3
&\text{~and~} V(\Ann (L(\lam)))=\overline{\mathcal{O}}({D}_5(a_1)+A_1).
\end{align*}



 \item[(22)] $\gkd L(\lambda)=99$ if and only if \begin{align*}
w_{\lambda}&\sim_{LR}
%s_3s_4s_2s_5s_4s_3s_6s_5s_4s_2s_6s_5s_4s_3s_6s_5s_4s_6s_5s_6s_8=
w_{\{\alpha_2,\alpha_3,\alpha_4,\alpha_5,\alpha_6,\alpha_8\}}=[1, 2, 3, 1, 2, 3, 4, 3, 1, 2, 3, 4, 5, 4, 3, 1, 2, 3, 4, 5,7]\\
& \text{~and~} V(\Ann (L(\lam)))=\overline{\mathcal{O}}({E}_6(a_3)),
 \end{align*}
 or $$w_{\lambda}\sim_{LR}
 %s_7s_6s_7s_5s_6s_7s_4s_5s_6s_7s_3s_4s_5s_6s_7s_1s_3s_4s_5s_6s_7=
 w_{\{\alpha_1,\alpha_3,\alpha_4,\alpha_5,\alpha_6,\alpha_7\}} \text{~and~} V(\Ann (L(\lam)))=\overline{\mathcal{O}}(D_4+A_2).$$
      
        \item[(23)] $\gkd L(\lambda)=100$ if and only if \begin{align*}w_{\lambda}&\sim_{LR}
        %s_3s_4s_2s_5s_4s_3s_6s_5s_4s_2s_6s_5s_4s_3s_6s_5s_4s_6s_5s_6=
        w_{\{\alpha_2,\alpha_3,\alpha_4,\alpha_5,\alpha_6\}}=[1, 2, 3, 1, 2, 3, 4, 3, 1, 2, 3, 4, 5, 4, 3, 1, 2, 3, 4, 5]\\
        &\text{~and~} V(\Ann (L(\lam)))=\overline{\mathcal{O}}(D_5).
        \end{align*}
  
    \item[(24)] $\gkd L(\lambda)=104$ if and only if \begin{align*}
          w_{\lambda}&\sim_{LR}w_{\{\alpha_1,\alpha_2,\alpha_4,\alpha_5,\alpha_6,\alpha_7\}}
          % &=s_1s_7s_6s_7s_5s_6s_7s_4s_5s_6s_7s_2s_4s_5s_6s_7
          \text{~and~} V(\Ann (L(\lam)))=\overline{\mathcal{O}}({E}_8(a_7)).
      \end{align*}
      
      
        \item[(25)] $\gkd L(\lambda)=105$ if and only if $$w_{\lambda}\sim_{LR}
       % s_6s_5s_6s_4s_5s_6s_3s_4s_5s_6s_1s_3s_4s_5s_6=
        w_{\{\alpha_1,\alpha_3,\alpha_4,\alpha_5,\alpha_6\}}
        \text{~and~} V(\Ann (L(\lam)))=\overline{\mathcal{O}}({D}_6(a_1)),$$
        or
\begin{align*}w_{\lambda}&\sim_{LR}
%s_2s_4s_3s_5s_4s_2s_5s_4s_3s_5s_4s_5s_8s_7s_8=
w_{\{\alpha_2,\alpha_3,\alpha_4,\alpha_5,\alpha_7,\alpha_8\}}=[1, 2, 3, 1, 2, 3, 4, 3, 1, 2, 3, 4,7,6,7]\\
 &\text{~and~} V(\Ann (L(\lam)))=\overline{\mathcal{O}}(A_6).
 \end{align*}


 
     \item[(26)] $\gkd L(\lambda)=106$ if and only if \begin{align*}
         w_{\lambda}&\sim_{LR}
     [0, 1, 2, 0, 3, 1, 2, 0, 3, 2, 4, 3, 1, 2, 0, 3, 4, 5, 4, 3, 6, 5, 4, 3, 1, 2, 0, 3]\\
     &\text{~and~} V(\Ann (L(\lam)))=\overline{\mathcal{O}}({E}_7(a_4)),
     \end{align*}
     % s_1s_2s_3s_1s_4s_2s_3s_1s_4s_3s_5s_4s_2s_3s_1s_4s_5s_6s_5s_4s_7s_6s_5s_4s_2s_3s_1s_4
      or $$w_{\lambda}\sim_{LR}
     %s_3s_1s_3s_2s_8s_7s_8s_6s_7s_8s_5s_6s_7s_8=
     w_{\{\alpha_1,\alpha_3,\alpha_2,\alpha_5,\alpha_6,\alpha_7,\alpha_8\}}\text{~and~} V(\Ann (L(\lam)))=\overline{\mathcal{O}}(A_6+A_1).$$
    
\item[(27)] $\gkd L(\lambda)=107$ if and only if \begin{align*}
    w_{\lambda}&\sim_{LR}
    %s_2s_4s_3s_5s_4s_2s_5s_4s_3s_5s_4s_5s_7=
    w_{\{\alpha_2,\alpha_3,\alpha_4,\alpha_5,\alpha_7\}}=[1, 2, 3, 1, 2, 3, 4, 3, 1, 2, 3, 4,6]\\
    &\text{~and~} V(\Ann (L(\lam)))=\overline{\mathcal{O}}({E}_6(a_1)), 
\end{align*}
        or
$$w_{\lambda}\sim_{LR}
%s_2s_4s_2s_3s_4s_2s_1s_3s_4s_2s_7s_6s_7=
w_{\{\alpha_1,\alpha_2,\alpha_3,\alpha_4,\alpha_6,\alpha_7\}} \text{~and~} V(\Ann (L(\lam)))=\overline{\mathcal{O}}(D_5+A_2).$$

      
\item[(28)] $\gkd L(\lambda)=108$ if and only if \begin{align*}
    w_{\lambda}&\sim_{LR}
%s_2s_4s_3s_5s_4s_2s_5s_4s_3s_5s_4s_5=
w_{\{\alpha_2,\alpha_3,\alpha_4,\alpha_5\}}=[1, 2, 3, 1, 2, 3, 4, 3, 1, 2, 3, 4]\\
&\text{~and~} V(\Ann (L(\lam)))=\overline{\mathcal{O}}({D}_7(a_2)), 
\end{align*}
or $$w_{\lambda}\sim_{LR}
%s_4s_3s_4s_1s_3s_4s_8s_7s_8s_6s_7s_8=
w_{\{\alpha_1,\alpha_3,\alpha_4,\alpha_6,\alpha_7,\alpha_8\}} 
 \text{~and~} V(\Ann (L(\lam)))=\overline{\mathcal{O}}({E}_6).$$
      
        \item[(29)] $\gkd L(\lambda)=109$ if and only if $$w_{\lambda}\sim_{LR}
       % s_1s_6s_5s_6s_4s_5s_6s_2s_4s_5s_6=
        w_{\{\alpha_1,\alpha_2,\alpha_4,\alpha_5,\alpha_6\}} \text{~and~} V(\Ann (L(\lam)))=\overline{\mathcal{O}}({E}_6(a_1)+A_1).$$
       
        \item[(30)] $\gkd L(\lambda)=110$ if and only if $$w_{\lambda}\sim_{LR}w_{\{\alpha_3,\alpha_4,\alpha_5,\alpha_6\}}
       \text{~and~} V(\Ann (L(\lam)))=\overline{\mathcal{O}}({E}_7(a_3)),$$ or
$$w_{\lambda}\sim_{LR}
%s_3s_1s_3s_2s_7s_6s_7s_5s_6s_7=
w_{\{\alpha_1,\alpha_3,\alpha_2,\alpha_5,\alpha_6,\alpha_7\}} 
 \text{~and~} V(\Ann (L(\lam)))=\overline{\mathcal{O}}({E}_8(b_6)).$$
 
  \item[(31)] $\gkd L(\lambda)=111$ if and only if $$w_{\lambda}\sim_{LR}w_{\{\alpha_1,\alpha_3,\alpha_5,\alpha_6,\alpha_7\}} \text{~and~} V(\Ann (L(\lam)))=\overline{\mathcal{O}}({D}_7(a_1)).$$
   
    \item[(32)] $\gkd L(\lambda)=112$ if and only if $$w_{\lambda}\sim_{LR}w_{\{\alpha_1,\alpha_2,\alpha_5,\alpha_6,\alpha_7\}}\text{~and~} V(\Ann (L(\lam)))=\overline{\mathcal{O}}({E}_8(a_6)).$$
      
        
        \item[(33)] $\gkd L(\lambda)=113$ if and only if $$w_{\lambda}\sim_{LR}w_{\{\alpha_1,\alpha_4,\alpha_5,\alpha_6\}}
        \text{~and~} V(\Ann (L(\lam)))=\overline{\mathcal{O}}({E}_8(b_5)).$$
      
       \item[(34)] $\gkd L(\lambda)=114$ if and only if $$w_{\lambda}\sim_{LR}w_{\{\alpha_4,\alpha_5,\alpha_6\}} \text{~and~} V(\Ann (L(\lam)))=\overline{\mathcal{O}}({E}_7(a_1)),$$ or $$w_{\lambda}\sim_{LR}w_{\{\alpha_1,\alpha_3,\alpha_5,\alpha_6\}} \text{~and~} V(\Ann (L(\lam)))=\overline{\mathcal{O}}({E}_8(a_5)).$$
        
        \item[(35)] $\gkd L(\lambda)=115$ if and only if $$w_{\lambda}\sim_{LR}w_{\{\alpha_1,\alpha_2,\alpha_5,\alpha_6\}}.$$
        \item[(36)] $\gkd L(\lambda)=116$ if and only if $$w_{\lambda}\sim_{LR}w_{\{\alpha_1,\alpha_5,\alpha_6\}} \text{~and~} V(\Ann (L(\lam)))=\overline{\mathcal{O}}({E}_8(a_4)).$$
        
        \item[(37)] $\gkd L(\lambda)=117$ if and only if $$w_{\lambda}\sim_{LR}w_{\{\alpha_7,\alpha_8\}}.$$ 
        \item[(38)] $\gkd L(\lambda)=118$ if and only if $w_{\lambda}\sim_{LR}w_{\{\alpha_1,\alpha_8\}}$.




      
\end{itemize}
\end{Prop}
 
We also have the following result for nonintegral highest weight modules of $E_8$.

\begin{Prop} Let $L(\lam)$ be a simple nonintegral highest weight module of $E_8$. Assume that $\Phi_{[\lam]}^\vee$ is pseudo-maximal.
	Then
    one of the following holds: 
    \begin{itemize}
	\item[(1)]  $\Phi_{[\lam]}\simeq D_8$. Then 
 $ \gkd L(\lam) $ can achieve the following values:
$$64,77,78,86,89,90,91,92,96,97,98,99,100,[102,120].$$













\item[(2)] $\Phi_{[\lam]}\simeq A_8$. Then
 $ \gkd L(\lam) $ can achieve the following values:
$$84,92,98,99,102,104, 105, [107,120].$$








\item[(3)] $\Phi_{[\lam]}\simeq E_7\times A_1$.
 Then 
 $ \gkd L(\lam) $ can achieve the following values:
$$56,57,73,74,82,83,84,89,90,94,95,97,98,99,100,[103,120].$$



\item[(4)] $\Phi_{[\lam]}\simeq A_7\times A_1$.
 Then 
 $ \gkd L(\lam) $ can achieve the following values:
$$91,92,98,99,[103,120].$$









\item[(5)] $\Phi_{[\lam]}\simeq A_5\times A_2\times A_1$ and 
 $$\Delta_{[\lam]}=\{\ep_8+\ep_3,\alpha_5,\alpha_6,\alpha_7,\alpha_8,\alpha_1,\alpha_3,\alpha_2\}.$$ Then 
 $ \gkd L(\lam) $ can achieve the following values:
$$[101,120].$$







\item[(6)] $\Phi_{[\lam]}\simeq  A_4\times A_4$ and 
 $$\Delta_{[\lam]}=\{\alpha_2,\alpha_4,\alpha_3,\alpha_1,\ep_8+\ep_4,\alpha_6,\alpha_7,\alpha_8\}.$$ Then 
 $ \gkd L(\lam) $ can achieve the following values:
$$100,104,[106,120].$$





\item[(7)] $\Phi_{[\lam]}\simeq D_5\times A_3$. Then 
 $ \gkd L(\lam) $ can achieve the following values:
 $$94,[97,102],[104,120].$$









\item[(8)] $\Phi_{[\lam]}\simeq E_6\times A_2$. Then 
 $ \gkd L(\lam) $ can achieve the following values:
$$81,83,84,92,94,95,97,99,100,102,[104,120].$$






\end{itemize}
\end{Prop}





% \begin{Rem}\label{e8-gou} From now on we denote $\beta_1=\ep_1-\ep_2, \beta_2=\ep_2-\ep_3, \beta_3=\ep_3-\ep_4$,  $\beta_4=\ep_4-\ep_5$, $\beta_5=\ep_5-\ep_6$, $\beta_6=\ep_6-\ep_7$, $\beta_7=\ep_7-\ep_8$  and $\beta_8=\ep_7+\ep_8$. Then the simple root system of $D_8$ is $\Delta=\{\beta_1, \beta_2, \cdots,\beta_8\}$. In general, by Theorem \ref{transf} and Corollary \ref{transf-2} we have the following:
%     \begin{enumerate}
%         \item In the case of $\mathfrak{l}_1$, we define an isomorphism map $\phi: \mathbb{C}^8\rightarrow \mathbb{C}^8$ such that 
%  $\phi|_{\Phi_{[\lambda]}}: \Phi_{[\lambda]} \rightarrow D_8$ is an isomorphism and  $\phi(\ep_8-\ep_7)=\beta_1, \phi(\alpha_8)=\beta_2, \phi(\alpha_7)=\beta_3, \phi(\alpha_6)=\beta_4,\phi(\alpha_5)=\beta_5,\phi(\alpha_4)=\beta_6,\phi(\alpha_3)=\beta_7,\phi(\alpha_2)=\beta_8$. Then we can extend it to an isomorphism and still denote it by $\phi: \mathfrak{h}_{\Phi_{[\lambda]}}^*\rightarrow \mathfrak{h}_{D_8}^*$.
%  Thus for $\lam\in \mathfrak{h}_{E_8}^* $, it will be an integral weight  in  $\mathfrak{h}_{\Phi_{[\lambda]}}^*$ and $$\phi({\lam})=\lam'=(\lam_8,\lam_7,\lam_6,\lam_5,\lam_4,\lam_3,\lam_2,\lam_1)$$  will  be an integral weight of type $D_8$.
      
%      \item In the case of $\mathfrak{l}_2$, we have two cases:
%      \begin{enumerate}
%          \item   we define an isomorphism map $\phi: \mathbb{C}^8\rightarrow \mathbb{C}^8$ such that 
%  $\phi|_{\Phi_{[\lambda]}}: \Phi_{[\lambda]} \rightarrow D_8$ is an isomorphism and $\phi(\alpha_1)=\beta_1, \phi(\alpha_3)=\beta_2, \phi(\alpha_4)=\beta_3, \phi(\alpha_5)=\beta_4,\phi(\alpha_6)=\beta_5,$
% $\phi(\alpha_7)=\beta_6,\phi(\alpha_8)=\beta_7,\phi(\gamma_2)=\beta_8$. 
%        Then we can extend it to an isomorphism and still denote it by $\phi: \mathfrak{h}_{\Phi_{[\lambda]}}^*\rightarrow \mathfrak{h}_{D_8}^*$.   Thus \begin{align*}
%            &\phi(\lam)=\lambda'\\
%            =&(\frac{-\lam_1-\cdots-\lam_7+3\lam_8}{4},
%                \frac{-3\lam_1+\lam_2+\cdots+\lam_8}{4},\frac{\lam_1-3\lam_2+\lam_3+\cdots+\lam_8}{4},\\
%            & \frac{\lam_1+\lam_2-3\lam_3+\lam_4+\cdots+\lam_8}{4},\frac{\lam_1+\lam_2+\lam_3-3\lam_4+\lam_5+\cdots+\lam_8}{4},\\
%            & \frac{\lam_1+\lam_2+\lam_3+\lam_4-3\lam_5+\lam_6+\lam_7+\lam_8}{4},\\
%            &\frac{\lam_1+\lam_2+\lam_3+\lam_4+\lam_5-3\lam_6+\lam_7+\lam_8}{4},\\
%            &
%            \frac{\lam_1+\lam_2+\cdots+\lam_6-3\lam_7+\lam_8}{4}) 
%         \end{align*} will be an integral weight of type $D_8$.
%         \item   we define an isomorphism map $\phi: \mathbb{C}^8\rightarrow \mathbb{C}^8$ such that 
%  $\phi|_{\Phi_{[\lambda]}}: \Phi_{[\lambda]} \rightarrow A_8$ is an isomorphism and $\phi(\gamma'_2)=\beta_1, \phi(\alpha_1)=\beta_2, \phi(\alpha_3)=\beta_3, \phi(\alpha_4)=\beta_4,\phi(\alpha_5)=\beta_5,$
% $\phi(\alpha_6)=\beta_6,\phi(\alpha_7)=\beta_7,\phi(\alpha_8)=\ep_7-\ep_8$. 
%        Then we can extend it to an isomorphism and still denote it by $\phi: \mathfrak{h}_{\Phi_{[\lambda]}}^*\rightarrow \mathfrak{h}_{A_8}^*$.   Thus \begin{align*}
%            &\phi(\lam)=\lambda'\\
%            =&(\frac{\lam_1+\cdots+\lam_7+5\lam_8}{6},
%                -\frac{\lam_1+\cdots+\lam_7-\lam_8}{3},\frac{-5\lam_1+\lam_2+\cdots+\lam_7-\lam_8}{6},\\
%            & \frac{\lam_1-5\lam_2+\lam_3+\cdots+\lam_7-\lam_8}{6},\frac{\lam_1+\lam_2-5\lam_3+\lam_4+\cdots+\lam_7-\lam_8}{6},\\
%            & \frac{\lam_1+\lam_2+\lam_3-5\lam_4+\lam_5+\lam_6+\lam_7-\lam_8}{6},\\
%            &\frac{\lam_1+\lam_2+\lam_3+\lam_4-5\lam_5+\lam_6+\lam_7-\lam_8}{6},\\
%            &
%            \frac{\lam_1+\cdots+\lam_5-5\lam_6+\lam_7-\lam_8}{6}, \frac{\lam_1+\cdots+\lam_6-5\lam_7-\lam_8}{6}) 
%         \end{align*} will be an integral weight of type $A_8$.
%      \end{enumerate}
     
     
     
   
 
%  \item In the case of $\mathfrak{l}_3$, we have two cases：
%  \begin{enumerate}
%      \item  
%  we define an isomorphism map $\phi: \mathbb{C}^8\rightarrow \mathbb{C}^8$ such that 
%  $\phi|_{\Phi_{[\lambda]}}: \Phi_{[\lambda]} \rightarrow E_7\times A_1$ is an isomorphism and  $\phi(\alpha_8)=\alpha_1, \phi(\alpha_7)=\alpha_3, \phi(\alpha_6)=\alpha_4, \phi(\gamma_3)=\alpha_2,\phi(\alpha_5)=\alpha_5,\phi(\alpha_4)=\alpha_6,\phi(\alpha_2)=\alpha_7, \phi(\alpha_1)=\ep_7+\ep_8$. Then we can extend it to an isomorphism and still denote it by $\phi: \mathfrak{h}_{\Phi_{[\lambda]}}^*\rightarrow \mathfrak{h}_{E_7\times A_1}^*$.
%  Thus for $\lam\in \mathfrak{h}_{E_8}^* $, it will be an integral weight  in  $\mathfrak{h}_{\Phi_{[\lambda]}}^*$. So \begin{align*}
%            &\phi(\lam)|_{\mathfrak{h}_{E_7}^*}\\
%            =&(\frac{-\lam_1+\lam_2+\cdots+\lam_5-3\lam_6-\lam_7+\lam_8}{4},
%                \frac{-\lam_1+\lam_2+\lam_3+\lam_4-3\lam_5+\lam_6-\lam_7+\lam_8}{4},\\
%            & \frac{-\lam_1+\lam_2+\lam_3-3\lam_4+\lam_5+\lam_6-\lam_7+\lam_8}{4},\frac{-\lam_1+\lam_2-3\lam_3+\lam_4+\lam_5+\lam_6-\lam_7+\lam_8}{4},\\
%            &\frac{-\lam_1-3\lam_2+\lam_3+\cdots+\lam_6-\lam_7+\lam_8}{4},\frac{3\lam_1+\lam_2+\cdots+\lam_6-\lam_7+\lam_8}{4},\\
%            &-\frac{\lam_7+\lam_8}{4},
%            \frac{\lam_7+\lam_8}{4}) 
%         \end{align*} will  be an integral weight of type $E_7$. 
        
%         And $$\phi(\lam)|_{\mathfrak{h}_{A_1}^*}=\frac{\lam_1-\lam_2-\cdots-\lam_7+\lam_8}{4}(\ep_7+\ep_8):=x_8(\ep_7+\ep_8).$$ The value of its $\aff$ function is $1$ if $x_8>0$  and $0$ if $x_8\lest 0$. 
% \item 

% \end{enumerate}
%  \item In the case of $\mathfrak{l}_4$, we define an isomorphism map $\phi: \mathbb{C}^8\rightarrow \mathbb{C}^8$ such that 
%  $\phi|_{\Phi_{[\lambda]}}: \Phi_{[\lambda]} \rightarrow E_7\times A_1$ is an isomorphism and  $\phi(\alpha_8)=\alpha_1, \phi(\alpha_7)=\alpha_3, \phi(\alpha_6)=\alpha_4, \phi(\alpha_5)=\alpha_2,\phi(\ep_3+\ep_4)=\alpha_5,\phi(\alpha_1)=\alpha_6,\phi(\alpha_3)=\alpha_7, \phi(\alpha_2)=\ep_7+\ep_8$. Then we can extend it to an isomorphism and still denote it by $\phi: \mathfrak{h}_{\Phi_{[\lambda]}}^*\rightarrow \mathfrak{h}_{E_7\times A_1}^*$.
%  Thus for $\lam\in \mathfrak{h}_{E_8}^* $, it will be an integral weight  in  $\mathfrak{h}_{\Phi_{[\lambda]}}^*$. So \begin{align*}
%            &\phi(\lam)|_{\mathfrak{h}_{E_7}^*}\\
%            =&(\frac{-\lam_3+\lam_4+\lam_5-\lam_6}{2},
%                \frac{-\lam_3+\lam_4-\lam_5+\lam_6}{2},\\
%            & \frac{-\lam_3-\lam_4+\lam_5+\lam_6}{2},\frac{\lam_3+\lam_4+\lam_5+\lam_6}{2},\frac{\lam_1-\lam_2-\lam_7+\lam_8}{2},\\
%            &\frac{-\lam_1+\lam_2-\lam_7+\lam_8}{2},-\frac{\lam_7+\lam_8}{2},
%            \frac{\lam_7+\lam_8}{2}) 
%         \end{align*} will  be an integral weight of type $E_7$. 
        
%         And $$\phi(\lam)|_{\mathfrak{h}_{A_1}^*}=\frac{\lam_1+\lam_2}{2}(\ep_7+\ep_8):=x_8(\ep_7+\ep_8).$$ The value of its $\aff$ function is $1$ if $x_8>0$  and $0$ if $x_8\lest 0$.
     
%       \item In the case of $\mathfrak{l}_5$, we define an isomorphism map $\phi: \mathbb{C}^8\rightarrow \mathbb{C}^8$ such that 
%  $\phi|_{\Phi_{[\lambda]}}: \Phi_{[\lambda]} \rightarrow D_8$ is an isomorphism and $\phi(\alpha_2)=\beta_1, \phi(\alpha_4)=\beta_2, \phi(\alpha_3)=\beta_3, \phi(\alpha_1)=\beta_4,\phi(\ep_4+\ep_5)=\beta_5,$
% $\phi(\alpha_7)=\beta_6,\phi(\alpha_6)=\beta_7,\phi(\alpha_8)=\beta_8$. 
%        Then we can extend it to an isomorphism and still denote it by $\phi: \mathfrak{h}_{\Phi_{[\lambda]}}^*\rightarrow \mathfrak{h}_{D_8}^*$.   Thus \begin{align*}
%            &\phi(\lam)=\lambda'\\
%            =&(\frac{\lam_1+\lam_2+\lam_3+\lam_8}{2},
%                \frac{-\lam_1\lam_2+\lam_3+\lam_8}{2},\frac{-\lam_1+\lam_2-\lam_3+\lam_8}{2},\\
%            &\frac{\lam_1-\lam_2-\lam_3+\lam_8}{2},\frac{\lam_4+\lam_5+\lam_6+\lam_7}{2},\frac{-\lam_4-\lam_5+\lam_6+\lam_7}{2},\\
%            &\frac{-\lam_4+\lam_5-\lam_6+\lam_7}{2},
%            \frac{\lam_4-\lam_5-\lam_6+\lam_7}{2}) 
%         \end{align*} will be an integral weight of type $D_8$.

%  \item In the case of $\mathfrak{l}_6$, we define an isomorphism map $\phi: \mathbb{C}^8\rightarrow \mathbb{C}^8$ such that 
%  $\phi|_{\Phi_{[\lambda]}}: \Phi_{[\lambda]} \rightarrow D_8$ is an isomorphism and $\phi(\alpha_7)=\beta_1, \phi(\alpha_8)=\beta_2, \phi(\ep_5+\ep_6)=\beta_3, \phi(\alpha_1)=\beta_4,\phi(\alpha_3)=\beta_5,$
% $\phi(\alpha_4)=\beta_6,\phi(\alpha_5)=\beta_7,\phi(\alpha_2)=\beta_8$. 
%        Then we can extend it to an isomorphism and still denote it by $\phi: \mathfrak{h}_{\Phi_{[\lambda]}}^*\rightarrow \mathfrak{h}_{D_8}^*$.   Thus \begin{align*}
%            &\phi(\lam)=\lambda'\\
%            =&(\frac{-\lam_5+\lam_6+\lam_7+\lam_8}{2},
%               \frac{\lam_5-\lam_6+\lam_7+\lam_8}{2}, \frac{\lam_5+\lam_6-\lam_7+\lam_8}{2},\\
%            &\frac{-\lam_5-\lam_6-\lam_7+\lam_8}{2},\frac{-\lam_1+\lam_2+\lam_3+\lam_4}{2},\frac{\lam_1-\lam_2+\lam_3+\lam_4}{2},\\
%            &\frac{\lam_1+\lam_2-\lam_3+\lam_4}{2},
%            \frac{\lam_1+\lam_2+\lam_3-\lam_4}{2}) 
%         \end{align*} will be an integral weight of type $D_8$.

        
% \item In the case of $\mathfrak{l}_7$, we define an isomorphism map $\phi: \mathbb{C}^8\rightarrow \mathbb{C}^8$ such that 
%  $\phi|_{\Phi_{[\lambda]}}: \Phi_{[\lambda]} \rightarrow E_7\times A_1$ is an isomorphism and  $\phi(\alpha_6)=\alpha_1, \phi(\alpha_5)=\alpha_3, \phi(\alpha_4)=\alpha_4, \phi(\alpha_2)=\alpha_2,\phi(\alpha_3)=\alpha_5,\phi(\alpha_1)=\alpha_6,\phi(\ep_7+\ep_6)=\alpha_7, \phi(\alpha_8)=\ep_7+\ep_8$. Then we can extend it to an isomorphism and still denote it by $\phi: \mathfrak{h}_{\Phi_{[\lambda]}}^*\rightarrow \mathfrak{h}_{E_7\times A_1}^*$.
%  Thus for $\lam\in \mathfrak{h}_{E_8}^* $, it will be an integral weight  in  $\mathfrak{h}_{\Phi_{[\lambda]}}^*$. So \begin{align*}
%            &\phi(\lam)|_{\mathfrak{h}_{E_7}^*}\\
%            =&(\frac{\lam_1+\lam_2+\lam_3-\lam_4}{2},
%                \frac{\lam_1+\lam_2-\lam_3+\lam_4}{2}, \frac{\lam_1-\lam_2+\lam_3+\lam_4}{2},\\
%            &\frac{-\lam_1+\lam_2+\lam_3+\lam_4}{2},\frac{-\lam_5-\lam_6-\lam_7+\lam_8}{2},\frac{-\lam_5+\lam_6+\lam_7+\lam_8}{2},\\
%            &-\frac{\lam_5+\lam_8}{2},
%            \frac{\lam_5+\lam_8}{2}) 
%         \end{align*} will  be an integral weight of type $E_7$. 
        
%         And $$\phi(\lam)|_{\mathfrak{h}_{A_1}^*}=\frac{\lam_7-\lam_6}{2}(\ep_1+\ep_2):=x_8(\ep_7+\ep_8).$$ The value of its $\aff$ function is $1$ if $x_8>0$  and $0$ if $x_8\lest 0$.

%  \item In the case of $\mathfrak{l}_8$, we define an isomorphism map $\phi: \mathbb{C}^8\rightarrow \mathbb{C}^8$ such that 
%  $\phi|_{\Phi_{[\lambda]}}: \Phi_{[\lambda]} \rightarrow E_7\times A_1$ is an isomorphism and $\phi(\alpha_1)=\alpha_1, \phi(\alpha_2)=\alpha_2, \phi(\alpha_3)=\alpha_3, \phi(\alpha_4)=\alpha_4,\phi(\alpha_5)=\alpha_5,\phi(\alpha_6)=\alpha_6,\phi(\alpha_7)=\alpha_7,\phi(\ep_7+\ep_8)=\ep_7+\ep_8$. Then we can extend it to an isomorphism and still denote it by $\phi: \mathfrak{h}_{\Phi_{[\lambda]}}^*\rightarrow \mathfrak{h}_{E_7\times A_1}^*$.
%  Thus for $\lam\in \mathfrak{h}_{E_8}^* $, it will be an integral weight  in  $\mathfrak{h}_{\Phi_{[\lambda]}}^*$. So \begin{align*}
%            &\phi(\lam)|_{\mathfrak{h}_{E_7}^*}
%            =(\lam_1,\cdots,\lam_6,\frac{\lam_7-\lam_8}{2},\frac{-\lam_7+\lam_8}{2}).
%         \end{align*} will  be an integral weight of type $E_7$. 
        
%         And $$\phi(\lam)|_{\mathfrak{h}_{A_1}^*}=\frac{\lam_7+\lam_8}{2}(\ep_7+\ep_8):=x_8(\ep_7+\ep_8).$$ The value of its $\aff$ function is $1$ if $x_8>0$  and $0$ if $x_8\lest 0$.
%     \end{enumerate}
% Then the value of $\aff(w_{\lam})=\aff(w_{\lam'})$ (or $\aff(w_{\lam'})$) can be computed by the algorithm in \cite{BXX} (or Theorem \ref{E6}).    
% \end{Rem}
\begin{proof}
The argument is similar to that in the proof of Proposition \ref{e6-non}.


\end{proof}
\begin{ex}
	Let $\mathfrak{g}=E_8$ and let $L(\lam)$ be the highest weight module of $\mathfrak{g}$ with  $\lambda=(\frac{1}{4},\frac{5}{4},\frac{9}{4},\frac{13}{4},\frac{9}{4},\frac{1}{4},
 \frac{5}{4},\frac{9}{4}).$ It is easy to verify that $\Phi_{[\lambda]}$ is a subsystem with simple roots $\{\alpha_1,\alpha_3,\alpha_4,\alpha_5,\alpha_6,\alpha_7,\alpha_8,\gamma\}$ where $\gamma=\frac{1}{2}(\ep_1+\ep_2+\ep_3+\ep_4+\ep_5-\ep_6-\ep_7+\ep_8)$. Therefore $\Phi_{[\lambda]}\simeq D_8$.  We define the isomorphism  $\phi:  \Phi_{[\lambda]} \rightarrow D_8$ such that $\phi(\alpha_1)=\beta_1, \phi(\alpha_3)=\beta_2, \phi(\alpha_4)=\beta_3, \phi(\alpha_5)=\beta_4,\phi(\alpha_6)=\beta_5,$
$\phi(\alpha_7)=\beta_6,\phi(\alpha_8)=\beta_7,\phi(\gamma)=\beta_8$.
Then  by the algorithm in \S \ref{nonint}, 
\begin{align*}
           \phi(\lam)=\lambda'&=-4\phi(\omega_1)+\phi(\omega_2)+\phi(\omega_3)+\phi(\omega_4)-\phi(\omega_5)-2\phi(\omega_6)+\phi(\omega_7)+5\phi(\omega_8)\\
           &=(-1,3,2,1,0,1,3,2)
        \end{align*} is an integral weight of type $D_8$.
  So $\aff(w_{\lambda})=\aff(w_{\lambda'})=17$ by the RS algorithm in \cite{BXX} and Proposition \ref{integral}.
 By Proposition \ref{pr:main1},  we have $\gkd L(\lambda)=|\Phi^+|-\aff(w_{\lambda})=120-17=103$.
\end{ex}



\begin{ex}
	Let $\mathfrak{g}=E_8$ and $L(\lam)$ be the highest weight module of $\mathfrak{g}$ with  $\lambda=(\frac{1}{2},-\frac{3}{2},-3,-2,-1,-4,-5,-19)$. It is easy to verify that $\Phi_{[\lambda]}$ is a subsystem with simple roots $\{\alpha_3,\alpha_1,\ep_3+\ep_4,\alpha_6,\alpha_5,\alpha_7,\alpha_8,\alpha_2\}$. Therefore $\Phi_{[\lambda]}\simeq E_7\times A_1$. Suppose the simple root system of $E_7\times A_1$ is $\Delta=\{\alpha_1, \alpha_2, \alpha_3,\alpha_4,\alpha_5,\alpha_6,\alpha_7\}\times \{\gamma_1=\ep_1-\ep_2\}$.
  We define the  isomorphism $\phi:  \Phi_{[\lambda]} \rightarrow E_7\times A_1$ such that  $\phi(\alpha_8)=\alpha_1, \phi(\alpha_7)=\alpha_3, \phi(\alpha_6)=\alpha_4, \phi(\alpha_5)=\alpha_2,\phi(\ep_3+\ep_4)=\alpha_5,\phi(\alpha_1)=\alpha_6,\phi(\alpha_3)=\alpha_7, \phi(\alpha_2)=\gamma_1$.

  Denote ${\lam}'=\phi(\lam|_{\mathfrak{h}_{\Phi_{[\lambda]}}})$. 
 Thus  by the algorithm in \S \ref{nonint},
 \begin{align*}
           &\lam'|_{{E_7}}=(2,-1,0,-5,-6,-8,12,-12)
        \end{align*} is an integral weight of type $E_7$, and 
$\lam'|_{A_1}=(-\frac{1}{2},\frac{1}{2})$ is an integral weight of type $A_1$ with $\aff(w_{\lam'|_{{A_1}}})=0$ by using  the RS algorithm in \cite{BXX} and Proposition \ref{integral}.


Denote $\gamma=\phi({\lam})|_{E_7}=(2,-1,0,-5,-6,-8,12,-12)$.
We can write $\gamma=-\omega_1+\omega_2-3\omega_3+\omega_4-5\omega_5-\omega_6-2\omega_7:=[-1,1,-3,1,-5,-1,-2]$. By using Lemma \ref{find-w-lambda}, we find that  $$s_4s_2s_4\gamma=[-1,-1,-1,-1,-3,-1,-2],$$ which is  antidominant. Thus we have $w_{\gamma}=s_4s_2s_4$.
Note that $w_{\gamma}=[3,1,3]$ in PyCox. Then by using PyCox, we find that  $\aff(w_{\gamma})=3$.
So $\aff(w_{\lambda})=\aff(w_{\lambda'})=3+0=3$.
 By Proposition \ref{pr:main1},  we have $\gkd L(\lambda)=|\Phi^+|-\aff(w_{\lambda})=120-3=117$.
\end{ex}


\section{Annihilator varieties of highest weight modules}\label{AVHWM}

Recall that one can compute the annihilator variety $\mathcal{O}_{{\rm Ann}(L(\lambda))}$ by applying the result of Joseph \cite{Jo85}. The goal of this section is to show that one method of computation is provided by the Sommers duality as defined in \cite{Sommers}. Note that in many cases, the dimension of a nilpotent orbit already determines uniquely the orbit itself, and thus such a computation can be avoided in those cases.


Joseph's result (see Proposition \ref{2dim}) asserts that
\begin{equation} \label{E:jos}
\mathcal{O}_{{\rm Ann}(L(\lambda))} = \mathcal{O}_{\rm Spr}^G \circ j_{W_{[\lambda]}}^{W(G)} (\pi_{w_\lambda}).
\end{equation}
Hence, in principle, one can use \eqref{E:jos} to compute the annihilator variety of $L(\lambda)$. However, explicating the $j$-induction in each case is non-trivial. Instead, we show that the right hand side of \eqref{E:jos} is intimately related to the Sommers duality studied in \cite{Sommers}, which entails a more direct formula in computing $\mathcal{O}_{{\rm Ann}(L(\lambda))}$.

We use notation as in \S \ref{sec:cell}. Let $G$ be an algebraic group over $\mathbb{C}$ with Lie algebra $\mathfrak{g}$. Recall that a pseudo-Levi subgroup $M \subseteq G$ is the connected  centralizer subgroup of a semisimple element of $G$. Such a group is essentially characterized by the fact that the Dynkin diagram of its root system is a subset of the extended Dynkin diagram of the root system of $G$. We consider pairs of the form $(M, \mathcal{O}_M)$, where $M$ is a pseudo-Levi subgroup and $\mathcal{O}_M \in \mathcal{N}(M)$ is a distinguished orbit. Let $G^\vee$ be the Langlands dual group of $G$. Then Sommers defined a natural map
$$d_{\rm Som}^G: \{(M, \mathcal{O}_M) \} \longrightarrow \mathcal{N}(G^\vee)$$
that extends the Barbasch--Vogan dual map $d_{\rm BV}: \mathcal{N}(G) \to \mathcal{N}^{\rm spe}(G^\vee)$, where $\mathcal{N}^{\rm spe}(G) \subset \mathcal{N}(G)$ denotes the subset of special nilpotent orbits of $G$. This we explain in more detail now.

First, one has the Lusztig--Spaltenstein map $d_{\rm LS}^G: \mathcal{N}(G) \longrightarrow \mathcal{N}^{\rm spe}(G)$. Moreover, one has a natural isomorphism 
$$\tau_{G, G^\vee}: W(G) \longrightarrow W(G^\vee)$$
of the two Weyl groups, which induces a bijection (using the same notation)
$$\tau_{G, G^\vee}: {\rm Irr}(W(G)) \longrightarrow {\rm Irr}(W(G^\vee)).$$
It furthermore induces a bijection (by abuse of notation again)
$$\tau_{G, G^\vee}: \mathcal{N}^{\rm spe}(G) \longrightarrow \mathcal{N}^{\rm spe}(G^\vee).$$
One has $d_{\rm BV}:= \tau_{G, G^\vee} \circ d_{\rm LS}^G$. All these maps mentioned above fit into the following commutative diagram:
$$\begin{tikzcd}
\mathcal{N}^{\rm spe}(G) \ar[r, "{\tau_{G,G^\vee}}"] & \mathcal{N}^{\rm spe}(G^\vee)  \ar[r, hook] & \mathcal{N}(G^\vee) \\
& \mathcal{N}(G) \ar[lu, "{d_{\rm LS}^G}"] \ar[u, "{d_{\rm BV}^G}"] \ar[r, hook] & {\{(M, \mathcal{O}_M) \} } \ar[u, "{d_{\rm Som}^G}"] .
\end{tikzcd}
$$
Here the bottom injection is provided by the Bala--Carter classification of nilpotent orbits, asserting that every $\mathcal{O} \in \mathcal{N}(G)$ corresponds uniquely to a $G$-conjugacy class of pair $(L, \mathcal{O}_L)$, where $L \subseteq G$ is a Levi subgroup and $\mathcal{O}_L$ a distinguished orbit for $L$. In this case, we write
$${\rm sat}_L^G(\mathcal{O}_L):= G \cdot \mathcal{O}_L \in \mathcal{N}(G)$$ for the nilpotent orbit thus obtained. We remark also that Achar actually gave a further extension of $d_{\rm Som}^G$ in \cite{ac03}.


Let $M \subset G$ be a pseudo-Levi subgroup and $\mathcal{O}_{M}$ a distinguished orbit for $M$. Then it was shown by Sommers \cite[\S 6]{Sommers} that 
\begin{equation} \label{E:dSom}
d_{\rm Som}^{G}(M, \mathcal{O}_{M}) = \mathcal{O}_{\rm Spr}^{G^\vee} \circ \tau_{G, G^\vee} \circ j_{W(M)}^{W(G)} \circ {\rm Spr}_{M, \mathbf{1}}^{-1} \circ d_{\rm LS}^{M}(\mathcal{O}_{M}).
\end{equation}

Note that the dual group $M^\vee$ of a pseudo-Levi $M$ of $G$ in general may not be a pseudo-Levi subgroup of $G^\vee$. However, one has a natural isomorphism $W(M) \simeq W(M^\vee)$. Also, for every $\sigma \in {\rm Irr}(W(M))$, one  $\sigma^\vee:= \tau_{M, M^\vee}(\sigma) \in {\rm Irr}(W(M^\vee))$ satisfying  
\begin{equation} \label{E:transp}
j_{W(M^\vee)}^{W(G^\vee)} (\sigma^\vee) = \tau_{G, G^\vee} \circ j_{W(M)}^{W(G)} (\sigma).
\end{equation}

\begin{Prop} \label{P:Som}
Keep notations as above. Assume $\sigma \in {\rm Irr}(W(M))$ is a special representation of $W(M)$. Furthermore, let $L\subset M$ be a Levi subgroup of $M$ and $\mathcal{O}_L$ be a distinguished orbit of $L$ such that $d_{\rm LS}^M \circ \mathcal{O}_{\rm Spr}^M(\sigma) = {\rm sat}_L^M(\mathcal{O}_L)$.
Then 
$$\mathcal{O}_{\rm Spr}^{G^\vee} \circ j_{W(M^\vee)}^{W(G^\vee)}(\sigma^\vee) = d_{\rm Som}^G(L, \mathcal{O}_L).$$
\end{Prop}
\begin{proof}
Note that $L$ is a Levi subgroup of $M$ and $M$ is a pseudo-Levi subgroup of $G$. Thus, $L$ is a pseudo-Levi subgroup of $G$ and the right hand side of the above equality is well-defined.

Now we have from \eqref{E:transp} that
$$\mathcal{O}_{\rm Spr}^{G^\vee} \circ j_{W(M^\vee)}^{W(G^\vee)}(\sigma^\vee) = \mathcal{O}_{\rm Spr}^{G^\vee} \circ \tau_{G, G^\vee} \circ j_{W(M)}^{W(G)} (\sigma).$$
Since $\sigma$ is a special representation of $W(M)$ and $d_{\rm LS}^M$ is a bijection when restricted to the special orbits of $M$, it follows that
$$\sigma = {\rm Spr}_{M, \mathbf{1}}^{-1} \circ d_{\rm LS}^M \circ {\rm sat}_{L}^M(\mathcal{O}_L).$$
Consider the induction of nilpotent orbits ${\rm ind}_L^M: \mathcal{N}(L) \longrightarrow \mathcal{N}(M)$. One has
$$d_{\rm LS}^M \circ {\rm sat}_L^M(\mathcal{O}_L) = {\rm ind}_L^M \circ d_{\rm LS}^L(\mathcal{O}_L),$$
see \cite[Theorem 8.3.1]{CM}. On the other hand, the Springer correspondence intertwins the maps  $j_{W(L)}^{W(M)}$ and ${\rm ind}_L^M$ (see \cite{LuSp1}), i.e.,
$${\rm Spr}_{M, \mathbf{1}}^{-1} \circ {\rm ind}_L^M = j_{W(L)}^{W(M)} \circ {\rm Spr}_{L, \mathbf{1}}^{-1}.$$
Combining the above four equalities and using  $j_{W(M)}^{W(G)} \circ j_{W(L)}^{W(M)} = j_{W(L)}^{W(G)}$, we get that
$$
\begin{aligned}
    \mathcal{O}_{\rm Spr}^{G^\vee} \circ j_{W(M^\vee)}^{W(G^\vee)}(\sigma^\vee) & = \mathcal{O}_{\rm Spr}^{G^\vee} \circ \tau_{G, G^\vee} \circ j_{W(L)}^{W(G)} \circ {\rm Spr}_{L, \mathbf{1}}^{-1} \circ d_{\rm LS}^L (\mathcal{O}_L) \\
     & = d_{\rm Som}^G(L, \mathcal{O}_L),
\end{aligned}$$
where the last equality follows from \eqref{E:transp}. This completes the proof.
\end{proof}



Now Proposition \ref{P:Som} coupled with \eqref{E:jos} enable us to compute $\mathcal{O}_{{\rm Ann}(L(\lambda))}$ more efficiently, by utilizing the tabulated data for the Sommers duality given in \cite{Sommers, ac03}. More precisely, we apply Proposition \ref{P:Som} with the role of $G$ and $G^\vee$ switched, while $\pi_{w_\lambda}$ substitutes the input in the left hand side of the equality in Proposition \ref{P:Som}.  This will enable us to completely determine the column of $\mathcal{O}_{{\rm Ann}(L(\lambda))}$ in the tables given as in the Appendix. In the remaining of this section, we give an elaboration on this.


% For  a totally ordered set $ \Gamma $, we  denote by $ \mathrm{Seq}_n (\Gamma)$ the set of sequences $ x=(x_1,x_2,\cdots, x_n) $   of length $ n $ with $ x_i\in\Gamma $. In this paper, we usually take $\Gamma$ to be $\mathbb{Z}$ or a coset of $\mathbb{Z}$ in $\mathbb{C}$. 

In \cite{BXX}, by using the RS algorithm for $x\in  \mathrm{Seq}_n (\Gamma)$, we can get a Young tableau $P(x)$ and a partition ${\bf p}(x)=\mathrm{sh}(P(x))=p(x)=[p_1,p_2,...,p_N]$, where $p_i$ is the number of boxes in the $i$-th
row of $P(x)$.  
% For  $ x=(x_1,x_2,\cdots,x_n)\in \mathrm{Seq}_n (\Gamma) $, set
% \begin{equation*}
% 	\begin{aligned}
% 		{x}^-=&(x_1,x_2,\cdots,x_{n-1}, x_n,-x_n,-x_{n-1},\cdots,-x_2,-x_1).
% 		%{}^-{x}=&(-x_n,-x_{n-1},\cdots, -x_2,-x_1,x_1,x_2,\cdots, x_{n-1}, x_n).
% 	\end{aligned}
% \end{equation*}
Recall that in \cite{BMW}, we can get a special partition $H(p(x^-))$ from the given Young tableau $P(x^-)$ or the partition ${\bf p}(x^-)=p(x^-)$ by using the H-algorithm defined in \cite{BMW}. We denote it by $p_{X}(x^-)^s$ when it is a special partition of type $X$ (for $X=B,C$ or $D$). Roughly speaking, $H(p(x^-))=p_{X}(x^-)^s$ is the unique  special partition which has the same odd (resp. even) boxes for type $B$
and $C$ (resp. type $D$) with the partition $p(x^-)$.



Recall that  $\mathfrak{g}$ is the Lie algebra of $G$. Let $\mathfrak{h}$ be a Cartan subalgebra.
	Let $\Phi^+\subset\Phi$ be the set of positive roots determined by a Borel subalgebra $\mathfrak{b}$ of $\mathfrak{g}$. We have a Cartan decomposition $\mathfrak{g}=\mathfrak{n}\oplus\mathfrak{h}\oplus\mathfrak{n}^-$ with $\mathfrak{b}=\mathfrak{n}\oplus \mathfrak{h}$. Denote by $\Delta$ the set of simple roots in $\Phi^+$.
By the results in \S \ref{nonint}, we may write $\Phi_{[\lambda]}\simeq \Phi_1\times \Phi_2\times \cdots\times \Phi_k$ for any $\lambda\in \mathfrak{h}^*$ and $$\lam'=\phi(\lam|_{\mathfrak{h}_{\Phi_{[\lambda]}}})=\prod\limits_{1\lest i\lest k} \phi(\lambda|_{\mathfrak{h}_{\Phi_{[\lambda]}}})|_{{\Phi_i}}=\prod\limits_{1\lest i\lest k} \lam'|_{{\Phi_i}}.$$ Thus ${\lam}'|_{\Phi_i}$ is an integral weight of type $\Phi_i$. When some $\Phi_i$ is of type $A$, we can get a special partition ${\bf p}_{i}:=p(\lam'|_{{\Phi_i}})$ of type $A$ which comes from the Young tableau $P(\lam'|_{{\Phi_i}})$.
When some $\Phi_i$ is of classical type $X_i$ ($X_i=B,C$ or $D$), by using the H-algorithm in \cite{BMW}, we can get two special partitions ${\bf p}_{i}:=p_{X_i}((\lam'|_{{\Phi_i}})^-)^s$ of type $X_i$ and  ${\bf p}^{\vee}_{i}:=p_{X_i^{\vee}}((\lam'|_{{\Phi_i}})^-)^s$ of type $X_i^{\vee}$, both of which come from the Young tableau $P((\lam'|_{{\Phi_i}})^-)$. When some $\Phi_i$ is of exceptional type, we can use  Lemma \ref{find-w-lambda} and PyCox to get  the corresponding character $\chi_i$ of $w_{\lam'|_{{\Phi_i}}}$ and the corresponding special  nilpotent orbit (Bala--Carter label) $\mathcal{O}_{\chi_i}$. 

For classical types, we identify a special partition ${\bf p}$ (corresponding to a special nilpotent orbit $\mathcal{O}_w$) and its corresponding special representation $\pi_w=\pi_{\bf p}$ of the Weyl group by the Springer correspondence. We write ${\bf p}^t$ for the transpose of a partition ${\bf p}$. Also, for a partition ${\bf p}$ of classical $X$-type, one has $d_{\rm LS}({\bf p}) = {\bf p}^t_X$.

From now on, we identify a root subsystem  of $\Phi$ with the corresponding analytic subgroup of the given Lie group $G$. Thus $\Phi$ is identified with $G$. We may write $L\subseteq \Phi$ or $\mathcal{O}_L\subseteq \Phi$ to mean that $L$ is a Levi subgroup of $G$ or $\mathcal{O}_L$ is the distinguished orbit for $L$.

%Then we have the following result.


\begin{Cor}
    Keep notations as above. Suppose the root system of $\mathfrak{g}$ is $\Phi$. Let $L(\lam)$ be a highest weight module of $\mathfrak{g}$. Denote
$H:=\Phi_{[\lambda]}\simeq \Phi_1 \times \Phi_2 \times ... \times \Phi_k$. Write $\lam'=\phi(\lam|_{\mathfrak{h}_{\Phi_{[\lambda]}}})=\prod\limits_{1\lest i\lest k} \phi(\lambda|_{\mathfrak{h}_{\Phi_{[\lambda]}}})|_{{\Phi_i}}=\prod\limits_{1\lest i\lest k} \lam'|_{{\Phi_i}}$. We have the following possibilities:
\begin{enumerate}
    \item Suppose $\Phi_i$ is of classical type $X_i$ and ${\bf p}_{i}$ is not a  very even partition of type $D$  for all $1\lest i \lest k$.   Then $\pi_{w_{\lam}}=\prod\limits_{1\lest i\lest k}{\bf p}_{i} \in {\rm Irr}(W(H)) $, $\pi^{\vee}_{w_{\lam}}=\prod\limits_{1\lest i\lest k}{\bf p}^{\vee}_{i} \in {\rm Irr}(W(H^{\vee})) $, and 
$$d_{\rm LS}^{H^\vee}(\mathcal{O}(\pi^{\vee}_{w_{\lam}}))=\prod\limits_{1\lest i\lest k}({\bf p}^{\vee}_{i})^t_{X_i} = {\rm sat}_{L^\vee}^{H^\vee} (\mathcal{O}_{L^\vee}),$$
where $L^\vee \subset \Phi^\vee$ and $\mathcal{O}_{L^\vee}$ a distinguished orbit of $L^\vee$.
% It follows from Proposition \ref{P:Som} that 
Also we have
$$\mathcal{O}_{{\rm Ann}(L(\lambda))} = d_{\rm Som}^{G^\vee}(L^{\vee}, \mathcal{O}_{L^\vee}).$$

\item Suppose $\Phi_i$ is of classical type $X_i$ and simply-laced  for all $1\lest i \lest k$ and one unique ${\bf p}_{i_0}$ is  a  very even partition of type $D$. Then the numeral of $\mathcal{O}_L$ is the same with the numeral of $\pi_{{\bf p}_{i_0}}$.
% Then the numeral $Y$ ($I$  or $II$) of ${\bf p}_{i_0}$ or $\pi_{w_{\lam'|_{{\Phi_{i_0}}}}}$ is determined by the numeral of the Kazhdan-Lusztig two-sided cell containing $w_{\lam'|_{{\Phi_{i_0}}}}$. The numeral of $\pi^t_{w_{\lam}}=(\pi^{\vee}_{w_{\lam}})^t$ is the same with $Y$.


\item Suppose there is a unique $\Phi_{i_0}$ of exceptional type $E$ and all other $\Phi_i$ is of classical type $X_i$ and simply-laced.
Suppose that $\chi_{i_0}$  is the corresponding character for $w_{\lam'|_{{\Phi_{i_0}}}}$ with Bala--Carter label $L_{i_0}$. Then $\pi^{\vee}_{w_{\lam}} \simeq \pi_{w_{\lam}}=\chi_{i_0}\times \prod\limits_{1\lest i\neq i_0\lest k}{\bf p}_{i}$ and $\mathcal{O}(\pi_{w_\lambda}^\vee)=L_{i_0}\times \prod\limits_{1\lest i\neq i_0\lest k}{\bf p}_{i}$. 
Also we have
$d_{\rm LS}^{H^\vee}(\mathcal{O}(\pi_{w_\lambda}^\vee)) = {\rm sat}_{L^\vee}^{H^\vee}(\mathcal{O}_{L^\vee})$ 
for a distinguished orbit $\mathcal{O}_{L^\vee}$ of $L^\vee$, 
and thus 
$$\mathcal{O}_{{\rm Ann}(L(\lambda))} = d_{\rm Som}^{G^\vee}(L^{\vee}, \mathcal{O}_{L^\vee}).$$

\end{enumerate}
\end{Cor}




% Here for illustration purpose, we give two examples extracted from Table \ref{} and \ref{}.



Recall that we use $\Delta=\{\alpha_1,\dots,\alpha_n\} $ to denote the simple system of the root system $\Phi$ for a simple Lie algebra $\mathfrak{g}$. Let $\tilde{\Delta}=\Delta\cup \{\alpha_0\}$, where  $-\alpha_0=\beta$ is the highest root of $\Phi$. We write $\theta=\sum_{1\leq i\leq n}c_i\alpha_i$ and set $c_0=1$. For any proper subset $J\subset \tilde{\Delta}$, let $d_J$ be the greatest common divisor of those $c_i$ for which $\alpha_i\in \tilde{\Delta}-J$.




The following result without proof will be used in the determination of the Bala--Carter label for a given $(L,\mathcal{O}_L)$. Some details can be found in \cite{Ca85, som-98}.

\begin{Lem}\label{prime1-2}
Let ${\bf p} $ be a partition corresponding to a nilpotent orbit of classical type $(L,\mathcal{O}_L)$. Then some correspondences are given in   Table \ref{label}.


For type $E_7$ and $E_8$, it is possible that there are two in-equivalence classes $L', L''$ (both of type $A_r$) associated with $J', J'' \subseteq \tilde{\Delta}$ respectively. Since such $L$ is of type $A_r$, a distinguished $\mathcal{O}_L$ is the regular orbit of $L$. Here $L'$ in 
the Bala--Carter label indicates that $J'$ is conjugate to a subset of $A_k\subset E_k$ ($k=7,8$); otherwise  we have $L''$. Moreover, for $E_8$ we have that $d_J=1$ if and only if the decoration is given by $L'$.

    
\end{Lem}

\begin{Rem}
    In Table \ref{primes} we elaborate on the different conventions and notation used in literature regarding very even orbits and their associated Weyl group representations. The notation is either self-explanatory, or a precise reference is given. 
    %For example, we use $([\bf u],[\bf u])$ to denote the bipartition corresponding to the partition ${\bf p}$ and $\pi$ to denote the representation corresponding to $([\bf u],[\bf u])$.


\begin{table}[htbp]\caption{Conventions for very even orbits and representations of $\mathfrak{so}(2n,\mathbb{C})$}
	\label{primes}
	\centering
	\renewcommand{\arraystretch}{1.5}
	\setlength\tabcolsep{5pt}
	
	\begin{tabular}{|c|c|c|}
		%\begin{tabular}{|l|l|l|}
		\hline {Our notation}  &${\bf p}^I$ & ${\bf p}^{II}$ \\ \hline
       
        \cite[p. 89]{lusztig1984char} & $\pi'$ & $\pi''$ \\
        \hline

        \cite[p. 170]{GP} & $([\bf u],+)$ & $([\bf u],-)$ \\
        \hline

 PyCox & $([\bf u],-)$ & $([\bf u],+)$ \\
        \hline
 \cite[p. 83]{CM} & $\mathcal{O}^I$ if $n\equiv 0 ~(\mathrm{mod}~ 4)$ & $\mathcal{O}^{II}$ if $n\equiv 0 ~(\mathrm{mod}~ 4)$\\
 & $\mathcal{O}^{II}$ if $n\equiv 2 ~(\mathrm{mod}~ 4)$ & $\mathcal{O}^{I}$ if $n\equiv 2 ~(\mathrm{mod}~ 4)$\\
\hline

\cite{Alv} & $([\bf u],[\bf u])^+$ if $n\equiv 0 ~(\mathrm{mod}~ 4)$ & $([\bf u],[\bf u])^-$ if $n\equiv 0 ~(\mathrm{mod}~ 4)$\\
 & $([\bf u],[\bf u])^-$ if $n\equiv 2 ~(\mathrm{mod}~ 4)$ & $([\bf u],[\bf u])^+$ if $n\equiv 2 ~(\mathrm{mod}~ 4)$\\
\hline

 
 \end{tabular}
	\bigskip
	
\end{table}


















\end{Rem}

The annihilator variety $V(\Ann(L(\lam)))$ can be determined by $\gkd L(\lam)$ if there is only one nilpotent orbit with dimension equal to $2\gkd L(\lam)$. Thus we only need to consider the cases when there is more than one nilpotent orbit with dimension equal to $2\gkd L(\lam)$.
For $L(\lam)$ with $\lambda$ integral, the results are given 
 in Propositions \ref{f4}, \ref{g2}, \ref{E6}, \ref{E7} and \ref{E8}. 
 
\subsection{Type \texorpdfstring{$G_2$}{}}
 
 If $L(\lam)$ is a simple  highest weight module of $G_2$, the annihilator variety $V(\Ann(L(\lam)))$ can be determined uniquely by $\gkd L(\lam)$ since all the nilpotent orbits have different dimensions, see \cite[\S 8.4]{CM}. Thus the annihilator variety $V(\Ann(L(\lam)))$ for type $G_2$ can be determined by our algorithm of GK dimensions in \S \ref{a-value-compu}.


\subsection{Type \texorpdfstring{$F_4$}{}}
 We only need to consider the nonintegral cases and the results are as follows.


\begin{Thm}\label{f4-av}
Let $L(\lam)$ be a simple nonintegral highest weight module of $F_4$  such that $\Phi_{[\lam]}^\vee$ is pseudo-maximal. Write $\lam'=\phi(\lam|_{\mathfrak{h}_{\Phi_{[\lambda]}}})=\prod\limits_{1\lest i\lest k} \phi(\lambda|_{\mathfrak{h}_{\Phi_{[\lambda]}}})|_{{\Phi_i}}=\prod\limits_{1\lest i\lest k} \lam'|_{{\Phi_i}}$.
Then the following holds:
    \begin{enumerate}
	\item Suppose $H=\Phi_{[\lam]}\simeq C_4$.
We only need to consider the following three cases:
$$\gkd L(\lam)=15, 18, 21.$$
In the above three cases, we have $H^{\vee}=B_4$,  
$\pi_{w_{\lam}}={\bf p}_{1}=p_{C}((\lam')^-)^s$ and $\pi^{\vee}_{w_{\lam}}={\bf p}^{\vee}_{1}=p_{B}((\lam')^-)^s=({\bf p}_{1}+[1,0,\cdots,0])_B$. Then $$d_{\rm LS}^{H^\vee}(\mathcal{O}(\pi^{\vee}_{w_{\lam}}))=({\bf p}^{\vee}_{1})^t_{B_4} = {\rm sat}_{L^\vee}^{H^\vee} (\mathcal{O}_{L^\vee}),$$
where $L^\vee \subset \Phi^\vee$ and $\mathcal{O}_{L^\vee}$ is a distinguished orbit of $L^\vee$. The explicit information about $\mathcal{O}_{{\rm Ann}(L(\lambda))}$ can be found in Table \ref{tabc4}.


% When $ \gkd L(\lam)=15 $, we have $V(\Ann (L(\lam)))=\overline{\mathcal{O}}(A_2)$ if and only if $\pi_{w_{\lam}}=([1^3],[1])$ $(\text{equivalently~} \mathcal{O}(\pi_{w_{\lam}})=[2,2,1^4])$.



% When $ \gkd L(\lam)=18 $, we have $V(\Ann (L(\lam)))=\overline{\mathcal{O}}(B_2)$ if and only if $\pi_{w_{\lam}}=([1,1],[1,1])$ $(\text{equivalently~} \mathcal{O}(\pi_{w_{\lam}})=[2^4])$.

% When $ \gkd L(\lam)=21 $, we have $V(\Ann (L(\lam)))=\overline{\mathcal{O}}(B_3)$ if and only if $\pi_{w_{\lam}}=([2,1],[1])$ $(\text{equivalently~} \mathcal{O}(\pi_{w_{\lam}})=[4,2,2])$.


\item Suppose $H=\Phi_{[\lam]}\simeq A_2\times \tilde{A}_2$. We only need to consider the following two cases:
$$\gkd L(\lam)=18, 21.$$
In the above two cases, we have $H^{\vee}=\tilde{A}_2 \times A_2$ and 
$\pi^{\vee}_{w_{\lam}}={\bf p}_{1}\times {\bf p}_{2}=p_{A}(\lam'|_{A_2})\times p_{A}(\lam'|_{\tilde{A}_2})$. Then $$\mathcal{O}(\pi^{\vee}_{w_{\lam}})^t={\bf p}^t_{1}\times {\bf p}^t_{2} = {\rm sat}_{L^\vee}^{H^\vee} (\mathcal{O}_{L^\vee}),$$
where $L^\vee \subset \Phi^\vee$ and $\mathcal{O}_{L^\vee}$ a distinguished orbit of $L^\vee$. The explicit information about $\mathcal{O}_{{\rm Ann}(L(\lambda))}$ can be found in Table \ref{taba2a2}.
  
% When $ \gkd L(\lam)=18 $, we have $V(\Ann (L(\lam)))=\overline{\mathcal{O}}(\tilde{A}_2+A_1)$ if and only if $\mathcal{O}(\pi_{w_{\lam}})=[1^3]\times [1^3]$.

% When $ \gkd L(\lam)=21 $, we have $V(\Ann (L(\lam)))=\overline{\mathcal{O}}(B_3)$ if and only if $\mathcal{O}(\pi_{w_{\lam}})=[1^3]\times [3]$, and $V(\Ann (L(\lam)))=\overline{\mathcal{O}}(C_3)$ if and only if $\mathcal{O}(\pi_{w_{\lam}})=[3]\times [1^3]$.





  \item Suppose $H=\Phi_{[\lam]}\simeq A_1\times \tilde{A}_3$. We only need to consider the following two cases:
$$\gkd L(\lam)=18, 21.$$
In the above two cases, we have $H^{\vee}=\tilde{A}_1 \times A_3$ and 
$\pi^{\vee}_{w_{\lam}}={\bf p}_{1}\times {\bf p}_{2}=p_{A}(\lam'|_{A_1})\times p_{A}(\lam'|_{\tilde{A}_3})$. Then $$\mathcal{O}(\pi^{\vee}_{w_{\lam}})^t={\bf p}^t_{1}\times {\bf p}^t_{2} = {\rm sat}_{L^\vee}^{H^\vee} (\mathcal{O}_{L^\vee}),$$
where $L^\vee \subset \Phi^\vee$ and $\mathcal{O}_{L^\vee}$ a distinguished orbit of $L^\vee$. The explicit information about $\mathcal{O}_{{\rm Ann}(L(\lambda))}$ can be found in Table \ref{taba1a3}.



% When $ \gkd L(\lam)=18 $, we have $V(\Ann (L(\lam)))=\overline{\mathcal{O}}(B_2)$ if and only if $\mathcal{O}(\pi_{w_{\lam}})=[2]\times [1^4]$.

% When $ \gkd L(\lam)=21 $, we have  $V(\Ann (L(\lam)))=\overline{\mathcal{O}}(C_3)$ if and only if $\mathcal{O}(\pi_{w_{\lam}})=[2]\times [2,1,1]$.

    
  







    
    
    \item Suppose $H=\Phi_{[\lam]}\simeq B_3\times \tilde{A}_1$. We only need to consider the following two cases:
$$\gkd L(\lam)=15, 21.$$
In the above two cases, we have $H^{\vee}={C}_3 \times A_1$, 
$\pi_{w_{\lam}}={\bf p}_{1}\times {\bf p}_{2}=p_{B}((\lam'|_{B_3})^-)^s\times p_{A}(\lam'|_{\tilde{A}_1})$ and $\pi^{\vee}_{w_{\lam}}={\bf p}^{\vee}_{1}\times {\bf p}_{2}=p_{C}((\lam'|_{B_3})^-)^s\times p_{A}(\lam'|_{\tilde{A}_1})$. Then $$d_{\rm LS}^{H^\vee}(\mathcal{O}(\pi^{\vee}_{w_{\lam}}))=({\bf p}^{\vee}_{1})^t\times {\bf p}^t_{2} = {\rm sat}_{L^\vee}^{H^\vee} (\mathcal{O}_{L^\vee}),$$
where $L^\vee \subset \Phi^\vee$ and $\mathcal{O}_{L^\vee}$ a distinguished orbit of $L^\vee$. The explicit information about $\mathcal{O}_{{\rm Ann}(L(\lambda))}$ can be found in Table \ref{tabb3a1}.

% When $ \gkd L(\lam)=15 $, we have $V(\Ann (L(\lam)))=\overline{\mathcal{O}}(\tilde{A}_2)$ if and only if $\pi_{w_{\lam}}=(\emptyset,[1^3])\times [2]$ $(\text{equivalently~}\mathcal{O}(\pi_{w_{\lam}})=[1^7]\times [2])$.

% When $ \gkd L(\lam)=21 $, we have $V(\Ann (L(\lam)))=\overline{\mathcal{O}}(C_3)$ if and only if $\pi_{w_{\lam}}=([1,1],[1])\times [2]$ $(\text{equivalently~}\mathcal{O}(\pi_{w_{\lam}})=[3,2,2]\times [2])$, and $V(\Ann (L(\lam)))=\overline{\mathcal{O}}(B_3)$ if and only if $\pi_{w_{\lam}}=([1],[2])\times [1,1]$ $(\text{equivalently~}\mathcal{O}(\pi_{w_{\lam}})=[3,3,1]\times [1,1])$.

   



 
\end{enumerate}
\end{Thm}
\begin{proof}
We give some details for (1), (4) and omit the argument for  (2) and (3) since the latter is similar and easier.

For (1), we have
$H=\Phi_{[\lambda]}\simeq C_4$ with $H^\vee=B_4$ being a pseudo-Levi subgroup of $F_4$. From $\lam'$, by using the H-algorithm in \cite{BMW}, we can get two special partitions $\pi_{w_{\lam}}={\bf p}_{1}=p_C((\lam')^-)^s$ of type $C_4$ and  $\pi^{\vee}_{w_{\lam}}={\bf p}^{\vee}_{1}=p_B((\lam')^-)^s$ of type $B_4$, which both come from the Young tableau $P((\lam')^-)$. Then $p_C((\lam')^-)^s$ will give us the special bipartition $\pi_{w_\lambda} = ({\bf d}, {\bf f})   \in {\rm Irr}(W(H))$. 
% We see that $H^\vee = B_4$ is a pseudo-Levi subgroup of $F_4^\vee$. 
Consider the representation $\sigma^\vee= \tau_{H, H^\vee}(\sigma) = ({\bf d}, {\bf f}) \in {\rm Irr}(W(H^\vee))$. A simple computation gives that
$$d_{\rm LS}^{H^\vee} \circ \mathcal{O}_{\rm Spr}^{H^\vee}(\sigma^\vee) = (p_B((\lam'))^-)^s)^t = {\rm sat}_{L^\vee}^{H^\vee} (\mathcal{O}_{L^\vee}),$$
where $L^\vee \subset F_4^\vee$ and $\mathcal{O}_{L^\vee}$ is a distinguished orbit of $L^\vee$.
It follows from Proposition \ref{P:Som} that $$\mathcal{O}_{{\rm Ann}(L(\lambda))} = d_{\rm Som}^{F_4^\vee}(L^{\vee}, \mathcal{O}_{L^\vee}).$$


  When $ \gkd L(\lam)=15, 18 $ or $21$,  there are more than one nilpotent orbits with
dimension equal to $2 \gkd L(\lam)$.
% the Bala--Carter label of the nilpotent orbit in $V(\Ann (L(\lam)))$ is determined by \cite[Table D.3]{GP}. When $ \gkd L(\lam)=18 $, the Bala--Carter label of the nilpotent orbit in $V(\Ann (L(\lam)))$ is determined by Alvis \cite{Alv}.
For example, when $ \gkd L(\lam)=15$, we should have $\aff(w_{\lam})=9$. By Lemma \ref{a-value}, there is only one special partition of type $C_4$ whose $\aff$ value is $9$: ${\bf p}=[2,2,1^4]$.
  The corresponding special partition of type $B_4$ is ${\bf p}^{\vee}=({\bf p}+[1,0,\cdots,0])_B=[3,1^6]$. Thus $({\bf p}^{\vee})^t=[7,1,1]={\rm sat}_{L^\vee}^{H^\vee} (\mathcal{O}_{L^\vee})$, where $\mathcal{O}_{L^\vee}=B_3$.  
  Thus from Proposition \ref{P:Som} and the LS duality in \cite[p. 440]{Ca85} or Table \ref{tabcharf4}, we have $V(\Ann (L(\lam)))=\overline{\mathcal{O}}(A_2)$ since $ d_{\rm Som}^{F_4^\vee}(L^{\vee}, \mathcal{O}_{L^\vee})=d_{\rm LS}^{F_4^\vee}(L^{\vee}, \mathcal{O}_{L^\vee})=A_2$ when $\mathcal{O}_{L^\vee}=B_3$.

The  argument for $\gkd L(\lam)=18$ and $21$ is similar. 



For (4) one has $H=\Phi_{[\lam]}\simeq B_3\times \tilde{A}_1$ and $H^{\vee}={C}_3 \times A_1$. By Lemma \ref{a-value}, we only need to consider  that $ \gkd L(\lam)=15 $ or $21$ since $\aff(w_{\lam})$ can not take the value $6$. For example, when $ \gkd L(\lam)=21$, we should have $\aff(w_{\lam})=3$. By Lemma \ref{a-value}, there are two special partitions of type $B_3\times {A}_1$ whose $\aff$ value is $3$: ${\bf p}=[3,3,1] \times [1,1]$ and ${\bf q}=[3,2,2] \times [2]$.
  The corresponding two special partitions of type $C_3\times A_1$ are:  ${\bf p}^{\vee}=[3,3] \times [1,1]$ and ${\bf q}^{\vee}=[2^3] \times [2]$. Thus $({\bf p}^{\vee})^t=[2^3]\times [2]={\rm sat}_{L^\vee}^{H^\vee} (\mathcal{O}_{L^\vee})$ where $\mathcal{O}_{L^\vee}=2A_1+\tilde{A}_1$, and $({\bf q}^{\vee})^t=[3,3]\times [1,1]={\rm sat}_{L^\vee}^{H^\vee} (\mathcal{O}_{L^\vee})$ where $\mathcal{O}_{L^\vee}=\tilde{A}_2$.  
  Thus from Proposition \ref{P:Som}, the LS duality in \cite[p. 440]{Ca85} and the Sommers duality in \cite[Table VII]{Sommers} or Table \ref{tabcharf4}, we have $V(\Ann (L(\lam)))=\overline{\mathcal{O}}(B_3)$ for $\pi_{\lam}={\bf p}$ since $ d_{\rm Som}^{F_4^\vee}(L^{\vee}, \mathcal{O}_{L^\vee})=B_3$ when $\mathcal{O}_{L^\vee}=2A_1+\tilde{A}_1$, and $V(\Ann (L(\lam)))=\overline{\mathcal{O}}(C_3)$ for $\pi_{\lam}={\bf q}$ since $ d_{\rm Som}^{F_4^\vee}(L^{\vee}, \mathcal{O}_{L^\vee})=d_{\rm LS}^{F_4^\vee}(L^{\vee}, \mathcal{O}_{L^\vee})=C_3$ when $\mathcal{O}_{L^\vee}=\tilde{A}_2$.

  The  argument for $\gkd L(\lam)=15$ is similar.




%     When $ \gkd L(\lam)=15 $ or $21$, the Bala--Carter label of the nilpotent orbit in $V(\Ann (L(\lam)))$ is determined by \cite[Table D.3]{GP}. When $ \gkd L(\lam)=18 $, the Bala--Carter label of the nilpotent orbit in $V(\Ann (L(\lam)))$ is determined by Alvis \cite{Alv}.
% For example, when $ \gkd L(\lam)=15$, we should have $\aff(w_{\lam})=9$. There is only one special $B$-symbol whose $\aff$ value is $9$: $\Lambda^s =\begin{pmatrix}
% 	0~1~2~4\\1~2~3
% 	\end{pmatrix} $, which gives  a bipartition
%  $({\bf p},{\bf q})=([1]],[1,1,1])$.
%  From \cite[ (4.5.4)]{lusztig1984char}, we have $\Lambda^s=\pi_{({\bf p},{\bf q})}=j_{W_3}^{W_4}\Lambda'$ with $\Lambda' =\begin{pmatrix}
% 	0~1~2~3\\1~2~3
% 	\end{pmatrix}.$ The corresponding bipartition of $\Lambda'$ is $({\bf p}',{\bf q}')=(\emptyset,[1,1,1])$. Thus from \cite[Table C.3 \& Table D.3]{GP} and \cite[p.428]{Ca85}, we have $V(\Ann (L(\lam)))=\overline{\mathcal{O}}(A_2)$.

\end{proof}

\begin{ex}Recall that in Example \ref{exf4},
	 $L(\lam)$ is a highest weight module of $F_4$ with  $\lambda=(4, 5, \frac{3}{2}, \frac{1}{2})$.  $\Phi_{[\lambda]}$ is a subsystem with simple roots $\{\alpha_2, \alpha_3, \alpha_4, \ep_2\}$ and $H=\Phi_{[\lambda]}\simeq C_4$. Thus  $H^\vee \simeq B_4$. Suppose the simple root system of $C_4$ is $\Delta=\{\beta_1=\ep_1-\ep_2, \beta_2=\ep_2-\ep_3, \beta_3=\ep_3-\ep_4,\beta_4=2\ep_4\}$. We define a map $\phi: \Phi_{[\lambda]} \rightarrow C_4$ such that $\phi(\varepsilon_2)=\beta_1, \phi(\alpha_4)=\beta_2, \phi(\alpha_3)=\beta_3, \phi(\alpha_2)=\beta_4$. Then we have $\phi(\lam)=\lam'=(9,-1,2,1)$. 
     % We give two methods of computing $\mathcal{O}_{{\rm Ann}(L(\lambda))}$.
 
 

% By Proposition \ref{pr:main1},  $L(\lambda)$ has $\gkd L(\lambda)=|\Phi^+|-\aff(w_{\lambda})=24-9=15$.
% and $V(\Ann (L(\lam)))=\overline{\mathcal{O}}(A_2)$.

% The first method is as follows. Note that $\lam'=(9,-1,2,1)$ will give us a $B$-symbol
% $\Lambda =\begin{pmatrix}
% 	0~1~2~3\\1~2~4
% 	\end{pmatrix} $. The corresponding special symbol is $\Lambda^s =\begin{pmatrix}
% 	0~1~2~4\\1~2~3
% 	\end{pmatrix} $, which gives  a bipartition
%  $({\bf p},{\bf q})=([1],[1,1,1])$.
%  From \cite[ (4.5.4)]{lusztig1984char}, we have $\Lambda=\pi_{({\bf p},{\bf q})}=j_{W_3}^{W_4}\Lambda'$ with $\Lambda' =\begin{pmatrix}
% 	0~1~2~3\\1~2~3
% 	\end{pmatrix}.$ The corresponding bipartition of $\Lambda'$ is $({\bf p}',{\bf q}')=(\emptyset,[1,1,1])$. Thus from \cite[Table C.3 \& Table D.3]{GP} and \cite[p. 428]{Ca85}, we have $V(\Ann (L(\lam)))=\overline{\mathcal{O}}(A_2)$.

% The second method is as follows. 
By using the RS algorithm in \cite{BXX} and the H-algorithm in \cite{BMW}, we have  $\pi_{w_\lambda}=p_C((\lam')^-)^s=[2,2,1^4]$ and $\pi^{\vee}_{w_\lambda}=p_B((\lam')^-)^s=[3,1^6]$.   A simple computation gives that
$$d_{\rm LS}^{H^\vee}(\mathcal{O}(\pi^{\vee}_{w_\lambda}))=(p_B((\lam')^-)^s)^t= [7,1,1] = {\rm sat}_{L^\vee}^{H^\vee} (\mathcal{O}_{L^\vee}^{\rm reg}),$$
where $\mathcal{O}_{L^\vee} = B_3$.
It follows from Proposition \ref{P:Som} and the LS duality in \cite[p. 440]{Ca85} or Table \ref{tabcharf4} that $\mathcal{O}_{{\rm Ann}(L(\lambda))} = d_{\rm Som}^{F_4^\vee}(L^{\vee}, \mathcal{O}_{L^\vee}) = d_{\rm LS}^{F_4^\vee}(L^{\vee}, \mathcal{O}_{L^\vee}) =\mathcal{O}(A_2)$ when $\mathcal{O}_{L^\vee}=B_3$.
\end{ex}





% \begin{ex}
% Consider $F_4$ and $\lambda$ such that $\Phi_{[\lambda]} = B_3 \times \tilde{A}_1:=H$, corresponding to the pseudo-Levi $H^\vee=C_3 \times A_1 \subseteq F_4^\vee$. We further assume that $\sigma:=\pi_{w_\lambda} = ([1, 1], [1]) \times [2]  \in {\rm Irr}(W(H))$. Consider the representation $\sigma^\vee= \tau_{H, H^\vee}(\sigma) = ([1, 1], [1]) \times [2] \in {\rm Irr}(W(H^\vee))$. A simple computation gives that
% $$d_{\rm LS}^{H^\vee} \circ \mathcal{O}_{\rm Spr}^{H^\vee}(\sigma^\vee) = [3,3] \times [1,1] = {\rm sat}_{L^\vee}^{H^\vee} (\mathcal{O}_{L^\vee}^{\rm reg}),$$
% where $L^\vee = \tilde{A}_2 \subset F_4^\vee$.
% It follows from Proposition \ref{P:Som} and \cite[Figure 2]{ac03} that $\mathcal{O}_{{\rm Ann}(L(\lambda))} = d_{\rm Som}^{F_4^\vee}(L^{\vee}, \mathcal{O}_{L^\vee}^{\rm reg}) = \mathcal{O}(C_3)$.
% \end{ex}

\subsection{Type \texorpdfstring{$E_6$}{}} We only need to consider the nonintegral cases and the results
are as following.
\begin{Thm}\label{e6-av}
Let $L(\lam)$ be a simple nonintegral highest weight module of $E_6$ such that $\Phi_{[\lam]}^\vee$ is pseudo-maximal. Write $\lam'=\phi(\lam|_{\mathfrak{h}_{\Phi_{[\lambda]}}})=\prod\limits_{1\lest i\lest k} \phi(\lambda|_{\mathfrak{h}_{\Phi_{[\lambda]}}})|_{{\Phi_i}}=\prod\limits_{1\lest i\lest k} \lam'|_{{\Phi_i}}$.
	Then
    the following holds:
    \begin{enumerate}
	\item Suppose $H=\Phi_{[\lam]}\simeq D_5$. We only need to consider the following two cases:
$$\gkd L(\lam)=30, 32.$$
In the above two cases, we have $H^{\vee}=D_5$ and   
$\pi^{\vee}_{w_{\lam}}={\bf p}_{1}=p_{D}((\lam')^-)^s$. Then $$d_{\rm LS}^{H^\vee} (\mathcal{O}(\pi_{w_{\lam}}^\vee))=(({\bf p}_{1})^t)_D = {\rm sat}_{L^\vee}^{H^\vee} (\mathcal{O}_{L^\vee}),$$
where $L^\vee \subset \Phi^\vee$. The explicit information about $\mathcal{O}_{{\rm Ann}(L(\lambda))}$ can be found in Table \ref{tabd5-1}.

% When $ \gkd L(\lam)=30 $,  we have $V(\Ann (L(\lam)))=\overline{\mathcal{O}}(A_4)$ if and only if $\pi_{w_{\lam}}=([1,1],[2,1])$ $(\text{equivalently~} \mathcal{O}(\pi_{w_{\lam}})=[3,3,2,2])$ or $\pi_{w_{\lam}}=([3,1,1],\emptyset)$ $(\text{equivalently~} \mathcal{O}(\pi_{w_{\lam}})=[5,1^5])$.





% When $ \gkd L(\lam)=32 $, we have $V(\Ann (L(\lam)))=\overline{\mathcal{O}}(D_5(a_1)$ if and only if $\pi_{w_{\lam}}=([2,1],[2])$ $(\text{equivalently~} \mathcal{O}(\pi_{w_{\lam}})=[4,4,1,1])$.

\item Suppose $H=\Phi_{[\lam]}\simeq A_5\times A_1$. 
We only need to consider the following two cases:
$$\gkd L(\lam)=30, 32.$$
In the above two cases, we have $H^{\vee}=A_5\times A_1$ and   
$\pi^{\vee}_{w_{\lam}}={\bf p}_{1}\times {\bf p}_{2}=p_{A}(\lam'|_{A_5})\times p_{A}(\lam'|_{{A}_1})$. Then $$\mathcal{O}(\pi^{\vee}_{w_{\lam}})^t={\bf p}^t_{1}\times {\bf p}^t_{2} = {\rm sat}_{L^\vee}^{H^\vee} (\mathcal{O}_{L^\vee}^{\rm reg})={\rm sat}_{L}^{H} (\mathcal{O}_{L}^{\rm reg}),$$
where $L^\vee \subset \Phi^\vee$. The explicit information about $\mathcal{O}_{{\rm Ann}(L(\lambda))}$ can be found in Table \ref{taba5a1}.




% When $ \gkd L(\lam)=30 $, we have $V(\Ann (L(\lam)))=\overline{\mathcal{O}}({A}_4)$ if and only if $\pi_{w_{\lambda}}=[3,1^3]\times [2]$, and $V(\Ann (L(\lam)))=\overline{\mathcal{O}}({D}_4)$ if and only if $\pi_{w_{\lambda}}=[2^3]\times[2]$.




% When $ \gkd L(\lam)=32 $, we have $V(\Ann (L(\lam)))=\overline{\mathcal{O}}(D_5(a_1))$ if and only if 
% $\pi_{w_{\lambda}}=[4,1,1]\times[1,1]$ or $\pi_{w_{\lambda}}=[3,2,1]\times[2]$, and $V(\Ann (L(\lam)))=\overline{\mathcal{O}}(A_5)$ if and only if $\pi_{w_{\lambda}}=[3,3]\times[1,1]$.


  \item Suppose $H=\Phi_{[\lam]}\simeq A_2\times A_2\times A_2$. 
We only need to consider the following two cases:
$$\gkd L(\lam)=30, 32.$$
In the above two cases, we have $H^{\vee}=A_2\times A_2\times A_2$ and   
$\pi^{\vee}_{w_{\lam}}={\bf p}_{1}\times {\bf p}_{2}\times {\bf p}_{3}=p_{A}(\lam'|_{A_2(1)})\times p_{A}(\lam'|_{{A}_2(2)}) \times p_{A}(\lam'|_{{A}_2(3)})$, where $A_2(i)$ denotes the $i$-th $A_2$ component  of $\Phi_{[\lam]}$. Then $$\mathcal{O}(\pi^{\vee}_{w_{\lam}})^t={\bf p}^t_{1}\times {\bf p}^t_{2}\times {\bf p}^t_{3} = {\rm sat}_{L^\vee}^{H^\vee} (\mathcal{O}_{L^\vee}^{\rm reg}),$$
where $L^\vee \subset \Phi^\vee$. The explicit information about $\mathcal{O}_{{\rm Ann}(L(\lambda))}$ can be found in Table \ref{taba2a2a2}.



            
% When $ \gkd L(\lam)=30 $, we have  $V(\Ann (L(\lam)))=\overline{\mathcal{O}}({D}_4)$ if and only if 
% $\pi_{w_{\lambda}}\simeq [1^3]\times[1^3]\times [3]$.

% When $ \gkd L(\lam)=32 $, we have $V(\Ann (L(\lam)))=\overline{\mathcal{O}}(D_5(a_1))$ if and only if 
% $\pi_{w_{\lambda}}\simeq [1^3]\times[2,1]\times [3]$.


\end{enumerate}
\end{Thm}

\begin{proof}
% (1) In this case
% $\Phi_{[\lambda]}\simeq D_5:=H$. Then the dual root system $H^\vee=D_5$ is that of a pseudo-Levi subgroup of $E_6$. From $\lam$, by using the H-algorithm in \cite{BMW}, we can get a special partition $p_D(\lam^-)^s$ of type $D_5$, which comes from the Young tableau $P(\lam^-)$. Then $p_D(\lam^-)^s$ will give us the special representation $\pi_{w_\lambda} = \{{\bf d}, {\bf f}\}   \in {\rm Irr}(W(H))$. We see that $H^\vee = D_5$ is a pseudo-Levi subgroup of $E_6^\vee=E_6$. Consider the representation $\sigma^\vee= \tau_{H, H^\vee}(\sigma) = \{{\bf d}, {\bf f}\} \in {\rm Irr}(W(H^\vee))$. A simple computation gives that
% $$d_{\rm LS}^{H^\vee} \circ \mathcal{O}_{\rm Spr}^{H^\vee}(\sigma^\vee) = ((p_D(\lam^-)^s)^t)_D = {\rm sat}_{L^\vee}^{H^\vee} (\mathcal{O}_{L^\vee}^{\rm reg}),$$
% where $L^\vee \subset E_6^\vee$.
% It follows from Proposition \ref{P:Som} that $$\mathcal{O}_{{\rm Ann}(L(\lambda))} = d_{\rm Som}^{E_6^\vee}(L, \mathcal{O}_{L^\vee}^{\rm reg}).$$


%   When $ \gkd L(\lam)=30 $ or $32$, the Bala--Carter label of the nilpotent orbit in $V(\Ann (L(\lam)))$ is determined by \cite[Table D.4]{GP}. 
% For example, when $ \gkd L(\lam)=30$, we should have $\aff(w_{\lam})=6$. By Lemma \ref{a-value}, there are two special partitions of type $D_5$ whose $\aff$ value is $9$: ${\bf p}=[5,1^5]$ and ${\bf q}=[3,3,2,2]$.
%   For ${\bf p}$, we have $({\bf p}^t)_D=[6,1^4]_D=[5,1^5]={\rm sat}_{L^\vee}^{H^\vee} (\mathcal{O}_{L^\vee}^{\rm reg})$, where $L^\vee=A_3 \subset F_4^\vee$.  
%   Thus from Proposition \ref{P:Som}, we have $V(\Ann (L(\lam)))=\overline{\mathcal{O}}(A_4)$ since $\mathcal{O}_{{\rm Ann}(L(\lambda))} = d_{\rm Som}^{F_4^\vee}(L, \mathcal{O}_{L^\vee}^{\rm reg})$.
% For ${\bf q}$, we have $({\bf q}^t)_D=[4,4,2]_D=[4,4,1,1]={\rm sat}_{L^\vee}^{H^\vee} (\mathcal{O}_{L^\vee}^{\rm reg})$, where $L^\vee=A_3 \subset F_4^\vee$.  
%   Thus from Proposition \ref{P:Som}, we have $V(\Ann (L(\lam)))=\overline{\mathcal{O}}(A_4)$ since $\mathcal{O}_{{\rm Ann}(L(\lambda))} = d_{\rm Som}^{F_4^\vee}(L, \mathcal{O}_{L^\vee}^{\rm reg})$.

The  argument is similar with that in Theorem \ref{f4-av}, and thus we omit the details.



\end{proof}


\subsection{Type \texorpdfstring{$E_7$}{}} 
Recall that in the case of type $D_{2m}$, a nilpotent orbit $\mathcal{O}$  may be associated with a very even partition ${\bf p}$ with the  numeral being I or II, see Table \ref{primes} for our convention. The reader is also referred to \cite{BMW} for more details. %In the case of type $E_7$,
%we only need to consider the nonintegral case and the results
%are as follows.

\begin{Thm}\label{e7-av}
Let $L(\lam)$ be a simple nonintegral highest weight module of $E_7$ such that $\Phi_{[\lam]}^\vee$ is pseudo-maximal. Write $\lam'=\phi(\lam|_{\mathfrak{h}_{\Phi_{[\lambda]}}})=\prod\limits_{1\lest i\lest k} \phi(\lambda|_{\mathfrak{h}_{\Phi_{[\lambda]}}})|_{{\Phi_i}}=\prod\limits_{1\lest i\lest k} \lam'|_{{\Phi_i}}$.
	Then
    the following holds:
    \begin{enumerate}
	

\item Suppose $H=\Phi_{[\lam]}\simeq A_7$. We only need to consider the following $12$ cases:
$$\gkd L(\lam)=42,47,48,50,51,53,54,55,56,57,59,60.$$
In the above $12$ cases, we have $H^{\vee}= A_7$ and   
$\pi^{\vee}_{w_{\lam}}={\bf p}_{1}=p_{A}(\lam')$. Then $$\mathcal{O}(\pi^{\vee}_{w_{\lam}})^t={\bf p}_{1}^t = {\rm sat}_{L^\vee}^{H^\vee} (\mathcal{O}_{L^\vee}^{\rm reg}),$$
where $L^\vee \subset \Phi^\vee$. The explicit information about $\mathcal{O}_{{\rm Ann}(L(\lambda))}$ can be found in Table \ref{taba7}. 

\item Suppose $H=\Phi_{[\lam]}\simeq A_5\times A_2$. We only need to consider the following $10$ cases:
$$\gkd L(\lam)=47,48,50,53,54,55,56,57,59,60.$$
In the above $10$ cases, we have $H^{\vee}=A_5\times A_2$ and   
$\pi^{\vee}_{w_{\lam}}={\bf p}_{1}\times {\bf p}_{2}=p_{A}(\lam'|_{A_5})\times p_{A}(\lam'|_{A_2})$. Then $$\mathcal{O}(\pi^{\vee}_{w_{\lam}})^t={\bf p}_{1}^t \times {\bf p}_{2}^t = {\rm sat}_{L^\vee}^{H^\vee} (\mathcal{O}_{L^\vee}^{\rm reg}),$$
where $L^\vee \subset \Phi^\vee$. The explicit information about $\mathcal{O}_{{\rm Ann}(L(\lambda))}$ can be found in Table \ref{taba5a2}.

\item Suppose $H=\Phi_{[\lam]}\simeq A_3\times A_3\times A_1$. We only need to consider the following $9$ cases:
$$\gkd L(\lam)=50,51,53,54,55,56,57,59,60.$$
In the above $9$ cases, we have $H^{\vee}=A_3\times A_3\times A_1$ and   
$\pi^{\vee}_{w_{\lam}}={\bf p}_{1}\times {\bf p}_{2}\times {\bf p}_{3}=p_{A}(\lam'|_{A_3(1)})\times p_{A}(\lam'|_{A_3(2)})\times p_{A}(\lam'|_{A_1})$. Then $$\mathcal{O}(\pi^{\vee}_{w_{\lam}})^t={\bf p}_{1}^t \times {\bf p}_{2}^t \times {\bf p}_{3}^t = {\rm sat}_{L^\vee}^{H^\vee} (\mathcal{O}_{L^\vee}^{\rm reg}),$$
where $L^\vee \subset \Phi^\vee$. The explicit information about $\mathcal{O}_{{\rm Ann}(L(\lambda))}$ can be found in Table \ref{taba3a3a1}.

\item Suppose $H=\Phi_{[\lam]}\simeq D_6\times A_1$. 
We only need to consider the following $12$ cases:
$$\gkd L(\lam)=42,47,48,50,51,53,54,55,56,57,59,60.$$
In the above $12$ cases, we have $H^{\vee}=D_6\times A_1$ and   
$\pi^{\vee}_{w_{\lam}}={\bf p}_{1}\times {\bf p}_{2}=p_{D}((\lam'|_{D_6})^-)^s\times p_{A}(\lam'|_{A_1})$. Then 
$$d_{\rm LS}^{H^\vee} (\mathcal{O}(\pi_{w_{\lam}}^\vee))=({\bf p}_{1}^t)_D \times {\bf p}^t_{2} = {\rm sat}_{L^\vee}^{H^\vee} (\mathcal{O}_{L^\vee}),$$
where $L^\vee \subset \Phi^\vee$ and $\mathcal{O}_{L^\vee}$ a distinguished orbit of $L^\vee$. The explicit information about $\mathcal{O}_{{\rm Ann}(L(\lambda))}$ can be found in Table \ref{tabd6a1} and Table \ref{tabd6a1-2}. When ${\bf p}_{1}$ is a very even partition, the numeral of $\mathcal{O}_{L^{\vee}}$ is the same with the numeral of ${\bf p}_{1}$.

\item Suppose $H=\Phi_{[\lam]}\simeq E_6$. We only need to consider the following $8$ cases:
$$\gkd L(\lam)=48,50,51,53,56,57,59,60.$$
In the above $8$ cases, we have $H^{\vee}=E_6$ and
$\pi^{\vee}_{w_{\lam}}={\chi}_{1}$. %Suppose that the Bala--Carter label of $\mathcal{O}({\chi}_{1})$ is $L_{1}$. Then $\mathcal{O}(\pi_{w_\lambda})=L_{1}$. 
Also we have
$d_{\rm LS}^{H^\vee}(\mathcal{O}(\pi_{w_\lambda}^\vee)) = {\rm sat}_{L^\vee}^{H^\vee}(\mathcal{O}_{L^\vee})$,
and $$\mathcal{O}_{{\rm Ann}(L(\lambda))} = d_{\rm Som}^{G^\vee}(L^{\vee}, \mathcal{O}_{L^\vee}),$$
where $L^\vee \subset \Phi^\vee$ and $\mathcal{O}_{L^\vee}$ a distinguished orbit of $L^\vee$. The explicit information about $\mathcal{O}_{{\rm Ann}(L(\lambda))}$ can be found in Table \ref{tabe6}.





















\end{enumerate}
\end{Thm}

%  From Table \ref{taba3a3a1}, we have the following result.
% \begin{Cor}Keep notations as above.
%     There is a special case:
%     $$J=\{-\beta, \alpha_7,\alpha_2\}\subset \tilde{\Delta}$$
%      such that  $J\subset \tilde{\Delta}$ is not a subset of $A_7\subset E_7$, but $J$ is conjugate a subset of $A_7\subset E_7$. 
% \end{Cor}
% \begin{proof}
%   From Tables \ref{taba7}--\ref{tabe6}, we can see that for all cases the Bala--Carter label of $L$ has one prime if the subset $J\subset \tilde{\Delta}$ is a subset of $A_7\subset E_7$, and for all cases it has two primes if $J\subset \tilde{\Delta}$ is not a subset of $A_7\subset E_7$ except for the case that $L$ corresponds to $[2,1,1]\times [2,1,1]\times [2]$.  In this case, we may choose $J=\{-\beta, \alpha_7,\alpha_2\}\subset \tilde{\Delta}$. The corresponding subsystem $\Phi_J$ has type $(3A_1)'$ or $(3A_1)''$. Our $L$ corresponds to this $J$. Denote $w=s_{\ep_4-\ep_2}s_{\ep_1+\ep_3}\in W(E_7)$. Then we have $w(J)=\{-\beta, \alpha_7,\alpha_5\}$, which is a subset of $A_7\subset E_7$. Therefore we must have 
%   $L=(3A_1)'$ by Lemma \ref{prime1-2}.
  
  
  
  
  
   
% \end{proof}






\begin{ex}
    Let $\mathfrak{g}=E_7$ and $L(\lam)$ be a highest weight module of $\mathfrak{g}$ with  $\lambda=(1,3,-5,-7,-9,-11, -\frac{1}{2}, \frac{1}{2})$. It is easy to verify that $\Phi_{[\lambda]}$ is a subsystem with simple roots $\{ \alpha_2,\alpha_3,\alpha_4,\alpha_5,\alpha_6,\alpha_7,\ep_8-\ep_7\}$. Therefore $H=\Phi_{[\lambda]}\simeq D_6\times A_1$. Suppose the simple root system of $D_6\times A_1$ is $\Delta=\{\beta_1=\ep_1-\ep_2, \beta_2=\ep_2-\ep_3, \beta_3=\ep_3-\ep_4,  \beta_4=\ep_4-\ep_5, \beta_5=\ep_5-\ep_6, \beta_6=\ep_5+\ep_6\}\times \{\gamma_1=e_1-e_2\}$.
    We define a map $\phi: \Phi_{[\lambda]} \rightarrow D_6\times A_1$ such that $\phi(\alpha_7)=\beta_1, \phi(\alpha_6)=\beta_2, \phi(\alpha_5)=\beta_3, \phi(\alpha_4)=\beta_4,\phi(\alpha_3)=\beta_5,\phi(\alpha_2)=\beta_6,\phi(\ep_8-\ep_7)=\gamma_1$.

   Thus by the algorithm in \S \ref{nonint},
   we have $$\lam'|_{D_6}=(-11,-9,-7,-5,3,1),$$
   and $$\lam'|_{A_1}= (\frac{1}{2},-\frac{1}{2}).$$
By using the RS algorithm in \cite{BXX} and the H-algorithm in \cite{BMW}, we have
$$\pi^{\vee}_{w_{\lam}}={\bf p}_{1}\times {\bf p}_{2}=p_{D}((\lam'|_{D_6})^-)^s\times p_{A}(\lam'|_{A_1})=[9,1^3]\times [1,1].$$ Then 
$$d_{\rm LS}^{H^\vee} (\mathcal{O}(\pi_{w_{\lam}}^\vee))=({\bf p}_{1}^t)_D \times {\bf p}^t_{2} =[3,1^9]\times [2]= {\rm sat}_{L^\vee}^{H^\vee} (\mathcal{O}_{L^\vee}),$$
where $\mathcal{O}_{L^\vee}=(3A_1)'' $ by Lemma \ref{prime1-2}. 
It then follows from Proposition \ref{P:Som} and the LS duality in \cite[p. 443]{Ca85} or Table \ref{tabchare7} that $\mathcal{O}_{{\rm Ann}(L(\lambda))} = d_{\rm Som}^{E_7^\vee}(L^\vee, \mathcal{O}_{L^\vee}) = d_{\rm LS}^{E_7^\vee}(L^\vee, \mathcal{O}_{L^\vee})= \mathcal{O}(E_{6})$ when $\mathcal{O}_{L^\vee}=(3A_1)''$.


 % So $\aff(w_{\lambda})=\aff(w_{\lambda'})=2+1=3$ by the algorithm in \cite{BXX}.
 % By Proposition \ref{pr:main1},  $L(\lambda)$ has $\gkd L(\lambda)=|\Phi^+|-\aff(w_{\lambda})=63-3=60$. 
% Thus from \cite[p.128]{CM}, we have $V(\Ann (L(\lam)))=\overline{\mathcal{O}}(C_3(a_1))$.
\end{ex}

















\subsection{Type \texorpdfstring{$E_8$}{}} We only need to consider the nonintegral cases and the results
are as follows.

\begin{Thm}\label{e8-av}
Let $L(\lam)$ be a simple nonintegral highest weight module of $E_8$ such that $\Phi_{[\lam]}^\vee$ is pseudo-maximal. Write $\lam'=\phi(\lam|_{\mathfrak{h}_{\Phi_{[\lambda]}}})=\prod\limits_{1\lest i\lest k} \phi(\lambda|_{\mathfrak{h}_{\Phi_{[\lambda]}}})|_{{\Phi_i}}=\prod\limits_{1\lest i\lest k} \lam'|_{{\Phi_i}}$.
	Then
    the following holds:
    \begin{enumerate}
	

\item Suppose $H=\Phi_{[\lam]}\simeq A_8$. We only need to consider the following $15$ cases:
$$\gkd L(\lam)=84,98,99,102,104,105,[107,114],116.$$
In the above $15$ cases, we have $H^{\vee}= A_8$ and   
$\pi^{\vee}_{w_{\lam}}={\bf p}_{1}=p_{A}(\lam')$. Then $$\mathcal{O}(\pi^{\vee}_{w_{\lam}})^t={\bf p}_{1}^t = {\rm sat}_{L^\vee}^{H^\vee} (\mathcal{O}_{L^\vee}^{\rm reg}),$$
where $L^\vee \subset \Phi^\vee$. The explicit information about $\mathcal{O}_{{\rm Ann}(L(\lambda))}$ can be found in Table \ref{taba8}.

\item Suppose $H=\Phi_{[\lam]}\simeq A_7\times A_1$. We only need to consider the following $15$ cases:
$$\gkd L(\lam)=92,98,99,102,[104,114],116.$$
In the above $15$ cases, we have $H^{\vee}= A_7\times A_1$ and   
$\pi^{\vee}_{w_{\lam}}={\bf p}_{1}\times {\bf p}_{2}=p_{A}(\lam'|_{A_7})\times p_{A}(\lam'|_{A_1})$. Then $$\mathcal{O}(\pi^{\vee}_{w_{\lam}})^t={\bf p}_{1}^t\times {\bf p}_{2}^t = {\rm sat}_{L^\vee}^{H^\vee} (\mathcal{O}_{L^\vee}^{\rm reg}),$$
where $L^\vee \subset \Phi^\vee$. The explicit information about $\mathcal{O}_{{\rm Ann}(L(\lambda))}$ can be found in Table \ref{taba7a1}.

\item Suppose $H=\Phi_{[\lam]}\simeq A_5\times A_2 \times A_1$. We only need to consider the following $14$ cases:
$$\gkd L(\lam)=101,102,[104,114],116.$$
In the above $14$ cases, we have $H^{\vee}= A_5\times A_2\times A_1$ and   
$\pi^{\vee}_{w_{\lam}}={\bf p}_{1}\times {\bf p}_{2}\times {\bf p}_{3}=p_{A}(\lam'|_{A_5})\times p_{A}(\lam'|_{A_2})\times p_{A}(\lam'|_{A_1})$. Then $$\mathcal{O}(\pi^{\vee}_{w_{\lam}})^t={\bf p}_{1}^t\times {\bf p}_{2}^t \times {\bf p}_{3}^t = {\rm sat}_{L^\vee}^{H^\vee} (\mathcal{O}_{L^\vee}^{\rm reg}),$$
where $L^\vee \subset \Phi^\vee$. The explicit information about $\mathcal{O}_{{\rm Ann}(L(\lambda))}$ can be found in Table \ref{taba5a2a1} and Table \ref{taba5a2a1-2}.


\item Suppose $H=\Phi_{[\lam]}\simeq A_4\times A_4$. We only need to consider the following $12$ cases:
$$\gkd L(\lam)=100,104,[106,114],116.$$
In the above $12$ cases, we have $H^{\vee}= A_4\times A_4$ and   
$\pi^{\vee}_{w_{\lam}}={\bf p}_{1}\times {\bf p}_{2}=p_{A}(\lam'|_{A_4(1)}\times p_{A}(\lam'|_{A_4(2)})$. Then $$\mathcal{O}(\pi^{\vee}_{w_{\lam}})^t={\bf p}_{1}^t\times {\bf p}_{2}^t = {\rm sat}_{L^\vee}^{H^\vee} (\mathcal{O}_{L^\vee}^{\rm reg}),$$
where $L^\vee \subset \Phi^\vee$. The explicit information about $\mathcal{O}_{{\rm Ann}(L(\lambda))}$ can be found in Table \ref{taba4a4}.


\item Suppose $H=\Phi_{[\lam]}\simeq D_8$. 
We only need to consider the following $17$ cases:
$$\gkd L(\lam)=92,98,99,100,102,[104,114],116.$$
In the above $17$ cases, we have $H^{\vee}=D_8$ and   
$\pi^{\vee}_{w_{\lam}}={\bf p}_{1}=p_{D}((\lam')^-)^s$. Then $$d_{\rm LS}^{H^\vee} (\mathcal{O}(\pi_{w_{\lam}}^\vee))=({\bf p}_{1}^t)_D  = {\rm sat}_{L^\vee}^{H^\vee} (\mathcal{O}_{L^\vee}),$$
where $L^\vee \subset \Phi^\vee$ and $\mathcal{O}_{L^\vee}$ a distinguished orbit of $L^\vee$. The explicit information about $\mathcal{O}_{{\rm Ann}(L(\lambda))}$ can be found in Table \ref{tabd8-1} and Table \ref{tabd8-2}. When ${\bf p}_{1}$ is a very even partition, the numeral of $\mathcal{O}_{L^{\vee}}$ is the same with the numeral of ${\bf p}_{1}$.



\item Suppose $H=\Phi_{[\lam]}\simeq D_5\times A_3$. 
We only need to consider the following $18$ cases:
$$\gkd L(\lam)=94,98,99,100,101,102,[104,114],116.$$
In the above $18$ cases, we have $H^{\vee}=D_5\times A_3$ and   
$\pi^{\vee}_{w_{\lam}}={\bf p}_{1}\times {\bf p}_{2}=p_{D}((\lam'|_{D_5})^-)^s\times p_{A}(\lam'|_{A_3})$. Then $$d_{\rm LS}^{H^\vee} (\mathcal{O}(\pi_{w_{\lam}}^\vee))=({\bf p}_{1}^t)_D\times {\bf p}^t_{2} = {\rm sat}_{L^\vee}^{H^\vee} (\mathcal{O}_{L^\vee}),$$
where $L^\vee \subset \Phi^\vee$ and $\mathcal{O}_{L^\vee}$ a distinguished orbit of $L^\vee$. The explicit information about $\mathcal{O}_{{\rm Ann}(L(\lambda))}$ can be found in Table \ref{tabd5a3} and Table \ref{tabd5a3-2}. 


\item Suppose $H=\Phi_{[\lam]}\simeq E_6\times A_2$. We only need to consider the following $18$ cases:
$$\gkd L(\lam)=84,92,94,99,100,102,[104,114],116.$$
In the above $18$ cases, we have $H^{\vee}=E_6\times A_2$ and
$\pi^{\vee}_{w_{\lam}}={\chi}_{1}\times {\bf p}_{2}={\chi}_{1}\times p_{A}(\lam'|_{A_2})$. 
Suppose the orbit $\mathcal{O}({\chi}_{1})$ is $L_{1}$. Then $\mathcal{O}(\pi_{w_\lambda}^\vee)=L_{1}\times {\bf p}_{2}$. 
Also we have
$$d_{\rm LS}^{H^\vee}(\mathcal{O}(\pi_{w_\lambda})) =d_{\rm LS}^{E_6}(L_1)\times {\bf p}^t_{2}= {\rm sat}_{L^\vee}^{H^\vee}(\mathcal{O}_{L^\vee}),$$
and $$\mathcal{O}_{{\rm Ann}(L(\lambda))} = d_{\rm Som}^{E_8^\vee}(L^{\vee}, \mathcal{O}_{L^\vee}),$$
where $L^\vee \subset \Phi^\vee$ and $\mathcal{O}_{L^\vee}$ a distinguished orbit of $L^\vee$. The explicit information about $\mathcal{O}_{{\rm Ann}(L(\lambda))}$ can be found in Table \ref{tabe6a2} and Table \ref{tabe6a2-2}.


\item Suppose $H=\Phi_{[\lam]}\simeq E_7\times A_1$. We only need to consider the following $17$ cases:
$$\gkd L(\lam)=84,94,98,99,100,[104,114],116.$$
In the above $17$ cases, we have $H^{\vee}=E_7\times A_1$ and
$\pi^{\vee}_{w_{\lam}}={\chi}_{1}\times {\bf p}_{2}={\chi}_{1}\times p_{A}(\lam'|_{A_1})$. Suppose that the orbit $\mathcal{O}({\chi}_{1})$ is $L_{1}$. Then $\mathcal{O}(\pi_{w_\lambda}^\vee)=L_{1}\times {\bf p}_{2}$. Also we have
$$d_{\rm LS}^{H^\vee}(\mathcal{O}(\pi_{w_\lambda}^\vee)) =d_{\rm LS}^{E_7}(L_1)\times {\bf p}^t_{2}= {\rm sat}_{L^\vee}^{H^\vee}(\mathcal{O}_{L^\vee}),$$
and $$\mathcal{O}_{{\rm Ann}(L(\lambda))} = d_{\rm Som}^{E_8^\vee}(L^{\vee}, \mathcal{O}_{L^\vee}),$$
where $L^\vee \subset \Phi^\vee$ and $\mathcal{O}_{L^\vee}$ a distinguished orbit of $L^\vee$. The explicit information about $\mathcal{O}_{{\rm Ann}(L(\lambda))}$ can be found in Table \ref{tabe7a1} and  Table \ref{tabe7a1-2}.













\end{enumerate}

\end{Thm}

% From Tables \ref{taba5a2a1}--\ref{taba5a2a1-2} and Table \ref{tabd8-2}, we have the following result.
% \begin{Cor}Keep notations as above.
%     There are three special cases:
%     $$J_1=\{\alpha_1,\alpha_2, \alpha_5,\alpha_6,\alpha_7\}\subset \tilde{\Delta}, J_2=\{\alpha_1,\alpha_2, \alpha_5,\alpha_7\}\subset \tilde{\Delta}, J_3=\{\alpha_2,\alpha_3, \alpha_5,\alpha_6,\alpha_7\}\subset \tilde{\Delta}$$
%     such that  $J_i\subset \tilde{\Delta}$ is not a subset of $A_8\subset E_8$, but $J_i$ is conjugate a subset of $A_8\subset E_8$. 
% \end{Cor}
% \begin{proof}
%   From Tables \ref{taba8}-\ref{tabe7a1}, we can see that for most cases the Bala--Carter label of $L$ has one prime if the subset $J\subset \tilde{\Delta}$ is a subset of $A_8\subset E_8$, and for most cases it has two primes if $J\subset \tilde{\Delta}$ is not a subset of $A_8\subset E_8$ except for the following three cases such that $L$ corresponds to $$[4,1,1]\times [2,1]\times [2], \text{~or~} [2,2,1,1]\times [2,1]\times [2], \text{~or~}[4,4,3,1^5].$$
%   Note that the first two cases come from Table \ref{taba5a2a1} and Table \ref{taba5a2a1-2}, and the third one comes from Table \ref{tabd8-2}. 
  
%   In the first case, we may choose $J_1=\{\alpha_1,\alpha_2, \alpha_5,\alpha_6,\alpha_7\}\subset \tilde{\Delta}$. The corresponding subsystem $\Phi_{J_1}$ has type $(A_3+2A_1)'$ or $(A_3+2A_1)''$. 
%   Denote $w_1=s_{\ep_2-\ep_1}s_{\ep_8-\ep_7}s_{\ep_7+\ep_2}s_{\ep_8+\ep_1}\in W(E_8)$. Then we have $w_1(J_1)=\{\alpha_1,-\beta, \alpha_5,\alpha_6,\alpha_7\}$, which is a subset of $A_8\subset E_8$. Therefore $\Phi_{J_1}$ corresponds to 
%   $L=(A_3+2A_1)'$ by Lemma \ref{prime1-2}.
  
  
%   In the second case, we may choose $J_2=\{\alpha_1,\alpha_2, \alpha_5,\alpha_7\}\subset \tilde{\Delta}$. The corresponding subsystem $\Phi_{J_2}$ has type $(4A_1)'$ or $(4A_1)''$. 
%   Then we have $w_1(J_2)=\{\alpha_1,-\beta, \alpha_5,\alpha_7\}$, which is a subset of $A_8\subset E_8$. Therefore $\Phi_{J_2}$ corresponds to 
%   $L=(4A_1)'$ by Lemma \ref{prime1-2}.
  
  
%  In the third case, we may choose $J_3=\{\alpha_2, \alpha_3,\alpha_5,\alpha_6, \alpha_7\}\subset \tilde{\Delta}$. The corresponding subsystem $\Phi_{J_3}$ has type $(A_3+2A_1)'$ or $(A_3+2A_1)''$. 
%  From \cite[\S 2.2]{som-98}, we know that $d_J$ is an invariant of the equivalence class of $J\subset \tilde{\Delta}$. For $J_3$, we have $d_{J_3}=\mathrm{gcd}(c_0,c_1,c_4,c_8)=\mathrm{gcd}(1,2,6,2)=1$. Thus $J_3$ is conjugate a subset of $A_8\subset E_8$ by \cite[\S 2.2]{som-98}. 
%  Therefore $\Phi_{J_3}$ corresponds to 
%   $L=(A_3+2A_1)'$ by Lemma \ref{prime1-2}.
 
 
%  % Denote $w_3=?$. Then we have $w_3(J_3)=\{\alpha_3, \alpha_5,\alpha_6,\alpha_7,-\beta\}$, which is a subset of $A_8\subset E_8$. Therefore $\Phi_{J_3}$ corresponds to 
%  %  $L=(A_3+2A_1)'$ by Lemma \ref{prime1-2}.
   
% \end{proof}

\begin{Rem}
    In $E_7$ or $E_8$, the partition  ${\bf p}=[3,1]$ associated with $D_2=2A_1$ corresponds to $\{\alpha_2, \alpha_3\} \subset \tilde{\Delta}_{E_7}$ and $\{\alpha_2, \alpha_5\} \subset \tilde{\Delta}_{E_8}$ respectively. Similarly, the partition  ${\bf p}=[5,1]$ associated with $D_3 = A_3$ corresponds to $\{\alpha_2, \alpha_4,\alpha_5\}\subset \tilde{\Delta}_{E_7}$ and $\{\alpha_2, \alpha_4,\alpha_5\}\subset \tilde{\Delta}_{E_8}$ for $E_7$ and $E_8$, respectively. 
\end{Rem}





\begin{ex}

Let $\mathfrak{g}=E_8$ and $L(\lam)$ be a highest weight module of $\mathfrak{g}$ with  $\lambda=(1,1,1,1,1,1, \frac{1}{2}, \frac{5}{2})$. It is easy to verify that $\Phi_{[\lambda]}$ is a subsystem with simple roots $\{ \alpha_1, \alpha_2,\alpha_3,\alpha_4,\alpha_5,\alpha_6,\alpha_7,\ep_7+\ep_8\}$. Therefore $H=\Phi_{[\lambda]}\simeq E_7\times A_1$. Suppose the simple root system of $E_7\times A_1$ is $\Delta=\{\alpha_1, \alpha_2,\alpha_3,\alpha_4,\alpha_5,\alpha_6,\alpha_7\}\times \{\gamma_1=e_1-e_2\}$.
    We define a map $\phi: \Phi_{[\lambda]} \rightarrow E_7\times A_1$ such that $\phi(\alpha_1)=\alpha_1, \phi(\alpha_2)=\alpha_2, \phi(\alpha_3)=\alpha_3, \phi(\alpha_4)=\alpha_4,\phi(\alpha_5)=\alpha_5,\phi(\alpha_6)=\alpha_6,\phi(\alpha_7)=\alpha_7,\phi(\ep_7+\ep_8)=\gamma_1$.
   Thus for $\lam\in \mathfrak{h}_{E_8}^* $, it will be an integral weight  of ${\Phi_{[\lambda]}}$. By the algorithm in \S \ref{nonint}, $$\lam'|_{{E_7}}
    =(1,1,1,1,1,1,-1,1)$$ is an integral weight of type $E_7$,         and $$\lam'|_{{A_1}}=\left(\frac{3}{2},-\frac{3}{2}\right)$$ is an integral weight of type $A_1$. 
   
  
   Denote $\gamma=\lam'|_{E_7}=(1,1,1,1,1,1,-1,1)$.
Then $\gamma=-\omega_1+2\omega_2:=[-1,2,0,0,0,0,0]$. By using Lemma \ref{find-w-lambda}, we can find a $w_{\gamma}$ such that  $$\gamma=w_{\gamma}\mu,$$ 
where  $\mu=-\omega_2-\omega_7:=[0,-1,0,0,0,0,-1]$ is antidominant and  
\begin{align*}
   w_{\gamma}=&[1,3,4,5,6,2,3,4,5,1,3,4,2,3,1,0,2,3,4,5,6,\\
   &1,3,4,5,2,3,4,1,3,2,0,2,3,4,5,1,3,4,2,3,1] 
\end{align*} in PyCox. Then by using PyCox, we find that the special character corresponding to $w_{\gamma}$ is
$105_b$ and its $\aff$ value is $6$. From \cite[p. 430]{Ca85}, the Bala--Carter label of $\mathcal{O}(105_b)$ is $A_6$.




By using the RS algorithm in \cite{BX}, we have
$${\bf p}_{2}= p_{A}(\lam'|_{A_1})=[1,1].$$
Now we have
$$\pi^{\vee}_{w_{\lam}}=\chi_{1}\times {\bf p}_{2}=105_b\times p_{A}(\lam'|_{A_1})=105_b\times [1,1].$$ 
Note that from  the LS duality in \cite[p. 443]{Ca85} , we have $d_{\rm LS}^{E_7}(A_6)=A_2+3A_1$. 
Then \begin{align*}
    d_{\rm LS}^{H^\vee} (\mathcal{O}(\pi_{w_{\lam}}^\vee))&=d_{\rm LS}^{E_7}(A_6) \times {\bf p}^t_{2} \\
    &=d_{\rm LS}^{E_7}(A_6) \times [2]\\
    &=(A_2+3A_1)\times [2]\\
    &= {\rm sat}_{L^\vee}^{H^\vee} (\mathcal{O}_{L^\vee}),
\end{align*}
where $\mathcal{O}_{L^\vee}=A_2+4A_1 $ by Lemma \ref{prime1-2}. 
It follows from Proposition \ref{P:Som} and \cite[Table X]{Sommers} that $\mathcal{O}_{{\rm Ann}(L(\lambda))} = d_{\rm Som}^{E_8}(L^\vee, \mathcal{O}_{L^\vee}) = \mathcal{O}(D_7)$.
\end{ex}











% \begin{ex}
% Consider $E_8$, for which every pseudo-Levi $M$ and its dual $M^\vee$ are naturally identified. Take $\lambda$ be such that $W_{[\lambda]} = A_7 \times A_1 =: M$, and that $\pi_{w_\lambda}$ and $\mathcal{O}(\pi_{w_\lambda})$ are associated with the partition $[2,2,1^4] \times [2]$ of $M$. In this case, we see that $d_{\rm LS}^{M}(\mathcal{O}(\pi_{w_\lambda})) = {\rm sat}_L^{M}(\mathcal{O}_L^{\rm reg})$,
% where $L:=(A_5 + A_1)''$ and $\mathcal{O}_L^{\rm reg}$ means the regular orbit of $L$. It then follows from Proposition \ref{P:Som} and \cite[Table X]{Sommers} that $\mathcal{O}_{{\rm Ann}(L(\lambda))} = \mathcal{O}(D_5 +A_1)$.
% \end{ex}
\section{The general case } \label{S:gen}


  For a given highest weight module $L(\lambda)$, the integral dual root system  $\Phi_{[\lam]}^\vee$ may not be pseudo-maximal. There are two methods of computing $\mathcal{O}_{{\rm Ann}(L(\lambda))}$. First, since $j$-induction satisfies the property of induction by stages, see \cite[Proposition 11.2.4]{Ca85}, from  $\pi_{w_{\lambda}}$ we can use $j$-induction to get a representation of $W(\Phi_1)$, where $\Phi_1^\vee \subset \Phi^\vee$ is pseudo-maximal and $W(\Phi_1)$ contains $W_{[\lambda]}$ as parabolic Weyl subgroup. In this case, one can still utilize the tables in Appendix to compute $\mathcal{O}_{{\rm Ann}(L(\lambda))}$. The second method is to apply Proposition \ref{P:Som} directly.
\begin{ex}
    Let $\mathfrak{g}=E_8$. We assume $H:=\Phi_{[\lambda]}\simeq A_7$ and $\pi_{w_{\lambda}}=[1^8]\in \mathrm{Irr} (W(A_7))=\mathrm{Irr}(S_8)$. If we consider $A_7$ as a root subsystem of $A_8$, then from \cite{BMW} the $j$-induction gives the representation $\pi_{\bf p}=j_{S_8\times S_1}^{S_9}(\pi_{[1^8]}\otimes\pi_{[1]})$ of $S_9=W(A_8)$ with ${\bf p}=[2,1^7]$. It then follows from Table \ref{taba8} that 
 $\mathcal{O}_{{\rm Ann}(L(\lambda))} = \mathcal{O}(D_4(a_1)+A_2)$.
 
 On the other hand, we see that $d_{\rm LS}^{M^\vee}(\mathcal{O}(\pi_{w_\lambda}^\vee)) = {\rm sat}_{L^\vee}^{M^\vee}(\mathcal{O}_{L^\vee}^{\rm reg})$,
where $\mathcal{O}_{L^\vee}:=(A_7)'$. It then also follows from Proposition \ref{P:Som} and \cite[Table X]{Sommers} or Table \ref{tabchare8} that $\mathcal{O}_{{\rm Ann}(L(\lambda))} = \mathcal{O}(D_4(a_1)+A_2)$.
\end{ex} 
 
 
 Let $J=H:=\Phi_{[\lambda]}\subset \tilde{\Delta}$ be a proper subset of $\tilde{\Delta}$ and $d_J$ be the greatest common divisor of those $c_i$ for which $c_i\in \tilde{\Delta}-J$. Let $\Phi_J$ be the root system generated by $J$.
As mentioned in \cite{som-98}, one needs to execute with care for the following cases of $\Phi_J$.
\begin{enumerate}
    \item For type $E_7$, $\Phi_J$ is one the following: \begin{equation}\label{E7-case}
        A_5,A_5+A_1,4A_1,3A_1,A_3+2A_1 \text{~or~} A_3+A_1.
    \end{equation}
When $J$ is conjugate to a subset of $A_7\subset E_7$ in the above cases, we give the corresponding root system  one prime, otherwise two primes.
\item For type $E_8$, $\Phi_J$ is one the following: \begin{equation}\label{E8-case}
    A_7,A_5+A_1,4A_1,2A_3 \text{~or~} A_3+2A_1.
\end{equation}
When $J$ is conjugate to a subset of $A_8\subset E_8$, we will have  $d_J=1$ in the above cases. Then we give the corresponding root system  one prime. When  $J$ is not conjugate to a subset of $A_8\subset E_8$, we will have $d_J=2$ and give the corresponding root system two primes.
\end{enumerate}


\begin{ex}
% Keep notations as above.
%     There are three special cases:
%     $$J_1=\{\alpha_1,\alpha_2, \alpha_5,\alpha_6,\alpha_7\}\subset \tilde{\Delta}, J_2=\{\alpha_1,\alpha_2, \alpha_5,\alpha_7\}\subset \tilde{\Delta}, J_3=\{\alpha_2,\alpha_3, \alpha_5,\alpha_6,\alpha_7\}\subset \tilde{\Delta}$$
%     such that  $J_i\subset \tilde{\Delta}$ is not a subset of $A_8\subset E_8$, but $J_i$ is conjugate a subset of $A_8\subset E_8$. 


  From Tables \ref{taba8}--\ref{tabe7a1}, we see that for most cases the Bala--Carter label of $L$ has one prime if the subset $J\subset \tilde{\Delta}$ is a subset of $A_8\subset E_8$, and for most cases it has two primes if $J\subset \tilde{\Delta}$ is not a subset of $A_8\subset E_8$ except for the following three cases such that $\mathcal{O}_{L}$ corresponds to $$[4,1,1]\times [2,1]\times [2], \text{~or~} [2,2,1,1]\times [2,1]\times [2], \text{~or~}[4,4,3,1^5].$$
  Note that the first two cases come from Table \ref{taba5a2a1} and Table \ref{taba5a2a1-2}, and the third one comes from Table \ref{tabd8-2}. 
  
  In the first case, we may choose $J_1=\{\alpha_1,\alpha_2, \alpha_5,\alpha_6,\alpha_7\}\subset \tilde{\Delta}$. The corresponding subsystem $\Phi_{J_1}$ has type $(A_3+2A_1)'$ or $(A_3+2A_1)''$. 
  Denote $w_1=s_{\ep_2-\ep_1}s_{\ep_8-\ep_7}s_{\ep_7+\ep_2}s_{\ep_8+\ep_1}\in W(E_8)$. Then we have $w_1(J_1)=\{\alpha_1,-\beta, \alpha_5,\alpha_6,\alpha_7\}$, which is a subset of $A_8\subset E_8$. Therefore $\Phi_{J_1}$ corresponds to 
  $\mathcal{O}_{L}=(A_3+2A_1)'$ by Lemma \ref{prime1-2}.
  
  
  In the second case, we may choose $J_2=\{\alpha_1,\alpha_2, \alpha_5,\alpha_7\}\subset \tilde{\Delta}$. The corresponding subsystem $\Phi_{J_2}$ has type $(4A_1)'$ or $(4A_1)''$. 
  Then we have $w_1(J_2)=\{\alpha_1,-\beta, \alpha_5,\alpha_7\}$, which is a subset of $A_8\subset E_8$. Therefore $\Phi_{J_2}$ corresponds to 
  $\mathcal{O}_{L}=(4A_1)'$ by Lemma \ref{prime1-2}.
  
  
 In the third case, we may choose $J_3=\{\alpha_2, \alpha_3,\alpha_5,\alpha_6, \alpha_7\}\subset \tilde{\Delta}$. The corresponding subsystem $\Phi_{J_3}$ has type $(A_3+2A_1)'$ or $(A_3+2A_1)''$. 
 From \cite[\S 2.2]{som-98}, we know that $d_J$ is an invariant of the equivalence class of $J\subset \tilde{\Delta}$. For $J_3$, we have $d_{J_3}=\mathrm{gcd}(c_0,c_1,c_4,c_8)=\mathrm{gcd}(1,2,6,2)=1$. Thus, $J_3$ conjugates to a subset of $A_8\subset E_8$ by \cite[\S 2.2]{som-98}. 
 Therefore, $\Phi_{J_3}$ corresponds to $\mathcal{O}_{L}=(A_3+2A_1)'$ by Lemma \ref{prime1-2}.
 
 

   
\end{ex}


For some cases, it is not easy to determine whether $J$ is conjugate to a subset of $A_7\subset E_7$ (resp. to a subset $A_8\subset E_8$). For these cases, we will need to compute the weighted Dynkin diagram of the corresponding nilpotent orbit of the saturation of $d_{\rm LS}^{H^\vee}(\mathcal{O}(\pi_{w_\lambda}^\vee))$. The algorithm is as follows:
\begin{enumerate}
    \item Find the simple system $\Delta_{[\lambda]}$ of $H:=\Phi_{[\lambda]}$;
    \item Suppose $d_{\rm LS}^{H^\vee}(\mathcal{O}(\pi_{w_\lambda}^\vee))=a_kA_k+a_1A_1$ for some integer $k>1$ and $a_1,a_k\in \mathbb{Z}_{\geq 0}$;
    \item Find  the neutral element of the $\mathfrak{sl}_2$-triple associated with the orbit $a_kA_k+a_1A_1$ from $\Delta_{[\lambda]}$, denoted it by $h:=[b_1,b_2,\dots,b_l]=b_1\omega_1+b_2\omega_2+\cdots+b_l\omega_l=c_1e_1+c_2e_2+\cdots+c_8e_8$, where $l=7$ for $E_7$ and $l=8$ for $E_8$;
    \item Apply the positive index reduction algorithm in Lemma \ref{find-w-lambda} to obtain an antidominant weight and then multiply by $-1$ to get a dominant weight $h^*$ in the Weyl-orbit of $h$;
    \item Then  
    $$(\langle h^*, \alpha_i \rangle)_{1\lest i \lest 7}$$
    or $$(\langle h^*, \alpha_i \rangle)_{1\lest i \lest 8}$$ is the weighted Dynkin diagram for the saturation of the orbit $d_{\rm LS}^{H^\vee}(\mathcal{O}(\pi_{w_\lambda}^\vee))$ in $E_7$ or $E_8$. Then by the tables in \cite{som-98} or Table \ref{tabchare7} and Table \ref{tabchare8}, we can determine whether the orbit $a_kA_k+a_1A_1$ has a decoration of one prime or two primes.
\end{enumerate}

\begin{Rem}
    For each case in (\ref{E7-case}) and (\ref{E8-case}), we may choose a fixed subset $J \subset \tilde{\Delta}$ such that $J$ is not conjugate to a subset of $A_7\subset E_7$ (resp. to a subset $A_8\subset E_8$). Then the corresponding root system $\Phi_J$ will have two primes.

    For type $E_7$, we have the following:
    \begin{enumerate}
        \item When $\Phi_J=(A_5)''$, we may choose $J=\{\alpha_2,\alpha_4,\alpha_5,\alpha_6,\alpha_7\}$. Then $h=5\alpha_2+8\alpha_4+9\alpha_5+8\alpha_6+5\alpha_7$.
        \item When $\Phi_J=(A_5+A_1)''$, we may choose $J=\{\alpha_1,\alpha_2,\alpha_4,\alpha_5,\alpha_6,\alpha_7\}.$ Then $h=\alpha_1+5\alpha_2+8\alpha_4+9\alpha_5+8\alpha_6+5\alpha_7$.
        \item When $\Phi_J=(4A_1)''$, we may choose $J=\{\alpha_2,\alpha_3,\alpha_5,\alpha_7\}.$ Then $h=\alpha_2+\alpha_3+\alpha_5+\alpha_7$.
        \item When $\Phi_J=(3A_1)''$, we may choose $J=\{\alpha_2,\alpha_5,\alpha_7\}.$ Then $h=\alpha_2+\alpha_5+\alpha_7$.
        \item When $\Phi_J=(A_3+2A_1)''$, we may choose $J=\{\alpha_2,\alpha_3,\alpha_5,\alpha_6,\alpha_7\}.$  Then $h=\alpha_2+\alpha_3+3\alpha_5+4\alpha_6+3\alpha_7$.
        \item When $\Phi_J=(A_3+A_1)''$, we may choose $J=\{\alpha_2,\alpha_4,\alpha_5,\alpha_7\}.$ Then $h=3\alpha_2+4\alpha_4+3\alpha_5+\alpha_7$.
    \end{enumerate}


For type $E_8$, we have the following:
    \begin{enumerate}
        \item When $\Phi_J=(A_7)''$, we may choose $J=\{\alpha_2,\alpha_4,\alpha_5,\alpha_6,\alpha_7,\alpha_8,\alpha_0\}.$
        \item When $\Phi_J=(A_5+A_1)''$, we may choose $J=\{\alpha_2,\alpha_5,\alpha_6,\alpha_7,\alpha_8,\alpha_0\}.$
        \item When $\Phi_J=(4A_1)''$, we may choose $J=\{\alpha_2,\alpha_5,\alpha_7,\alpha_0\}.$ 
        \item When $\Phi_J=(2A_3)''$, we may choose $J=\{\alpha_2,\alpha_4,\alpha_5,\alpha_7,\alpha_8,\alpha_0\}.$
        \item When $\Phi_J=(A_3+2A_1)''$, we may choose $J=\{\alpha_2,\alpha_4,\alpha_5,\alpha_7,\alpha_0\}.$
       
    \end{enumerate}

An explicit form of the neutral element $h$ can be found from \cite[\S 3.6]{CM}.
\end{Rem}

\begin{ex}\label{e7-ex}
Consider $\mathfrak{g}=E_7$ and $$\lambda=(0,1,1,2,3,3,-7,7)/4.$$
The root system $H:=\Phi_{[\lambda]} \simeq D_4 \times A_1 \times A_1$ has simple roots
$$ \left\{\alpha:=\frac{1}{2}(-e_1-e_2-e_3+e_4-e_5-e_6-e_7+e_8), e_2 + e_5, e_6 - e_5, e_3- e_2 \right\}$$
for $D_4$ and $\{\beta:=\frac{1}{2}(e_1+e_2+e_3-e_4-e_5-e_6-e_7+e_8)\}$ for the first $A_1$ and $\{\gamma:=\frac{1}{2}(e_1-e_2-e_3-e_4+e_5+e_6-e_7+e_8)\}$  for the second $A_1$. 

% Suppose the simple root system of $E_7\times A_1$ is $\Delta=\{\alpha_1, \alpha_2,\alpha_3,\alpha_4,\alpha_5,\alpha_6,\alpha_7\}\times \{\gamma_1=e_1-e_2\}$.
%     We define a map $\phi: \Phi_{[\lambda]} \rightarrow E_7\times A_1$ such that $\phi(\alpha_1)=\alpha_1, \phi(\alpha_2)=\alpha_2, \phi(\alpha_3)=\alpha_3, \phi(\alpha_4)=\alpha_4,\phi(\alpha_5)=\alpha_5,\phi(\alpha_6)=\alpha_6,\phi(\alpha_7)=\alpha_7,\phi(\ep_7+\ep_8)=\gamma_1$.
   Thus for $\lam\in \mathfrak{h}_{E_7}^* $, it will be an integral weight  of ${\Phi_{[\lambda]}}$. By the algorithm in \S \ref{nonint}, $$\lam'|_{{D_4}}
    =(2,1,0,0)$$ is an integral weight of type $D_4$,         $$\lam'|_{{A_1(1)}}=\left(\frac{1}{2},-\frac{1}{2}\right)$$ is an integral weight of type $A_1$ and $$\lam'|_{{A_1(2)}}=(1,-1)$$ is an integral weight of type $A_1$. 
   
  
%    Denote $\gamma=\lam'|_{E_7}=(1,1,1,1,1,1,-1,1)$.
% Then $\gamma=-\omega_1+2\omega_2:=[-1,2,0,0,0,0,0]$. By using Lemma \ref{find-w-lambda}, we can find a $w_{\gamma}$ such that  $$\gamma=w_{\gamma}\mu,$$ 
% where  $\mu=-\omega_2-\omega_7:=[0,-1,0,0,0,0,-1]$ is antidominant and  
% \begin{align*}
%    w_{\gamma}=&[1,3,4,5,6,2,3,4,5,1,3,4,2,3,1,0,2,3,4,5,6,\\
%    &1,3,4,5,2,3,4,1,3,2,0,2,3,4,5,1,3,4,2,3,1] 
% \end{align*} in PyCox. Then by using PyCox, we find that the special character corresponding to $w_{\gamma}$ is
% $105_b$ and its $\aff$ value is $6$. From \cite[p. 430]{Ca85}, the Bala--Carter label of $\mathcal{O}(105_b)$ is $A_6$.




By using the RS algorithm in \cite{BXX} and the H-algorithm in \cite{BMW},  we have 
$${\bf p}_{1}=p_{D}((\lam'|_{D_4})^-)^s=[3,1^5],$$
$${\bf p}_{2}= p_{A}(\lam'|_{A_1(1)})=[1,1],$$
 and $${\bf p}_{3}= p_{A}(\lam'|_{A_1(2)})=[1,1].$$
 
Thus
$$\pi^{\vee}_{w_{\lam}}={\bf p}_{1}\times {\bf p}_{2}\times {\bf p}_{3}=[3,1^5]\times [1,1]\times [1,1].$$ 

So \begin{align*}
    d_{\rm LS}^{H^\vee} (\mathcal{O}(\pi_{w_{\lam}}^\vee))&= ({\bf p}_{1}^t)_D\times {\bf p}^t_{2}\times {\bf p}^t_{3} \\
    &=[5,1^3] \times [2]\times [2]\\
    &= {\rm sat}_{L^\vee}^{H^\vee} (\mathcal{O}_{L^\vee}).
\end{align*}
% where $\mathcal{O}_{L^\vee}=A_3+2A_1 \subset \Phi$ by Lemma \ref{prime1-2}. 
Moreover, we have
$$\mathcal{O}_{L^\vee} \simeq D_3 + A_1+A_1 \simeq A_3+2A_1.$$
To determine whether it is $(A_3+2A_1)'$ or $(A_3+2A_1)''$, we note that the neutral element of the $\mathfrak{sl}_2$-triple associated with the orbit $A_3+2A_1$ is of the form
$$h:=3(e_3 - e_2) + 4(e_2+e_5) + 3(e_6 - e_5)+\beta+\gamma=e_1+e_2+3e_3-e_4+e_5+3e_6-e_7+e_8.$$
This element is clearly not dominant. We apply the positive index reduction algorithm in Lemma \ref{find-w-lambda} to obtain a dominant $h^*$ in the Weyl-orbit of $h$; also, by \cite[\S 3.5 ]{CM} the element $h^*$ gives rise to a weighted Dynkin diagram decorated by $$(\langle h^*, \alpha_i \rangle)_{1\lest i \lest 7} = (0,0,2,0,0,0,0).$$
This shows that the saturation of the orbit $\mathcal{O}_{L^\vee}$ in $E_7$ is $D_4(a_1)$ by the Table in \cite[p. 130]{CM}. It then follows (for example, from \cite[p. 559]{som-98} or Table \ref{tabchare7}) that $\mathcal{O}_{L^\vee} = (A_3+2A_1)'$. It follows from Proposition \ref{P:Som} and \cite[Table IX]{Sommers}  that $\mathcal{O}_{{\rm Ann}(L(\lambda))} = D_6(a_2)$.
\end{ex}

\begin{ex} 
Consider $\mathfrak{g}=E_8$ and $$\lambda=(\omega_2 + \omega_7 + 2\omega_8)/4 =(1,1,1,1,1,3,7,13)/8.$$
We obtain that the root system $H:=\Phi_{[\lambda]} \simeq D_6 \times A_1$ has simple roots
$$\left\{\frac{1}{2}(e_1-e_2-e_3-e_4-e_5-e_6-e_7+e_8), e_1 + e_7, -e_1 + e_2, -e_2 + e_3, -e_3 + e_4, -e_4 + e_5 \right\}$$
for $D_6$ and $\{e_6 + e_8\}$ for $A_1$. Moreover, by the similar process with Example \ref{e7-ex}, we have
$$\mathcal{O}_{L^\vee} \simeq (3A_1) + A_1 \simeq 4A_1.$$
To determine whether it is $(4A_1)'$ or $(4A_1)''$, we note that the neutral element of the $\mathfrak{sl}_2$-triple associated with orbit $3A_1 + A_1$ is of the form
$$h:=(e_5 - e_4) + (e_3 - e_2) + (e_1 + e_7) + (e_6 + e_8).$$
This element is clearly not dominant. We apply the positive index reduction algorithm in Lemma \ref{find-w-lambda} to obtain a dominant $h^*$ in the Weyl-orbit of $h$; also, by \cite[\S 3.5 ]{CM} $h^*$ gives rise to a weighted Dynkin diagram decorated by $$(\langle h^*, \alpha_i \rangle)_{1\lest i \lest 8} = (0,0,0,0,0,0,0,2).$$
This shows that the saturation of the orbit $\mathcal{O}_{L^\vee}$ in $E_8$ is $A_2$ by the Table in \cite[p. 130]{CM}. It then follows (for example, from \cite[p. 559]{som-98} or Table \ref{tabchare8}) that $\mathcal{O}_{L^\vee} = (4A_1)''$. It follows from Proposition \ref{P:Som} and \cite[Table X]{Sommers} that $\mathcal{O}_{{\rm Ann}(L(\lambda))} = E_7$.
\end{ex}







\section{Appendix}
We tabulate all the relevant data when $\Phi_{[\lambda]}^\vee$ is a pseudo-maximal root subsystem of $\Phi^\vee$.
We explain the notation used in the tables:
\begin{enumerate}
    \item[$\bullet$] $\mathrm{GKd} L(\lambda)$ means the Gelfand--Kirillov dimension of $L(\lam)$;
    \item[$\bullet$] cha$(\mathcal{O}):={\rm Spr}_{G,\mathbf{1}}^{-1}(\mathcal{O})$ means the character associated with the nilpotent orbit $\mathcal{O}$ via the Springer correspondence;
    \item[$\bullet$] the orbit $\mathcal{O}_{L^\vee}$ (as in Tables \ref{tabcharf4}--\ref{tabchare8}) means a distinguished orbit of $L^\vee$, where $L^\vee \subset M^\vee$ is a Levi group of $M^\vee$ and $M^\vee \subseteq G^\vee$ is a pseudo-Levi subgroup (thus $L^\vee \subseteq G^\vee$ is a pseudo-Levi subgroup of $G^\vee$), as in Proposition \ref{P:Som}; 
    %
    \item[$\bullet$] $\mathcal{O}_{L^\vee}^{\rm BC}:=G^\vee \cdot \mathcal{O}_{L^\vee}$ (as in Tables \ref{tabcharf4}--\ref{tabchare8}) means the saturation of the orbit $\mathcal{O}_{L^\vee}$, written using the Bala--Carter label;
    \item[$\bullet$] the orbit $\mathcal{O}_{{\rm Ann}(L(\lambda))}$ is the annihilator variety of $L(\lambda)$, and is equal to $d_{\rm Som}^{G^\vee}(L^{\vee}, \mathcal{O}_{L^\vee}).$
    %
    \item[$\bullet$] $\mathcal{O}(\pi):=\mathcal{O}_{\rm Spr}^H(\pi)$ in Tables \ref{tabc4}--\ref{tabb3a1}, where $H^\vee \subseteq G^\vee$ is pseudo-Levi subgroup with root system being  a fixed pseudo-maximal root subsystem of $G^\vee$, and $\pi$ a special representation of $W(H)$; thus $\pi^\vee \in {\rm Irr}(W(H^\vee))$ is special and $\mathcal{O}^\vee(\pi) = \mathcal{O}_{\rm Spr}^{H^\vee}(\pi^\vee)$. 
\end{enumerate}
For $E_6, E_7, E_8$ which are simply-laced, the root system of a Levi and its dual system are canonically identified and of the same type, thus from Table \ref{tabd5-1} onwards, we make such an identification and omit the column of $\mathcal{O}^\vee(\pi)$ as in Tables \ref{tabc4}--\ref{tabb3a1}.



\begin{table}[htbp]\caption{ Bala--Carter label for some nilpotent orbits }
	\label{label}
	\centering
	\renewcommand{\arraystretch}{1.5}
	\setlength\tabcolsep{5pt}
	
	%\begin{tabular}{|c|c|c|}
		\begin{tabular}{|c|c|c|}
		\hline \text{Type}  &${\bf p}$ & \text{Bala--Carter label} \\ \hline
        $A$ & $[k+1]$ & $A_k$ \\
        \hline

        $B$ & $[2k+1]$ & $B_k$ \\
        &$[5,3,1]$ & $B_4(a_2)$\\
        & $[k+1,k+1]$ & $A_k$ \\
        \hline
$C$ & $[2k]$ & $C_k$ \\
& $[4,2]$ & $C_3(a_1)$ \\
        & $[k+1,k+1]$ & $\tilde{A}_k$ \\
        & $[2^3]$ & $A_1+\tilde{A}_1$\\
        \hline
     

      $D$  &  $ [7,5, 3, 1] $  &   $ D_8(a_5)$      \\ 

        &   $ [9, 3] $  &  $D_6(a_1)$      \\
         
       &   $ [7, 5] $  &  $D_6(a_2)$      \\
         
       &   $ [7, 3] $  &  $D_5(a_1)$      \\
         
       &   $ [5, 3] $  &  $D_4(a_1)$      \\
          
        &  $ [5, 1] $  &  $D_3$      \\
         &  $ [5,3,3, 1] $  &  $A_3+A_2$      \\ 
%&  $ [5, 1^7] $  &  $(A_3)''$      \\


         
       &   $ [3, 1] $  &  $D_2=2A_1$      \\
      &   $ [2k-2i-1, 2i+1] $  &  $D_k(a_i)$      \\
        
        
        
        &$[k+1,k+1]$ & $A_k$ \\
        &$[2k-1,1]$ & $D_k$ \\
        
        
        
        
        
        \hline

    
 
 
 \end{tabular}
	\bigskip
	
\end{table}




% \begin{table}[htbp]\caption{ Characters in PyCox and nilpotent orbits for $ F_4$}
% 	\label{tabcharef4-1}
% 	\centering
% 	\renewcommand{\arraystretch}{1.3}
% 	\setlength\tabcolsep{5pt}
% 	%\setlength{\leftmargin}{15cm}
% {	\begin{tabular}{|l|l|l|l|}
% 		%\begin{tabular}{|l|l|l|l|}
% 		\hline  cha$(\mathcal{O})$ in PyCox & Nilpotent orbit $\mathcal{O}$ & $\dim \mathcal{O}$ \\ \hline
 
%   $1_1$  &  $F_4$  & $48$ \\\hline
%   $4_2$  &  $F_4(a_1)$  & $46$ \\
% \hline
%   $9_1$  &  $F_4(a_2)$  & $44$ \\
% \hline
%   $8_1$  &  $B_3$  & $42$ \\
%   $8_3$  &  $C_3$  & $42$ \\
% \hline
%  $12$  &  $F_4(a_3)$  & $40$ \\
% \hline
%  $8_2$  &  $\tilde{A}_2$  & $30$ \\
%  $8_4$  &  ${A}_2$  & $30$ \\
% \hline
%  $9_4$  &  $A_1+\tilde{A}_1$  & $28$ \\
% \hline
%  $4_5$  &  $\tilde{A}_1$  & $22$ \\
% \hline
%  $1_4$  &  $0$  & $0$ \\
% \hline



% \end{tabular}}
% 	\bigskip
	

% \end{table}

\begin{table}[htbp]\caption{  Nilpotent orbits and characters in PyCox  for $ F_4$}
	\label{tabcharef4-1}
	\centering
	\renewcommand{\arraystretch}{1.3}
	\setlength\tabcolsep{5pt}
{	\begin{tabular}{|c|c|c|}
		\hline Nilpotent orbit $\mathcal{O}$ & cha$(\mathcal{O})$ in PyCox & $\dim \mathcal{O}$ \\ \hline
 
  $F_4$  &  $1_1$  & $48$ \\\hline
  $F_4(a_1)$  &  $4_2$  & $46$ \\
\hline
  $F_4(a_2)$  &  $9_1$  & $44$ \\
\hline
  $B_3$  &  $8_1$  & $42$ \\
  $C_3$  &  $8_3$  & $42$ \\
\hline
 $F_4(a_3)$  &  $12$  & $40$ \\
\hline
 $\tilde{A}_2$  &  $8_2$  & $30$ \\
 ${A}_2$  &  $8_4$  & $30$ \\
\hline
 $A_1+\tilde{A}_1$  &  $9_4$  & $28$ \\
\hline
 $\tilde{A}_1$  &  $4_5$  & $22$ \\
\hline
 $0$  &  $1_4$  & $0$ \\
\hline

\end{tabular}}
	\bigskip
	
\end{table}

% \begin{table}[htbp]\caption{ Characters in PyCox and nilpotent orbits for $ E_6$ }
% 	\label{tabchare6-0}
% 	\centering
% 	\renewcommand{\arraystretch}{1.3}
% 	\setlength\tabcolsep{5pt}
% 	%\setlength{\leftmargin}{15cm}
% {	\begin{tabular}{|l|l|l|l|}
% 		%\begin{tabular}{|l|l|l|l|}
% 		\hline   cha$(\mathcal{O})$ in PyCox & Nilpotent orbit $\mathcal{O}$ & $\dim \mathcal{O}$ \\ \hline
%  $1_p$& $E_6$& $72$\\\hline
%  $6_p$& $E_{6}(a_1)$& $70$\\\hline
%  $20_p$& $D_5$& $68$\\\hline
%  $30_p$& $E_{6}(a_3)$& $66$\\\hline
%  $64_p$& $D_{5}(a_1)$& $64$\\\hline
 
%  $60_p$& $A_4+A_1$& $62$\\\hline
%  $81_p$& $A_4$& $60$\\
%  $24_p$& $D_4$& $60$\\\hline
%  $80_s$& $D_4(a_1)$& $58$\\\hline
% $81_p'$& $A_3$& $52$\\\hline
% $60_p'$& $A_2+2A_1$& $50$\\\hline
% $24_p'$& $2A_2$& $48$\\\hline
% $64_p'$& $A_2+A_1$& $46$\\\hline
% $30_p'$& $A_2$& $42$\\\hline
% $20_p'$& $2A_1$& $32$\\\hline
% $6_p'$& $A_1$& $22$\\\hline
% $1_p'$& $0$& $0$\\\hline



% \end{tabular}}
% 	\bigskip
	

% \end{table}


\begin{table}[htbp]\caption{Nilpotent orbits and characters in PyCox  for $ E_6$ }
\label{tabchare6-0}
\centering
\renewcommand{\arraystretch}{1.3}
\setlength\tabcolsep{5pt}
{ \begin{tabular}{|c|c|c|}
\hline Nilpotent orbit $\mathcal{O}$ & cha$(\mathcal{O})$ in PyCox & $\dim \mathcal{O}$ \\ \hline
$E_6$& $1_p$& $72$\\\hline
$E_{6}(a_1)$& $6_p$& $70$\\\hline
$D_5$& $20_p$& $68$\\\hline
$E_{6}(a_3)$& $30_p$& $66$\\\hline
$D_{5}(a_1)$& $64_p$& $64$\\\hline

$A_4+A_1$& $60_p$& $62$\\\hline
$A_4$& $81_p$& $60$\\
$D_4$& $24_p$& $60$\\\hline
$D_4(a_1)$& $80_s$& $58$\\\hline
$A_3$& $81_p'$& $52$\\\hline
$A_2+2A_1$& $60_p'$& $50$\\\hline
$2A_2$& $24_p'$& $48$\\\hline
$A_2+A_1$& $64_p'$& $46$\\\hline
$A_2$& $30_p'$& $42$\\\hline
$2A_1$& $20_p'$& $32$\\\hline
$A_1$& $6_p'$& $22$\\\hline
$0$& $1_p'$& $0$\\\hline

\end{tabular}}
\bigskip

\end{table}







% \begin{table}[htbp]\caption{ Characters in PyCox and nilpotent orbits for $ E_7$ }
% 	\label{tabchare7-0}
% 	\centering
% 	\renewcommand{\arraystretch}{1.3}
% 	\setlength\tabcolsep{5pt}
% 	%\setlength{\leftmargin}{15cm}
% {	\begin{tabular}{|l|l|l|}
% 		%\begin{tabular}{|l|l|l|l|}
% 		\hline   cha$(\mathcal{O})$ in PyCox & Nilpotent orbit $\mathcal{O}$ & $\dim \mathcal{O}$ \\ \hline
%  $1_a$& $E_7$& $126$\\\hline
%  $7_a'$& $E_{7}(a_1)$& $124$\\\hline
%  $27_a$& $E_{7}(a_2)$& $122$\\\hline
%  $21_b'$& $E_{6}$& $120$\\
%  $56_a'$& $E_{7}(a_3)$& $120$\\\hline
%  $120_a$& $E_{6}$& $118$\\\hline
%  $189_b'$& $E_{7}(a_4)$& $116$\\\hline
%  $168_a$& $D_5+A_1$& $114$\\
%  $105_b$& $A_6$& $114$\\
%  $210_a$& $D_6(a_1)$& $114$\\\hline
%   $ 189_c'$  &  $D_5$  & $112$ \\
%   $ 315_a'$  &  $E_7(a_5)$  & $112$ \\\hline
%   $ 405_a$  &  $E_6(a_3)$  & $110$ \\
% \hline
%   $ 378_a'$  &  $D_5(a_1)+A_1$  & $108$ \\
% \hline
%   $ 420_a$  &  $D_5(a_1)$  & $106$ \\
%   $ 210_b$  &  $A_4+A_2$  & $106$ \\
% \hline
%   $ 512_a'$  &  $A_4+A_1$  & $104$ \\
% \hline
%   $ 105_c$  &  $(A_5)''$  & $102$ \\
% \hline
%   $ 210_b'$  &  $A_3+A_2+A_1$  & $100$ \\
%   $ 420_a'$  &  $A_4$  & $100$ \\
% \hline
%   $ 378_a$  &  $A_3+A_2$  & $98$ \\
% \hline
%   $ 405_a'$  &  $D_4(a_1)+A_1$  & $96$ \\

%   $ 105_c'$  &  $D_4$  & $96$ \\
% \hline
%   $315_a$  &  $D_4(a_1)$  & $94$ \\
% \hline
%   $189_c$  &  $(A_3+A_1)''$  & $86$ \\
% \hline
%   $168_a'$  &  $2A_2$  & $84$ \\
%   $210_a'$  &  $A_3$  & $84$ \\
%   $105_b'$  &  $A_3+3A_1$  & $84$ \\
% \hline
%   $189_b$  &  $A_2+2A_1$  & $82$ \\
% \hline
%   $120_a'$  &  $A_2+A_1$  & $76$ \\
% \hline
%   $56_a$  &  $A_2$  & $66$ \\
% \hline

%   $21_b$  &  $(3A_1)''$  & $54$ \\
% \hline
%   $27_a'$  &  $2A_1$  & $52$ \\
% \hline
%   $7_a$  &  $A_1$  & $34$ \\
% \hline
%   $1_a'$  &  $0$  & $0$ \\
% \hline

% \end{tabular}}
% 	\bigskip
% 	\end{table}
\begin{table}[htbp]\caption{Nilpotent orbits and characters in PyCox for $ E_7$ }
	\label{tabchare7-0}
	\centering
	\renewcommand{\arraystretch}{1.3}
	\setlength\tabcolsep{5pt}
	%\setlength{\leftmargin}{15cm}
{	\begin{tabular}{|c|c|c|c|}
		%\begin{tabular}{|l|l|l|l|}
		\hline   Nilpotent orbit $\mathcal{O}$ & cha$(\mathcal{O})$ in PyCox & $\dim \mathcal{O}$ \\ \hline
 $E_7$& $1_a$& $126$\\\hline
 $E_{7}(a_1)$& $7_a'$& $124$\\\hline
 $E_{7}(a_2)$& $27_a$& $122$\\\hline
 $E_{6}$& $21_b'$& $120$\\
 $E_{7}(a_3)$& $56_a'$& $120$\\\hline
 $E_{6}(a_1)$& $120_a$& $118$\\\hline
 $E_{7}(a_4)$& $189_b'$& $116$\\\hline
 $D_5+A_1$& $168_a$& $114$\\
 $A_6$& $105_b$& $114$\\
 $D_6(a_1)$& $210_a$& $114$\\\hline
  $D_5$  &  $ 189_c'$  & $112$ \\
  $E_7(a_5)$  &  $ 315_a'$  & $112$ \\\hline
  $E_6(a_3)$  &  $ 405_a$  & $110$ \\
\hline
  $D_5(a_1)+A_1$  &  $ 378_a'$  & $108$ \\
\hline
  $D_5(a_1)$  &  $ 420_a$  & $106$ \\
  $A_4+A_2$  &  $ 210_b$  & $106$ \\
\hline
  $A_4+A_1$  &  $ 512_a'$  & $104$ \\
\hline
  $(A_5)''$  &  $ 105_c$  & $102$ \\
\hline
  $A_3+A_2+A_1$  &  $ 210_b'$  & $100$ \\
  $A_4$  &  $ 420_a'$  & $100$ \\
\hline
  $A_3+A_2$  &  $ 378_a$  & $98$ \\
\hline
  $D_4(a_1)+A_1$  &  $ 405_a'$  & $96$ \\

  $D_4$  &  $ 105_c'$  & $96$ \\
\hline
  $D_4(a_1)$  &  $315_a$  & $94$ \\
\hline
  $(A_3+A_1)''$  &  $189_c$  & $86$ \\
\hline
  $2A_2$  &  $168_a'$  & $84$ \\
  $A_3$  &  $210_a'$  & $84$ \\
  $A_3+3A_1$  &  $105_b'$  & $84$ \\
\hline
  $A_2+2A_1$  &  $189_b$  & $82$ \\
\hline
  $A_2+A_1$  &  $120_a'$  & $76$ \\
\hline
  $A_2$  &  $56_a$  & $66$ \\
\hline

  $(3A_1)''$  &  $21_b$  & $54$ \\
\hline
  $2A_1$  &  $27_a'$  & $52$ \\
\hline
  $A_1$  &  $7_a$  & $34$ \\
\hline
  $0$  &  $1_a'$  & $0$ \\
\hline

\end{tabular}}
	\bigskip
	\end{table}

% \begin{table}[htbp]\caption{ Characters in PyCox and nilpotent orbits for $ E_8$ }
% 	\label{tabchare8-0}
% 	\centering
% 	\renewcommand{\arraystretch}{1.3}
% 	\setlength\tabcolsep{5pt}
% 	%\setlength{\leftmargin}{15cm}
% {	\begin{tabular}{|l|l|l|}
% 		%\begin{tabular}{|l|l|l|l|}
% 		\hline   cha$(\mathcal{O})$ in PyCox & Nilpotent orbit $\mathcal{O}$ & $\dim \mathcal{O}$ \\ \hline
% % $D_6(a_1)$ 
% % &  $ 210_a'$  &  $A_3$ & $84$  \\

% % $D_5+A_1$ 
% % &  $ 168_a'$  &  $2A_2$ & $84$  \\

% % $A_6$ 
% % &  $ 105_b'$  &  $A_2+3A_1$  & $84$ \\
% % \hline
%  $1_x$& $E_8$& $240$\\\hline
%  $8_z$& $E_{8}(a_1)$& $238$\\\hline
%  $35_x$& $E_{8}(a_2)$& $236$\\\hline
%  $112_z$& $E_{8}(a_3)$& $234$\\\hline
%  $210_x$& $E_{8}(a_4)$& $232$\\\hline
%  $560_z$& $E_{8}(b_4)$& $230$\\\hline
%   $567_x$& $E_{7}(a_1)$& $228$\\
%     $700_x$& $E_{8}(a_5)$& $228$\\\hline
%       $1400_z$& $E_{8}(b_5)$& $226$\\\hline
%       $1400_x$& $E_{8}(a_6)$& $224$\\\hline
%       $3240_z$& $D_{7}(a_1)$& $222$\\\hline

%       $2268_x$& $E_{7}(a_3)$& $220$\\
%       $2240_x$& $E_{8}(b_6)$& $220$\\\hline
% $4096_z$& $E_{6}(a_1)+A_1$& $218$\\\hline
% $525_x$& $E_{6}$& $216$\\
% $4200_x$& $D_{7}(a_2)$& $216$\\\hline
% $4536_z$& $D_5+A_2$& $214$\\
% $2800_z$& $E_6(a_1)$& $214$\\\hline
% $2835_x$& $A_6+A_1$& $212$\\
% $6075_x$& $E_7(a_4)$& $212$\\\hline
% $4200_z$& $A_6$& $210$\\
% $5600_z$& $D_6(a_1)$& $210$\\\hline
% $4480_y$& $E_8(a_7)$& $208$\\\hline
% $2100_y$& $D_5$& $200$\\\hline

% $4200_z'$& $D_6+A_2$& $198$\\
% $5600_z'$& $E_6(a_3)$& $198$\\\hline
% $6075_x'$& $D_5(a_1)+A_1$& $196$\\
% $2835_x'$& $A_4+A_2+A_1$& $196$\\\hline
% $4536_z'$& $A_4+A_2$& $194$\\\hline
% $4200_x'$& $A_4+2A_1$& $192$\\\hline
% $2800_z'$& $D_5(a_1)$& $190$\\\hline
% $4096_x'$& $A_4+A_1$& $188$\\\hline
% $2240_x'$& $D_4(a_1)+A_2$& $184$\\\hline

% \end{tabular}}
% 	\bigskip
	

% \end{table}
\begin{table}[htbp]\caption{Nilpotent orbits and characters in PyCox for $ E_8$ }
	\label{tabchare8-0}
	\centering
	\renewcommand{\arraystretch}{1.3}
	\setlength\tabcolsep{5pt}
{	\begin{tabular}{|c|c|c|}
		\hline   Nilpotent orbit $\mathcal{O}$ & cha$(\mathcal{O})$ in PyCox & $\dim \mathcal{O}$ \\ \hline
 $E_8$& $1_x$& $240$\\\hline
 $E_{8}(a_1)$& $8_z$& $238$\\\hline
 $E_{8}(a_2)$& $35_x$& $236$\\\hline
 $E_{8}(a_3)$& $112_z$& $234$\\\hline
 $E_{8}(a_4)$& $210_x$& $232$\\\hline
 $E_{8}(b_4)$& $560_z$& $230$\\\hline
 $E_{7}(a_1)$& $567_x$& $228$\\
 $E_{8}(a_5)$& $700_x$& $228$\\\hline
 $E_{8}(b_5)$& $1400_z$& $226$\\\hline
 $E_{8}(a_6)$& $1400_x$& $224$\\\hline
 $D_{7}(a_1)$& $3240_z$& $222$\\\hline
 $E_{7}(a_3)$& $2268_x$& $220$\\
 $E_{8}(b_6)$& $2240_x$& $220$\\\hline
 $E_{6}(a_1)+A_1$& $4096_z$& $218$\\\hline
 $E_{6}$& $525_x$& $216$\\
 $D_{7}(a_2)$& $4200_x$& $216$\\\hline
 $D_5+A_2$& $4536_z$& $214$\\
 $E_6(a_1)$& $2800_z$& $214$\\\hline
 $A_6+A_1$& $2835_x$& $212$\\
 $E_7(a_4)$& $6075_x$& $212$\\\hline
 $A_6$& $4200_z$& $210$\\
 $D_6(a_1)$& $5600_z$& $210$\\\hline
 $E_8(a_7)$& $4480_y$& $208$\\\hline
 $D_5$& $2100_y$& $200$\\\hline
 $D_6+A_2$& $4200_z'$& $198$\\
 $E_6(a_3)$& $5600_z'$& $198$\\\hline
 $D_5(a_1)+A_1$& $6075_x'$& $196$\\
 $A_4+A_2+A_1$& $2835_x'$& $196$\\\hline
 $A_4+A_2$& $4536_z'$& $194$\\\hline
 $A_4+2A_1$& $4200_x'$& $192$\\\hline
 $D_5(a_1)$& $2800_z'$& $190$\\\hline
 $A_4+A_1$& $4096_x'$& $188$\\\hline
 $D_4(a_1)+A_2$& $2240_x'$& $184$\\\hline
\end{tabular}}
	\bigskip
\end{table}

% \begin{table}[htbp]\caption{ Characters in PyCox and nilpotent orbits for $ E_8$(continued) }
% 	\label{tabchare8-1}
% 	\centering
% 	\renewcommand{\arraystretch}{1.3}
% 	\setlength\tabcolsep{5pt}
% 	%\setlength{\leftmargin}{15cm}
% {	\begin{tabular}{|l|l|l|}
% 		%\begin{tabular}{|l|l|l|l|}
% 		\hline   cha$(\mathcal{O})$ in PyCox & Nilpotent orbit $\mathcal{O}$ & $\dim \mathcal{O}$ \\ \hline
 
%   $2268_x'$  &  $A_4$  & $180$ \\\hline
%   $3240_z'$  &  $A_3+A_2$  & $178$ \\
% \hline
%   $1400_x'$  &  $D_4(a_1)+A_1$  & $176$ \\
% \hline
%   $525_x'$  &  $D_4$  & $168$ \\
% \hline
%  $1400_z'$  &  $D_4(a_1)$  & $166$ \\
% \hline
%  $700_x'$  &  $2A_2$  & $156$ \\
% \hline
% $567_x'$  &  $A_3$  & $148$ \\
% \hline
% $560_z'$  &  $A_2+2A_1$  & $146$ \\
% \hline
%   $210_x'$  &  $A_2+A_1$  & $136$ \\
% \hline
% $112_z'$  &  $A_2$  & $114$ \\
% \hline
%  $35_x'$  &  $2A_1$  & $92$ \\
% \hline
  
  
  
%   $8_z'$  &  $A_1$  & $58$ \\
% \hline
%   $1_x'$  &  $0$  & $0$ \\
% \hline

% \end{tabular}}
% 	\bigskip
	

% \end{table}

\begin{table}[htbp]\caption{Nilpotent orbits and characters in PyCox for $ E_8$ (continued) }
\label{tabchare8-1}
\centering
\renewcommand{\arraystretch}{1.3}
\setlength\tabcolsep{5pt}
{ \begin{tabular}{|c|c|c|}
\hline Nilpotent orbit $\mathcal{O}$ & cha$(\mathcal{O})$ in PyCox & $\dim \mathcal{O}$ \\ \hline

$A_4$ & $2268_x'$ & $180$ \\\hline
$A_3+A_2$ & $3240_z'$ & $178$ \\
\hline
$D_4(a_1)+A_1$ & $1400_x'$ & $176$ \\
\hline
$D_4$ & $525_x'$ & $168$ \\
\hline
$D_4(a_1)$ & $1400_z'$ & $166$ \\
\hline
$2A_2$ & $700_x'$ & $156$ \\
\hline
$A_3$ & $567_x'$ & $148$ \\
\hline
$A_2+2A_1$ & $560_z'$ & $146$ \\
\hline
$A_2+A_1$ & $210_x'$ & $136$ \\
\hline
$A_2$ & $112_z'$ & $114$ \\
\hline
$2A_1$ & $35_x'$ & $92$ \\
\hline
$A_1$ & $8_z'$ & $58$ \\
\hline
$0$ & $1_x'$ & $0$ \\
\hline

\end{tabular}}
\bigskip
\end{table}


















\begin{table}[htbp]\caption{Sommers duality and weighted diagrams for $ F_4$ }
	\label{tabcharf4}
	\centering
	\renewcommand{\arraystretch}{1.3}
	\setlength\tabcolsep{5pt}
	%\setlength{\leftmargin}{15cm}
\small{	\begin{tabular}{|c|c|c|c|c|c|}
		%\begin{tabular}{|c|c|c|c|}
		\hline $\mathcal{O}_{L^\vee}$  & $\mathcal{O}_{L^\vee}^{\rm BC} $ & diagram of $\mathcal{O}_{L^\vee}^{\rm BC} $ &$\mathcal{O}:=d_{\rm Som}^{G^\vee}(L^\vee, \mathcal{O}_{L^\vee})$ &  cha$(\mathcal{O})$ in Alvis/PyCox &  $\dim \mathcal{O}$ \\ \hline

$B_3$
& $B_3$ & $[2,2,0,0]$& $A_2$  & $\chi_{8,4},8_4$  &  $30$ \\
$C_3$
&$C_3$ & $[1,0,1,2]$&$\tilde{A}_2$  & $\chi_{8,2},8_2$  &   $30$ \\
\hline

$A_3$
& $B_2$ &$[2,0,0,1] $&$B_2$  &  $\chi_{9,2}$  &   $36ns$ \\
$B_4(a_2)$
& $F_4(a_3)$ & $[0,2,0,0]$  & $B_2$  & $\chi_{9,2}$  &   $36ns$ \\
$A_2+\tilde{A}_2$
& $F_4(a_3)$ & $[0,2,0,0]$  & $\tilde{A}_2+A_1$  & $\chi_{6,1}$  &   $36ns$ \\
\hline


$A_2$
& $A_2$ & $[2,0,0,0]$ & $B_3$  & $\chi_{8,1},8_1$  &   $42$ \\
$2A_1+\tilde{A}_1$
&$A_2$ & $[2,0,0,0]$&  $B_3$  &$\chi_{8,1},8_1$  &   $42$ \\
$\tilde{A}_2$
& $\tilde{A}_2$ & $[0,0,0,2]$& $C_3$  &  $\chi_{8,3},8_3$  & $42$ \\


\hline

\hline

\end{tabular}}
	\bigskip
 
 In the last column,    `$ns$' means that the nilpotent orbit is not special.		

\end{table}









\begin{table}[htbp]\caption{Sommers duality and weighted diagrams for $ E_7$ }
	\label{tabchare7}
	\centering
	\renewcommand{\arraystretch}{1.3}
	\setlength\tabcolsep{5pt}
	%\setlength{\leftmargin}{15cm}
{	\begin{tabular}{|c|c|c|c|c|c|}
		%\begin{tabular}{|c|c|c|c|}
		\hline $\mathcal{O}_{L^\vee}$  & $\mathcal{O}_{L^\vee}^{\rm BC} $ & diagram of $\mathcal{O}_{L^\vee}^{\rm BC} $ &$\mathcal{O}:=d_{\rm Som}^{G^\vee}(L^\vee, \mathcal{O}_{L^\vee})$ &  cha$(\mathcal{O})$&  $\dim \mathcal{O}$ \\ \hline

$(A_5+A_1)'$ &$E_6(a_3)$
&$[0,0,2,0,0,2,0]$& $A_3+2A_1$  & $ 216_a$  &   $94ns$ \\
$(A_5+A_1)''$ 
&$(A_5+A_1)''$ 
&$[1,0,0,1,0,1,2]$&$D_4(a_1)$  &  $315_a$  &   $94$ \\
$E_7(a_5)$ 
&$E_7(a_5)$ 
&$[0,0,0,2,0,0,2]$& $D_4(a_1)$  & $315_a$  &   $94$ \\
$D_6(a_2)$ 
&$D_6(a_2)$ 
&$[0,1,1,0,1,0,2]$& $D_4(a_1)$  & $315_a$  &   $94$ \\
\hline


$E_6(a_3)$ 
& $E_6(a_3)$ 
&$[0,0,2,0,0,2,0]$&$D_4(a_1)+A_1$  & $ 405_a'$  &   $96$ \\
$(A_5)'$ 
&$(A_5)'$ 
&$[1,0,0,1,0,2,0]$& $D_4(a_1)+A_1$  &  $ 405_a'$  &  $96$ \\
$(A_5)''$ 
&$(A_5)''$ 
&$[2,0,0,0,0,2,2]$& $D_4$  & $ 105_c'$  &   $96$ \\
\hline

% $D_5(a_1)$ 
% &  $ 420_a'$  &  $A_4$  & $100$ \\
% $D_4+A_1$ 
% &  $ 420_a'$  &  $A_4$  & $100$ \\
% $A_4+A_2$ 
% &  $ 210_b'$  &  $A_3+A_2+A_1$  & $100$ \\
% $A_1+2A_3$ 
% &  $ 210_b'$  &  $A_3+A_2+A_1$  & $100$ \\
% \hline


$2A_3$ 
&$A_4$&$[2,0,0,0,0,2,0]$& $D_4+A_1$  & $ 84_a$  &   $102ns$ \\
$D_4$ 
& $D_4$ 
&$[2,0,2,0,0,0,0]$&$(A_5)''$  &  $ 105_c$  &   $102$ \\
\hline

% $A_4$ 
% &  $ 420_a$  &  $D_5(a_1)$  & $106$ \\
% $A_3+A_2+A_1$ 
% &  $ 210_b$  &  $A_4+A_2$  & $106$ \\
% \hline


$3A_2$ 
&$D_4(a_1)$  &$[0,0,2,0,0,0,0]$& $(A_5+A_1)''$  &  $ 70_a'$  &   $108ns$ \\
$A_3+3A_1$ &$D_4(a_1)+A_1$
&$[0,1,1,0,0,0,1]$& $(A_5)'$  & $ 216_a'$  &   $108ns$ \\
$A_3+A_2$ 
& $A_3+A_2$ 
&$[0,0,0,1,0,1,0]$&$D_5(a_1)+A_1$  &  $ 378_a'$  &  $108$ \\
$D_4(a_1)+2A_1$ 
&  $A_3+A_2$ 
&$[0,0,0,1,0,1,0]$&   $D_5(a_1)+A_1$  &   $ 378_a'$  &  $108$ \\
\hline





$(A_3+2A_1)'$&$D_4(a_1)$  
&$[0,0,2,0,0,0,0]$& $D_6(a_2)$  & $ 280_b$  &   $110ns$ \\
$(A_3+2A_1)''$
&$(A_3+2A_1)''$
&$[1,0,0,0,1,0,1]$& $E_6(a_3)$  & $ 405_a$  &   $110$ \\
$D_4(a_1)+A_1$
&$D_4(a_1)+A_1$
&$[0,1,1,0,0,0,1]$& $E_6(a_3)$  & $ 405_a$  &   $110$ \\
\hline


$(A_3+A_1)''$
& $(A_3+A_1)''$
&$[2,0,0,0,0,0,2]$& $D_5$  & $ 189_c'$  &  $112$ \\

$(A_3+A_1)'$
& $(A_3+A_1)'$
&$[1,0,0,1,0,0,0]$&$E_7(a_5)$  & $ 315_a'$  &   $112$ \\
$2A_2+A_1$
&$2A_2+A_1$
&$[0,0,1,0,0,1,0]$& $E_7(a_5)$  & $ 315_a'$  &   $112$ \\
$D_4(a_1)$
&$D_4(a_1)$
&$[0,0,2,0,0,0,0]$&$E_7(a_5)$  & $ 315_a'$  &   $112$ \\
\hline

% $A_3+2A_1$
% &  $ 105_b$  &  $A_6$  & $114$ \\
% $2A_2$
% &  $ 168_a$  &  $D_5+A_1$  & $114$ \\
% $A_3$
% &  $ 210_a$  &  $D_6(a_1)$  & $114$ \\
% \hline



$(4A_1)'$ &$A_2$
&$[2,0,0,0,0,0,0]$& $D_6$  & $ 35_b$  &   $118ns$ \\
$(4A_1)''$
&  $(4A_1)''$
&$[0,1,0,0,0,0,1]$&  $E_6(a_1)$  &  $ 120_a$  &   $118$ \\
$A_2+A_1$
& $A_2+A_1$
&$[1,0,0,0,0,1,0]$&$E_6(a_1)$  & $ 120_a$  &   $118$ \\
\hline

$A_2$
& $A_2$
&$[2,0,0,0,0,0,0]$&  $E_7(a_3)$  &  $ 56_a'$  &   $120$ \\
$(3A_1)'$
& $(3A_1)'$
&$[0,0,1,0,0,0,0]$&   $E_7(a_3)$  &  $ 56_a'$  &   $120$ \\
$(3A_1)''$
& $(3A_1)''$
&$[0,0,0,0,0,0,2]$& $E_6$  & $ 21_b'$  &   $120$ \\
\hline



\end{tabular}}
	\bigskip
	
We list all the possible $\mathcal{O}_{L^\vee}$ when $\mathcal{O}_{L^\vee}$ or $\mathcal{O}$ contains one prime or two primes.
\end{table}


\begin{table}[htbp]\caption{Sommers duality and weighted diagrams for $ E_8$ }
	\label{tabchare8}
	\centering

	\renewcommand{\arraystretch}{1.3}
	\setlength\tabcolsep{5pt}
	%\setlength{\leftmargin}{15cm}
\small{	\begin{tabular}{|c|c|c|c|c|c|}
		%\begin{tabular}{|c|c|c|c|}
		\hline $\mathcal{O}_{L^\vee}$  & $\mathcal{O}_{L^\vee}^{\rm BC} $ & diagram of $\mathcal{O}_{L^\vee}^{\rm BC} $ &$\mathcal{O}:=d_{\rm Som}^{G^\vee}(L^\vee, \mathcal{O}_{L^\vee})$ &  cha$(\mathcal{O})$ &  $\dim \mathcal{O}$  \\ \hline

$(A_7)''$
&    $E_6(a_1)$ &$[2,0,0,0,0,2,0,2]$&$D_4+A_1$  &  $ 700_{xx}'$  &   $184ns$ \\
% $E_6(a_1)$
% & $D_4+A_1$  &  $ 700_{xx}'$  &  $184ns$ \\

$(A_7)'$
& $(A_7)'$&$[1,0,0,1,0,1,1,0]$& $D_4(a_1)+A_2$  &  $ 2240_{x}'$  &  $184$ \\
$E_8(b_6)$ &$E_8(b_6)$
&$[0,0,0,2,0,0,0,2]$& $D_4(a_1)+A_2$  & $ 2240_{x}'$  &   $184$ \\
$E_6(a_1)+A_2$
& $E_8(b_6)$&$[0,0,0,2,0,0,0,2]$& $D_4(a_1)+A_2$  & $ 2240_{x}'$  &   $184$ \\
\hline

$(A_5+A_1)''$
& $E_6(a_3)$ &$[2,0,0,0,0,0,2,0]$& $D_5+A_1$  &  $ 3200_{x}$  &  $208ns$ \\
$(A_5+A_1)'$
& $(A_5+A_1)'$
&$[1,0,0,1,0,0,0,1]$& $E_8(a_7)$  & $ 4480_y$  &  $208$ \\
$E_7(a_5)$
& $E_7(a_5)$
&$[0,0,0,1,0,1,0,0]$&$E_8(a_7)$  & $ 4480_y$  &   $208$ \\
$E_6(a_3)+A_1$
& $E_6(a_3)+A_1$
&$[1,0,0,0,1,0,1,0]$&$E_8(a_7)$  &  $ 4480_y$  &  $208$ \\
$D_5(a_1)+A_2$
&$D_5(a_1)+A_2$
&$[0,0,1,0,0,1,0,1]$ &$E_8(a_7)$  & $ 4480_y$  &   $208$ \\
$D_6(a_2)$
& $D_6(a_2)$
&$[0,1,1,0,0,0,1,0]$&$E_8(a_7)$  & $ 4480_y$  &   $208$ \\
$A_4+A_3$
&$A_4+A_3$
&$[0,0,0,1,0,0,1,0]$& $E_8(a_7)$  &  $ 4480_y$  &  $208$ \\
 $E_8(a_7)$  & $E_8(a_7)$&$[0,0,0,0,2,0,0,0]$&$E_8(a_7)$  &  $ 4480_y$  &  $208$ \\
\hline

$(2A_3)''$& $A_4$
&$[2,0,0,0,0,0,0,2]$&$D_6$  &  $ 972_{x}$  &   $216ns$ \\
$(2A_3)'$
&$(2A_3)'$
&$[1,0,0,0,1,0,0,0]$& $D_7(a_2)$  &  $ 4200_x$  &   $216$ \\
$A_4+2A_1$
& $A_4+2A_1$
&$[0,0,0,1,0,0,0,1]$&$D_7(a_2)$  &  $ 4200_x$  &  $216$ \\
$D_4$
&$D_4$
&$[0,0,0,0,0,0,2,2]$& $E_6$  &  $ 525_x$  &  $216$ \\
\hline

$(A_3+2A_1)''$ & $D_4(a_1)$
&$[0,0,0,0,0,0,2,0]$& $E_7(a_2)$  &  $ 1344_{x}$  &  $224ns$ \\
$(A_3+2A_1)'$
& $(A_3+2A_1)'$
&$[0,1,0,0,0,0,0,1]$& $E_8(a_6)$  & $ 1400_{x}$  &  $224$ \\
$2A_2+2A_1$
& $2A_2+2A_1$
& $[0,0,0,0,1,0,0,0]$&  $E_8(a_6)$  &  $ 1400_{x}$  &   $224$ \\
$D_4(a_1)+A_1$
&$D_4(a_1)+A_1$
&$[0,1,0,2,0,0,1,0]$ &$E_8(a_6)$  &  $ 1400_{x}$  &  $224$ \\
\hline



% $A_2$
% & $E_7$  & $84_{x}$  &   $232ns$ \\
$(4A_1)''$
& $A_2$ &$[0,0,0,0,0,0,0,2]$& $E_7$  & $84_{x}$  &   $232ns$ \\

$(4A_1)'$
&$(4A_1)'$
&$[0,1,0,0,0,0,0,0]$& $E_8(a_4)$  & $ 210_{x}$  &   $232$ \\
$A_2+A_1$
&$A_2+A_1$
&$[1,0,0,0,0,0,0,1]$& $E_8(a_4)$  & $ 210_{x}$  &   $232$ \\


\hline








\end{tabular}}
\vspace{0.2cm}
 


	


\end{table}






\begin{table}[htbp]\caption{ GKdim and AV for $H:=C_4\subset F_4$}
	\label{tabc4}
	\centering
	\renewcommand{\arraystretch}{1.5}
	\setlength\tabcolsep{5pt}
	
{	\begin{tabular}{|c|c|c|c|c|}
		%\begin{tabular}{|c|c|c|c|c|}
		\hline $\mathrm{GKd} L(\lambda)$ &    $\mathcal{O}({\pi_{w_{\lam}}})$  &    $\mathcal{O}^{\vee}({\pi_{w_{\lam}}})$  & $d_{\rm LS}^{H^{\vee}}(\mathcal{O}^{\vee}(\pi_{w_\lambda}))\rightarrow    O_{L^\vee}$ & $\mathcal{O}_{{\rm Ann}(L(\lambda))}$ \\ \hline
$15$ 
&  $ [2,2,1^4]$ & $ [3,1^6]$  & \small{$[7,1,1]\rightarrow B_3$}&   $A_2$ \\
\hline

$18$ 
&  $ [2^4]$ & $ [3,2,2,1,1]$ & \small{$[5,3,1]\rightarrow B_4(a_2)$ }&   $B_2$ \\
\hline

$21$ &  $ [4,2,2]$ &  $ [5,2,2]$ & \small{$[3,3,1^3]\rightarrow A_2$}&   $B_3$ \\
\hline

\end{tabular}}
	\bigskip
	

\end{table}

\begin{table}[htbp]\caption{ GKdim and AV for $H:=A_2\times \tilde{A}_2\subset F_4$ }
	\label{taba2a2}
	\centering
	\renewcommand{\arraystretch}{1.5}
	\setlength\tabcolsep{5pt}
	
{	\begin{tabular}{|c|c|c|c|c|}
		%\begin{tabular}{|c|c|c|c|c|}
		\hline $\mathrm{GKd} L(\lambda)$ &     $\mathcal{O}({\pi_{w_{\lam}}})$  &    $\mathcal{O}^{\vee}({\pi_{w_{\lam}}})$   & $d_{\rm LS}^{H^{\vee}}(\mathcal{O}^{\vee}(\pi_{w_\lambda}))\rightarrow    O_{L^\vee}$& $\mathcal{O}_{{\rm Ann}(L(\lambda))}$ \\ \hline


$18$ 
&  $ [1^3]\times [1^3]$ & $ [1^3]\times [1^3]$ & \small{$[3]\times [3]\rightarrow A_2+\tilde{A}_2$}&   $\tilde{A}_2+A_1$ \\
\hline

$21$&  $ [1^3]\times [3]$ &  $ [1^3]\times [3]$ & \small{$[3]\times [1^3]\rightarrow \tilde{A}_2$}&   $C_3$ \\
&  $ [3]\times [1^3]$ &  $ [3]\times [1^3]$ & \small{$[1^3]\times [3]\rightarrow {A}_2$}&   $B_3$ \\
\hline

\end{tabular}}
	\bigskip
	

\end{table}

\begin{table}[htbp]\caption{ GKdim and AV for $H:=A_1\times \tilde{A}_3\subset F_4$}
	\label{taba1a3}
	\centering
	\renewcommand{\arraystretch}{1.5}
	\setlength\tabcolsep{5pt}
	
{	\begin{tabular}{|c|c|c|c|c|}
		%\begin{tabular}{|c|c|c|c|c|}
		\hline $\mathrm{GKd} L(\lambda)$ &     $\mathcal{O}({\pi_{w_{\lam}}})$  &    $\mathcal{O}^{\vee}({\pi_{w_{\lam}}})$ &$d_{\rm LS}^{H^{\vee}}(\mathcal{O}^{\vee}(\pi_{w_\lambda}))\rightarrow    O_{L^\vee}$& $\mathcal{O}_{{\rm Ann}(L(\lambda))}$\\ \hline


$18$ 
&  $ [2]\times [1^4]$ &  $ [2]\times [1^4]$ & \small{$[1,1]\times [4]\rightarrow A_3$}&   $B_2$ \\
\hline

$21$&  $ [2]\times [2,1,1]$ &  $ [2]\times [2,1,1]$ & \small{$[1,1]\times [3,1]\rightarrow {A}_2$}&   $B_3$ \\
\hline

\end{tabular}}
	\bigskip
	

\end{table}



\begin{table}[htbp]\caption{ GKdim and AV for $H:=B_3\times \tilde{A}_1\subset F_4$}
	\label{tabb3a1}
	\centering
	\renewcommand{\arraystretch}{1.5}
	\setlength\tabcolsep{5pt}
	
{	\begin{tabular}{|c|c|c|c|c|}
		%\begin{tabular}{|c|c|c|c|c|}
		\hline $\mathrm{GKd} L(\lambda)$ &    $\mathcal{O}({\pi_{w_{\lam}}})$  &    $\mathcal{O}^{\vee}({\pi_{w_{\lam}}})$ &$d_{\rm LS}^{H^{\vee}}(\mathcal{O}^{\vee}(\pi_{w_\lambda}))\rightarrow    O_{L^\vee}$& $\mathcal{O}_{{\rm Ann}(L(\lambda))}$ \\ \hline


$15$ 
&  $ [1^7]\times [2]$ &  $ [1^6]\times [2]$ & \small{$[6]\times [1,1]\rightarrow C_3$ }&   $\tilde{A}_2$ \\
\hline

$21$&  $ [3,3,1]\times [1,1]$ &  $ [3,3]\times [1,1]$ & \small{$[2^3]\times [2]\rightarrow 2A_1+\tilde{A}_1$ }&   $B_3$ \\
&  $ [3,2,2]\times [2]$ &  $ [2,2,2]\times [2]$ & \small{$[3,3]\times [1,1]\rightarrow \tilde{A}_2$  }&   $C_3$ \\
\hline

\end{tabular}}
	\bigskip
	

\end{table}





\begin{table}[htbp]\caption{ GKdim and AV for $H:=D_5\subset E_6$}
	\label{tabd5-1}
	\centering
	\renewcommand{\arraystretch}{1.5}
	\setlength\tabcolsep{5pt}
	
	\begin{tabular}{|c|c|c|c|c|}
		%\begin{tabular}{|c|c|c|c|}
		\hline $\mathrm{GKd} L(\lambda)$ &     $\mathcal{O}({\pi_{w_{\lam}}})$& $d_{\rm LS}^{H}(\mathcal{O}(\pi_{w_\lambda}))\rightarrow O_{L}$ & $\mathcal{O}_{{\rm Ann}(L(\lambda))}$ \\ \hline
     $30$  &  $ [3,3,2,2] $  & $[4,4,1,1]\rightarrow  A_3 $     &   $ A_4$      \\
      &  $ [5,1^5] $  & $[5,1^5]\rightarrow  A_3 $     &   $ A_4$      \\
      \hline
        
         $ 32 $    &  $ [4,4,1,1] $  & $[3,3,2,2]\rightarrow A_2+A_1 $     &   $ D_5(a_1)$      \\\hline
        
\end{tabular}
	\bigskip
	
\end{table}


\begin{table}[htbp]\caption{ GKdim and AV for $H:=A_5\times A_1 \subset E_6$}
	\label{taba5a1}
	\centering
	\renewcommand{\arraystretch}{1.5}
	\setlength\tabcolsep{5pt}
	
	\begin{tabular}{|c|c|c|c|c|}
		%\begin{tabular}{|c|c|c|c|}
		\hline $\mathrm{GKd} L(\lambda)$ &     $\mathcal{O}({\pi_{w_{\lam}}})$& $d_{\rm LS}^{H}(\mathcal{O}(\pi_{w_\lambda}))\rightarrow O_{L}$ & $\mathcal{O}_{{\rm Ann}(L(\lambda))}$ \\ \hline
     $30$  &  $ [3,1^3]\times [2] $  & $[4,1,1]\times [1,1]\rightarrow A_3 $     &   $ A_4$      \\
    &  $ [2^3]\times [2]$  & $[3,3]\times [1,1]\rightarrow 2A_2 $     &   $ D_4$      \\
      \hline
        
         $ 32 $    &  $ [4,1,1]\times [1,1] $  & $[3,1^3]\times [2]\rightarrow A_2+A_1$     &   $ D_5(a_1)$      \\
         &  $ [3,2,1]\times [2] $  & $[3,2,1]\times [1,1]\rightarrow A_2+A_1$     &   $ D_5(a_1)$\\
         &  $ [3,3]\times [1,1] $  & $[2^3]\times [2]\rightarrow4A_1$     &   $ A_5$\\
         \hline
        
\end{tabular}
	\bigskip
	
\end{table}


\begin{table}[htbp]\caption{ GKdim and AV for $H:=A_2\times A_2\times A_2 \subset E_6$}
	\label{taba2a2a2}
	\centering
	\renewcommand{\arraystretch}{1.5}
	\setlength\tabcolsep{5pt}
	
	\begin{tabular}{|c|c|c|c|c|}
		%\begin{tabular}{|c|c|c|c|}
		\hline $\mathrm{GKd} L(\lambda)$ &   $\mathcal{O}({\pi_{w_{\lam}}})$& $d_{\rm LS}^{H}(\mathcal{O}(\pi_{w_\lambda}))\rightarrow O_{L}$ & $\mathcal{O}_{{\rm Ann}(L(\lambda))}$ \\ \hline
     $30$  &  $ [1^3]\times[1^3]\times [3] $  & $[3]\times [3]\times [1^3]\rightarrow 2A_2 $     &   $ D_4$      \\
      \hline
        
         $ 32 $   &  $ [1^3]\times[2,1]\times [3] $  & $[3]\times [2,1]\times [1^3]\rightarrow A_2+A_1 $     &   $ D_5(a_1)$      \\
      \hline
        
\end{tabular}
	\bigskip
	
\end{table}












\begin{table}[htbp]\caption{ GKdim and AV for $H:=A_7\subset E_7$}
	\label{taba7}
	\centering
	\renewcommand{\arraystretch}{1.5}
	\setlength\tabcolsep{5pt}
	
	\begin{tabular}{|c|c|c|c|c|}
		%\begin{tabular}{|c|c|c|c|}
		\hline $\mathrm{GKd} L(\lambda)$ &     $\mathcal{O}({\pi_{w_{\lam}}})$& $d_{\rm LS}^{H}(\mathcal{O}(\pi_{w_\lambda}))\rightarrow O_{L}$ & $\mathcal{O}_{{\rm Ann}(L(\lambda))}$ \\ \hline
$42$ 
&  $ [1^7]$  &{$[7]\rightarrow A_6$  }&   $A_2+3A_1$ \\
\hline

$47$ 
&  $ [2,2,1^4]$  &{$[6,2]\rightarrow (A_5+A_1)'$  }&   $A_3+2A_1$ \\
\hline

$48$ 
&  $ [3,1^5]$  &{$[6,1,1]\rightarrow (A_5)'$ }&   $D_4(a_1)+A_1$ \\
\hline


$50$ 
&  $ [2^3,1,1]$  &$[5,3]\rightarrow A_4+A_2$  &   $A_3+A_2+A_1$ \\
\hline


$51$ 
&  $ [2^4]$  &$[4,4]\rightarrow 2A_3$  &   $D_4+A_1$ \\
\hline

$53$ 
&  $ [4,1^4]$  &$[5,1^3]\rightarrow A_4$  &   $D_5(a_1)$ \\
\hline

$54$ 
&  $ [3,2,2,1]$  &$[4,3,1]\rightarrow A_3+A_2$  &   $D_5(a_1)+A_1$ \\
\hline

$55$ 
&  $ [3,3,1,1]$  &$[4,2,2]\rightarrow (A_3+2A_1)'$  &   $D_6(a_2)$ \\
\hline

$56$ 
&  $ [4,2,1,1]$  &$[4,2,1,1]\rightarrow (A_3+A_1)'$  &   $E_7(a_5)$ \\
& $ [3,3,2]$  &$[3,3,2]\rightarrow 2A_2+A_1$  &   $E_7(a_5)$ \\
\hline

$57$ 
&  $ [5,1^3]$  &$[4,1^4]\rightarrow A_3$  &   $D_6(a_1)$ \\
& $ [4,2,2]$  &$[3,3,1,1]\rightarrow 2A_2$  &   $D_5+A_1$ \\
\hline

$59$ 
&  $ [5,2,1]$  &$[3,2,1^3]\rightarrow A_2+A_1$  &   $E_6(a_1)$ \\
& $ [4,4]$  &$[2^4]\rightarrow (4A_1)'$  &   $D_6$ \\
\hline

$60$ 
&  $ [6,1,1]$  &$[3,1^5]\rightarrow A_2$  &   $E_7(a_3)$ \\
&  $ [5,3]$  &$[2^3,1,1]\rightarrow (3A_1)'$   &   $E_7(a_3)$  \\
\hline



\end{tabular}
	\bigskip
	
\end{table}


\begin{table}[htbp]\caption{ GKdim and AV for $H:=A_5\times A_2\subset E_7$}
	\label{taba5a2}
	\centering
	\renewcommand{\arraystretch}{1.5}
	\setlength\tabcolsep{5pt}
	
	\begin{tabular}{|c|c|c|c|c|}
		%\begin{tabular}{|c|c|c|c|}
		\hline $\mathrm{GKd} L(\lambda)$ &   $\mathcal{O}({\pi_{w_{\lam}}})$& $d_{\rm LS}^{H}(\mathcal{O}(\pi_{w_\lambda}))\rightarrow O_{L}$ & $\mathcal{O}_{{\rm Ann}(L(\lambda))}$ \\ \hline

$47$ 
&  $ [1^6]\times [2,1]$  &$[6]\times [2,1]\rightarrow (A_5+A_1)''$ &   $D_4(a_1)$ \\
\hline

$48$ 
&  $ [1^6]\times [3]$  &{$[6]\times [1^3]\rightarrow (A_5)''$  }&   $D_4$ \\
\hline


$50$ 
&  $ [2,1^4]\times [1^3]$  &$[5,1]\times [3]\rightarrow A_4+A_2$  &   $A_3+A_2+A_1$ \\
\hline




$53$ 
&  $ [2,1^4]\times [3]$  &$[5,1]\times [1^3]\rightarrow A_4$  &   $D_5(a_1)$ \\
\hline

$54$ 
&  $ [2^3]\times [1^3]$  &$[3,3]\times [3]\rightarrow 3A_2$  &   $A_5+A_1$ \\
&  $ [3,1^3]\times [1^3]$  &$[4,1,1]\times [3]\rightarrow A_3+A_2$  &   $D_5(a_1)+A_1$ \\
\hline

$55$ 
&  $ [2,2,1,1]\times [2,1]$  &$[4,2]\times [2,1]\rightarrow (A_3+2A_1)''$  &   $E_6(a_3)$ \\
\hline

$56$ 
&  $ [2,2,1,1]\times [3]$  &  $[4,2]\times [1^3]\rightarrow (A_3+A_1)''$ &   $D_5$ \\

& $ [3,1^3]\times [2,1]$  &$[4,1,1]\times [2,1]\rightarrow (A_3+A_1)'$ &   $E_7(a_5)$ \\

& $ [3,2,1]\times [1^3]$  &$[3,2,1]\times [3]\rightarrow 2A_2+A_1$  &   $E_7(a_5)$ \\
& $ [2^3]\times [2,1]$  &$[3,3]\times [2,1]\rightarrow 2A_2+A_1$  &   $E_7(a_5)$ \\\hline

$57$ 
& $ [3,1^3]\times [3]$  &$[4,1,1]\times [1^3]\rightarrow A_3$  &   $D_6(a_1)$ \\
& $ [2^3]\times [3]$  &$[3,3]\times [1^3]\rightarrow 2A_2$  &   $D_5+A_1$ \\
& $ [4,1,1]\times [1^3]$  &$[3,1^3]\times [3]\rightarrow 2A_2$  &   $D_5+A_1$ \\
& $ [3,3]\times [1^3]$  &$[2^3]\times [3]\rightarrow A_2+3A_1$  &   $A_6$ \\
\hline

$59$ 
& $ [4,1,1]\times [2,1]$  &$[3,1^3]\times [2,1]\rightarrow A_2+A_1$  &   $E_6(a_1)$ \\
& $ [3,3]\times [2,1]$  &$[2^3]\times [2,1]\rightarrow (4A_1)''$ &   $E_6(a_1)$ \\
& $ [5,1]\times [1^3]$  &$[2,1^4]\times [3]\rightarrow A_2+A_1$  &   $E_6(a_1)$ \\
& $ [3,2,1]\times [3]$  &$[3,2,1]\times [1^3]\rightarrow A_2+A_1$  &   $E_6(a_1)$ \\
\hline

$60$ 
& $ [4,1,1]\times [3]$  &$[3,1^3]\times [1^3]\rightarrow A_2$  &   $E_7(a_3)$ \\
& $ [3,3]\times [3]$  &$[2^3]\times [1^3]\rightarrow (3A_1)''$  &   $E_6$ \\
& $ [4,2]\times [2,1]$  &$[2,2,1,1]\times [2,1]\rightarrow (3A_1)'$  &   $E_7(a_3)$ \\
& $ [6]\times [1^3]$  &$[1^6]\times [3]\rightarrow A_2$  &   $E_7(a_3)$ \\
\hline


\end{tabular}
	\bigskip
	
\end{table}

\hspace{3cm}

\begin{table}[htbp]\caption{ GKdim and AV for $H:=A_3\times A_3\times A_1\subset E_7$}
	\label{taba3a3a1}
	
	\renewcommand{\arraystretch}{1.5}
	\setlength\tabcolsep{5pt}
	
\footnotesize{	\begin{tabular}{|c|c|c|c|c|}
		%\begin{tabular}{|c|c|c|c|}
		\hline $\mathrm{GKd} L(\lambda)$ & $\mathcal{O}({\pi_{w_{\lam}}})$& $d_{\rm LS}^{H}(\mathcal{O}(\pi_{w_\lambda}))\rightarrow O_{L}$ & $\mathcal{O}_{{\rm Ann}(L(\lambda))}$ \\ \hline
$50$ 
&  $ [1^4]\times [1^4]\times [1,1]$  &$[4]\times [4]\times [2]\rightarrow 2A_3+A_1$  &   $A_3+A_2+A_1$ \\
\hline

$51$ 
&  $ [1^4]\times [1^4]\times [2]$  &$[4]\times [4]\times [1,1]\rightarrow 2A_3$  &   $D_4+A_1$ \\
\hline


$53$ 
&  $ [1^4]\times [2,1,1]\times [1,1]$  &$[4]\times [3,1]\times [2]\rightarrow A_3+A_2+A_1$  &   $A_4+A_1$ \\
\hline

$54$ 
&  $ [1^4]\times [2,1,1]\times [2]$  &$[4]\times [3,1]\times [1,1]\rightarrow A_3+A_2$  &   $D_5(a_1)+A_1$ \\
\hline

$55$ 
& $ [1^4]\times [2,2]\times [2]$  &$[4]\times [2,2]\times [1,1]\rightarrow (A_3+2A_1)'$    &   $D_6(a_2)$ \\

& $ [1^4]\times [3,1]\times [1,1]$  &$[4]\times [2,1,1]\times [2]\rightarrow (A_3+2A_1)''$    &   $E_6(a_3)$ \\
\hline


$56$ 
& $ [1^4]\times [3,1]\times [2]$  & \small{ $[4]\times [2,1,1]\times [1,1]\rightarrow (A_3+A_1)'$ } &   $E_7(a_5)$ \\

& $ [1^4]\times [4]\times [1,1]$  &$[4]\times [1^4]\times [2]\rightarrow   (A_3+A_1)''$    &   $D_5$ \\

& \small{$ [2,1,1]\times [2,1,1]\times [1,1]$ } &$[3,1]\times [3,1]\times [2]\rightarrow   2A_2+A_1$   &   $E_7(a_5)$ \\
\hline

$57$ 
& $ [1^4]\times [4]\times [2]$  &$[4]\times [1^4]\times [1,1]\rightarrow A_3$   &   $D_6(a_1)$ \\

& $ [2,1,1]\times [2,1,1]\times [2]$  &$[3,1]\times [3,1]\times [1,1]\rightarrow   2A_2$   &   $D_5+A_1$ \\

& $ [2,2]\times [2,1,1]\times [1,1]$  &$[2,2]\times [3,1]\times [2]\rightarrow   A_2+3A_1$   &   $A_6$ \\
\hline

$59$ 
& $ [2,1,1]\times [3,1]\times [2]$  &\small{ $[3,1]\times [2,1,1]\times [1,1]\rightarrow   A_2+A_1$  } &   $E_6(a_1)$ \\

& $ [2,1,1]\times [4]\times [1,1]$  &$[3,1]\times [1^4]\times [2]\rightarrow   A_2+A_1$   &   $E_6(a_1)$ \\
&\small{ $ [3,1]\times [2,2]\times [1,1]$  } &$[2,1,1]\times [2,2]\times [2]\rightarrow   (4A_1)''$  &   $E_6(a_1)$ \\
& $[2,2]\times [2,2]\times [2]$ &$[2,2]\times [2,2]\times [1,1]\rightarrow   (4A_1)'$   &   $D_6$ \\
\hline

$60$ 
& $ [2,1,1]\times [4]\times [2]$  & $[3,1]\times [1^4]\times [1,1]\rightarrow   A_2$   &   $E_6(a_1)$ \\

& $ [3,1]\times [2,2]\times [2]$  &$[2,1,1]\times [2,2]\times [1,1]\rightarrow   (3A_1)'$ &   $E_7(a_3)$ \\
& $ [4]\times [2,2]\times [1,1]$  &$[1^4]\times [2,2]\times [2]\rightarrow   (3A_1)''$   &   $E_6$ \\
& $[3,1]\times [3,1]\times [1,1]$ &$[2,1,1]\times [2,1,1]\times [2]\rightarrow   (3A_1)'$  &   $E_7(a_3)$ \\
\hline


\end{tabular}}
	
 \bigskip
	

By symmetry, we only list one representation of $A_3\times A_3$ in the above table. Such as $[2,1,1]\times [4]$ is listed in the table, but $[4]\times [2,1,1]$ is not.
\end{table}












\begin{table}[htbp]\caption{ GKdim and AV for $H:=D_6\times A_1\subset E_7$}
	\label{tabd6a1}
	\centering
	\renewcommand{\arraystretch}{1.5}
	\setlength\tabcolsep{3pt}
{	
	%\begin{tabular}{|c|c|c|c|c|}
		\begin{tabular}{|c|c|c|c|c|}
		\hline $\mathrm{GKd} L(\lambda)$ &    $\mathcal{O}({\pi_{w_{\lam}}})$& $d_{\rm LS}^{H}(\mathcal{O}(\pi_{w_\lambda}))\rightarrow O_{L}$ & $\mathcal{O}_{{\rm Ann}(L(\lambda))}$ \\ \hline
$ 42$ 
&  $ [3,1^9]\times [1,1]$  &$[9,1^3]\times [2]\rightarrow  D_5+A_1$  &    $2A_2$       \\

&  $ [2,2,1^8]\times [2]$  &$[9,3]\times [1,1]\rightarrow  D_6(a_1)$  &   $A_3$\\
\hline

$ 47$ 
&  $ [2^4,1^4]\times [2]$  &$[7,5]\times [1,1]\rightarrow  D_6(a_2)$  &    $D_4(a_1)$       \\

&  \small{  $ [2^6]^{I}\times [1,1]$ }&$[6,6]^I\times [2]\rightarrow  (A_5+A_1)'$  &   $A_3+2A_1$\\
& \small{   $ [2^6]^{II}\times [1,1]$ } &$[6,6]^{II}\times [2]\rightarrow  (A_5+A_1)''$  &   $D_4(a_1)$\\

\hline

$ 48$ 
&  $ [2^6]^{I}\times [2]$ &$[6,6]^I\times [1,1]\rightarrow  (A_5)'$  &   $D_4(a_1)+A_1$\\
&   $ [2^6]^{II}\times [2]$  &$[6,6]^{II}\times [1,1]\rightarrow (A_5)''$ &   $D_4$\\
\hline

$ 50$ 
&  $ [3,3,1^6]\times [2]$  &$[7,3,1,1]\times [1,1]\rightarrow  D_5(a_1)$  &    $A_4$       \\

&  $ [5,1^7]\times [1,1]$  &$[7,1^5]\times [2]\rightarrow  D_4+A_1$  &   $A_4$\\
\hline
$ 51$ 
&  $ [5,1^7]\times [2]$  & $[7,1^5]\times [1,1]\rightarrow  D_4$  &   $(A_5)''$\\
\hline

$ 53$ 
&  $ [3,3,2,2,1,1]\times [2]$  & $[5,5,1,1]\times [1,1]\rightarrow  A_4$  &   $D_5(a_1)$ \\
&  $ [3^3,1^3]\times [1,1]$  & \small{$[5,3,3,1]\times [2]\rightarrow  A_3+A_2+A_1$  }&   $A_4+A_2$ \\
\hline

$54$ 
&  $ [3^3,1^3]\times [2]$  & \small{$[5,3,3,1]\times [1,1]\rightarrow  A_3+A_2$  }&   $D_5(a_1)+A_1$ \\
&  $ [4,4,1^4]\times [1,1]$  & \small{$[5,3,2,2]\times [2]\rightarrow  D_4(a_1)+2A_1$  }&   $D_5(a_1)+A_1$ \\
&  $ [3^4]\times [1,1]$  & \small{$[4,4,3,1]\times [2]\rightarrow  A_3+3A_1$  }&   $(A_5)'$ \\
\hline

$55$ 
&  $ [4,4,1^4]\times [2]$  & \small{$[5,3,2,2]\times [1,1]\rightarrow  D_4(a_1)+A_1$ }&   $E_6(a_3)$ \\
&  $ [3^4]\times [2]$  & \small{$[4,4,3,1]\times [1,1]\rightarrow  (A_3+2A_1)''$  }&   $E_6(a_3)$ \\
&  $ [5,3,1^4]\times [1,1]$  & \small{$[5,3,1^4]\times [2]\rightarrow  D_4(a_1)+A_1$  }&   $E_6(a_3)$ \\

& \footnotesize{   $ [4,4,2,2]^{I}\times [1,1]$} & \small{$[4,4,2,2]^{I}\times [2]\rightarrow  (A_3+2A_1)'$ }&   $D_6(a_2)$ \\
& \footnotesize{   $ [4,4,2,2]^{II}\times [1,1]$ }& \small{$[4,4,2,2]^{II}\times [2]\rightarrow  (A_3+2A_1)''$ }&   $E_6(a_3)$ \\
\hline

$56$ 
&  $ [5,3,1^4]\times [2]$  & \small{$[5,3,1^4]\times [1,1]\rightarrow  D_4(a_1)$  }&   $E_7(a_5)$ \\
& \footnotesize{  $ [4,4,2,2]^{I}\times [2]$ }& \small{$[4,4,2,2]^I\times [1,1]\rightarrow  (A_3+A_1)'$ }&   $E_7(a_5)$ \\
& \footnotesize{  $ [4,4,2,2]^{II}\times [2]$ } & \small{$[4,4,2,2]^{II}\times [1,1]\rightarrow  (A_3+A_1)''$ }&   $D_5$ \\
&  $ [7,1^5]\times [1,1]$  & \small{$[5,1^7]\times [2]\rightarrow  (A_3+A_1)''$  }&   $D_5$ \\
&  $ [5,3,2,2]\times [1,1]$  & \small{$[4,4,1^4]\times [2]\rightarrow (A_3+A_1)'$  }&   $E_7(a_5)$ \\
&  $ [4,4,3,1]\times [1,1]$  & \small{$[3^4]\times [2]\rightarrow 2A_2+A_1$}&   $E_7(a_5)$ \\
\hline





\end{tabular}}
	\bigskip
	
\end{table}

\begin{table}[htbp]\caption{ GKdim and AV for $H:=D_6\times A_1\subset E_7$ (continued)}
	\centering
	\renewcommand{\arraystretch}{1.5}
	\setlength\tabcolsep{5pt}
	%\setlength{\leftmargin}{15cm}
{	\begin{tabular}{|c|c|c|c|c|}
		%\begin{tabular}{|c|c|c|c|}
		\hline $\mathrm{GKd} L(\lambda)$ &   $\mathcal{O}({\pi_{w_{\lam}}})$& $d_{\rm LS}^{H}(\mathcal{O}(\pi_{w_\lambda}))\rightarrow O_{L}$ & $\mathcal{O}_{{\rm Ann}(L(\lambda))}$ \\ \hline

$57$ 
&  $ [7,1^5]\times [2]$  & \small{$[5,1^7]\times [1,1]\rightarrow  A_3$  }&   $D_6(a_1)$ \\
&  $ [5,3,2,2]\times [2]$  & \small{$[4,4,1^4]\times [1,1]\rightarrow A_3$  }&   $D_6(a_1)$ \\
&  $ [4,4,3,1]\times [2]$  & \small{$[3^4]\times [1,1]\rightarrow 2A_2$}&   $D_5+A_1$ \\

&  $ [5,3,3,1]\times [1,1]$  & \small{$[3^3,1^3]\times [2]\rightarrow A_2+3A_1$  }&   $A_6$ \\
\hline


  
$59$ 
&  $ [7,3,1,1]\times [1,1]$  & \small{$[3,3,1^6]\times [2]\rightarrow A_2+A_1$  }&   $E_6(a_1)$ \\
&  $ [5,5,1,1]\times [2]$  & \small{$[3,3,2,2,1,1]\times [1,1]\rightarrow A_2+A_1$  }&   $E_6(a_1)$ \\
& \small{  $ [6,6]^{I}\times [1,1]$} & \small{$[2^6]^I\times [2]\rightarrow (4A_1)'$}&   $D_6$ \\
& \small{   $ [6,6]^{II}\times [1,1]$ } & \small{$[2^6]^{II}\times [2]\rightarrow (4A_1)''$}&   $E_6(a_1)$ \\

\hline 

$60$ 
&  $ [7,3,1,1]\times [2]$  & \small{$[3,3,1^6]\times [1,1]\rightarrow A_2$  }&   $E_7(a_3)$ \\


& \small{ $ [6,6]^{I}\times [2]$} & \small{$[2^6]^{I}\times [1,1]\rightarrow (3A_1)'$ }&   $E_7(a_3)$ \\
& \small{   $ [6,6]^{II}\times [2]$} & \small{$[2^6]^{II}\times [1,1]\rightarrow (3A_1)''$}  &   $E_6$ \\

&  $ [7,5]\times [1,1]$  & \small{$[2^4,1^4]\times [2]\rightarrow (3A_1)'$  }&   $E_7(a_3)$ \\

&  $ [9,1^3]\times [1,1]$  & \small{$[3,1^9]\times [2]\rightarrow (3A_1)''$  }&   $E_6$ \\
\hline





\end{tabular}}
	\bigskip
	
	\label{tabd6a1-2}

\end{table}

\begin{table}[htbp]\caption{ GKdim and AV for $H:=E_6\subset E_7$}\label{tabe6}
	\centering
	\renewcommand{\arraystretch}{1.5}
	\setlength\tabcolsep{5pt}
	
{	\begin{tabular}{|c|c|c|c|c|}
		%\begin{tabular}{|c|c|c|c|}
		\hline $\mathrm{GKd} L(\lambda)$ & ${\pi_{w_{\lam}}}$ or    $\mathcal{O}({\pi_{w_{\lam}}})$& $d_{\rm LS}^{H}(\mathcal{O}(\pi_{w_\lambda}))\rightarrow O_{L}$ & $\mathcal{O}_{{\rm Ann}(L(\lambda))}$ \\ \hline

$48$ 
&  $ 30_p'$ or $A_2$ & $30_p \rightarrow E_6(a_3)$  &   $D_4(a_1)+A_1$ \\
% &  $ 15_p'$ or $A_2$ & $E_6(a_3)$  &   $D_4(a_1)+A_1$ \\
% &  $ 15_p'$ or $A_2$ & $E_6(a_3)$  &   $D_4(a_1)+A_1$ \\
\hline



$50$ 
&  $ 64_p'$ or $A_2+A_1$ & $64_p \rightarrow D_5(a_1)$  &   $A_4$ \\
\hline



$51$ 
&  $ 24_p'$ or $2A_2$ & $24_p \rightarrow D_4$  &   $(A_5)''$ \\
\hline

$53$ 
&  $ 81_p'$ or $A_3$ & $81_p \rightarrow A_4$  &   $D_5(a_1)$ \\
\hline

$56$ 
% &  $ 90_s$ or $D_4(a_1)$ & $D_4(a_1)$  &   $E_7(a_5)$ \\
&  $ 80_s$ or $D_4(a_1)$ & $80_s \rightarrow D_4(a_1)$  &   $E_7(a_5)$ \\
% &  $ 60_s$ or $D_4(a_1)$ & $D_4(a_1)$  &   $E_7(a_5)$ \\24_p \rightarrow 
% &  $ 20_s$ or $D_4(a_1)$ & $D_4(a_1)$  &   $E_7(a_5)$ \\
% &  $ 10_s$ or $D_4(a_1)$ & $D_4(a_1)$  &   $E_7(a_5)$ \\
\hline




$57$ 
&  $ 81_p$ or $A_4$ & $81_{p}'\rightarrow A_3$  &   $D_6(a_1)$ \\
&  $ 24_p$ or $D_4$ & $24_{p}'\rightarrow 2A_2$  &   $D_5+A_1$ \\
\hline

$59$ 
&  $ 64_p$ or $D_5(a_1)$ & $64_{p}'\rightarrow A_2+A_1$  &   $E_6(a_1)$ \\
\hline

$60$ 
&  $ 30_p$ or $E_6(a_3)$ & $30_{p}'\rightarrow A_2$  &   $E_7(a_3)$ \\
\hline



\end{tabular}}
	\bigskip
	
	
\end{table}





\begin{table}[htbp]\caption{ GKdim and AV for $H:=A_8\subset E_8$}
	\label{taba8}
	\centering
	\renewcommand{\arraystretch}{1.5}
	\setlength\tabcolsep{5pt}
	
	\begin{tabular}{|c|c|c|c|c|}
		%\begin{tabular}{|c|c|c|c|}
		\hline $\mathrm{GKd} L(\lambda)$ &     $\mathcal{O}({\pi_{w_{\lam}}})$& $d_{\rm LS}^{H}(\mathcal{O}(\pi_{w_\lambda}))\rightarrow    O_{L}$ & $\mathcal{O}_{{\rm Ann}(L(\lambda))}$ \\ \hline
		 $ 84$    &  $ [1^9]$  &  $ [9]\rightarrow A_8$     &   $ 2A_2+2A_1$      \\ \hline
	 $ 92 $    &  $ [2,1^7]$  &  $ [8,1]\rightarrow (A_7)'$     &   $ D_4(a_1)+A_2$      \\ 
 \hline	
  
  $ 98$    &  $ [2,2,1^5]$  &$[7,2]\rightarrow      A_6+A_1 $  &   $A_4+A_2+A_1 $        \\ \hline

$ 99$    &  $ [3, 1^6]$  &$[7,1,1]\rightarrow      A_6 $  &   $D_4+A_2 $        \\ \hline
$ 102$    &  $ [2^3, 3]$  &$[6,3]\rightarrow      A_5+A_2 $  &   $E_6(a_3)+A_1 $        \\ \hline
$ 104$    &  $ [3,2,1^4]$  &$[6,2,1]\rightarrow      (A_5+A_1)' $  &   $E_8(a_7) $        \\ 
% &  $ [3,2,1^4]$  &$[6,2,1]\rightarrow      (A_5+A_1)'' $  &   $D_5+A_1 $        \\
   &  $ [2^4,1]$  &$[5,4]\rightarrow      A_4+A_3 $  &   $E_8(a_7) $        \\ \hline

  $ 105$    &  $ [4, 1^5]$  &$[6,1^3]\rightarrow      A_5 $  &   $D_6(a_1)$        \\ \hline
$ 107$    &  $ [3,2,2, 1,1]$  &$[5,3,1]\rightarrow      A_4+A_2 $  &   $D_5+A_2$        \\ \hline
$ 108$    &  $ [3,3,1^3]$  &$[5,2,2]\rightarrow      A_4+2A_1 $  &   $D_7(a_2) $        \\ 
   &  $ [3,2^3]$  &$[4,4,1]\rightarrow      (2A_3)' $  &   $D_7(a_2) $        \\ 
  \hline
$ 109$    &  $ [4, 2, 1^3]$  &$[5,2,1,1]\rightarrow      A_4+A_1 $  &   $E_6(a_1)+A_1$        \\ \hline

$ 110$    &  $ [5,1^4]$  &$[5,1^4]\rightarrow      A_4 $  &   $E_7(a_3) $        \\ 
   &  $ [3,3,2,1]$  &$[4,3,2]\rightarrow      A_3+A_2+A_1 $  &   $E_8(b_6) $        \\ \hline

$ 111$    &  $ [4,2,2,1]$  &$[4,3,1,1]\rightarrow      A_3+A_2 $  &   $D_7(a_1) $        \\ 
   &  $ [3^3]$  &$[3^3]\rightarrow      3A_2 $  &   $E_6+A_1 $        \\ \hline

$ 112$    &  $ [4,3,1,1]$  &$[4,2,2,1]\rightarrow      (A_3+2A_1)' $  &   $E_8(a_6) $        \\ 
\hline
$ 113$    &  $ [5,2,1,1]$  &$[4,2,1,1]\rightarrow      A_3+A_1 $  &   $E_8(b_5) $        \\ 
   &  $ [4,3,2]$  &$[3,3,2,1]\rightarrow      2A_2+A_1 $  &   $E_8(b_5) $        \\ \hline

$ 114$    &  $ [6,1^3]$  &$[4,1^5]\rightarrow      A_3 $  &   $E_7(a_1) $        \\ 
   &  $ [5,2,2]$  &$[3,3,1^3]\rightarrow      2A_2 $  &   $E_8(a_5) $        \\ 
   &  $ [4,4,1]$  &$[3,2^3]\rightarrow      A_2+3A_1 $  &   $E_8(a_5) $        \\ \hline
$ 116$    &  $ [5,4]$  &$[2^4,1]\rightarrow      (4A_1)' $  &   $E_8(a_4) $        \\
% &  $ [5,4]$  &$[2^4,1]\rightarrow      (4A_1)'' $  &   $E_7 $        \\
\hline
  
	\end{tabular}
	\bigskip
	
\end{table}

\begin{table}[htbp]\caption{ GKdim and AV for $H:=A_7\times A_1\subset E_8$}
	\label{taba7a1}
	\centering
	\renewcommand{\arraystretch}{1.5}
	\setlength\tabcolsep{5pt}
	
\small	\begin{tabular}{|c|c|c|c|c|}
		%\begin{tabular}{|c|c|c|c|}
		\hline $\mathrm{GKd} L(\lambda)$ &    $\mathcal{O}({\pi_{w_{\lam}}})$& $d_{\rm LS}^{H}(\mathcal{O}(\pi_{w_\lambda}))\rightarrow   O_{L}$ & $\mathcal{O}_{{\rm Ann}(L(\lambda))}$ \\ \hline
		 $ 92 $    &  $ [1^8]\times [2] $  &  $ [8]\times [1,1]\rightarrow (A_7)''$     &   $ D_4+A_1$      \\ 
 \hline
		$ 98$    &  $ [2,1^6]\times [1^2]$  &$[7,1]\times [2]\rightarrow      A_6+A_1 $  &   $A_4+A_2+A_1 $        \\ \hline

$ 99$    &  $ [2,1^6]\times [2] $  &$[7,1]\times [1,1]\rightarrow      A_6$  &   $D_4+A_2 $        \\ \hline

$ 104$    &  $ [2,2,1^4]\times [2]$  &$[6,2]\times [1,1]\rightarrow      (A_5+A_1)''$ &   $D_5+A_1 $        \\

&  $ [3,1^5]\times [1,1]$  &$[6,1,1]\times [2]\rightarrow      (A_5+A_1)'$  &   $E_8(a_7) $        \\
\hline
$ 105$ 
&  $ [3,1^5]\times [2]$  &$[6,1,1]\times [1,1]\rightarrow  A_5$  &   $D_6(a_1) $       \\\hline

$ 106$ 
&  $ [2^3,1,1]\times [1,1]$  &$[5,3]\times [2]\rightarrow  A_4+A_2+A_1$  &   $A_6+A_1 $       \\\hline

$ 107$ 
&  $ [2^3,1,1]\times [2]$  &$[5,3]\times [1,1]\rightarrow  A_4+A_2$  &   $D_5+A_2 $       \\
&  $ [2^4]\times [1,1]$  &$[4,4]\times [2]\rightarrow  2A_3+A_1$  &   $D_5+A_2 $       \\
\hline

$ 108$ 
&  $ [2^4]\times [2]$  &$[4,4]\times [1,1]\rightarrow  (2A_3)''$ &   $D_6 $       \\

&  $ [3,2,1^3]\times [1,1]$  &$[5,2,1]\times [2]\rightarrow  A_4+2A_1$  &   $D_7(a_2) $       \\
\hline

$ 109$ 
&  $ [3,2,1^3]\times [2]$  &$[5,2,1]\times [1,1]\rightarrow  A_4+A_1$  &   $E_6(a_1)+A_1 $       \\
&  $ [4,1^4]\times [1,1]$  &$[5,1^3]\times [2]\rightarrow  A_4+A_1$  &   $E_6(a_1)+A_1$       \\
\hline

$ 110$ 
&  $ [3,2,2,1]\times [1,1]$  &$[4,3,1]\times [2]\rightarrow  A_3+A_2+A_1$  &   $E_8(b_6) $       \\
&  $ [4,1^4]\times [2]$  &$[5,1^3]\times [1,1]\rightarrow  A_4$  &   $E_7(a_3)$       \\
\hline

$ 111$ 
&  $ [3,2,2,1]\times [2]$  &$[4,3,1]\times [1,1]\rightarrow  A_3+A_2$  &   $D_7(a_1) $       \\
&  $ [3,3,1,1]\times [1,1]$  &$[4,2,2]\times [2]\rightarrow  A_3+3A_1$  &   $D_7(a_1)$       \\
\hline

$ 112$ 
&  $ [3,3,1,1]\times [2]$  &$[4,2,2]\times [1,1]\rightarrow  (A_3+2A_1)''$  &   $E_7(a_2)$\\

&  $ [4,2,1,1]\times [1,1]$  &$[4,2,1,1]\times [2]\rightarrow  (A_3+2A_1)'$  &   $E_8(a_6)$\\

&  $ [3,3,2]\times [1,1]$  &$[3,3,2]\times [2]\rightarrow  2A_2+2A_1$  &    $E_8(a_6)$       \\
\hline


$ 113$ 
&  $ [4,2,1,1]\times [2]$  &$[4,2,1,1]\times [1,1]\rightarrow  A_3+A_1$  &   $E_8(b_5)$\\
&  $ [5,1^3]\times [1,1]$  &$[4,1^4]\times [2]\rightarrow  A_3+A_1$  &   $E_8(b_5)$\\
&  $ [3,3,2]\times [2]$  &$[3,3,2]\times [1,1]\rightarrow  2A_2+A_1$  &    $E_8(b_5)$       \\
&  $ [4,2,2]\times [1,1]$  &$[3,3,1,1]\times [2]\rightarrow  2A_2+A_1$  &   $E_8(b_5)$\\
\hline
$ 114$ 
&  $ [4,3,1]\times [1,1]$  &$[3,2,2,1]\times [2]\rightarrow  A_2+3A_1$  &    $E_8(a_5)$       \\

&  $ [4,2,2]\times [2]$  &$[3,3,1,1]\times [1,1]\rightarrow  2A_2$  &   $E_8(a_5)$\\
&  $ [5,1^3]\times [2]$  &$[4,1^4]\times [1,1]\rightarrow  A_3$  &   $E_7(a_1)$\\
\hline
$ 116$ 
&  $ [5,2,1]\times [2]$  &$[3,2,1^3]\times [1,1]\rightarrow  A_2+A_1$  &    $E_8(a_4)$       \\
&  $ [6,1,1]\times [1,1]$  &$[3,1^5]\times [2]\rightarrow  A_2+A_1$  &   $E_8(a_4)$\\
&  $ [4,4]\times [2]$  &$[2^4]\times [1,1]\rightarrow  (4A_1)''$  &   $E_7$\\


&  $ [5,3]\times [1,1]$  &$[2^3,1,1]\times [2]\rightarrow  (4A_1)'$  &   $E_8(a_4)$\\
\hline


\end{tabular}
	\bigskip
	
\end{table}

% \begin{table}[htbp]
% 	\centering
% 	\renewcommand{\arraystretch}{1.5}
% 	\setlength\tabcolsep{5pt}
	
% 	\begin{tabular}{|c|c|c|c|c|}
% 		%\begin{tabular}{|c|c|c|c|}
% 		\hline $\gkd L(\lambda)$ & ${\pi_{w_{\lam}}}$ or    $\mathcal{O}({\pi_{w_{\lam}}})$& $d_{\rm LS}^{H}(\mathcal{O}(\pi_{w_\lambda}))\rightarrow    L$ & $\mathcal{O}_{{\rm Ann}(L(\lambda))}$ \\ \hline




% \end{tabular}
% 	\bigskip
% 	\caption{ GKdim and AV for $H:=A_7\times A_1\subset E_8$ (continued)}
% 	\label{taba7a1-2}
% \end{table}

% \newpage

\begin{table}[htbp]\caption{ GKdim and AV for $H:=A_5\times A_2\times A_1\subset E_8$}
	\label{taba5a2a1}
	\centering
	\renewcommand{\arraystretch}{1.5}
	\setlength\tabcolsep{5pt}
	
\footnotesize{	\begin{tabular}{|c|c|c|c|c|}
		%\begin{tabular}{|c|c|c|c|}
		\hline $\mathrm{GKd} L(\lambda)$ &   $\mathcal{O}({\pi_{w_{\lam}}})$& $d_{\rm LS}^{H}(\mathcal{O}(\pi_{w_\lambda}))\rightarrow    O_{L}$ & $\mathcal{O}_{{\rm Ann}(L(\lambda))}$ \\ \hline
		 $ 101 $    &  $ [1^6]\times [1^3]\times [1,1] $  &  $ [6]\times [3]\times [2] \rightarrow A_5+A_2+A_1$     &   $ A_5+A_1$      \\ \hline
		$ 102$    &  $ [1^6]\times [1^3]\times [2] $  &$[6]\times [3]\times [1,1]\rightarrow      A_5+A_2 $  &   $E_6(a_3)+A_1 $        \\ \hline

$ 104$    &  $ [1^6]\times [2,1]\times [2] $  &$[6]\times [2,1]\times [1,1]\rightarrow      (A_5+A_1)' $  &   $E_8(a_7) $        \\ 
&  $ [1^6]\times [3]\times [1,1] $  &$[6]\times [1^3]\times [2]\rightarrow      (A_5+A_1)'' $ &   $D_5+A_1$        \\
\hline
$ 105$    &  $ [1^6]\times [3]\times [2] $  &$[6]\times [1^3]\times [1,1]\rightarrow      A_5$  &   $D_6(a_1)$        \\ \hline

$ 106$    &  $ [2,1^4]\times [1^3]\times [1,1] $  &$[5,1]\times [3]\times [2]\rightarrow      A_4+A_2 +A_1$  &   $A_6+A_1 $        \\ \hline

$ 107$    &  $ [2,1^4]\times [1^3]\times [2] $  &$[5,1]\times [3]\times [1,1]\rightarrow      A_4+A_2 $  &   $D_5+A_2 $  \\ \hline

  $ 108$    &  $ [2,1^4]\times [2,1]\times [1,1] $  &$[5,1]\times [2,1]\times [2]\rightarrow      A_4+2A_1 $  &   $D_7(a_2) $  \\ \hline

$ 109$    &  $ [2,1^4]\times [2,1]\times [2] $  &$[5,1]\times [2,1]\times [1,1]\rightarrow      A_4+A_1 $  &   $E_6(a_1)+A_1 $  \\ 
&  $ [2,1^4]\times [3]\times [1,1] $  &$[5,1]\times [1^3]\times [2]\rightarrow      A_4+A_1 $  &   $E_6(a_1)+A_1 $  \\
&  $ [2,2,1,1]\times [1^3]\times [1,1] $  &$[4,2]\times [3]\times [2]\rightarrow      A_3+A_2+2A_1 $  &   $A_7 $  \\
\hline

$ 110$   &  $ [2,2,1,1]\times [1^3]\times [2] $  &$[4,2]\times [3]\times [1,1]\rightarrow      A_3+A_2+A_1 $  &   $E_8(b_6) $  \\
&  $ [3,1^3]\times [1^3]\times [1,1] $  &$[4,1,1]\times [3]\times [2]\rightarrow      A_3+A_2+A_1 $  &   $E_8(b_6) $  \\
&  $ [2^3]\times [1^3]\times [1,1] $  &$[3,3]\times [3]\times [2]\rightarrow      3A_2+A_1 $  &   $E_8(b_6) $  \\
\hline

$ 111$   &  $ [2,2,1,1]\times [2,1]\times [1,1] $  &$[4,2]\times [2,1]\times [2]\rightarrow      A_3+3A_1 $  &   $D_7(a_1) $  \\
&  $ [3,1^3]\times [1^3]\times [2] $  &$[4,1,1]\times [3]\times [1,1]\rightarrow      A_3+A_2 $  &   $D_7(a_1) $  \\
&  $ [2^3]\times [1^3]\times [2] $  &$[3,3]\times [3]\times [1,1]\rightarrow      3A_2 $  &   $E_6+A_1 $  \\
\hline

$ 112$   &  $ [2,2,1,1]\times [2,1]\times [2] $  &$[4,2]\times [2,1]\times [1,1]\rightarrow      (A_3+2A_1)' $  &   $E_8(a_6) $  \\

&  $ [3,1^3]\times [2,1]\times [1,1] $  &$[4,1,1]\times [2,1]\times [2]\rightarrow      (A_3+2A_1)' $  &   $E_8(a_6) $  \\

&  $ [2^3]\times [2,1]\times [1,1] $  &$[3,3]\times [2,1]\times [2]\rightarrow      2A_2+2A_1 $  &   $E_8(a_6) $  \\
&  $ [3,2,1]\times [1^3]\times [1,1] $  &$[3,2,1]\times [3]\times [2]\rightarrow      2A_2+2A_1 $  &   $E_8(a_6) $  \\
\hline

$ 113$   &  $ [2,2,1,1]\times [3]\times [2] $  &$[4,2]\times [1^3]\times [1,1]\rightarrow      A_3+A_1 $  &   $E_8(b_5) $  \\
&  $ [3,1^3]\times [2,1]\times [2] $  &$[4,1,1]\times [2,1]\times [1,1]\rightarrow      A_3+A_1 $  &   $E_8(b_5) $  \\
&  $ [3,1^3]\times [3]\times [1,1] $  &$[4,1,1]\times [1^3]\times [2]\rightarrow      A_3+A_1 $  &   $E_8(b_5) $  \\
&  $ [2^3]\times [2,1]\times [2] $  &$[3,3]\times [2,1]\times [1,1]\rightarrow      2A_2+A_1 $  &   $E_8(b_5) $  \\
&  $ [2^3]\times [3]\times [1,1] $  &$[3,3]\times [1^3]\times [2]\rightarrow      2A_2+A_1 $  &   $E_8(b_5) $  \\
  &  $ [3,2,1]\times [1^3]\times [2] $  &$[3,2,1]\times [3]\times [1,1]\rightarrow      2A_2+A_1 $  &   $E_8(b_5) $  \\
&  $ [4,1,1]\times [1^3]\times [1,1] $  &$[3,1^3]\times [3]\times [2]\rightarrow      2A_2+A_1 $  &   $E_8(b_5) $  \\
&  $ [3,3]\times [1^3]\times [1,1] $  &$[2,2,2]\times [3]\times [2]\rightarrow      A_2+4A_1 $  &   $D_7 $  \\
\hline



	\end{tabular}}
	\bigskip
	
\end{table}


\begin{table}[htbp]\caption{ GKdim and AV for $H:=A_5\times A_2\times A_1\subset E_8$ (continued)}
	\label{taba5a2a1-2}
	\centering
	\renewcommand{\arraystretch}{1.5}
	\setlength\tabcolsep{5pt}
	
\small	\begin{tabular}{|c|c|c|c|c|}
		%\begin{tabular}{|c|c|c|c|}
		\hline $\mathrm{GKd} L(\lambda)$ &    $\mathcal{O}({\pi_{w_{\lam}}})$& $d_{\rm LS}^{H}(\mathcal{O}(\pi_{w_\lambda}))\rightarrow   O_{L}$ & $\mathcal{O}_{{\rm Ann}(L(\lambda))}$ \\ \hline

$114$ &    $ [3,1^3]\times [3]\times [2] $  &$[4,1,1]\times [1^3]\times [1,1]\rightarrow      A_3 $  &   $E_7(a_1) $  \\
&  $ [4,1,1]\times [1^3]\times [2] $  &$[3,1^3]\times [3]\times [1,1]\rightarrow      2A_2 $  &   $E_8(a_5) $  \\
&  $ [2^3]\times [3]\times [2] $  &$[3,3]\times [1^3]\times [1,1]\rightarrow      2A_2 $  &   $E_8(a_5) $  \\
&  $ [3,2,1]\times [2,1]\times [1,1] $  &$[3,2,1]\times [2,1]\times [2]\rightarrow      A_2+3A_1 $  &   $E_8(a_5) $  \\


&  $ [3,3]\times [1^3]\times [2] $  &$[2^3]\times [3]\times [1,1]\rightarrow      A_2+3A_1 $  &   $E_8(a_5) $  \\
&  $ [4,2]\times [1^3]\times [1,1] $  &$[2,2,1,1]\times [3]\times [2]\rightarrow      A_2+3A_1 $  &   $E_8(a_5) $  \\
\hline




$116$ &     $ [3,2,1]\times [3]\times [2] $  &$[3,2,1]\times [1^3]\times [1,1]\rightarrow      A_2+A_1 $  &   $E_8(a_4) $  \\
&  $ [4,1,1]\times [2,1]\times [2] $  &$[3,1^3]\times [2,1]\times [1,1]\rightarrow      A_2+A_1 $  &   $E_8(a_4) $  \\
&  $ [4,1,1]\times [3]\times [1,1] $  &$[3,1^3]\times [1^3]\times [2]\rightarrow      A_2+A_1 $  &   $E_8(a_4) $  \\
&  $ [5,1]\times [1^3]\times [2] $  &$[2,1^4]\times [3]\times [1,1]\rightarrow      A_2 +A_1$  &   $E_8(a_4) $  \\
% &  $ [3,3]\times [3]\times [1,1] $  &$[2^3]\times [1^3]\times [2]\rightarrow      4A_1 $  &   $E_8(a_4) $  \\
&  $ [3,3]\times [3]\times [1,1] $  &$[2^3]\times [1^3]\times [2]\rightarrow      (4A_1)'' $  &   $E_7 $  \\
&  $ [3,3]\times [2,1]\times [2] $  &$[2^3]\times [2,1]\times [1,1]\rightarrow      (4A_1)' $ &   $E_8(a_4) $  \\
% &  $ [3,3]\times [2,1]\times [2] $  &$[2^3]\times [2,1]\times [1,1]\rightarrow      (4A_1)'' $  &   $E_7 $  \\
&  $ [4,2]\times [2,1]\times [1,1] $  &$[2,2,1,1]\times [2,1]\times [2]\rightarrow      (4A_1)' $  &   $E_8(a_4) $  \\
% &  $ [4,2]\times [2,1]\times [1,1] $  &$[2,2,1,1]\times [2,1]\times [2]\rightarrow      (4A_1)'' $  &   $E_7$  \\
\hline
\end{tabular}
	\bigskip
	
\end{table}
\newpage 






\begin{table}[htbp]\caption{ GKdim and AV for $H:=A_4\times A_4\subset E_8$}
	\label{taba4a4}
	\centering
	\renewcommand{\arraystretch}{1.5}
	\setlength\tabcolsep{5pt}
	
	\begin{tabular}{|c|c|c|c|c|}
		%\begin{tabular}{|c|c|c|c|}
		\hline $\mathrm{GKd} L(\lambda)$ &    $\mathcal{O}({\pi_{w_{\lam}}})$& $d_{\rm LS}^{H}(\mathcal{O}(\pi_{w_\lambda}))\rightarrow   O_{L}$ & $\mathcal{O}_{{\rm Ann}(L(\lambda))}$ \\ \hline
		 $ 100 $    &  $ [1^5]\times [1^5]$  &  $ [5]\times [5] \rightarrow 2A_4$     &   $ A_4+A_3$      \\ \hline
		$ 104$   
  &  $ [2,1^3]\times [1^5]$  &$[4,1]\times [5]\rightarrow      A_4+A_3 $  &   $E_8(a_7) $        \\  \hline

$ 106$    &  $ [1^5]\times [2,2,1] $  & $[5]\times [3,2]\rightarrow      A_4+A_2+A_1 $  &   $A_6+A_1 $        \\
\hline

$ 107$    &  $ [1^5]\times [3,1,1] $  & $[5]\times [3,1,1]\rightarrow      A_4+A_2 $  &   $D_5+A_2 $        \\
\hline
$ 108$    &  $ [1^5]\times [3,2] $  & $[5]\times [2,2,1]\rightarrow      A_4+2A_1 $  &   $D_7(a_2) $        \\
&  $ [2,1^3]\times [2,1^3] $  & $[4,1]\times [4,1]\rightarrow      (2A_3)' $&   $D_7(a_2) $        \\
% &  $ [2,1^3]\times [2,1^3] $  & $[4,1]\times [4,1]\rightarrow      (2A_3)'' $  &   $D_6 $        \\
\hline
$ 109$    &  $ [1^5]\times [4,1] $  & $[5]\times [2,1^3]\rightarrow      A_4+A_1 $  &   $E_6(a_1)+A_1 $        \\
\hline

$ 110$    &  $ [1^5]\times [5] $  & $[5]\times [1^5]\rightarrow      A_4 $  &   $E_7(a_3)$        \\
&  $ [2,1^3]\times [2,2,1] $  & $[4,1]\times [3,2]\rightarrow      A_3+A_2+A_1 $  &   $E_8(b_6)$        \\
\hline

$ 111$    &  $ [2,1^3]\times [3,1,1] $  & $[4,1]\times [3,1,1]\rightarrow      A_3+A_2 $  &   $D_7(a_1) $        \\
\hline

$ 112$    &  $ [2,2,1]\times [2,2,1] $  & $[3,2]\times [3,2]\rightarrow      2A_2+2A_1 $  &   $E_8(a_6)$        \\

&  $ [2,1^3]\times [3,2] $  & $[4,1]\times [2,2,1]\rightarrow      (A_3+2A_1)' $  &   $E_8(a_6)$        \\
% &  $ [2,1^3]\times [3,2] $  & $[4,1]\times [2,2,1]\rightarrow      (A_3+2A_1)'' $  &   $E_7(a_2)$        \\
\hline

$ 113$    
&  $ [2,1^3]\times [4,1] $  & $[4,1]\times [2,1^3]\rightarrow      A_3+A_1 $  &   $E_8(b_5)$        \\
&  $ [2,2,1]\times [3,1,1] $  & $[3,2]\times [3,1,1]\rightarrow      2A_2+A_1 $  &   $E_8(b_5)$        \\
\hline

$ 114$    
&  $ [2,1^3]\times [5] $  & $[4,1]\times [1^5]\rightarrow      A_3 $  &   $E_7(a_1)$        \\
&  $ [2,2,1]\times [3,2] $  & $[3,2]\times [2,2,1]\rightarrow      A_2+3A_1 $  &   $E_8(a_5)$        \\
&  $ [3,1,1]\times [3,1,1] $  & $[3,1,1]\times [3,1,1]\rightarrow      2A_2$  &   $E_8(a_5)$        \\
\hline

$ 116$    
&  $ [2,2,1]\times [5] $  & $[3,2]\times [1^5]\rightarrow      A_2+A_1 $  &   $E_8(a_4)$        \\
&  $ [3,1,1]\times [4,1] $  & $[3,1,1]\times [2,1^3]\rightarrow      A_2+A_1 $  &   $E_8(a_4)$        \\
&  $ [3,2]\times [3,2] $  & $[2,2,1]\times [2,2,1]\rightarrow      (4A_1)'$  &   $E_8(a_4)$        \\
% &  $ [3,2]\times [3,2] $  & $[2,2,1]\times [2,2,1]\rightarrow      (4A_1)''$  &   $E_7$        \\
\hline
	\end{tabular}
	\bigskip
	
\end{table}

\hspace{15cm}

\newpage




\begin{table}[htbp]\caption{ GKdim and AV for $H:=D_8\subset E_8$}
	\label{tabd8-1}
	\centering
	\renewcommand{\arraystretch}{1.5}
	\setlength\tabcolsep{5pt}
	
{	\begin{tabular}{|c|c|c|c|c|}
		%\begin{tabular}{|c|c|c|c|}
		\hline $\mathrm{GKd} L(\lambda)$ & $\mathcal{O}(\pi_{w_{\lam}})$& $d_{\rm LS}^{H}(\mathcal{O}(\pi_{w_\lambda}))\rightarrow    O_{L}$& $\mathcal{O}_{{\rm Ann}(L(\lambda))}$ \\ \hline
		 $ 92 $    &   $[2^8]^I$ & $[8,8]^I\rightarrow    (A_7)' $     &   $ D_4(a_1)+A_2$      \\ 
		   &   $[2^8]^{II}$ &    $[8,8]^{II}\rightarrow     (A_7)'' $   &   $D_4+A_1 $        \\ \hline
     $ 98 $    &  $ [4,4,1^8] $  & $[9,3,2,2]\rightarrow     D_6(a_1)+A_1 $     &   $ D_5(a_1)+A_1$      \\
     \hline
      $ 99 $    &  $ [5, 3, 1^8] $  & $[9,3,1^4]\rightarrow     D_6(a_1) $     &   $ E_6(a_3)$      \\ 
          % &  $ [5,2,2,1^7] $  & $[9,3,1^4]\rightarrow     ? $     &   $ E_6(a_3)$      \\
          % &  $ [3,3, 2^4, 1,1] $  & $[7,7,1,1]\rightarrow     A_6 $     &   $ D_4+A_2$      \\
          &  $ [3, 3, 2^4, 1, 1] $  & $[7,7,1,1]\rightarrow     A_6$     &   $ D_4+A_2$      \\ \hline
       $ 100 $    &  $ [7,1^9] $  & $[9,1^7]\rightarrow     D_5 $     &   $ D_5$      \\ \hline
        $ 102 $    &  $ [3^4,1^4] $  & $[7,5,3,1]\rightarrow    D_8(a_5)$     &   $ D_6(a_2)$     
        % \\ 
        % &  $ [3^3,2^2,1^3] $  & $[7,5,3,1]\rightarrow     ? $     &   $ ?$ 
        \\\hline
         $ 104 $    &  $ [3^5,1] $  & $[5^3,1]\rightarrow     A_4+A_3$     &   $ E_8(a_7)$      \\ 
         &  $ [5,3,2,2,1^4] $  & $[7,5,1^4]\rightarrow    D_6(a_2)  $     &   $ E_8(a_7)$      \\
         % &  $ [5, 2^4, 1^3] $  & $[7,5,1^4]\rightarrow     ? $     &   $ E_8(a_7)$      \\
         &  $ [4,4,3,1^5] $  & $[7,3^3]\rightarrow     D_5(a_1)+A_2 $     &   $ E_8(a_7)$      \\
         &     $ [4,4,2^4]^I $ & $[6,6,2,2]^I\rightarrow    (A_5+A_1)' $     &   $ E_8(a_7)$      \\
         &  $ [4,4,2^4]^{II} $  & $[6,6,2,2]^{II}\rightarrow    (A_5+A_1)'' $     &   $ D_5+A_1$      \\\hline
         $ 105 $    &  $ [5, 3, 3, 1^5] $  & $[7,3,3,1^3]\rightarrow     D_4+A_2 $     &   $ A_6$      \\ 
        &  $ [5, 3, 2^4] $  & $[6,6,1^4]\rightarrow     A_5 $     &   $ D_6(a_1)$      \\\hline
        $ 106 $    &  $ [5, 5, 1^6] $  & $[7,3,2,2,1,1]\rightarrow     D_5(a_1)+A_1 $     &   $ E_7(a_4)$        \\\hline
        $ 107 $    &  $ [4,4, 3, 3,1,1] $  & $[5,5,3,3]\rightarrow     A_4+A_2 $     &   $ D_5+A_2$      \\ 
         
          % &  $ [7,2,2, 1^5] $  & $[7,3,1^6]\rightarrow     ? $     &   $ E_6(a_1)$      \\
          &  $ [7,3,1^6] $  & $[7,3,1^6]\rightarrow     D_5(a_1)$     &   $ E_6(a_1)$      \\ \hline

          $ 108 $    &   $ [4^4]^{I} $  & $[4^4]^I\rightarrow    (2A_3)' $     &   $ D_7(a_2)$      \\ 
          &  $ [4^4]^{II} $  & $ [4^4]^{II} \rightarrow   (2A_3)'' $     &   $ D_6$      \\
         &  $ [5,3^3,1,1] $  & $[5,5,3,1^3]\rightarrow     A_4+2A_1 $     &   $ D_7(a_2)$      \\
         % &  $ [5, 3,3, 2,2,1] $  & $[5,5,3,1^3]\rightarrow     ? $     &   $ D_7(a_2)$      \\
         &  $ [9,1^7] $  & $[7,1^9]\rightarrow     D_4 $     &   $ E_6$      \\\hline
         $ 109 $    &  $ [5,5,2,2,1,1] $  & $[5,5,2,2,1,1]\rightarrow    A_4+A_1 $     &   $ E_6(a_1)+A_1$      \\\hline

         $ 110 $    &  $ [5,5, 3, 1^3] $  & $[5,3^3,1,1]\rightarrow     D_4(a_1)+A_2 $     &   $ E_8(b_6)$      \\ 
          &  $ [7,3,2,2,1,1] $  & $[5,5,1^6]\rightarrow     A_4 $     &   $ E_7(a_3)$      \\ \hline
  $ 111 $    &  $ [7, 3, 3, 1^3] $  & $[5,3,3,1^5]\rightarrow     A_3+A_2 $     &   $ D_7(a_1)$      \\ 
          &  $ [5,5,3,3] $  & $[4,4,3,3,1,1]\rightarrow     A_3+A_2 $     &   $D_7(a_1)$      \\
          &  $ [6,6,1^4] $  & $[5,3,2^4]\rightarrow     D_4(a_1)+2A_1$     &   $D_7(a_1)$      \\ \hline
        
        
\end{tabular}}
	\bigskip
	
\end{table}

\begin{table}[htbp]\caption{GKdim and AV for $H:=D_8 \subset E_8$ (continued)}
	\label{tabd8-2}
	\centering
	\renewcommand{\arraystretch}{1.5}
	\setlength\tabcolsep{5pt}
	
	\begin{tabular}{|c|c|c|c|c|}
		%\begin{tabular}{|c|c|c|c|}
		\hline $\mathrm{GKd} L(\lambda)$ &  $\mathcal{O}(\pi_{w_{\lam}})$& $d_{\rm LS}^{H}(\mathcal{O}(\pi_{w_\lambda}))\rightarrow    O_{L}$ & $\mathcal{O}_{{\rm Ann}(L(\lambda))}$ \\ \hline
         
          $ 112 $    &  $ [5^3,1] $  & $[3^5,1]\rightarrow     2A_2+2A_1$     &   $ E_8(a_6)$      \\ 
         &  $ [7,5,1^4] $  & $[5,3,2,2,1^4]\rightarrow    D_4(a_1)+A_1$     &   $ E_8(a_6)$      \\
         % &  $ [5, 2^4, 1^3] $  & $[7,5,1^4]\rightarrow     ? $     &   $ E_8(a_7)$      \\
         &  $ [7,3^3] $  & $[4,4,3,1^5]\rightarrow     (A_3+2A_1)' $   &   $ E_8(a_6)$      \\
         % &  $ [7,3^3] $  & $[4,4,3,1^5]\rightarrow     (A_3+2A_1)''? $     &   $ E_7(a_2)$      \\
         &   $ [6,6,2,2]^I $  & $[4,4,2^4]^I\rightarrow     (A_3+2A_1)' $     &   $ E_8(a_6)$      \\
         &    $ [6,6,2,2]^{II} $  & $[4,4,2^4]^{II}\rightarrow    (A_3+2A_1)'' $     &   $ E_7(a_2)$   \\\hline 

         $ 113 $    &  $ [6, 6, 3, 1] $  & $[3^4,2,2]\rightarrow     2A_2+A_1 $     &   $ E_8(b_5)$      \\ 
          &  $ [9,3,1^4] $  & $[5,3,1^8]\rightarrow     D_4(a_1) $     &   $E_8(b_5)$      \\
          &  $ [7,5,2,2] $  & $[4,4,2,2,1^4]\rightarrow     A_3+A_1$     &   $E_8(b_5)$      \\ \hline

          $ 114 $    &  $ [7, 5, 3, 1] $  & $[3^4,1^4]\rightarrow     2A_2 $     &   $ E_8(a_5)$      \\ 
          &  $ [9,3,2,2] $  & $[4,4,1^8]\rightarrow     A_3 $     &   $E_7(a_1)$      \\
          &  $ [11,1^5] $  & $[5,1^{11}]\rightarrow     A_3$     &   $E_7(a_1)$      \\ 
          % &  $ [11,5] $  & $[2^4,1^{8}]\rightarrow     2A_2\color{red}{?}$     &   $E_8(a_5)\color{red}{?}$      \\
          \hline
	
 $ 116 $    &  $ [9,5,1,1] $  & $[3,3,2,2,1^6]\rightarrow     A_2+A_1$     &   $ E_8(a_4)$      \\ 
         &  $ [8,8]^I $  & $[2^8]^I\rightarrow    (4A_1)' $     &   $ E_8(a_4)$      \\
        % & $([4], [4])^+$ or  $ [8,8]^{II} $  & $[2^8]^{II}\rightarrow    (4A_1)'' $     &   $ E_7$      \\ 
        &   $ [8,8]^{II} $  & $[2^8]^{II}\rightarrow    (4A_1)'' $     &   $ E_7$      \\ 
        
        \hline
 
 
 
 
 \end{tabular}
	\bigskip
	
\end{table}



\newpage 

\begin{table}[htbp]\caption{ GKdim and AV for $H:=D_5\times A_3\subset E_8$}
	\label{tabd5a3}
	\centering
	\renewcommand{\arraystretch}{1.5}
	\setlength\tabcolsep{5pt}
	
	\begin{tabular}{|c|c|c|c|c|}
		%\begin{tabular}{|c|c|c|c|}
		\hline $\mathrm{GKd} L(\lambda)$ &     $\mathcal{O}({\pi_{w_{\lam}}})$& $d_{\rm LS}^{H}(\mathcal{O}(\pi_{w_\lambda}))\rightarrow    O_{L}$ & $\mathcal{O}_{{\rm Ann}(L(\lambda))}$ \\ \hline
		 $ 94 $    &  $ [1^{10}]\times [1^4]$  &  $ [9,1]\times [4] \rightarrow D_5+A_3$     &   $ 2A_3$      \\ \hline
	$ 98 $    &  $ [1^{10}]\times [2,2]$  &  $ [9,1]\times [2,2] \rightarrow D_5+2A_1$     &   $ A_5$      \\ \hline
  $ 99 $    &  $ [1^{10}]\times [3,1]$  &  $ [9,1]\times [2,1,1] \rightarrow D_5+A_1$     &   $ E_6(a_3)$      \\ \hline
  $ 100 $    &  $ [1^{10}]\times [4]$  &  $ [9,1]\times [1^4] \rightarrow D_5$     &   $ D_5$      \\ \hline

  $ 101 $    &  $ [2,2,1^{6}]\times [1^4]$  &  $ [7,3]\times [4] \rightarrow D_5(a_1)+A_3$     &   $ D_5(a_1)+A_2$      \\ \hline
  
  $ 102$    &  $ [3,1^7]\times [1^4] $  &$[7,1^3]\times [4]\rightarrow      D_4+A_3 $  &   $D_6(a_2) $        \\ \hline

$ 104$    &  $ [2^4,1,1]\times [1^4] $  & $[5,5]\times [4]\rightarrow      A_4+A_3 $  &   $E_8(a_7) $        \\
&  $ [2,2,1^6]\times [2,1,1] $  & $[7,3]\times [3,1]\rightarrow      D_5(a_1)+A_2$  &   $E_8(a_7) $        \\
\hline
$ 105$    &  $ [2,2,1^6]\times [2,2] $  & $[7,3]\times [2,2]\rightarrow      D_5(a_1)+2A_1$  &   $A_6 $        \\
&  $ [3,1^7]\times [2,1,1] $  & $[7,1^3]\times [3,1]\rightarrow      D_4+A_2 $  &   $A_6 $        \\
\hline

$ 106$    &  $ [2,2,1^6]\times [3,1] $  & $[7,3]\times [2,1,1]\rightarrow      D_5(a_1)+A_1$  &   $E_7(a_4) $        \\
&  $ [3,1^7]\times [2,2] $  & $[7,1^3]\times [2,2]\rightarrow      D_4+2A_1 $  &   $E_7(a_4) $        \\
\hline

$ 107$   &  $ [3,1^7]\times [3,1] $  & $[7,1^3]\times [2,1,1]\rightarrow      D_4+A_1 $  &   $E_6(a_1)$        \\ 
&  $ [2,2,1^6]\times [4] $  & $[7,3]\times [1^4]\rightarrow      D_5(a_1)$  &   $E_6(a_1) $        \\
&  $ [3,3,1^4]\times [1^4] $  & $[5,3,1,1]\times [4]\rightarrow      D_4(a_1)+A_3 $  &   $D_5+A_2$        \\

&  $ [2^4,1,1]\times [2,1,1] $  & $[5,5]\times [3,1]\rightarrow      A_4+A_2 $  &   $D_5+A_2$        \\
% &  $ [3,2,2,1^3]\times [1^4] $  & $[5,3,1,1]\times [4]\rightarrow      A_4+A_2 $  &   $D_5+A_2$        \\
\hline

$ 108$   &  $ [3,1^7]\times [4] $  & $[7,1^3]\times [1^4]\rightarrow      D_4 $  &   $E_6$        \\ 
&  $ [2^4,1,1]\times [2,2] $  & $[5,5]\times [2,2]\rightarrow      A_4+2A_1$  &   $D_7(a_2) $        \\
& \small{  $ [3,3,2,2]\times [1^4] $ } & $[4,4,1,1]\times [4]\rightarrow      (2A_3)' $ &   $D_7(a_2)$        \\

% &  $ [5,1^5]\times [1^4] $ or $ ([3,1,1],\emptyset)\times [1^4] $  & $[5,1^5]\times [4]\rightarrow      2A_3\color{red}{?}$  &   $D_7(a_2)\color{red}{?}$        \\
&  $ [5,1^5]\times [1^4] $   & $[5,1^5]\times [4]\rightarrow      (2A_3)''$   &   $D_6$        \\
\hline

$ 109$  
&  $ [2^4,1,1]\times [3,1] $  & $[5,5]\times [2,1,1]\rightarrow      A_4+A_1$  &   $E_6(a_1) +A_1$        \\
&  $ [3^3,1]\times [1^4] $ & $[3^3,1]\times [4]\rightarrow      A_3+A_2+2A_1 $  &   $A_7$        \\ \hline
$ 110$  
&  $ [2^4,1,1]\times [4] $  & $[5,5]\times [1^4]\rightarrow      A_4$  &   $E_7(a_3)$        \\
&  $ [3,3,1^4]\times [2,1,1] $ & $[5,3,1,1]\times [3,1]\rightarrow      D_4(a_1)+A_2$  &   $E_8(b_6)$        \\
&  $ [4,4,1,1]\times [1^4] $ & $[4,2^3]\times [4]\rightarrow      A_3+A_2+A_1$  &   $E_8(b_6)$        \\
\hline

$ 111$  
&  $ [3,3,1^4]\times [2,2] $ & \footnotesize{$[5,3,1,1]\times [2,2]\rightarrow      D_4(a_1)+2A_1$ } &   $D_7(a_1)$        \\
&  $ [3,3,2,2]\times [2,1,1] $ & $[4,4,1,1]\times [3,1]\rightarrow      A_3+A_2$  &   $D_7(a_1)$        \\
&  $ [5,1^5]\times [2,1,1] $ & $[5,1^5]\times [3,1]\rightarrow      A_3+A_2$  &   $D_7(a_1)$        \\
&  $ [5,3,1,1]\times [1^4] $ & $[3,3,1^4]\times [4]\rightarrow      A_3+A_2$  &   $D_7(a_1)$        \\
\hline


  
	\end{tabular}
	\bigskip
	
\end{table}

\begin{table}[htbp]\caption{ GKdim and AV for $H:=D_5\times A_3\subset E_8$ (continued)}
	\label{tabd5a3-2}
	\centering
	\renewcommand{\arraystretch}{1.5}
	\setlength\tabcolsep{5pt}
	
	\begin{tabular}{|c|c|c|c|c|}
		%\begin{tabular}{|c|c|c|c|}
		\hline $\mathrm{GKd} L(\lambda)$ & $\mathcal{O}(\pi_{w_{\lam}})$& $d_{\rm LS}^{H}(\mathcal{O}(\pi_{w_\lambda}))\rightarrow    O_{L}$& $\mathcal{O}_{{\rm Ann}(L(\lambda))}$ \\ \hline

$ 112$  
&  $ [3,3,1^4]\times [3,1] $ & \footnotesize{ $[5,3,1,1]\times [2,1,1]\rightarrow      D_4(a_1)+A_1$ } &   $E_8(a_6)$        \\
&  $ [3^3,1]\times [2,1,1] $ & $[3^3,1]\times [3,1]\rightarrow      2A_2+2A_1$  &   $E_8(a_6)$        \\

&  \footnotesize{$ [3,3,2,2]\times [2,2] $ } & $[4,4,1,1]\times [2,2]\rightarrow      (A_3+2A_1)'$ &   $E_8(a_6)$        \\
&\small{  $ [5,1^5]\times [2,2] $  }& $[5,1^5]\times [2,2]\rightarrow      (A_3+2A_1)''$   &   $E_7(a_2)$        \\


&  \small{$ [7,1^3]\times [1^4] $  }& $[3,1^7]\times [4]\rightarrow      (A_3+2A_1)''$  &   $E_7(a_2)$        \\
& \small{ $ [5,5]\times [1^4] $ }& $[2^4,1,1]\times [4]\rightarrow      (A_3+2A_1)'$  &   $E_8(a_6)$        \\
\hline
$ 113$  
&  $ [3,3,1^4]\times [4] $ & \small{ $[5,3,1,1]\times [1^4]\rightarrow      D_4(a_1)$ } &   $E_8(b_5)$        \\
&  $ [3,3,2,2]\times [3,1] $ & $[4,4,1,1]\times [2,1,1]\rightarrow      A_3+A_1$  &   $E_8(b_5)$        \\


&{  $ [5,1^5]\times [3,1] $}& $[5,1^5]\times [2,1,1]\rightarrow      A_3+A_1$ &   $E_8(b_5)$        \\


&  {$ [4,4,1,1]\times [2,1,1] $ }& $[3,3,2,2]\times [3,1]\rightarrow      2A_2+A_1$  &   $E_8(b_5)$        \\
&{ $ [7,3]\times [1^4] $ }& $[2,2,1^6]\times [4]\rightarrow      A_3+A_1$  &   $E_8(b_5)$        \\
&  $ [3,3,3,1]\times [2,2] $ & $[3,3,3,1]\times [2,2]\rightarrow      A_2+4A_1$  &   $D_7$        \\
\hline

$ 114$  
&  $ [3,3,2,2]\times [4] $ & { $[4,4,1,1]\times [1^4]\rightarrow     A_3$ } &   $E_7(a_1)$        \\

&{  $ [5,1^5]\times [4] $}& $[5,1^5]\times [1^4]\rightarrow      A_3$ &   $E_7(a_1)$        \\
&  $ [3,3,3,1]\times [3,1] $ & $[3,3,3,1]\times [2,1,1]\rightarrow      A_2+3A_1$   &   $E_8(a_5)$        \\

&  {$ [4,4,1,1]\times [2,2] $ }& $[3,3,2,2]\times [2,2]\rightarrow      A_2+3A_1$  &   $E_8(a_5)$        \\
&{ $ [5,3,1,1]\times [2,1,1] $ }& $[3,3,1^4]\times [3,1]\rightarrow      2A_2$  &   $E_8(a_5)$        \\
&  $ [9,1]\times [1^4] $ & $[1^{10}]\times [4]\rightarrow      A_3$  &   $E_7(a_1)$        \\
\hline

$ 116$  
&  {$ [4,4,1,1]\times [4] $ }& $[3,3,2,2]\times [1^4]\rightarrow      A_2+A_1$  &   $E_8(a_4)$        \\
&{ $ [5,3,1,1]\times [3,1] $ }& $[3,3,1^4]\times [2,1,1]\rightarrow      A_2+A_1$  &   $E_8(a_4)$        \\
% &  $ [7,1^3]\times [2,2] $ or $([4,1],\emptyset)\times[2,2]$ & $[3,1^{7}]\times [2,2]\rightarrow      4A_1$ \color{red}{?} &   $E_8(a_4)$        \\
&  $ [7,1^3]\times [2,2] $  & $[3,1^{7}]\times [2,2]\rightarrow      (4A_1)''$   &   $E_7$        \\
&  $ [5,5]\times [2,2] $  & $[2^{4},1,1]\times [2,2]\rightarrow      (4A_1)'$  &   $E_8(a_4)$        \\


&  $ [7,3]\times [2,1,1] $ & $[2^{2},1^6]\times [3,1]\rightarrow      A_2+A_1$  &   $E_8(a_4)$        \\
\hline

  
\end{tabular}
	\bigskip
	
\end{table}

\begin{table}[htbp]\caption{GKdim and AV for $H:=E_6 \times A_2 \subset E_8$ }
	\label{tabe6a2}
	\centering
	\renewcommand{\arraystretch}{1.5}
	\setlength\tabcolsep{5pt}

	\begin{tabular}{|c|c|c|c|c|}
		%\begin{tabular}{|c|c|c|c|}
		\hline $\mathrm{GKd} L(\lambda)$&${\pi_{w_{\lam}}}$ or $\mathcal{O}(\pi_{w_{\lam}})$& $d_{\rm LS}^{H}(\mathcal{O}(\pi_{w_\lambda}))\rightarrow O_{L}$& $\mathcal{O}_{{\rm Ann}(L(\lambda))}$ \\ \hline
          $ 84 $    &  $ 1_p'\times[3] $ or $1\times [3]$ &   $1_p\times[1^3]\rightarrow E_6 $     &   $ D_4$      \\ \hline
        $92$ &  $6_p'\times [1^3] $ or $A_1\times [1^3]$ &  $6_p\times[3]\rightarrow E_6(a_1)+A_2 $     &   $ D_4(a_1)+A_2$      \\ \hline
         % &  $ [5, 2^4, 1^3] $  & $[7,5,1^4]\rightarrow     ? $     &   $ E_8(a_7)$      \\
      $94$   &  $ 6_p'\times [2,1]$ or $A_1\times [2,1]$ & $6_p\times[2,1]\rightarrow E_6(a_1)+A_1$     &   $A_4+A_1$      \\ \hline
     $99$    &  $20_p'\times [2,1] $ or $2A_1\times [2,1]$ & $20_p\times[2,1]\rightarrow  D_5+A_1 $     &   $ E_6(a_3)$      \\ \hline
    $100$     & $20_p'\times [3] $ or $2A_1\times [3]$  & $ 20_p\times[1^3]\rightarrow  D_5 $     &   $ D_5$   \\\hline

         $ 102 $    &  $30_p'\times [1^3] $ or $A_2\times [1^3]$ &  $30_p\times[3]\rightarrow  E_6(a_3)+A_2 $     &   $ E_6(a_3)+A_1$      \\ 
          % &  $15_q'\times [1^3] $ or $A_2\times [1^3]$ &  $E_6(a_3)+A_2 $     &   $ E_6(a_3)+A_1$  \\
          % & $15_p'\times [1^3] $ or $A_2\times [1^3]$ &  $E_6(a_3)+A_2 $     &   $ E_6(a_3)+A_1$  \\
          \hline

          $ 104 $    & $30_p'\times [2,1] $ or $A_2\times [2,1]$ &  $30_p\times[2,1]\rightarrow E_6(a_3)+A_1 $     &   $ E_8(a_7)$   \\ 
          &  $64_p'\times [1^3] $ or $(A_2+A_1)\times [1^3]$ &  $64_p\times[3]\rightarrow D_5(a_1)+A_2 $     &   $ E_8(a_7)$   \\   \hline
	
 $ 105 $    & $30_p'\times [3] $ or $A_2\times [3]$ &  $30_p\times[1^3]\rightarrow E_6(a_3) $     &   $ D_6(a_1)$     \\ 
         &  $24_p'\times [1^3] $ or $2A_2\times [1^3]$ &  $24_p\times[3]\rightarrow D_4+A_2 $     &   $ A_6$    \\ \hline
 $ 106 $    & $64_p'\times [2,1] $ or $(A_2+A_1)\times [2,1]$ &  $64_p\times[2,1]\rightarrow D_5(a_1)+A_1 $     &   $ E_7(a_4)$     \\ 
         &  $60_p'\times [1^3] $ or $(A_2+2A_1)\times [1^3]$ &  $60_p\times[3]\rightarrow A_4+A_2+A_1 $     &   $ A_6+A_1$    \\ \hline
 $ 107 $    & $64_p'\times [3] $ or $A_2+A_1\times [3]$ &  $64_p\times[1^3]\rightarrow D_5(a_1) $     &   $ E_6(a_1)$     \\ 
         &  $24_p'\times [2,1] $ or $2A_2\times [2,1]$ &  $24_p\times[2,1]\rightarrow D_4+A_1 $     &   $ E_6(a_1)$    \\ 
      &  $81_p'\times [1^3] $ or $A_3\times [1^3]$ &  $81_p\times[3]\rightarrow A_4+A_2 $     &   $ D_5+A_2$    \\   
         \hline
 
 $ 108 $    & $24_p'\times [3] $ or $2A_2\times [3]$ &  $24_p\times[1^3]\rightarrow D_4 $     &   $ E_6$     \\ 
         &  $60_p'\times [2,1] $ or $(A_2+2A_1)\times [2,1]$ &  $60_p\times[2,1]\rightarrow A_4+2A_1 $     &   $ D_7(a_2)$    \\  
         \hline

  $ 109 $    & $60_p'\times [3] $ or $(A_2+2A_1)\times [3]$ &  $60_p\times[1^3]\rightarrow A_4+A_1 $     &   $ E_6(a_1)+A_1$     \\ 
         &  $81_p'\times [2,1] $ or $A_3\times [2,1]$ &  $81_p\times[2,1]\rightarrow A_4+A_1 $     &   $ E_6(a_1)+A_1$    \\  
         \hline  

$ 110 $    & $81_p'\times [3] $ or $A_3\times [3]$ &  $81_p\times[1^3]\rightarrow A_4 $     &   $ E_7(a_3)$     \\ 
         % &  $90_s\times [1^3] $ or $D_4(a_1)\times [1^3]$ &  $D_4(a_1)+A_2 $     &   $ E_8(b_6)$    \\  
         &  $80_s\times [1^3] $ or $D_4(a_1)\times [1^3]$ &  $80_s\times[3]\rightarrow D_4(a_1)+A_2 $     &   $ E_8(b_6)$    \\
         % &  $60_s\times [1^3] $ or $D_4(a_1)\times [1^3]$ &  $D_4(a_1)+A_2 $     &   $ E_8(b_6)$    \\
         % &  $20_s\times [1^3] $ or $D_4(a_1)\times [1^3]$ &  $D_4(a_1)+A_2 $     &   $ E_8(b_6)$    \\
         % &  $10_s\times [1^3] $ or $D_4(a_1)\times [1^3]$ &  $D_4(a_1)+A_2$     &   $ E_8(b_6)$    \\
         \hline

$ 111 $    & $81_p\times [1^3] $ or $A_4\times [1^3]$ &  $81_{p}'\times[3]\rightarrow A_3+A_2 $     &   $ D_7(a_1)$     \\ 
         &  $24_p\times [1^3] $ or $D_4\times [1^3]$ &  $24_{p}'\times[1^3]\rightarrow 3A_2 $     &   $ E_6+A_1$    \\  
         \hline
   $ 112 $  
    % &  $90_s\times [2,1] $ or $D_4(a_1)\times [2,1]$ &  $D_4(a_1)+A_1 $     &   $ E_8(a_6)$    \\  
         &  $80_s\times [2,1] $ or $D_4(a_1)\times [2,1]$ &  $80_s\times[2,1]\rightarrow D_4(a_1)+A_1 $     &   $ E_8(a_6)$  \\
         % &  $60_s\times [1^3] $ or $D_4(a_1)\times [2,1]$ &  $D_4(a_1)+A_1 $     &   $ E_8(a_6)$  \\
         % &  $20_s\times [1^3] $ or $D_4(a_1)\times [2,1]$ &  $D_4(a_1)+A_1 $     &   $ E_8(a_6)$  \\
         % &  $10_s\times [1^3] $ or $D_4(a_1)\times [2,1]$ &  $D_4(a_1)+A_1$     &   $ E_8(a_6)$ \\
         &  $60_p\times [1^3] $ or $(A_4+A_1)\times [1^3]$ &  $60_{p}'\times[3]\rightarrow 2A_2+2A_1 $     &   $ E_8(a_6)$ \\
         \hline
     
         
 \end{tabular}
	\bigskip
	
\end{table}


\begin{table}[htbp]\caption{GKdim and AV for $H:=E_6 \times A_2 \subset E_8$ (continued)}
	\label{tabe6a2-2}
	\centering
	\renewcommand{\arraystretch}{1.5}
	\setlength\tabcolsep{5pt}
	
	\begin{tabular}{|c|c|c|c|c|}
		%\begin{tabular}{|c|c|c|c|}
		\hline $\mathrm{GKd} L(\lambda)$&${\pi_{w_{\lam}}}$ or $\mathcal{O}(\pi_{w_{\lam}})$& $d_{\rm LS}^{H}(\mathcal{O}(\pi_{w_\lambda}))\rightarrow O_{L}$& $\mathcal{O}_{{\rm Ann}(L(\lambda))}$ \\ \hline
  
       
$ 113 $ 
% &  $90_s\times [3] $ or $D_4(a_1)\times [3]$ &  $D_4(a_1) $     &   $ E_8(b_5)$    \\  
         &  $80_s\times [3] $ or $D_4(a_1)\times [3]$ &  $80_s\times[1^3]\rightarrow D_4(a_1) $     &   $ E_8(b_5)$  \\
         % &  $60_s\times [3] $ or $D_4(a_1)\times [3]$ &  $D_4(a_1) $     &   $ E_8(b_5)$  \\
         % &  $20_s\times [3] $ or $D_4(a_1)\times [3]$ &  $D_4(a_1) $     &   $ E_8(b_5)$  \\
         % &  $10_s\times [3] $ or $D_4(a_1)\times [3]$ &  $D_4(a_1)$     &   $ E_8(b_5)$ \\
         &  $81_p\times [2,1] $ or $A_4\times [2,1]$ &  $81_{p}'\times[2,1]\rightarrow A_3+A_1 $     &   $ E_8(b_5)$ \\
        
         &  $24_p\times [2,1] $ or $D_4\times [2,1]$ &  $24_{p}'\times[2,1]\rightarrow 2A_2+A_1 $     &   $ E_8(b_5)$ \\
          &  $64_p\times [1^3] $ or $D_5(a_1)\times [1^3]$ &  $64_{p}'\times[3]\rightarrow 2A_2+A_1 $     &   $ E_8(b_5)$ \\
         \hline
$ 114 $  &  $81_p\times [3] $ or $A_4\times [3]$ &  $81_{p}'\times[1^3]\rightarrow A_3 $     &   $ E_7(a_1)$    \\  
         &  $24_p\times [3] $ or $D_4\times [3]$ &  $24_{p}'\times[1^3]\rightarrow 2A_2 $     &   $ E_8(a_5)$  \\
         &  $60_p\times [2,1] $ or $(A_4+A_1)\times [2,1]$ &  $60_{p}'\times[2,1]\rightarrow A_2+3A_1 $     &   $ E_8(a_5)$  \\
         &  $30_p\times [1^3] $ or $E_6(a_3)\times [1^3]$ &  $30_{p}'\times[3]\rightarrow 2A_2 $     &   $ E_8(a_5)$  \\
         % &  $15_q\times [1^3] $ or $E_6(a_3)\times [1^3]$ &  $2A_2 $     &   $ E_8(a_5)$\\
         % &  $15_p\times [1^3] $ or $E_6(a_3)\times [1^3]$ &  $2A_2 $     &   $ E_8(a_5)$\\
         \hline
        
       $ 116 $  &  $64_p\times [3] $ or $D_5(a_1)\times [3]$ &  $64_{p}'\times[1^3]\rightarrow A_2+A_1 $     &   $ E_8(a_4)$    \\  
        
         &  $30_p\times [2,1] $ or $E_6(a_3)\times [2,1]$ &  $30_{p}'\times[2,1]\rightarrow A_2 +A_1$     &   $ E_8(a_4)$  \\
         % &  $15_q\times [2,1] $ or $E_6(a_3)\times [2,1]$ &  $A_2 +A_1$     &   $ E_8(a_4)$\\
         % &  $15_p\times [2,1] $ or $E_6(a_3)\times [2,1]$ &  $A_2 +A_1$     &   $ E_8(a_4)$\\
         &  $6_p\times [1^3] $ or $E_6(a_1)\times [1^3]$ &  $6_{p}'\times[3]\rightarrow A_2 +A_1$     &   $ E_8(a_4)$\\
         \hline
\end{tabular}
	\bigskip
	
\end{table}



\begin{table}[htbp]\caption{ GKdim and AV for $H:=E_7\times A_1\subset E_8$} 
	\label{tabe7a1}
	\centering
	\renewcommand{\arraystretch}{1.5}
	\setlength\tabcolsep{5pt}
	
 \small  \begin{tabular}{|c|c|c|c|c|}
		%\begin{tabular}{|c|c|c|c|}
		\hline $\mathrm{GKd} L(\lambda)$ & ${\pi_{w_{\lam}}}$ or    $\mathcal{O}({\pi_{w_{\lam}}})$& $d_{\rm LS}^{H}(\mathcal{O}(\pi_{w_\lambda}))\rightarrow  O_{L}$ & $\mathcal{O}_{{\rm Ann}(L(\lambda))}$ \\ \hline
	$ 84 $    &  $21_b\times [2] $ or $(3A_1)''\times [2] $ &  $ 21_{b}'\times [1,1] \rightarrow E_6 $     &   $ D_4$      \\ \hline	
   
   
   $ 94$    &  $ 120_a'\times [1,1]$  or $(A_2+A_1)\times [1,1]$&   $120_a\times [2] \rightarrow E_6(a_1)+A_1$     &   $ A_4+A_1$      \\ 
   % &  $ 105_a\times [1,1]$  or $(A_2+A_1)\times [1,1]$&   $E_6(a_1)+A_1$     &   $ A_4+A_1$      \\
   % &  $ 15_a\times [1,1]$  or $(A_2+A_1)\times [1,1]$&    $E_6(a_1)+A_1$     &   $ A_4+A_1$      \\   
   \hline

$ 98$    &  $ 189_b\times [2]$  or $(A_2+2A_1)\times [2]$&     $189_{b}'\times [1,1] \rightarrow E_7(a_4)$     &   $ D_5(a_1)+A_1$      \\
&  $ 210_a'\times [1,1]$  or $A_3\times [1,1]$&   $210_a\times [2] \rightarrow D_6(a_1)+A_1$     &   $ D_5(a_1)+A_1$      \\ 
   &  $ 168_a'\times [1,1]$  or $2A_2\times [1,1]$&   $168_a\times [2] \rightarrow D_5+2A_1$     &   $ A_5$      \\
   &  $ 105_b'\times [1,1]$  or $(A_2+3A_1)\times [1,1]$&    $105_b\times [2] \rightarrow A_6+A_1$     &   $ A_4+A_2+A_1$      \\
  \hline

$ 99$    &  $ 189_c\times [1,1]$  or $(A_3+A_1)''\times [1,1]$&     $189_c'\times [2] \rightarrow D_5+A_1$     &   $ E_6(a_3)$      \\
&  $ 210_a'\times [2]$  or $A_3\times [2]$&   $210_a\times [1,1] \rightarrow D_6(a_1)$     &   $ E_6(a_3)$      \\ 
   &  $ 168_a'\times [2]$  or $2A_2\times [2]$&   $168_a\times [1,1] \rightarrow D_5+A_1$     &   $ E_6(a_3)$      \\
   &  $ 105_b'\times [2]$  or $(A_2+3A_1)\times [2]$&    $105_b\times [1,1] \rightarrow A_6$     &   $ D_4+A_2$      \\
  \hline

$ 100$    &  $ 189_c\times [2]$  or $(A_3+A_1)''\times [2]$&     $189_c'\times [1,1] \rightarrow D_5$     &   $ D_5$      \\
  \hline

$ 104$    
   &  $ 105_c'\times [1,1]$  or $D_4\times [1,1]$&    $105_c\times [2] \rightarrow (A_5+A_1)''$     &   $ D_5+A_1$      \\
   &  $ 405_a'\times [1,1]$  or $(D_4(a_1)+A_1)\times [1,1]$&    $405_a\times [2] \rightarrow E_6(a_3)+A_1$     &   $ E_8(a_7)$      \\
 % &  $ 189_a'\times [1,1]$  or $(D_4(a_1)+A_1)\times [1,1]$&    $E_6(a_3)+A_1$     &   $ E_8(a_7)$      \\
 % &  $ 216_a\times [1,1]$  or $(D_4(a_1)+A_1)\times [1,1]$&    $E_6(a_3)+A_1$     &   $ E_8(a_7)$      \\
  &  $ 315_a\times [2]$  or $D_4(a_1)\times [2]$&    $315_a'\times [1,1] \rightarrow E_7(a_5)$ &   $ E_8(a_7)$      \\
  % &  $ 280_a\times [2]$  or $D_4(a_1)\times [2]$&    $E_7(a_5)$  &   $ E_8(a_7)$      \\
 
  % &  $ 70_a\times [2]$  or $D_4(a_1)\times [2]$&    $E_7(a_5)$ &   $ E_8(a_7)$      \\
  %  &  $ 35_a\times [2]$  or $D_4(a_1)\times [2]$&   $E_7(a_5)$  &   $ E_8(a_7)$      \\
  % &  $ 280_b'\times [2]$  or $D_4(a_1)\times [2]$&    $E_7(a_5)$  &   $ E_8(a_7)$      \\
  \hline

$ 105$  
% &  $ 216_a\times [2]$  or $(D_4(a_1)+A_1)\times [2]$&    $E_6(a_3)$     &   $ D_6(a_1)$      \\
&  $ 405_a'\times [2]$  or $(D_4(a_1)+A_1)\times [2]$&    $405_a\times [1,1] \rightarrow E_6(a_3)$     &   $ D_6(a_1)$      \\

 % &  $ 189_a'\times [2]$  or $(D_4(a_1)+A_1)\times [2]$&    $E_6(a_3)$     &   $ D_6(a_1)$      \\
     &  $ 105_c'\times [2]$  or $(D_4(a_1)+A_1)\times [2]$&    $105_c\times [1,1] \rightarrow E_6(a_3)$     &   $ D_6(a_1)$      \\
  \hline

  $ 106$   
  % &  $ 336_a\times [1,1]$  or $A_4\times [1,1]$&     $D_5(a_1)+A_1$    &   $ E_7(a_4)$      \\
&  $ 420_a'\times [1,1]$ or $A_4\times [1,1]$&     $420_a\times [2] \rightarrow D_5(a_1)+A_1$    &   $ E_7(a_4)$      \\
% &  $ 84_a'\times [1,1]$ or $A_4\times [1,1]$&     $D_5(a_1)+A_1$    &   $ E_7(a_4)$      \\

 &  $ 210_b'\times [1,1]$  or $(A_3+A_2+A_1)\times [1,1]$&    $210_b\times [2] \rightarrow A_4+A_2+A_1$     &   $ A_6+A_1$      \\
  \hline
$ 107$   
% &  $ 336_a\times [2]$  or $A_4\times [2]$&     $D_5(a_1)$    &   $ E_6(a_1)$      \\
&  $ 420_a'\times [2]$ or $A_4\times [2]$&    $420_a\times [1,1] \rightarrow D_5(a_1)$  &   $ E_6(a_1)$      \\
% &  $ 84_a'\times [2]$ or $A_4\times [2]$&   $D_5(a_1)$    &   $ E_6(a_1)$      \\
&  $ 105_c\times [1,1]$ or $(A_5)''\times [1,1]$&    $105_{c'}\times [2] \rightarrow D_4+A_1$     &   $ E_6(a_1)$      \\

 &  $ 210_b'\times [2]$  or $(A_3+A_2+A_1)\times [2]$&    $210_b\times [1,1] \rightarrow A_4+A_2$     &   $ D_5+A_2$      \\
  \hline

$ 108$  
&  $ 105_c\times [2]$ or $(A_5)''\times [2]$&    $105_{c}'\times [1,1] \rightarrow D_4$     &   $ E_6$      \\
% &  $ 512_a\times [1,1]$  or $(A_4+A_1)\times [1,1]$&    $A_4+2A_1$     &   $ D_7(a_2)$      \\
 &  $ 512_a'\times [1,1]$  or $(A_4+A_1)\times [1,1]$&    $512_a'\times [2] \rightarrow A_4+2A_1$     &   $ D_7(a_2)$      \\
  \hline

$ 109$  
 &  $ 512_a'\times [2]$  or $(A_4+A_1)\times [2]$&    $512_a'\times [1,1] \rightarrow A_4+A_1$     &   $ E_6(a_1)+A_1$      \\
 &  $ 420_a\times [1,1]$ or $D_5(a_1)\times [1,1]$&    $420_a'\times [2] \rightarrow A_4+A_1$     &   $ E_6(a_1)+A_1$      \\
 % &  $ 336_a'\times [1,1]$ or $D_5(a_1)\times [1,1]$&    $A_4+A_1$     &   $ E_6(a_1)+A_1$      \\
 % &  $ 84_a\times [1,1]$ or $D_5(a_1)\times [1,1]$&    $A_4+A_1$     &   $ E_6(a_1)+A_1$      \\
 &  $ 210_b\times [1,1]$  or $(A_4+A_2)\times [1,1]$&    $210_b'\times [2] \rightarrow A_3+A_2+2A_1$     &   $ A_7$      \\
  \hline




	\end{tabular}
	\bigskip
	
\end{table}


\begin{table}[htbp]\caption{ GKdim and AV for $H:=E_7\times A_1\subset E_8$ (continued)} 
	\label{tabe7a1-2}
	\centering
	\renewcommand{\arraystretch}{1.5}
	\setlength\tabcolsep{5pt}
	
\small	\begin{tabular}{|c|c|c|c|c|}
		%\begin{tabular}{|c|c|c|c|}
		\hline $\mathrm{GKd} L(\lambda)$ & ${\pi_{w_{\lam}}}$ or    $\mathcal{O}({\pi_{w_{\lam}}})$& $d_{\rm LS}^{H}(\mathcal{O}(\pi_{w_\lambda}))\rightarrow O_{L}$ & $\mathcal{O}_{{\rm Ann}(L(\lambda))}$ \\ \hline





$ 110$  
 &  $ 420_a\times [2]$ or $D_5(a_1)\times [2]$&    $420_a'\times [1,1] \rightarrow A_4$     &   $ E_7(a_3)$      \\
 % &  $ 336_a'\times [2]$ or $D_5(a_1)\times [2]$&    $A_4$     &   $ E_7(a_3)$         \\
 % &  $ 84_a\times [2]$ or $D_5(a_1)\times [2]$&    $A_4$     &   $ E_7(a_3)$      \\
 
 &  $ 210_b\times [2]$  or $(A_4+A_2)\times [2]$&    $210_b'\times [1,1] \rightarrow A_3+A_2+A_1$     &   $ E_8(b_6)$      \\

 &  $ 378_a'\times [1,1]$  or $(D_5(a_1)+A_1)\times [1,1]$&    $378_a\times [2] \rightarrow A_3+A_2+A_1$     &   $ E_8(b_6)$      \\
  \hline

$ 111$  
 &  $ 378_a'\times [2]$  or $(D_5(a_1)+A_1)\times [2]$&    $378_a\times [1,1] \rightarrow A_3+A_2$     &   $ D_7(a_1)$      \\
 &  $ 405_a\times [1,1]$  or $E_6(a_3)\times [1,1]$&     $405_a'\times [2] \rightarrow D_4(a_1)+2A_1$  &   $ D_7(a_1)$      \\
 % &  $ 189_a\times [1,1]$  or $E_6(a_3)\times [1,1]$&     $D_4(a_1)+2A_1$   &   $ D_7(a_1)$      \\
 % &  $ 216_a'\times [1,1]$  or $E_6(a_3)\times [1,1]$&     $D_4(a_1)+2A_1$  &   $ D_7(a_1)$      \\
  \hline
$ 112$  
 &  $ 405_a\times [2]$  or $E_6(a_3)\times [2]$&     $405_a'\times [1,1] \rightarrow D_4(a_1)+A_1$  &   $ E_8(a_6)$      \\
 % &  $ 189_a\times [2]$  or $E_6(a_3)\times [2]$&     $D_4(a_1)+A_1$ &   $ E_8(a_6)$      \\
 % &  $ 216_a'\times [2]$  or $E_6(a_3)\times [2]$&    $D_4(a_1)+A_1$ &   $ E_8(a_6)$      \\
&  $ 315_a'\times [1,1]$  or $E_7(a_5)\times [1,1]$&     $315_a\times [2] \rightarrow D_4(a_1)+A_1$  &   $ E_8(a_6)$      \\
 % &  $ 280_a'\times [1,1]$  or $E_7(a_5)\times [1,1]$&     $D_4(a_1)+A_1$   &   $ E_8(a_6)$      \\
 % &  $ 70_a'\times [1,1]$  or $E_7(a_5)\times [1,1]$&     $D_4(a_1)+A_1$  &   $ E_8(a_6)$      \\
 % &  $ 35_a'\times [1,1]$  or $E_7(a_5)\times [1,1]$&     $D_4(a_1)+A_1$  &   $ E_8(a_6)$      \\
 % &  $ 280_b\times [1,1]$  or $E_7(a_5)\times [1,1]$&     $D_4(a_1)+A_1$  &   $ E_8(a_6)$      \\
 &  $ 189_c'\times [1,1]$  or $D_5\times [1,1]$&    $189_c\times [2] \rightarrow (A_3+2A_1)''$   &   $ E_7(a_2)$      \\
  \hline




% \end{tabular}
% 	\bigskip
% 	\caption{ GKdim and AV for $H:=E_7\times A_1\subset E_8$ (continued)} 
% 	\label{tabe7a1-2}
% \end{table}
% \begin{table}[htbp]
% 	\centering
% 	\renewcommand{\arraystretch}{1.5}
% 	\setlength\tabcolsep{5pt}
	
% 	\begin{tabular}{|c|c|c|c|c|}
% 		%\begin{tabular}{|c|c|c|c|}
% 		\hline $\mathrm{GKd} L(\lambda)$ & ${\pi_{w_{\lam}}}$ or    $\mathcal{O}({\pi_{w_{\lam}}})$& $d_{\rm LS}^{H}(\mathcal{O}(\pi_{w_\lambda}))=    L$ & $\mathcal{O}_{{\rm Ann}(L(\lambda))}$ \\ \hline

$ 113$  
&  $ 315_a'\times [2]$  or $E_7(a_5)\times [2]$&    $315_a\times [1,1] \rightarrow D_4(a_1)$  &   $ E_8(b_5)$      \\
 % &  $ 280_a'\times [2]$  or $E_7(a_5)\times [2]$&    $D_4(a_1)$  &   $ E_8(b_5)$      \\
 % &  $ 70_a'\times [2]$  or $E_7(a_5)\times [2]$&    $D_4(a_1)$   &   $ E_8(b_5)$      \\
 % &  $ 35_a'\times [2]$ or $E_7(a_5)\times [2]$&    $D_4(a_1)$   &   $ E_8(b_5)$      \\
 % &  $ 280_b\times [2]$ or $E_7(a_5)\times [2]$&    $D_4(a_1)$  &   $ E_8(b_5)$      \\
 &  $ 189_c'\times [2]$  or $D_5\times [2]$&    $189_c\times [1,1] \rightarrow (A_3+A_1)''$  &   $ E_8(b_5)$      \\

 &  $ 210_a\times [1,1]$  or $D_6(a_1)\times [1,1]$&    $210_a'\times [2] \rightarrow A_3+A_1$  &   $ E_8(b_5)$      \\
 &  $ 168_a\times [1,1]$  or $(D_5+A_1)\times [1,1]$&    $168_a'\times [2] \rightarrow 2A_2+A_1$  &   $ E_8(b_5)$      \\
 &  $ 105_b\times [1,1]$  or $A_6\times [1,1]$&    $105_b'\times [2] \rightarrow A_2+4A_1$  &   $ D_7$      \\
  \hline

$ 114$  
&  $ 210_a\times [2]$  or $D_6(a_1)\times [2]$&    $210_a'\times [1,1] \rightarrow A_3$  &   $ E_7(a_1)$      \\
 &  $ 168_a\times [2]$  or $(D_5+A_1)\times [2]$&    $168_a'\times [1,1] \rightarrow 2A_2$  &   $ E_8(a_5)$      \\
 &  $ 105_b\times [2]$  or $A_6\times [2]$&    $105_b'\times [1,1] \rightarrow A_2+3A_1$  &   $ E_8(a_5)$      \\

 &  $ 189_b'\times [1,1]$  or $E_7(a_4)\times [1,1]$&    $189_b\times [2] \rightarrow A_2+3A_1$  &   $ E_8(a_5)$      \\
  \hline

$ 116$  
&  $ 120_a\times [2]$  or $E_6(a_1)\times [2]$&    $120_a'\times [1,1] \rightarrow A_2+A_1$ &   $ E_8(a_4)$      \\
 % &  $ 105_a'\times [2]$  or $E_6(a_1)\times [2]$&    $A_2+A_1$ &   $ E_8(a_4)$      \\
 % &  $ 15_a'\times [2]$  or $E_6(a_1)\times [2]$&    $A_2+A_1$ &   $ E_8(a_4)$      \\

 &  $ 56_a'\times [1,1]$  or $E_7(a_3)\times [1,1]$&    $56_a\times [2] \rightarrow A_2+A_1$  &   $ E_8(a_4)$      \\
 % &  $ 35_b\times [1,1]$  or $E_7(a_3)\times [1,1]$&    $A_2+A_1$  &   $ E_8(a_4)$      \\
 % &  $ 21_a\times [1,1]$  or $E_7(a_3)\times [1,1]$&    $A_2+A_1$  &   $ E_8(a_4)$      \\
 &  $ 21_b'\times [1,1]$  or $E_6\times [1,1]$&    $21_b\times [2] \rightarrow (4A_1)''$  &   $ E_7$      \\
  \hline




  
\end{tabular}
	\bigskip
	
\end{table}



% \bibliographystyle{plain}
 \bibliographystyle{alpha}
 \bibliography{BGWX}
%\printbibliography


\end{document}

